
\documentclass[11pt,a4paper]{article}

% ================================================================
% TD de classe préparatoire scientifique
% Analyse 3 -- Fonctions réelles : limites, continuité, TVI, bijections
% Document autonome prêt à compiler
% ================================================================

\usepackage[T1]{fontenc}
\usepackage[utf8]{inputenc}
\usepackage[french]{babel}
\usepackage{lmodern}
\usepackage{microtype}
\usepackage{amsmath,amssymb,amsthm,mathtools}
\usepackage{enumitem}
\usepackage[most]{tcolorbox}
\usepackage{geometry}
\usepackage{fancyhdr}
\usepackage{lastpage}
\usepackage{xcolor}
\usepackage{hyperref}
\usepackage{array,booktabs,multicol}
\usepackage{tikz}
\usetikzlibrary{arrows.meta,calc}

\geometry{left=1.65cm,right=1.65cm,top=1.75cm,bottom=1.9cm}
\setlength{\parindent}{0pt}
\setlength{\parskip}{0.45em}
\setlength{\headheight}{14pt}
\setlist[enumerate]{label=\textbf{\arabic*.},leftmargin=*,itemsep=0.35em,topsep=0.25em}
\setlist[itemize]{leftmargin=*,itemsep=0.2em,topsep=0.15em}
\hypersetup{colorlinks=true,linkcolor=blue!50!black,urlcolor=blue!50!black,citecolor=blue!50!black}

\definecolor{bleuprepa}{RGB}{16,56,115}
\definecolor{vertprepa}{RGB}{20,105,70}
\definecolor{rougeprepa}{RGB}{150,36,36}
\definecolor{orangeprepa}{RGB}{190,105,20}
\definecolor{grisprepa}{RGB}{245,247,250}
\definecolor{violetprepa}{RGB}{90,55,140}

\pagestyle{fancy}
\fancyhf{}
\lhead{\textsc{TD -- Analyse 3}}
\chead{\textsc{Fonctions réelles}}
\rhead{\textsc{MPSI/PCSI/MP}}
\cfoot{Page \thepage\ sur \pageref{LastPage}}
\renewcommand{\headrulewidth}{0.4pt}
\renewcommand{\footrulewidth}{0.4pt}

\newcommand{\etoiles}[1]{\hfill\textcolor{orangeprepa}{\textbf{#1}}}
\newcommand{\R}{\mathbb{R}}
\newcommand{\N}{\mathbb{N}}
\newcommand{\e}{\mathrm{e}}
\newcommand{\dd}{\,\mathrm{d}}
\newcommand{\Id}{\mathrm{Id}}
\newcommand{\dom}{\mathcal D}
\newcommand{\indic}[1]{\mathbf 1_{#1}}
\newcommand{\suite}[1]{(#1_n)_{n\in\N}}
\newcommand{\limite}[2]{\lim_{#1} #2}
\newcommand{\fonction}[5]{#1:#2\longrightarrow #3,\quad #4\longmapsto #5}

\tcbset{enhanced,breakable,sharp corners,boxrule=0.7pt,fonttitle=\bfseries}
\newtcolorbox{objectifbox}{colback=blue!2!white,colframe=bleuprepa,title=Objectifs du TD}
\newtcolorbox{methodbox}{colback=green!2!white,colframe=vertprepa,title=Méthode à retenir}
\newtcolorbox{attentionbox}{colback=red!2!white,colframe=rougeprepa,title=Point de vigilance}
\newtcolorbox{exobox}[2][]{colback=white,colframe=bleuprepa,title={#2},#1}
\newtcolorbox{corrbox}[2][]{colback=grisprepa,colframe=vertprepa,title={#2},#1}
\newtcolorbox{problembox}[2][]{colback=white,colframe=violetprepa,title={#2},#1}
\newtcolorbox{fichebox}{colback=orange!3!white,colframe=orangeprepa,title=Fiche méthode}

\begin{document}

\begin{titlepage}
\begin{tikzpicture}[remember picture,overlay]
  \fill[bleuprepa] (current page.north west) rectangle ([yshift=-4.2cm]current page.north east);
  \fill[grisprepa] ([yshift=-4.2cm]current page.north west) rectangle (current page.south east);
\end{tikzpicture}
\vspace*{0.6cm}
\begin{center}
{\color{white}\Huge\bfseries Analyse 3 -- TD complet de fonctions réelles}\\[0.3cm]
{\color{white}\Large Limites, continuité, théorème des valeurs intermédiaires, bijections}\\[0.4cm]
{\color{white}\large Niveau MPSI -- PCSI -- MP}\\[2.1cm]
\begin{tcolorbox}[width=0.88\textwidth,colback=white,colframe=bleuprepa,arc=2mm,boxrule=1pt]
\centering
\Large \textbf{Polycopié d'exercices corrigés}\\[0.25cm]
\normalsize 50 exercices progressifs, problèmes longs type concours et corrigé entièrement rédigé.\\[0.2cm]
\normalsize Document original construit pour l'entraînement en classe préparatoire scientifique.
\end{tcolorbox}
\vfill
{\large \today}
\end{center}
\end{titlepage}

\tableofcontents
\newpage

\section*{Avant-propos}
\addcontentsline{toc}{section}{Avant-propos}
Ce document constitue un TD autonome sur les fonctions réelles d'une variable réelle. Il est centré sur les notions suivantes : domaine de définition, image directe, image réciproque d'un ensemble, limites finies ou infinies, limites latérales, caractérisation séquentielle, comparaison, continuité, prolongement par continuité, théorème des valeurs intermédiaires, image d'un intervalle, théorème des bornes atteintes, injectivité, surjectivité, bijectivité, théorème de la bijection continue et fonction réciproque.

\begin{attentionbox}
Le chapitre demandé est un chapitre d'analyse. Les notions d'application linéaire, de noyau, d'image, de rang, d'injectivité et de surjectivité au sens de l'algèbre linéaire n'interviennent pas dans ce polycopié, sauf pour rappeler que les mots injectif, surjectif et bijectif ont aussi un sens général pour les fonctions. Le TD reste donc strictement dans les fonctions réelles, conformément au chapitre \og Analyse 3 -- fonctions réelles \fg{}.
\end{attentionbox}

\begin{objectifbox}
À la fin du TD, l'étudiant doit savoir :
\begin{itemize}
\item déterminer proprement un domaine de définition et une image d'intervalle ;
\item prouver une limite avec les définitions, par opérations, par encadrement ou par suites ;
\item distinguer limite, limite latérale et valeur de la fonction ;
\item décider si une fonction est continue et construire un prolongement par continuité ;
\item utiliser le théorème des valeurs intermédiaires pour une existence ;
\item utiliser la stricte monotonie pour une unicité ;
\item appliquer le théorème de la bijection continue et déterminer une réciproque.
\end{itemize}
\end{objectifbox}

\begin{fichebox}
Dans toutes les corrections, la rédaction suit le schéma d'une copie de concours : on annonce la fonction, son domaine, la continuité ou les limites utilisées, puis le théorème appliqué. Les calculs sont séparés des arguments d'existence et d'unicité.
\end{fichebox}

\newpage

\section{Énoncés des exercices}
\begin{methodbox}
Les étoiles mesurent la difficulté globale : calcul, longueur, idée à trouver, niveau de rédaction attendu et proximité avec les oraux de concours.
\end{methodbox}
\section{Exercices d'application directe}
\begin{exobox}{Exercice 1 -- Domaines élémentaires \etoiles{$\star$}}

Soit
\[
f(x)=\frac{\sqrt{2x-1}}{x^2-4}.
\]
\begin{enumerate}
\item Déterminer le domaine de définition de $f$.
\item Préciser sur quels intervalles $f$ est continue.
\item Dire si $f$ peut être prolongée par continuité en $2$.
\end{enumerate}

\medskip
\textbf{Consignes de rédaction.} Identifier le théorème utilisé, vérifier les hypothèses et conclure par une phrase complète.
\end{exobox}


\begin{exobox}{Exercice 2 -- Image d'un segment par une fonction quadratique \etoiles{$\star$}}

On considère $f(x)=x^2-2x$.
\begin{enumerate}
\item Écrire $f$ sous forme canonique.
\item Déterminer $f([0,3])$.
\item Déterminer $f^{-1}([-1,0])$ où $f^{-1}$ désigne l'image réciproque d'un ensemble.
\end{enumerate}

\medskip
\textbf{Consignes de rédaction.} Identifier le théorème utilisé, vérifier les hypothèses et conclure par une phrase complète.
\end{exobox}


\begin{exobox}{Exercice 3 -- Parité et domaine \etoiles{$\star$}}

Soit
\[
f(x)=\frac{x^3}{1+x^2}.
\]
\begin{enumerate}
\item Déterminer le domaine de définition.
\item Étudier la parité de $f$.
\item En déduire une symétrie du graphe.
\end{enumerate}

\medskip
\textbf{Consignes de rédaction.} Identifier le théorème utilisé, vérifier les hypothèses et conclure par une phrase complète.
\end{exobox}


\begin{exobox}{Exercice 4 -- Limite par factorisation \etoiles{$\star$}}

Calculer
\[
\lim_{x\to2}\frac{x^2-4}{x-2}.
\]
Rédiger la justification en indiquant pourquoi la simplification est licite.

\medskip
\textbf{Consignes de rédaction.} Identifier le théorème utilisé, vérifier les hypothèses et conclure par une phrase complète.
\end{exobox}


\begin{exobox}{Exercice 5 -- Limites latérales simples \etoiles{$\star$}}

Soit $f(x)=\frac{|x|}{x}$ définie sur $\R^*$.
\begin{enumerate}
\item Calculer la limite à droite en $0$.
\item Calculer la limite à gauche en $0$.
\item Conclure sur l'existence de la limite en $0$.
\end{enumerate}

\medskip
\textbf{Consignes de rédaction.} Identifier le théorème utilisé, vérifier les hypothèses et conclure par une phrase complète.
\end{exobox}


\begin{exobox}{Exercice 6 -- Encadrement classique \etoiles{$\star$}}

Montrer que
\[
\lim_{x\to0}x^2\cos\left(\frac1x\right)=0.
\]
La preuve doit faire apparaître explicitement le théorème utilisé.

\medskip
\textbf{Consignes de rédaction.} Identifier le théorème utilisé, vérifier les hypothèses et conclure par une phrase complète.
\end{exobox}


\begin{exobox}{Exercice 7 -- Continuité en un point défini par morceaux \etoiles{$\star$}}

On définit
\[
f(x)=\begin{cases}
\dfrac{x^2-1}{x-1} & \text{si }x\neq1,\\
3 & \text{si }x=1.
\end{cases}
\]
\begin{enumerate}
\item Calculer $\lim_{x\to1}f(x)$.
\item La fonction est-elle continue en $1$ ?
\item Quelle valeur aurait-il fallu choisir en $1$ pour obtenir la continuité ?
\end{enumerate}

\medskip
\textbf{Consignes de rédaction.} Identifier le théorème utilisé, vérifier les hypothèses et conclure par une phrase complète.
\end{exobox}


\begin{exobox}{Exercice 8 -- Prolongement de $\sin x/x$ \etoiles{$\star$}}

Soit
\[
f(x)=\begin{cases}
\dfrac{\sin x}{x} & \text{si }x\neq0,\\
a & \text{si }x=0.
\end{cases}
\]
Déterminer la valeur du réel $a$ pour laquelle $f$ est continue en $0$.

\medskip
\textbf{Consignes de rédaction.} Identifier le théorème utilisé, vérifier les hypothèses et conclure par une phrase complète.
\end{exobox}


\begin{exobox}{Exercice 9 -- TVI : existence d'un point fixe du cosinus \etoiles{$\star$}}

Montrer que l'équation
\[
\cos x=x
\]
admet au moins une solution dans $[0,1]$.

\medskip
\textbf{Consignes de rédaction.} Identifier le théorème utilisé, vérifier les hypothèses et conclure par une phrase complète.
\end{exobox}


\begin{exobox}{Exercice 10 -- Bijection continue élémentaire \etoiles{$\star$}}

Montrer que $f:x\mapsto x^3+x$ réalise une bijection de $\R$ sur $\R$.

\medskip
\textbf{Consignes de rédaction.} Identifier le théorème utilisé, vérifier les hypothèses et conclure par une phrase complète.
\end{exobox}


\section{Exercices de méthode}
\begin{exobox}{Exercice 11 -- Limite avec quantité conjuguée \etoiles{$\star\star$}}

Calculer
\[
\lim_{x\to0}\frac{\sqrt{1+x}-1}{x}.
\]
On justifiera la transformation algébrique utilisée puis on interprétera le résultat comme un prolongement par continuité.

\medskip
\textbf{Consignes de rédaction.} Identifier le théorème utilisé, vérifier les hypothèses et conclure par une phrase complète.
\end{exobox}


\begin{exobox}{Exercice 12 -- Limites et signe du dénominateur \etoiles{$\star\star$}}

Calculer les deux limites latérales en $5$ de
\[
f(x)=\frac{x^2-11x+28}{x^2-25}.
\]
On précisera le signe des facteurs qui tendent vers $0$.

\medskip
\textbf{Consignes de rédaction.} Identifier le théorème utilisé, vérifier les hypothèses et conclure par une phrase complète.
\end{exobox}


\begin{exobox}{Exercice 13 -- Composition de limites \etoiles{$\star\star$}}

Calculer
\[
\lim_{x\to0^+}\sqrt{1+x\sin(1/x)}.
\]
On commencera par montrer que l'expression est bien définie au voisinage à droite de $0$.

\medskip
\textbf{Consignes de rédaction.} Identifier le théorème utilisé, vérifier les hypothèses et conclure par une phrase complète.
\end{exobox}


\begin{exobox}{Exercice 14 -- Non-existence par deux suites \etoiles{$\star\star$}}

Montrer que la fonction
\[
f(x)=\sin\left(\frac1x\right)
\]
n'admet pas de limite en $0$.

\medskip
\textbf{Consignes de rédaction.} Identifier le théorème utilisé, vérifier les hypothèses et conclure par une phrase complète.
\end{exobox}


\begin{exobox}{Exercice 15 -- Continuité d'une fonction à paramètre \etoiles{$\star\star$}}

Soit
\[
f_a(x)=\begin{cases}
\dfrac{\e^x-1}{x} & \text{si }x\neq0,\\
a & \text{si }x=0.
\end{cases}
\]
Déterminer les valeurs de $a$ pour lesquelles $f_a$ est continue sur $\R$.

\medskip
\textbf{Consignes de rédaction.} Identifier le théorème utilisé, vérifier les hypothèses et conclure par une phrase complète.
\end{exobox}


\begin{exobox}{Exercice 16 -- Image d'un intervalle non fermé \etoiles{$\star\star$}}

Déterminer l'image de l'intervalle $]0,+\infty[$ par
\[
f(x)=x+\frac1x.
\]
La solution ne doit pas utiliser de dérivée.

\medskip
\textbf{Consignes de rédaction.} Identifier le théorème utilisé, vérifier les hypothèses et conclure par une phrase complète.
\end{exobox}


\begin{exobox}{Exercice 17 -- Une équation avec racine \etoiles{$\star\star$}}

Montrer que l'équation
\[
\sqrt{x+2}=x
\]
admet une unique solution réelle, puis la déterminer.

\medskip
\textbf{Consignes de rédaction.} Identifier le théorème utilisé, vérifier les hypothèses et conclure par une phrase complète.
\end{exobox}


\begin{exobox}{Exercice 18 -- Bijection de $]0,+\infty[$ sur $\R$ \etoiles{$\star\star$}}

Soit $f(x)=x+\ln x$ sur $]0,+\infty[$.
\begin{enumerate}
\item Montrer que $f$ est continue.
\item Montrer que $f$ est strictement croissante sans utiliser la dérivée.
\item En déduire que $f$ réalise une bijection de $]0,+\infty[$ sur $\R$.
\end{enumerate}

\medskip
\textbf{Consignes de rédaction.} Identifier le théorème utilisé, vérifier les hypothèses et conclure par une phrase complète.
\end{exobox}


\begin{exobox}{Exercice 19 -- Image réciproque et fonction réciproque \etoiles{$\star\star$}}

Soit $f:\R\to\R$, $f(x)=x^2$.
\begin{enumerate}
\item Calculer $f^{-1}([1,4])$ au sens de l'image réciproque d'un ensemble.
\item Expliquer pourquoi $f$ n'admet pas de fonction réciproque de $\R$ dans $\R$.
\item Restreindre le domaine pour obtenir une bijection et donner la réciproque.
\end{enumerate}

\medskip
\textbf{Consignes de rédaction.} Identifier le théorème utilisé, vérifier les hypothèses et conclure par une phrase complète.
\end{exobox}


\begin{exobox}{Exercice 20 -- Périodicité et limite \etoiles{$\star\star$}}

Soit $f:\R\to\R$ une fonction périodique de période $T>0$. On suppose que $\lim_{x\to+\infty}f(x)=\ell\in\R$. Montrer que $f$ est constante égale à $\ell$.

\medskip
\textbf{Consignes de rédaction.} Identifier le théorème utilisé, vérifier les hypothèses et conclure par une phrase complète.
\end{exobox}


\begin{exobox}{Exercice 21 -- TVI avec paramètre \etoiles{$\star\star$}}

Soit $a>0$. Montrer que l'équation
\[
x^3+ax-1=0
\]
admet une unique solution réelle.

\medskip
\textbf{Consignes de rédaction.} Identifier le théorème utilisé, vérifier les hypothèses et conclure par une phrase complète.
\end{exobox}


\begin{exobox}{Exercice 22 -- Réciproque explicite \etoiles{$\star\star$}}

Soit
\[
f:[1,+\infty[\to[0,+\infty[,\qquad f(x)=x^2-2x+1.
\]
\begin{enumerate}
\item Montrer que $f$ est bijective.
\item Déterminer $f^{-1}$.
\item Vérifier algébriquement les identités de réciprocité.
\end{enumerate}

\medskip
\textbf{Consignes de rédaction.} Identifier le théorème utilisé, vérifier les hypothèses et conclure par une phrase complète.
\end{exobox}


\section{Exercices de démonstration}
\begin{exobox}{Exercice 23 -- Unicité de la limite \etoiles{$\star\star\star$}}

Soit $f:D\to\R$ et $a\in\R$. Montrer, avec la définition $\varepsilon$--$\eta$, que si $f$ admet une limite finie en $a$, alors cette limite est unique.

\medskip
\textbf{Consignes de rédaction.} Identifier le théorème utilisé, vérifier les hypothèses et conclure par une phrase complète.
\end{exobox}


\begin{exobox}{Exercice 24 -- Fonction ayant une limite finie : bornitude locale \etoiles{$\star\star\star$}}

Soit $f:D\to\R$ telle que $\lim_{x\to a}f(x)=\ell\in\R$. Montrer que $f$ est bornée au voisinage épointé de $a$.

\medskip
\textbf{Consignes de rédaction.} Identifier le théorème utilisé, vérifier les hypothèses et conclure par une phrase complète.
\end{exobox}


\begin{exobox}{Exercice 25 -- Somme de deux limites \etoiles{$\star\star\star$}}

Démontrer que si $f(x)\to\ell$ et $g(x)\to m$ lorsque $x\to a$, alors $(f+g)(x)\to\ell+m$.

\medskip
\textbf{Consignes de rédaction.} Identifier le théorème utilisé, vérifier les hypothèses et conclure par une phrase complète.
\end{exobox}


\begin{exobox}{Exercice 26 -- Produit de deux limites \etoiles{$\star\star\star$}}

Démontrer que si $f(x)\to\ell$ et $g(x)\to m$ lorsque $x\to a$, alors $(fg)(x)\to\ell m$.

\medskip
\textbf{Consignes de rédaction.} Identifier le théorème utilisé, vérifier les hypothèses et conclure par une phrase complète.
\end{exobox}


\begin{exobox}{Exercice 27 -- Caractérisation séquentielle : sens direct \etoiles{$\star\star\star$}}

Soit $f:D\to\R$, $a\in\R$ et $\ell\in\R$. Démontrer que si $f(x)\to\ell$ lorsque $x\to a$, alors pour toute suite $(x_n)$ d'éléments de $D\setminus\{a\}$ telle que $x_n\to a$, on a $f(x_n)\to\ell$.

\medskip
\textbf{Consignes de rédaction.} Identifier le théorème utilisé, vérifier les hypothèses et conclure par une phrase complète.
\end{exobox}


\begin{exobox}{Exercice 28 -- Caractérisation séquentielle : contraposée \etoiles{$\star\star\star$}}

Démontrer la réciproque de la caractérisation séquentielle des limites par contraposée.

\medskip
\textbf{Consignes de rédaction.} Identifier le théorème utilisé, vérifier les hypothèses et conclure par une phrase complète.
\end{exobox}


\begin{exobox}{Exercice 29 -- Image d'un intervalle par une fonction continue \etoiles{$\star\star\star$}}

Soit $I$ un intervalle de $\R$ et $f:I\to\R$ continue. Montrer que $f(I)$ est un intervalle.

\medskip
\textbf{Consignes de rédaction.} Identifier le théorème utilisé, vérifier les hypothèses et conclure par une phrase complète.
\end{exobox}


\begin{exobox}{Exercice 30 -- Fonction strictement monotone et injectivité \etoiles{$\star\star\star$}}

Soit $f:I\to\R$ strictement monotone sur un intervalle $I$. Montrer que $f$ est injective.

\medskip
\textbf{Consignes de rédaction.} Identifier le théorème utilisé, vérifier les hypothèses et conclure par une phrase complète.
\end{exobox}


\begin{exobox}{Exercice 31 -- Théorème de la bijection continue : existence et unicité \etoiles{$\star\star\star$}}

Soit $f:I\to\R$ continue et strictement monotone sur un intervalle $I$. Montrer que pour tout $y\in f(I)$, l'équation $f(x)=y$ admet une unique solution dans $I$.

\medskip
\textbf{Consignes de rédaction.} Identifier le théorème utilisé, vérifier les hypothèses et conclure par une phrase complète.
\end{exobox}


\begin{exobox}{Exercice 32 -- Continuité de la réciproque : cas croissant sur un segment \etoiles{$\star\star\star$}}

Soit $f:[a,b]\to\R$ continue et strictement croissante. On pose $J=f([a,b])$. Montrer que $f^{-1}:J\to[a,b]$ est continue.

\medskip
\textbf{Consignes de rédaction.} Identifier le théorème utilisé, vérifier les hypothèses et conclure par une phrase complète.
\end{exobox}


\section{Exercices de réflexion}
\begin{exobox}{Exercice 33 -- Une fonction continue injective est monotone sur un intervalle \etoiles{$\star\star\star\star$}}

Soit $I$ un intervalle et $f:I\to\R$ continue et injective. Montrer que $f$ est strictement monotone sur $I$.

\medskip
\textbf{Consignes de rédaction.} Identifier le théorème utilisé, vérifier les hypothèses et conclure par une phrase complète.
\end{exobox}


\begin{exobox}{Exercice 34 -- Équation $x+\sin x=a$ \etoiles{$\star\star\star\star$}}

Pour $a\in\R$, étudier l'existence et l'unicité des solutions de
\[
x+\sin x=a.
\]

\medskip
\textbf{Consignes de rédaction.} Identifier le théorème utilisé, vérifier les hypothèses et conclure par une phrase complète.
\end{exobox}


\begin{exobox}{Exercice 35 -- Point fixe d'une fonction de $[0,1]$ dans lui-même \etoiles{$\star\star\star\star$}}

Soit $f:[0,1]\to[0,1]$ continue. Montrer qu'il existe $c\in[0,1]$ tel que $f(c)=c$.

\medskip
\textbf{Consignes de rédaction.} Identifier le théorème utilisé, vérifier les hypothèses et conclure par une phrase complète.
\end{exobox}


\begin{exobox}{Exercice 36 -- Image d'un segment et bornes atteintes \etoiles{$\star\star\star\star$}}

Soit $f:[a,b]\to\R$ continue. Montrer que $f$ est bornée et atteint ses bornes. On admettra seulement la propriété de Bolzano-Weierstrass pour les suites réelles bornées.

\medskip
\textbf{Consignes de rédaction.} Identifier le théorème utilisé, vérifier les hypothèses et conclure par une phrase complète.
\end{exobox}


\begin{exobox}{Exercice 37 -- Une fonction continue sans zéros change-t-elle de signe ? \etoiles{$\star\star\star\star$}}

Soit $I$ un intervalle et $f:I\to\R$ continue. On suppose que $f(x)\neq0$ pour tout $x\in I$. Montrer que $f$ garde un signe constant sur $I$.

\medskip
\textbf{Consignes de rédaction.} Identifier le théorème utilisé, vérifier les hypothèses et conclure par une phrase complète.
\end{exobox}


\begin{exobox}{Exercice 38 -- Équation avec paramètre et intervalle d'image \etoiles{$\star\star\star\star$}}

Soit $f(x)=x^2+\frac1{x^2}$ sur $]0,+\infty[$.
\begin{enumerate}
\item Déterminer l'image de $]0,+\infty[$ par $f$.
\item Pour quels réels $a$ l'équation $x^2+\frac1{x^2}=a$ admet-elle au moins une solution positive ?
\item Préciser le nombre de solutions positives selon $a$.
\end{enumerate}

\medskip
\textbf{Consignes de rédaction.} Identifier le théorème utilisé, vérifier les hypothèses et conclure par une phrase complète.
\end{exobox}


\begin{exobox}{Exercice 39 -- Prolongement non trivial \etoiles{$\star\star\star\star$}}

Étudier le prolongement par continuité en $0$ de
\[
f(x)=\frac{\sqrt{1+x+x^2}-1}{x}
\]
définie au voisinage épointé de $0$.

\medskip
\textbf{Consignes de rédaction.} Identifier le théorème utilisé, vérifier les hypothèses et conclure par une phrase complète.
\end{exobox}


\begin{exobox}{Exercice 40 -- Continuité d'une fonction définie par image rationnelle \etoiles{$\star\star\star\star$}}

Soit
\[
f(x)=\begin{cases}
\dfrac{x^2\sin(1/x)}{|x|} & \text{si }x\neq0,\\
0 & \text{si }x=0.
\end{cases}
\]
Étudier la continuité de $f$ sur $\R$ et sa parité.

\medskip
\textbf{Consignes de rédaction.} Identifier le théorème utilisé, vérifier les hypothèses et conclure par une phrase complète.
\end{exobox}


\begin{exobox}{Exercice 41 -- Injectivité sans dérivée \etoiles{$\star\star\star\star$}}

Montrer que la fonction
\[
f(x)=x^5+x^3+x
\]
est une bijection de $\R$ sur $\R$, sans utiliser la dérivée.

\medskip
\textbf{Consignes de rédaction.} Identifier le théorème utilisé, vérifier les hypothèses et conclure par une phrase complète.
\end{exobox}


\begin{exobox}{Exercice 42 -- Fonction réciproque implicite \etoiles{$\star\star\star\star$}}

On considère $f(x)=x+\e^x$ sur $\R$.
\begin{enumerate}
\item Montrer que $f$ réalise une bijection de $\R$ sur $\R$.
\item On note $g=f^{-1}$. Déterminer $g(1)$.
\item Résoudre $g(y)\leq0$.
\end{enumerate}

\medskip
\textbf{Consignes de rédaction.} Identifier le théorème utilisé, vérifier les hypothèses et conclure par une phrase complète.
\end{exobox}


\section{Exercices type oraux et concours}
\begin{exobox}{Exercice 43 -- Oral CCP : racine unique dans un intervalle mobile \etoiles{$\star\star\star\star\star$}}

Soit $n\in\N^*$ et
\[
f_n(x)=x^n+x-1.
\]
\begin{enumerate}
\item Montrer que l'équation $f_n(x)=0$ admet une unique solution $\alpha_n$ dans $[0,1]$.
\item Montrer que $\alpha_n\in]0,1[$.
\item Étudier le sens de variation de la suite $(\alpha_n)$.
\item Montrer que $(\alpha_n)$ converge et déterminer sa limite.
\end{enumerate}

\medskip
\textbf{Consignes de rédaction.} Identifier le théorème utilisé, vérifier les hypothèses et conclure par une phrase complète.
\end{exobox}


\begin{exobox}{Exercice 44 -- Oral Mines : équation fonctionnelle de continuité \etoiles{$\star\star\star\star\star$}}

Soit $f:\R\to\R$ continue telle que
\[
\forall x\in\R,\qquad f(f(x))=x.
\]
Montrer que $f$ est bijective. Montrer ensuite que $f$ est strictement monotone. Que peut-on dire de son sens de variation si $f$ n'est pas l'identité ?

\medskip
\textbf{Consignes de rédaction.} Identifier le théorème utilisé, vérifier les hypothèses et conclure par une phrase complète.
\end{exobox}


\begin{exobox}{Exercice 45 -- Oral Centrale : image et antécédents d'une fonction oscillante amortie \etoiles{$\star\star\star\star\star$}}

Soit
\[
f(x)=\begin{cases}
x\sin(1/x)&\text{si }x\neq0,\\
0&\text{si }x=0.
\end{cases}
\]
\begin{enumerate}
\item Montrer que $f$ est continue sur $\R$.
\item Montrer que $f$ n'est monotone sur aucun voisinage de $0$.
\item Montrer que l'équation $f(x)=0$ possède une infinité de solutions dans tout voisinage de $0$.
\end{enumerate}

\medskip
\textbf{Consignes de rédaction.} Identifier le théorème utilisé, vérifier les hypothèses et conclure par une phrase complète.
\end{exobox}


\begin{exobox}{Exercice 46 -- Oral Mines : continuité et densité élémentaire \etoiles{$\star\star\star\star\star$}}

Soit $f:\R\to\R$ continue. On suppose que $f(x)=0$ pour tout $x\in\mathbb Q$. Montrer que $f$ est identiquement nulle. On pourra utiliser la caractérisation séquentielle de la continuité et le fait admis que tout réel est limite d'une suite de rationnels.

\medskip
\textbf{Consignes de rédaction.} Identifier le théorème utilisé, vérifier les hypothèses et conclure par une phrase complète.
\end{exobox}


\begin{exobox}{Exercice 47 -- Oral CCP : comparaison et fausse intuition graphique \etoiles{$\star\star\star\star\star$}}

Étudier l'existence de
\[
\lim_{x\to0}\frac{x^2\sin(1/x)}{|x|+x^2}.
\]

\medskip
\textbf{Consignes de rédaction.} Identifier le théorème utilisé, vérifier les hypothèses et conclure par une phrase complète.
\end{exobox}


\begin{exobox}{Exercice 48 -- Oral Centrale : suite d'antécédents \etoiles{$\star\star\star\star\star$}}

Soit $f:[0,1]\to[0,1]$ continue. Pour $n\geq1$, montrer qu'il existe $x_n\in[0,1]$ tel que
\[
f(x_n)=x_n^n.
\]

\medskip
\textbf{Consignes de rédaction.} Identifier le théorème utilisé, vérifier les hypothèses et conclure par une phrase complète.
\end{exobox}


\begin{exobox}{Exercice 49 -- Problème long : réciproque et équation implicite \etoiles{$\star\star\star\star\star$}}

On considère $f(x)=x+\ln(1+x)$ sur $]-1,+\infty[$.
\begin{enumerate}
\item Montrer que $f$ réalise une bijection de $]-1,+\infty[$ sur $\R$.
\item On note $g=f^{-1}$. Déterminer le signe de $g(y)$ selon $y$.
\item Résoudre l'inéquation $g(y)\geq y$.
\item Déterminer pour quels $y$ on a $g(y)\leq y+1$.
\end{enumerate}

\medskip
\textbf{Consignes de rédaction.} Identifier le théorème utilisé, vérifier les hypothèses et conclure par une phrase complète.
\end{exobox}


\begin{exobox}{Exercice 50 -- Problème long : existence, unicité, encadrement \etoiles{$\star\star\star\star\star$}}

Pour $a>0$, on considère l'équation
\[
\e^{-x}=ax.
\]
\begin{enumerate}
\item Montrer qu'elle admet une unique solution $x_a$ dans $]0,+\infty[$.
\item Montrer que $0<x_a<1/a$.
\item Montrer que si $0<a<b$, alors $x_b<x_a$.
\item Étudier les limites de $x_a$ lorsque $a\to0^+$ puis lorsque $a\to+\infty$.
\end{enumerate}

\medskip
\textbf{Consignes de rédaction.} Identifier le théorème utilisé, vérifier les hypothèses et conclure par une phrase complète.
\end{exobox}


\section{Problèmes longs type concours}
\begin{problembox}{Problème A -- Théorème des valeurs intermédiaires renforcé \etoiles{$\star\star\star\star\star$}}

Soit $f:[a,b]\to\R$ continue. On suppose $f(a)<0<f(b)$.
\begin{enumerate}
\item On pose $E=\{x\in[a,b]\mid f(x)\leq0\}$. Montrer que $E$ est non vide et majoré.
\item Soit $c=\sup E$. Montrer que $c\in[a,b]$.
\item Montrer que $f(c)=0$.
\item On suppose en plus que $f$ est strictement croissante. Montrer que $c$ est l'unique zéro de $f$ sur $[a,b]$.
\item Application : montrer que $x^3+x-1=0$ admet une unique racine dans $[0,1]$.
\end{enumerate}

\medskip
\textbf{Objectif concours.} Le but est de combiner continuité, image d'intervalle, monotonie et caractérisation séquentielle sans utiliser d'outils hors chapitre.
\end{problembox}
\begin{problembox}{Problème B -- Bijection continue et paramètres \etoiles{$\star\star\star\star\star$}}

Pour $\lambda>0$, on définit
\[
f_\lambda(x)=x+\lambda\ln x
\]
sur $]0,+\infty[$.
\begin{enumerate}
\item Montrer que $f_\lambda$ réalise une bijection de $]0,+\infty[$ sur $\R$.
\item Pour $y\in\R$, on note $g_\lambda(y)$ l'unique réel $x>0$ tel que $x+\lambda\ln x=y$. Déterminer $g_\lambda(1)$ lorsque $\lambda=1$.
\item Montrer que si $0<\lambda<\mu$, alors $g_\mu(y)>g_\lambda(y)$ pour tout $y$ tel que $g_\lambda(y)<1$, et préciser ce qui se passe si $g_\lambda(y)=1$ ou $g_\lambda(y)>1$.
\item Résoudre l'inéquation $g_\lambda(y)\geq1$.
\end{enumerate}

\medskip
\textbf{Objectif concours.} Le but est de combiner continuité, image d'intervalle, monotonie et caractérisation séquentielle sans utiliser d'outils hors chapitre.
\end{problembox}
\begin{problembox}{Problème C -- Continuité, zéros et oscillations \etoiles{$\star\star\star\star\star$}}

On définit
\[
f(x)=\begin{cases}
x\sin(1/x)&\text{si }x\neq0,\\
0&\text{si }x=0.
\end{cases}
\]
\begin{enumerate}
\item Montrer que $f$ est continue sur $\R$.
\item Déterminer une infinité de zéros de $f$ convergeant vers $0$.
\item Trouver deux suites $(u_n)$ et $(v_n)$ tendant vers $0$ telles que $f(u_n)>0$ et $f(v_n)<0$ pour tout $n$.
\item Montrer que $f$ n'est injective sur aucun intervalle $]-\alpha,\alpha[$ avec $\alpha>0$.
\item Comparer avec la fonction $g(x)=x^2\sin(1/x)$ prolongée par $g(0)=0$.
\end{enumerate}

\medskip
\textbf{Objectif concours.} Le but est de combiner continuité, image d'intervalle, monotonie et caractérisation séquentielle sans utiliser d'outils hors chapitre.
\end{problembox}
\newpage
\section{Corrigé complet et rédigé}
\begin{attentionbox}
Les corrections ci-dessous sont volontairement détaillées. Dans une copie, il faut garder les justifications principales : domaine, continuité, monotonie, théorème appliqué, conclusion.
\end{attentionbox}
\section{Corrigés -- Exercices d'application directe}
\begin{corrbox}{Correction de l'exercice 1 -- Domaines élémentaires}

La racine impose $2x-1\geq 0$, donc $x\geq \frac12$. Le dénominateur impose $x^2-4\neq0$, donc $x\neq -2$ et $x\neq2$. Comme $x\geq\frac12$, la condition $x\neq-2$ est inutile. Ainsi
\[
\dom_f=\left[\frac12,+\infty\right[\setminus\{2\}.
\]
Sur chacun des intervalles $\left[\frac12,2\right[$ et $]2,+\infty[$, la fonction est quotient de fonctions continues avec dénominateur non nul. Elle est donc continue sur ces deux intervalles.
En $2$, le numérateur tend vers $\sqrt3\neq0$ tandis que le dénominateur tend vers $0$. La limite n'est pas finie. Il n'y a donc pas de prolongement par continuité en $2$.

\medskip
\textbf{Vérification finale.} Les hypothèses du théorème utilisé ont été explicitées et la conclusion porte bien sur l'objet demandé dans l'énoncé.
\end{corrbox}


\begin{corrbox}{Correction de l'exercice 2 -- Image d'un segment par une fonction quadratique}

On a $f(x)=x^2-2x=(x-1)^2-1$. Sur $[0,3]$, le minimum est atteint en $1$ et vaut $-1$. Aux bornes, $f(0)=0$ et $f(3)=3$, donc le maximum vaut $3$. Par continuité de $f$ sur $[0,3]$, l'image est un intervalle :
\[
f([0,3])=[-1,3].
\]
Pour l'image réciproque,
\[
-1\leq (x-1)^2-1\leq0 \Longleftrightarrow 0\leq (x-1)^2\leq1.
\]
La première inégalité est automatique. La seconde donne $|x-1|\leq1$, donc $x\in[0,2]$. Ainsi $f^{-1}([-1,0])=[0,2]$.

\medskip
\textbf{Vérification finale.} Les hypothèses du théorème utilisé ont été explicitées et la conclusion porte bien sur l'objet demandé dans l'énoncé.
\end{corrbox}


\begin{corrbox}{Correction de l'exercice 3 -- Parité et domaine}

Le dénominateur $1+x^2$ ne s'annule jamais sur $\R$. Donc $\dom_f=\R$, qui est symétrique par rapport à $0$. Pour tout $x\in\R$,
\[
f(-x)=\frac{(-x)^3}{1+(-x)^2}=-\frac{x^3}{1+x^2}=-f(x).
\]
La fonction est donc impaire. Son graphe est symétrique par rapport à l'origine du repère.

\medskip
\textbf{Vérification finale.} Les hypothèses du théorème utilisé ont été explicitées et la conclusion porte bien sur l'objet demandé dans l'énoncé.
\end{corrbox}


\begin{corrbox}{Correction de l'exercice 4 -- Limite par factorisation}

Pour $x\neq2$,
\[
\frac{x^2-4}{x-2}=\frac{(x-2)(x+2)}{x-2}=x+2.
\]
La limite en $2$ ne dépend que des valeurs de la fonction au voisinage de $2$, avec $x\neq2$. On peut donc remplacer l'expression par $x+2$ pour le calcul de limite. Ainsi
\[
\lim_{x\to2}\frac{x^2-4}{x-2}=\lim_{x\to2}(x+2)=4.
\]

\medskip
\textbf{Vérification finale.} Les hypothèses du théorème utilisé ont été explicitées et la conclusion porte bien sur l'objet demandé dans l'énoncé.
\end{corrbox}


\begin{corrbox}{Correction de l'exercice 5 -- Limites latérales simples}

Si $x>0$, $|x|=x$, donc $f(x)=1$. Ainsi $\lim_{x\to0^+}f(x)=1$. Si $x<0$, $|x|=-x$, donc $f(x)=-1$. Ainsi $\lim_{x\to0^-}f(x)=-1$. Les deux limites latérales existent mais sont différentes. La limite bilatérale en $0$ n'existe donc pas.

\medskip
\textbf{Vérification finale.} Les hypothèses du théorème utilisé ont été explicitées et la conclusion porte bien sur l'objet demandé dans l'énoncé.
\end{corrbox}


\begin{corrbox}{Correction de l'exercice 6 -- Encadrement classique}

Pour tout $x\neq0$, $-1\leq\cos(1/x)\leq1$. Comme $x^2\geq0$, on obtient
\[
-x^2\leq x^2\cos\left(\frac1x\right)\leq x^2.
\]
Or $-x^2\to0$ et $x^2\to0$ lorsque $x\to0$. Par le théorème des gendarmes,
\[
x^2\cos\left(\frac1x\right)\to0.
\]

\medskip
\textbf{Vérification finale.} Les hypothèses du théorème utilisé ont été explicitées et la conclusion porte bien sur l'objet demandé dans l'énoncé.
\end{corrbox}


\begin{corrbox}{Correction de l'exercice 7 -- Continuité en un point défini par morceaux}

Pour $x\neq1$,
\[
\frac{x^2-1}{x-1}=\frac{(x-1)(x+1)}{x-1}=x+1.
\]
Donc $\lim_{x\to1}f(x)=2$. Or $f(1)=3$. La limite existe mais ne vaut pas la valeur de la fonction. Donc $f$ n'est pas continue en $1$. Il aurait fallu poser $f(1)=2$.

\medskip
\textbf{Vérification finale.} Les hypothèses du théorème utilisé ont été explicitées et la conclusion porte bien sur l'objet demandé dans l'énoncé.
\end{corrbox}


\begin{corrbox}{Correction de l'exercice 8 -- Prolongement de $\sin x/x$}

On utilise la limite classique
\[
\lim_{x\to0}\frac{\sin x}{x}=1.
\]
La continuité en $0$ équivaut à $\lim_{x\to0}f(x)=f(0)$, donc à $1=a$. La seule valeur possible est donc $a=1$.

\medskip
\textbf{Vérification finale.} Les hypothèses du théorème utilisé ont été explicitées et la conclusion porte bien sur l'objet demandé dans l'énoncé.
\end{corrbox}


\begin{corrbox}{Correction de l'exercice 9 -- TVI : existence d'un point fixe du cosinus}

On pose $g(x)=\cos x-x$. La fonction $g$ est continue sur $[0,1]$ comme somme de fonctions continues. On calcule
\[
g(0)=1>0,\qquad g(1)=\cos1-1<0.
\]
Comme $0$ est compris entre $g(0)$ et $g(1)$, le théorème des valeurs intermédiaires assure l'existence de $c\in[0,1]$ tel que $g(c)=0$. Ainsi $\cos c=c$.

\medskip
\textbf{Vérification finale.} Les hypothèses du théorème utilisé ont été explicitées et la conclusion porte bien sur l'objet demandé dans l'énoncé.
\end{corrbox}


\begin{corrbox}{Correction de l'exercice 10 -- Bijection continue élémentaire}

La fonction $f$ est polynomiale, donc continue sur $\R$. Soient $x<y$. Alors
\[
f(y)-f(x)=y^3-x^3+y-x=(y-x)(y^2+xy+x^2+1).
\]
Or $y^2+xy+x^2=\frac12(x^2+y^2)+\frac12(x+y)^2\geq0$, donc $y^2+xy+x^2+1>0$. Comme $y-x>0$, $f(y)-f(x)>0$. Ainsi $f$ est strictement croissante. Enfin
\[
\lim_{x\to-\infty}f(x)=-\infty,\qquad \lim_{x\to+\infty}f(x)=+\infty.
\]
Par le théorème de la bijection continue, $f$ est une bijection de $\R$ sur $\R$.

\medskip
\textbf{Vérification finale.} Les hypothèses du théorème utilisé ont été explicitées et la conclusion porte bien sur l'objet demandé dans l'énoncé.
\end{corrbox}


\section{Corrigés -- Exercices de méthode}
\begin{corrbox}{Correction de l'exercice 11 -- Limite avec quantité conjuguée}

Pour $x\neq0$ et $x>-1$,
\[
\frac{\sqrt{1+x}-1}{x}=\frac{(\sqrt{1+x}-1)(\sqrt{1+x}+1)}{x(\sqrt{1+x}+1)}=\frac{x}{x(\sqrt{1+x}+1)}=\frac1{\sqrt{1+x}+1}.
\]
Donc
\[
\lim_{x\to0}\frac{\sqrt{1+x}-1}{x}=\frac1{\sqrt1+1}=\frac12.
\]
La fonction initiale, définie pour $x>-1$ et $x\neq0$, est donc prolongeable par continuité en $0$ en posant $f(0)=1/2$.

\medskip
\textbf{Vérification finale.} Les hypothèses du théorème utilisé ont été explicitées et la conclusion porte bien sur l'objet demandé dans l'énoncé.
\end{corrbox}


\begin{corrbox}{Correction de l'exercice 12 -- Limites et signe du dénominateur}

On factorise
\[
x^2-11x+28=(x-4)(x-7),\qquad x^2-25=(x-5)(x+5).
\]
Au voisinage de $5$, $(x-4)>0$, $(x-7)<0$ et $(x+5)>0$. Le numérateur tend vers $-2$ et ne s'annule pas. Ainsi le signe et la divergence sont gouvernés par $x-5$. Lorsque $x\to5^+$, le dénominateur tend vers $0^+$ et le numérateur vers un nombre négatif ; la limite vaut $-\infty$. Lorsque $x\to5^-$, le dénominateur tend vers $0^-$ et le numérateur vers un nombre négatif ; la limite vaut $+\infty$.

\medskip
\textbf{Vérification finale.} Les hypothèses du théorème utilisé ont été explicitées et la conclusion porte bien sur l'objet demandé dans l'énoncé.
\end{corrbox}


\begin{corrbox}{Correction de l'exercice 13 -- Composition de limites}

Pour tout $x>0$, $-x\leq x\sin(1/x)\leq x$. Si $0<x<1$, alors $1+x\sin(1/x)\geq1-x>0$, donc la racine est bien définie. Par encadrement, $x\sin(1/x)\to0$. Donc $1+x\sin(1/x)\to1$. La fonction racine carrée est continue en $1$, donc
\[
\sqrt{1+x\sin(1/x)}\to\sqrt1=1.
\]

\medskip
\textbf{Vérification finale.} Les hypothèses du théorème utilisé ont été explicitées et la conclusion porte bien sur l'objet demandé dans l'énoncé.
\end{corrbox}


\begin{corrbox}{Correction de l'exercice 14 -- Non-existence par deux suites}

On utilise la caractérisation séquentielle. Posons
\[
x_n=\frac1{\frac\pi2+2n\pi},\qquad y_n=\frac1{-\frac\pi2+2n\pi}
\]
pour $n$ assez grand afin que les dénominateurs soient non nuls. Alors $x_n\to0$ et $y_n\to0$. Mais
\[
f(x_n)=\sin\left(\frac\pi2+2n\pi\right)=1,
\]
et
\[
f(y_n)=\sin\left(-\frac\pi2+2n\pi\right)=-1.
\]
Si une limite en $0$ existait, toutes les images de suites tendant vers $0$ auraient la même limite. Ce n'est pas le cas. Donc la limite n'existe pas.

\medskip
\textbf{Vérification finale.} Les hypothèses du théorème utilisé ont été explicitées et la conclusion porte bien sur l'objet demandé dans l'énoncé.
\end{corrbox}


\begin{corrbox}{Correction de l'exercice 15 -- Continuité d'une fonction à paramètre}

Sur $\R^*$, la fonction est quotient de fonctions continues avec dénominateur non nul. Il ne reste qu'à étudier $0$. On sait que
\[
\lim_{x\to0}\frac{\e^x-1}{x}=1.
\]
La continuité en $0$ impose et suffit que $a=1$. Donc $f_a$ est continue sur $\R$ si et seulement si $a=1$.

\medskip
\textbf{Vérification finale.} Les hypothèses du théorème utilisé ont été explicitées et la conclusion porte bien sur l'objet demandé dans l'énoncé.
\end{corrbox}


\begin{corrbox}{Correction de l'exercice 16 -- Image d'un intervalle non fermé}

Pour $x>0$, l'inégalité $x+1/x\geq2$ vient de $(x-1)^2\geq0$. L'égalité a lieu pour $x=1$. Donc $f(]0,+\infty[)\subset[2,+\infty[$ et $2\in f(]0,+\infty[)$. Réciproquement, soit $y\geq2$. Résoudre $x+1/x=y$ revient à résoudre
\[
x^2-yx+1=0.
\]
Son discriminant vaut $\Delta=y^2-4\geq0$, et les racines
\[
x=\frac{y\pm\sqrt{y^2-4}}2
\]
sont positives. Il existe donc $x>0$ tel que $f(x)=y$. Ainsi
\[
f(]0,+\infty[)=[2,+\infty[.
\]

\medskip
\textbf{Vérification finale.} Les hypothèses du théorème utilisé ont été explicitées et la conclusion porte bien sur l'objet demandé dans l'énoncé.
\end{corrbox}


\begin{corrbox}{Correction de l'exercice 17 -- Une équation avec racine}

Le domaine impose $x\geq-2$, et le membre de droite doit être positif car le membre de gauche l'est, donc toute solution vérifie $x\geq0$. Sur $[0,+\infty[$, les deux membres sont bien définis. En élevant au carré, opération licite pour $x\geq0$,
\[
x+2=x^2\Longleftrightarrow x^2-x-2=0\Longleftrightarrow (x-2)(x+1)=0.
\]
La seule solution compatible avec $x\geq0$ est $x=2$. Vérification : $\sqrt{2+2}=2$. Donc l'équation admet une unique solution réelle, $2$.

\medskip
\textbf{Vérification finale.} Les hypothèses du théorème utilisé ont été explicitées et la conclusion porte bien sur l'objet demandé dans l'énoncé.
\end{corrbox}


\begin{corrbox}{Correction de l'exercice 18 -- Bijection de $]0,+\infty[$ sur $\R$}

La fonction $x\mapsto x$ et la fonction $\ln$ sont continues sur $]0,+\infty[$, donc $f$ est continue sur cet intervalle. Si $0<x<y$, alors $x<y$ et, comme $\ln$ est strictement croissante, $\ln x<\ln y$. En additionnant, $x+\ln x<y+\ln y$. Donc $f$ est strictement croissante. Enfin
\[
\lim_{x\to0^+}(x+\ln x)=-\infty,
\qquad
\lim_{x\to+\infty}(x+\ln x)=+\infty.
\]
Par le théorème de la bijection continue, $f$ réalise une bijection de $]0,+\infty[$ sur $\R$.

\medskip
\textbf{Vérification finale.} Les hypothèses du théorème utilisé ont été explicitées et la conclusion porte bien sur l'objet demandé dans l'énoncé.
\end{corrbox}


\begin{corrbox}{Correction de l'exercice 19 -- Image réciproque et fonction réciproque}

On cherche $x\in\R$ tel que $x^2\in[1,4]$, c'est-à-dire $1\leq x^2\leq4$, ou encore $1\leq |x|\leq2$. Ainsi
\[
f^{-1}([1,4])=[-2,-1]\cup[1,2].
\]
La fonction $x\mapsto x^2$ n'est pas injective sur $\R$ car $f(1)=f(-1)$. Elle n'est donc pas bijective de $\R$ dans $\R$ et n'admet pas de fonction réciproque sur $\R$. En revanche, restreinte à $[0,+\infty[$, elle est continue, strictement croissante et d'image $[0,+\infty[$. Elle est bijective de $[0,+\infty[$ sur $[0,+\infty[$ et sa réciproque est $y\mapsto\sqrt y$.

\medskip
\textbf{Vérification finale.} Les hypothèses du théorème utilisé ont été explicitées et la conclusion porte bien sur l'objet demandé dans l'énoncé.
\end{corrbox}


\begin{corrbox}{Correction de l'exercice 20 -- Périodicité et limite}

Soit $a\in\R$. Pour tout $n\in\N$, $a+nT\to+\infty$. Par hypothèse, $f(a+nT)\to\ell$. Or, par périodicité, $f(a+nT)=f(a)$ pour tout $n$. La suite constante de valeur $f(a)$ converge donc vers $\ell$, ce qui impose $f(a)=\ell$. Comme $a$ est arbitraire, $f$ est constante égale à $\ell$ sur $\R$.

\medskip
\textbf{Vérification finale.} Les hypothèses du théorème utilisé ont été explicitées et la conclusion porte bien sur l'objet demandé dans l'énoncé.
\end{corrbox}


\begin{corrbox}{Correction de l'exercice 21 -- TVI avec paramètre}

Posons $f(x)=x^3+ax-1$. La fonction est polynomiale, donc continue sur $\R$. De plus
\[
\lim_{x\to-\infty}f(x)=-\infty,\qquad \lim_{x\to+\infty}f(x)=+\infty.
\]
Elle admet donc au moins une racine réelle par le TVI sur un segment assez grand. Pour l'unicité, soient $x<y$. Alors
\[
f(y)-f(x)=y^3-x^3+a(y-x)=(y-x)(y^2+xy+x^2+a).
\]
Comme $y^2+xy+x^2\geq0$ et $a>0$, le facteur entre parenthèses est strictement positif. Donc $f(y)>f(x)$. La fonction est strictement croissante, donc injective. L'équation admet une unique solution réelle.

\medskip
\textbf{Vérification finale.} Les hypothèses du théorème utilisé ont été explicitées et la conclusion porte bien sur l'objet demandé dans l'énoncé.
\end{corrbox}


\begin{corrbox}{Correction de l'exercice 22 -- Réciproque explicite}

Pour $x\geq1$, $f(x)=(x-1)^2$. La fonction $x\mapsto x-1$ est strictement croissante de $[1,+\infty[$ sur $[0,+\infty[$, et $u\mapsto u^2$ est strictement croissante sur $[0,+\infty[$. Donc $f$ est strictement croissante. Elle est continue et son image est $[0,+\infty[$. Elle est donc bijective de $[1,+\infty[$ sur $[0,+\infty[$. Pour $y\geq0$, $y=(x-1)^2$ avec $x\geq1$ donne $x=1+\sqrt y$. Ainsi
\[
f^{-1}(y)=1+\sqrt y.
\]
On vérifie : $f(f^{-1}(y))=(1+\sqrt y-1)^2=y$ et $f^{-1}(f(x))=1+\sqrt{(x-1)^2}=1+|x-1|=x$ car $x\geq1$.

\medskip
\textbf{Vérification finale.} Les hypothèses du théorème utilisé ont été explicitées et la conclusion porte bien sur l'objet demandé dans l'énoncé.
\end{corrbox}


\section{Corrigés -- Exercices de démonstration}
\begin{corrbox}{Correction de l'exercice 23 -- Unicité de la limite}

Supposons que $f(x)\to\ell$ et $f(x)\to m$ lorsque $x\to a$. On veut montrer $\ell=m$. Soit $\varepsilon>0$. Par la première limite, il existe $\eta_1>0$ tel que $0<|x-a|\leq\eta_1$ implique $|f(x)-\ell|\leq\varepsilon/2$. Par la seconde, il existe $\eta_2>0$ tel que $0<|x-a|\leq\eta_2$ implique $|f(x)-m|\leq\varepsilon/2$. Pour $\eta=\min(\eta_1,\eta_2)$ et pour un $x\in D$ vérifiant $0<|x-a|\leq\eta$, on obtient
\[
|\ell-m|\leq |\ell-f(x)|+|f(x)-m|\leq\varepsilon.
\]
Comme cette inégalité est vraie pour tout $\varepsilon>0$, on conclut $|\ell-m|=0$, donc $\ell=m$.

\medskip
\textbf{Vérification finale.} Les hypothèses du théorème utilisé ont été explicitées et la conclusion porte bien sur l'objet demandé dans l'énoncé.
\end{corrbox}


\begin{corrbox}{Correction de l'exercice 24 -- Fonction ayant une limite finie : bornitude locale}

On prend $\varepsilon=1$ dans la définition de la limite. Il existe $\eta>0$ tel que, pour tout $x\in D$, si $0<|x-a|\leq\eta$, alors
\[
|f(x)-\ell|\leq1.
\]
Par l'inégalité triangulaire,
\[
|f(x)|\leq |f(x)-\ell|+|\ell|\leq1+|\ell|.
\]
Donc $f$ est bornée sur le voisinage épointé $D\cap(]a-\eta,a+\eta[\setminus\{a\})$.

\medskip
\textbf{Vérification finale.} Les hypothèses du théorème utilisé ont été explicitées et la conclusion porte bien sur l'objet demandé dans l'énoncé.
\end{corrbox}


\begin{corrbox}{Correction de l'exercice 25 -- Somme de deux limites}

Soit $\varepsilon>0$. Comme $f(x)\to\ell$, il existe $\eta_1>0$ tel que $0<|x-a|\leq\eta_1$ entraîne $|f(x)-\ell|\leq\varepsilon/2$. Comme $g(x)\to m$, il existe $\eta_2>0$ tel que $0<|x-a|\leq\eta_2$ entraîne $|g(x)-m|\leq\varepsilon/2$. Posons $\eta=\min(\eta_1,\eta_2)$. Alors
\[
|(f(x)+g(x))-(\ell+m)|\leq |f(x)-\ell|+|g(x)-m|\leq\varepsilon.
\]
Cela prouve $(f+g)(x)\to\ell+m$.

\medskip
\textbf{Vérification finale.} Les hypothèses du théorème utilisé ont été explicitées et la conclusion porte bien sur l'objet demandé dans l'énoncé.
\end{corrbox}


\begin{corrbox}{Correction de l'exercice 26 -- Produit de deux limites}

Comme $f$ admet une limite finie, elle est bornée au voisinage épointé de $a$ : il existe $M>0$ et $\eta_0>0$ tels que $|f(x)|\leq M$ pour $0<|x-a|\leq\eta_0$. On écrit
\[
fg-\ell m=f(g-m)+m(f-\ell).
\]
Soit $\varepsilon>0$. Choisissons $\eta_1$ tel que $|g(x)-m|\leq \frac{\varepsilon}{2M}$ et $\eta_2$ tel que $|f(x)-\ell|\leq \frac{\varepsilon}{2(1+|m|)}$ dès que $x$ est assez proche de $a$. Pour $\eta=\min(\eta_0,\eta_1,\eta_2)$,
\[
|f(x)g(x)-\ell m|\leq M\frac{\varepsilon}{2M}+|m|\frac{\varepsilon}{2(1+|m|)}\leq\varepsilon.
\]
Donc $fg\to\ell m$.

\medskip
\textbf{Vérification finale.} Les hypothèses du théorème utilisé ont été explicitées et la conclusion porte bien sur l'objet demandé dans l'énoncé.
\end{corrbox}


\begin{corrbox}{Correction de l'exercice 27 -- Caractérisation séquentielle : sens direct}

Soit $\varepsilon>0$. Par définition de la limite de $f$ en $a$, il existe $\eta>0$ tel que
\[
0<|x-a|\leq\eta\Longrightarrow |f(x)-\ell|\leq\varepsilon.
\]
Comme $x_n\to a$, il existe $N\in\N$ tel que pour tout $n\geq N$, $|x_n-a|\leq\eta$. De plus $x_n\neq a$ par hypothèse, donc $0<|x_n-a|\leq\eta$. Par conséquent, pour $n\geq N$,
\[
|f(x_n)-\ell|\leq\varepsilon.
\]
Ainsi $f(x_n)\to\ell$.

\medskip
\textbf{Vérification finale.} Les hypothèses du théorème utilisé ont été explicitées et la conclusion porte bien sur l'objet demandé dans l'énoncé.
\end{corrbox}


\begin{corrbox}{Correction de l'exercice 28 -- Caractérisation séquentielle : contraposée}

Supposons que $f(x)$ ne tende pas vers $\ell$ lorsque $x\to a$. La négation de la définition donne : il existe $\varepsilon_0>0$ tel que pour tout $\eta>0$, il existe $x\in D$ vérifiant
\[
0<|x-a|\leq\eta\quad\text{et}\quad |f(x)-\ell|>\varepsilon_0.
\]
Pour chaque $n\geq1$, on applique cette propriété avec $\eta=1/n$. Il existe alors $x_n\in D$ tel que
\[
0<|x_n-a|\leq\frac1n\quad\text{et}\quad |f(x_n)-\ell|>\varepsilon_0.
\]
La première inégalité donne $x_n\to a$, avec $x_n\neq a$, tandis que la seconde empêche $f(x_n)$ de tendre vers $\ell$. La propriété séquentielle est donc fausse. Par contraposée, si la propriété séquentielle est vraie, alors $f(x)\to\ell$.

\medskip
\textbf{Vérification finale.} Les hypothèses du théorème utilisé ont été explicitées et la conclusion porte bien sur l'objet demandé dans l'énoncé.
\end{corrbox}


\begin{corrbox}{Correction de l'exercice 29 -- Image d'un intervalle par une fonction continue}

Soient $u,v\in f(I)$ avec $u\leq v$. Il existe $a,b\in I$ tels que $u=f(a)$ et $v=f(b)$. Soit $y\in[u,v]$. Comme $I$ est un intervalle, le segment $[\min(a,b),\max(a,b)]$ est inclus dans $I$. La restriction de $f$ à ce segment est continue. Par le théorème des valeurs intermédiaires, il existe $c$ entre $a$ et $b$ tel que $f(c)=y$. Donc $y\in f(I)$. Ainsi tout réel compris entre deux éléments de $f(I)$ appartient à $f(I)$ : $f(I)$ est un intervalle.

\medskip
\textbf{Vérification finale.} Les hypothèses du théorème utilisé ont été explicitées et la conclusion porte bien sur l'objet demandé dans l'énoncé.
\end{corrbox}


\begin{corrbox}{Correction de l'exercice 30 -- Fonction strictement monotone et injectivité}

Supposons par exemple $f$ strictement croissante. Soient $x,y\in I$ avec $x\neq y$. Si $x<y$, alors $f(x)<f(y)$ ; si $y<x$, alors $f(y)<f(x)$. Dans les deux cas $f(x)\neq f(y)$. Donc $f(x)=f(y)$ implique nécessairement $x=y$. La fonction est injective. Le cas strictement décroissant se traite de la même manière en inversant les inégalités.

\medskip
\textbf{Vérification finale.} Les hypothèses du théorème utilisé ont été explicitées et la conclusion porte bien sur l'objet demandé dans l'énoncé.
\end{corrbox}


\begin{corrbox}{Correction de l'exercice 31 -- Théorème de la bijection continue : existence et unicité}

Par définition de $f(I)$, si $y\in f(I)$, il existe au moins un $x\in I$ tel que $f(x)=y$. C'est l'existence. Pour l'unicité, si $x_1,x_2\in I$ vérifient $f(x_1)=f(x_2)=y$, alors l'injectivité de $f$, conséquence de sa stricte monotonie, donne $x_1=x_2$. Ainsi l'équation admet une unique solution dans $I$.

\medskip
\textbf{Vérification finale.} Les hypothèses du théorème utilisé ont été explicitées et la conclusion porte bien sur l'objet demandé dans l'énoncé.
\end{corrbox}


\begin{corrbox}{Correction de l'exercice 32 -- Continuité de la réciproque : cas croissant sur un segment}

Soit $y_0\in J$ et $x_0=f^{-1}(y_0)$. On montre la continuité en $y_0$. Si $x_0\in]a,b[$, prenons $\varepsilon>0$ tel que $[x_0-\varepsilon,x_0+\varepsilon]\subset[a,b]$. Par stricte croissance,
\[
f(x_0-\varepsilon)<f(x_0)=y_0<f(x_0+\varepsilon).
\]
Posons
\[
\alpha=\min\{y_0-f(x_0-\varepsilon),f(x_0+\varepsilon)-y_0\}>0.
\]
Si $|y-y_0|<\alpha$ et $y\in J$, alors $f(x_0-\varepsilon)<y<f(x_0+\varepsilon)$. En appliquant $f^{-1}$, qui est croissante, on obtient $x_0-\varepsilon<f^{-1}(y)<x_0+\varepsilon$. Donc $|f^{-1}(y)-x_0|<\varepsilon$. Aux bornes, on adapte avec des demi-intervalles. La réciproque est continue.

\medskip
\textbf{Vérification finale.} Les hypothèses du théorème utilisé ont été explicitées et la conclusion porte bien sur l'objet demandé dans l'énoncé.
\end{corrbox}


\section{Corrigés -- Exercices de réflexion}
\begin{corrbox}{Correction de l'exercice 33 -- Une fonction continue injective est monotone sur un intervalle}

On raisonne par contradiction. Supposons que $f$ ne soit pas monotone. Alors il existe $a<b<c$ dans $I$ tels que $f(b)$ ne soit pas compris de manière cohérente entre $f(a)$ et $f(c)$ : par exemple $f(b)>f(a)$ et $f(b)>f(c)$, ou bien $f(b)<f(a)$ et $f(b)<f(c)$. Traitons le premier cas. Si $f(a)=f(c)$, l'injectivité est déjà contredite. Sinon, supposons par exemple $f(a)<f(c)<f(b)$. Par continuité de $f$ sur $[a,b]$, le TVI donne un $u\in(a,b)$ tel que $f(u)=f(c)$. Comme $u\neq c$, cela contredit l'injectivité. Les autres configurations se traitent de même. Donc une fonction continue injective sur un intervalle est strictement monotone.

\medskip
\textbf{Vérification finale.} Les hypothèses du théorème utilisé ont été explicitées et la conclusion porte bien sur l'objet demandé dans l'énoncé.
\end{corrbox}


\begin{corrbox}{Correction de l'exercice 34 -- Équation $x+\sin x=a$}

Posons $f(x)=x+\sin x$. La fonction est continue sur $\R$. Pour $x<y$,
\[
f(y)-f(x)=(y-x)+(\sin y-\sin x).
\]
On utilise l'inégalité $|\sin y-\sin x|\leq |y-x|$. Elle ne suffit qu'à obtenir $f(y)-f(x)\geq0$. Pour l'unicité stricte, si $f(y)=f(x)$ avec $x<y$, alors il faudrait $\sin y-\sin x=-(y-x)$. Or l'égalité dans l'inégalité lipschitzienne du sinus sur un intervalle non réduit n'est pas possible ici, car $\sin$ ne coïncide pas avec une droite de pente $-1$ sur $[x,y]$. On peut l'établir par le TVI appliqué à la fonction cosinus ou par les formules trigonométriques. Ainsi $f$ est strictement croissante. De plus $f(x)\to-\infty$ quand $x\to-\infty$ et $f(x)\to+\infty$ quand $x\to+\infty$. Par bijection continue, pour tout $a\in\R$, l'équation admet une unique solution.

\medskip
\textbf{Vérification finale.} Les hypothèses du théorème utilisé ont été explicitées et la conclusion porte bien sur l'objet demandé dans l'énoncé.
\end{corrbox}


\begin{corrbox}{Correction de l'exercice 35 -- Point fixe d'une fonction de $[0,1]$ dans lui-même}

Posons $g(x)=f(x)-x$. La fonction $g$ est continue sur $[0,1]$. Comme $f(0)\in[0,1]$, on a $g(0)=f(0)\geq0$. Comme $f(1)\in[0,1]$, on a $g(1)=f(1)-1\leq0$. Si $g(0)=0$ ou $g(1)=0$, on a un point fixe. Sinon, $g(0)>0$ et $g(1)<0$, et le TVI donne $c\in]0,1[$ tel que $g(c)=0$. Donc $f(c)=c$.

\medskip
\textbf{Vérification finale.} Les hypothèses du théorème utilisé ont été explicitées et la conclusion porte bien sur l'objet demandé dans l'énoncé.
\end{corrbox}


\begin{corrbox}{Correction de l'exercice 36 -- Image d'un segment et bornes atteintes}

On démontre la bornitude par contradiction. Si $f$ n'est pas bornée, pour tout $n\geq1$, il existe $x_n\in[a,b]$ tel que $|f(x_n)|\geq n$. La suite $(x_n)$ est bornée, donc elle possède une sous-suite $(x_{\varphi(n)})$ convergeant vers un $c\in[a,b]$. Par continuité, $f(x_{\varphi(n)})\to f(c)$, donc cette suite d'images est bornée, contradiction avec $|f(x_{\varphi(n)})|\geq\varphi(n)\to+\infty$.
Pour l'atteinte du maximum, soit $M=\sup f([a,b])$, qui existe par bornitude. Pour tout $n$, il existe $x_n\in[a,b]$ tel que $M-1/n\leq f(x_n)\leq M$. Une sous-suite $x_{\varphi(n)}$ converge vers $c\in[a,b]$. Par continuité, $f(x_{\varphi(n)})\to f(c)$, et par encadrement $f(x_{\varphi(n)})\to M$. Donc $f(c)=M$. Le minimum se traite avec l'infimum.

\medskip
\textbf{Vérification finale.} Les hypothèses du théorème utilisé ont été explicitées et la conclusion porte bien sur l'objet demandé dans l'énoncé.
\end{corrbox}


\begin{corrbox}{Correction de l'exercice 37 -- Une fonction continue sans zéros change-t-elle de signe ?}

Supposons qu'il existe $a,b\in I$ tels que $f(a)>0$ et $f(b)<0$. Comme $I$ est un intervalle, le segment entre $a$ et $b$ est inclus dans $I$. La fonction $f$ est continue sur ce segment. Le nombre $0$ est compris entre $f(a)$ et $f(b)$. Par le TVI, il existe $c$ entre $a$ et $b$ tel que $f(c)=0$, contradiction. Donc il n'existe pas deux points où les signes sont opposés. Comme $f$ ne s'annule jamais, elle est partout strictement positive ou partout strictement négative.

\medskip
\textbf{Vérification finale.} Les hypothèses du théorème utilisé ont été explicitées et la conclusion porte bien sur l'objet demandé dans l'énoncé.
\end{corrbox}


\begin{corrbox}{Correction de l'exercice 38 -- Équation avec paramètre et intervalle d'image}

Posons $u=x^2>0$. L'équation devient $u+1/u$. Par l'inégalité $u+1/u\geq2$, avec égalité si et seulement si $u=1$, on a $f(x)\geq2$, égalité pour $x=1$. Réciproquement, pour $a\geq2$, résoudre $u+1/u=a$ revient à résoudre $u^2-au+1=0$, dont le discriminant $a^2-4$ est positif ou nul. Les racines sont positives, car leur somme vaut $a>0$ et leur produit vaut $1>0$. Il existe donc des $u>0$, puis des $x=\sqrt u>0$. L'image est $[2,+\infty[$. Si $a<2$, aucune solution. Si $a=2$, une seule solution $x=1$. Si $a>2$, deux racines positives distinctes en $u$, donc deux solutions positives distinctes en $x$.

\medskip
\textbf{Vérification finale.} Les hypothèses du théorème utilisé ont été explicitées et la conclusion porte bien sur l'objet demandé dans l'énoncé.
\end{corrbox}


\begin{corrbox}{Correction de l'exercice 39 -- Prolongement non trivial}
,
Pour $x\neq0$ suffisamment proche de $0$,
\[
f(x)=\frac{\sqrt{1+x+x^2}-1}{x}\cdot \frac{\sqrt{1+x+x^2}+1}{\sqrt{1+x+x^2}+1}.
\]
On obtient
\[
f(x)=\frac{x+x^2}{x(\sqrt{1+x+x^2}+1)}=\frac{1+x}{\sqrt{1+x+x^2}+1}.
\]
En faisant tendre $x$ vers $0$,
\[
\lim_{x\to0}f(x)=\frac{1}{1+1}=\frac12.
\]
La fonction est donc prolongeable par continuité en $0$ en posant $f(0)=1/2$.

\medskip
\textbf{Vérification finale.} Les hypothèses du théorème utilisé ont été explicitées et la conclusion porte bien sur l'objet demandé dans l'énoncé.
\end{corrbox}


\begin{corrbox}{Correction de l'exercice 40 -- Continuité d'une fonction définie par image rationnelle}

Pour $x\neq0$,
\[
f(x)=\frac{x^2}{|x|}\sin(1/x)=|x|\sin(1/x).
\]
On a donc $|f(x)|\leq |x|$, ce qui donne $f(x)\to0=f(0)$ lorsque $x\to0$. La fonction est continue en $0$. Sur $\R^*$, elle est quotient et composée de fonctions continues avec dénominateur non nul, donc continue. Elle est continue sur $\R$. Pour la parité, pour $x\neq0$,
\[
f(-x)=|-x|\sin(-1/x)=-|x|\sin(1/x)=-f(x),
\]
et $f(0)=0$. Donc $f$ est impaire.

\medskip
\textbf{Vérification finale.} Les hypothèses du théorème utilisé ont été explicitées et la conclusion porte bien sur l'objet demandé dans l'énoncé.
\end{corrbox}


\begin{corrbox}{Correction de l'exercice 41 -- Injectivité sans dérivée}

La fonction est polynomiale, donc continue sur $\R$. Si $x<y$, alors $y-x>0$. On a
\[
y^5-x^5=(y-x)(y^4+y^3x+y^2x^2+yx^3+x^4),
\]
mais le facteur peut être pénible à signer directement. On utilise plutôt que la fonction $t\mapsto t^5$ est strictement croissante sur $\R$ et $t\mapsto t^3$ aussi : si $x<y$, alors $x^5<y^5$, $x^3<y^3$ et $x<y$. En additionnant, $f(x)<f(y)$. Ainsi $f$ est strictement croissante. De plus $f(x)\to-\infty$ quand $x\to-\infty$ et $f(x)\to+\infty$ quand $x\to+\infty$. Par le théorème de la bijection continue, $f$ est une bijection de $\R$ sur $\R$.

\medskip
\textbf{Vérification finale.} Les hypothèses du théorème utilisé ont été explicitées et la conclusion porte bien sur l'objet demandé dans l'énoncé.
\end{corrbox}


\begin{corrbox}{Correction de l'exercice 42 -- Fonction réciproque implicite}

La fonction $f$ est continue sur $\R$. Si $x<y$, alors $x<y$ et $\e^x<\e^y$, donc $x+\e^x<y+\e^y$. Ainsi $f$ est strictement croissante. De plus $f(x)\to-\infty$ quand $x\to-\infty$ et $f(x)\to+\infty$ quand $x\to+\infty$. Elle est donc bijective de $\R$ sur $\R$.
Pour $g(1)$, on cherche $x$ tel que $x+\e^x=1$. On voit que $x=0$ convient, donc $g(1)=0$. Comme $g$ est strictement croissante, $g(y)\leq0$ équivaut à $y\leq f(0)=1$. L'ensemble des solutions est donc $]-\infty,1]$.

\medskip
\textbf{Vérification finale.} Les hypothèses du théorème utilisé ont été explicitées et la conclusion porte bien sur l'objet demandé dans l'énoncé.
\end{corrbox}


\section{Corrigés -- Exercices type oraux et concours}
\begin{corrbox}{Correction de l'exercice 43 -- Oral CCP : racine unique dans un intervalle mobile}

Sur $[0,1]$, $f_n$ est continue. On a $f_n(0)=-1$ et $f_n(1)=1$. Le TVI donne au moins une racine dans $[0,1]$. Si $0\leq x<y\leq1$, alors $x^n<y^n$ et $x<y$, donc $f_n(x)<f_n(y)$. La fonction est strictement croissante, donc la racine est unique. Comme $f_n(0)<0$ et $f_n(1)>0$, elle est dans $]0,1[$.
On a $\alpha_n^n+\alpha_n=1$. Pour $0<\alpha_n<1$, $\alpha_n^{n+1}<\alpha_n^n$, donc
\[
f_{n+1}(\alpha_n)=\alpha_n^{n+1}+\alpha_n-1<0.
\]
Comme $f_{n+1}$ est strictement croissante et s'annule en $\alpha_{n+1}$, on obtient $\alpha_n<\alpha_{n+1}$. La suite est croissante et majorée par $1$, donc elle converge vers un $\ell\in[0,1]$. Si $\ell<1$, alors pour un réel $q$ tel que $\ell<q<1$, on a $\alpha_n\leq q$ à partir d'un certain rang, donc $\alpha_n^n\leq q^n\to0$. En passant à la limite dans $\alpha_n^n+\alpha_n=1$, on obtient $\ell=1$, contradiction. Donc $\ell=1$.

\medskip
\textbf{Vérification finale.} Les hypothèses du théorème utilisé ont été explicitées et la conclusion porte bien sur l'objet demandé dans l'énoncé.
\end{corrbox}


\begin{corrbox}{Correction de l'exercice 44 -- Oral Mines : équation fonctionnelle de continuité}

L'identité $f(f(x))=x$ donne immédiatement l'injectivité : si $f(a)=f(b)$, alors en appliquant $f$ aux deux membres, $f(f(a))=f(f(b))$, donc $a=b$. Elle donne aussi la surjectivité : pour tout $y\in\R$, en posant $x=f(y)$, on a $f(x)=f(f(y))=y$. Donc $f$ est bijective.
Comme $f$ est continue et injective sur l'intervalle $\R$, elle est strictement monotone. Si elle était strictement décroissante, sa réciproque serait aussi strictement décroissante. Or $f=f^{-1}$ car $f\circ f=\Id$. C'est compatible. Si elle est strictement croissante et n'est pas l'identité, on obtient une contradiction : s'il existe $a$ tel que $f(a)>a$, alors par croissance $f(f(a))>f(a)>a$, contradiction avec $f(f(a))=a$. De même si $f(a)<a$. Donc une solution croissante est nécessairement l'identité. Ainsi, si $f$ n'est pas l'identité, elle est strictement décroissante.

\medskip
\textbf{Vérification finale.} Les hypothèses du théorème utilisé ont été explicitées et la conclusion porte bien sur l'objet demandé dans l'énoncé.
\end{corrbox}


\begin{corrbox}{Correction de l'exercice 45 -- Oral Centrale : image et antécédents d'une fonction oscillante amortie}

Sur $\R^*$, $f$ est continue. En $0$, on a $|x\sin(1/x)|\leq |x|$, donc $f(x)\to0=f(0)$ : $f$ est continue sur $\R$.
Les zéros non nuls sont les $x$ tels que $\sin(1/x)=0$, donc $1/x=k\pi$, soit $x=1/(k\pi)$. Ces zéros s'accumulent en $0$. De plus, les points où $\sin(1/x)=1$ et $\sin(1/x)=-1$ donnent des valeurs positives et négatives arbitrairement proches de $0$. Une fonction monotone sur un intervalle qui possède plusieurs zéros et change de signe de manière alternée ne peut présenter une telle oscillation. Plus précisément, entre deux zéros consécutifs on trouve un extremum de signe alterné, ce qui contredit la monotonie sur tout voisinage de $0$. Enfin, puisque $1/(k\pi)\to0$, tout voisinage de $0$ contient une infinité de zéros de $f$.

\medskip
\textbf{Vérification finale.} Les hypothèses du théorème utilisé ont été explicitées et la conclusion porte bien sur l'objet demandé dans l'énoncé.
\end{corrbox}


\begin{corrbox}{Correction de l'exercice 46 -- Oral Mines : continuité et densité élémentaire}
,
Soit $a\in\R$. Par densité de $\mathbb Q$ dans $\R$, il existe une suite $(r_n)$ de rationnels telle que $r_n\to a$. Pour tout $n$, $f(r_n)=0$. Comme $f$ est continue en $a$, la caractérisation séquentielle donne $f(r_n)\to f(a)$. Or la suite $f(r_n)$ est constante égale à $0$, donc sa limite vaut $0$. Ainsi $f(a)=0$. Comme $a$ est arbitraire, $f$ est identiquement nulle sur $\R$.

\medskip
\textbf{Vérification finale.} Les hypothèses du théorème utilisé ont été explicitées et la conclusion porte bien sur l'objet demandé dans l'énoncé.
\end{corrbox}


\begin{corrbox}{Correction de l'exercice 47 -- Oral CCP : comparaison et fausse intuition graphique}

Pour $x\neq0$,
\[
\left|\frac{x^2\sin(1/x)}{|x|+x^2}\right|\leq \frac{x^2}{|x|+x^2}.
\]
Comme $|x|+x^2\geq |x|$, on a
\[
0\leq \frac{x^2}{|x|+x^2}\leq \frac{x^2}{|x|}=|x|.
\]
Donc le majorant tend vers $0$. Par encadrement,
\[
\frac{x^2\sin(1/x)}{|x|+x^2}\to0.
\]
La limite existe et vaut $0$, malgré l'oscillation de $\sin(1/x)$.

\medskip
\textbf{Vérification finale.} Les hypothèses du théorème utilisé ont été explicitées et la conclusion porte bien sur l'objet demandé dans l'énoncé.
\end{corrbox}


\begin{corrbox}{Correction de l'exercice 48 -- Oral Centrale : suite d'antécédents}

Posons $g_n(x)=f(x)-x^n$. Cette fonction est continue sur $[0,1]$. On a $g_n(0)=f(0)\geq0$ puisque $f(0)\in[0,1]$. On a aussi $g_n(1)=f(1)-1\leq0$ puisque $f(1)\in[0,1]$. Si l'une des valeurs est nulle, on a une solution. Sinon, $g_n(0)>0$ et $g_n(1)<0$, donc le TVI donne $x_n\in]0,1[$ tel que $g_n(x_n)=0$, c'est-à-dire $f(x_n)=x_n^n$.

\medskip
\textbf{Vérification finale.} Les hypothèses du théorème utilisé ont été explicitées et la conclusion porte bien sur l'objet demandé dans l'énoncé.
\end{corrbox}


\begin{corrbox}{Correction de l'exercice 49 -- Problème long : réciproque et équation implicite}

La fonction $f$ est continue sur $]-1,+\infty[$. Si $-1<x<y$, alors $x<y$ et $1+x<1+y$, donc $\ln(1+x)<\ln(1+y)$. Ainsi $f(x)<f(y)$ : $f$ est strictement croissante. De plus
\[
\lim_{x\to-1^+}f(x)=-\infty,\qquad \lim_{x\to+\infty}f(x)=+\infty.
\]
Par le théorème de la bijection continue, $f$ réalise une bijection de $]-1,+\infty[$ sur $\R$.
Comme $f(0)=0$ et comme $f$ est strictement croissante, on a $g(y)<0$ si $y<0$, $g(0)=0$ et $g(y)>0$ si $y>0$.
Pour résoudre $g(y)\geq y$, posons $x=g(y)$. Alors $y=f(x)=x+\ln(1+x)$. L'inégalité $g(y)\geq y$ devient
\[
x\geq x+\ln(1+x),
\]
donc $\ln(1+x)\leq0$, soit $1+x\leq1$, donc $x\leq0$. Comme $f$ est croissante, cela équivaut à $y=f(x)\leq f(0)=0$. Ainsi l'ensemble des solutions est $]-\infty,0]$.
Enfin $g(y)\leq y+1$ équivaut, avec $x=g(y)$, à
\[
x\leq x+\ln(1+x)+1,
\]
donc à $\ln(1+x)\geq-1$, soit $1+x\geq e^{-1}$, c'est-à-dire $x\geq e^{-1}-1$. En revenant à $y=f(x)$ et en utilisant la croissance de $f$, cela équivaut à
\[
y\geq f(e^{-1}-1)=e^{-1}-1+\ln(e^{-1})=e^{-1}-2.
\]

\medskip
\textbf{Vérification finale.} Les hypothèses du théorème utilisé ont été explicitées et la conclusion porte bien sur l'objet demandé dans l'énoncé.
\end{corrbox}


\begin{corrbox}{Correction de l'exercice 50 -- Problème long : existence, unicité, encadrement}

Posons $f_a(x)=\e^{-x}-ax$ sur $[0,+\infty[$. La fonction est continue. On a $f_a(0)=1>0$ et $f_a(1/a)=\e^{-1/a}-1<0$. Par le TVI, il existe une solution dans $]0,1/a[$. Si $0<x<y$, alors $\e^{-y}<\e^{-x}$ et $-ay<-ax$, donc $f_a(y)<f_a(x)$ : $f_a$ est strictement décroissante. La solution est unique.
L'encadrement $0<x_a<1/a$ vient de la construction précédente. Si $0<a<b$, alors
\[
\e^{-x_a}=a x_a < b x_a,
\]
donc $f_b(x_a)=\e^{-x_a}-b x_a<0$. Comme $f_b$ est strictement décroissante et s'annule en $x_b$, on doit avoir $x_b<x_a$.
Lorsque $a\to+\infty$, $0<x_a<1/a$ donne $x_a\to0$. Lorsque $a\to0^+$, si $(x_a)$ ne tendait pas vers $+\infty$, on pourrait extraire une sous-famille bornée ; le membre $a x_a$ tendrait alors vers $0$, tandis que $\e^{-x_a}$ resterait positive le long d'une sous-suite convergente finie, contradiction. Donc $x_a\to+\infty$.

\medskip
\textbf{Vérification finale.} Les hypothèses du théorème utilisé ont été explicitées et la conclusion porte bien sur l'objet demandé dans l'énoncé.
\end{corrbox}


\section{Corrigés des problèmes longs}
\begin{corrbox}{Correction -- Problème A -- Théorème des valeurs intermédiaires renforcé}

L'ensemble $E$ est non vide car $a\in E$, puisque $f(a)<0$. Il est majoré par $b$, car $E\subset[a,b]$. Il admet donc une borne supérieure $c\in[a,b]$.
Montrons que $f(c)\leq0$. Par caractérisation de la borne supérieure, il existe une suite $(x_n)$ d'éléments de $E$ telle que $x_n\to c$. Pour tout $n$, $f(x_n)\leq0$. Par continuité de $f$ en $c$, $f(x_n)\to f(c)$, donc $f(c)\leq0$.
Montrons que $f(c)\geq0$. Si $c=b$, alors $f(c)=f(b)>0$, ce qui contredirait $f(c)\leq0$ ; en réalité le raisonnement précédent force alors une analyse fine, mais comme $f(b)>0$, on a nécessairement $c<b$. Pour $n$ assez grand, $c+1/n\in[a,b]$ et $c+1/n>c$. Par définition de la borne supérieure, $c+1/n\notin E$, donc $f(c+1/n)>0$. En passant à la limite, par continuité, $f(c)\geq0$. Ainsi $f(c)=0$.
Si $f$ est strictement croissante, elle est injective. Elle ne peut donc avoir deux zéros distincts. Le zéro $c$ est unique.
Pour $p(x)=x^3+x-1$, on a $p(0)=-1<0$ et $p(1)=1>0$. La fonction est continue. De plus, si $x<y$, alors $x^3<y^3$ et $x<y$, donc $p(x)<p(y)$. Elle est strictement croissante. L'équation admet donc une unique racine dans $[0,1]$.

\medskip
\textbf{Bilan méthode.} Une existence vient en général du TVI ou de la définition de l'image ; une unicité vient de l'injectivité, souvent obtenue par stricte monotonie.
\end{corrbox}
\begin{corrbox}{Correction -- Problème B -- Bijection continue et paramètres}

La fonction $f_\lambda$ est continue sur $]0,+\infty[$. Si $0<x<y$, alors $x<y$ et $\ln x<\ln y$, donc $x+\lambda\ln x<y+\lambda\ln y$. Elle est strictement croissante. De plus $f_\lambda(x)\to-\infty$ quand $x\to0^+$ et $f_\lambda(x)\to+\infty$ quand $x\to+\infty$. Le théorème de la bijection continue donne une bijection de $]0,+\infty[$ sur $\R$.
Pour $\lambda=1$ et $y=1$, on cherche $x+\ln x=1$. Comme $x=1$ convient, l'unicité donne $g_1(1)=1$.
Soit $x=g_\lambda(y)$. Alors $y=x+\lambda\ln x$. Pour $\mu>\lambda$,
\[
f_\mu(x)=x+\mu\ln x=y+(\mu-\lambda)\ln x.
\]
Si $x<1$, alors $\ln x<0$, donc $f_\mu(x)<y$. Comme $f_\mu$ est strictement croissante et s'annule au niveau $y$ en $g_\mu(y)$, on obtient $g_\mu(y)>x=g_\lambda(y)$. Si $x=1$, alors $f_\mu(1)=y$ et $g_\mu(y)=1$. Si $x>1$, alors $f_\mu(x)>y$, donc $g_\mu(y)<x$.
Enfin, $g_\lambda(y)\geq1$ équivaut, par croissance de $f_\lambda$, à $y\geq f_\lambda(1)=1$. Donc l'ensemble des $y$ cherchés est $[1,+\infty[$.

\medskip
\textbf{Bilan méthode.} Une existence vient en général du TVI ou de la définition de l'image ; une unicité vient de l'injectivité, souvent obtenue par stricte monotonie.
\end{corrbox}
\begin{corrbox}{Correction -- Problème C -- Continuité, zéros et oscillations}

La continuité sur $\R^*$ est immédiate par opérations et composition. En $0$, $|f(x)|\leq |x|$ pour $x\neq0$, donc $f(x)\to0=f(0)$ par encadrement.
Les zéros non nuls sont donnés par $\sin(1/x)=0$, donc $1/x=k\pi$, soit $x_k=1/(k\pi)$ pour $k\in\N^*$. Alors $x_k\to0$ et $f(x_k)=0$.
Pour obtenir des valeurs positives et négatives, on prend
\[
u_n=\frac1{\pi/2+2n\pi},\qquad v_n=\frac1{-\pi/2+2n\pi}
\]
pour $n$ assez grand. Alors $u_n\to0$, $v_n\to0$, $f(u_n)=u_n>0$ et $f(v_n)=-v_n$ ; selon le signe de $v_n$, on choisit au besoin une suite positive correspondant à $3\pi/2+2n\pi$ pour obtenir une valeur négative. Dans tout voisinage de $0$, il existe plusieurs zéros distincts. Une fonction injective ne peut pas prendre la même valeur $0$ en deux points distincts. Donc $f$ n'est injective sur aucun voisinage symétrique de $0$.
Pour $g(x)=x^2\sin(1/x)$, on a $|g(x)|\leq x^2$, donc la continuité en $0$ est encore plus forte. Néanmoins, les zéros $1/(k\pi)$ subsistent et empêchent aussi l'injectivité sur tout voisinage de $0$.

\medskip
\textbf{Bilan méthode.} Une existence vient en général du TVI ou de la définition de l'image ; une unicité vient de l'injectivité, souvent obtenue par stricte monotonie.
\end{corrbox}
\newpage
\section{Annexe -- Fiches méthodes détaillées}
\begin{fichebox}
\textbf{Calculer une limite finie en un point.}
\begin{enumerate}
\item Chercher d'abord le domaine de définition près du point.
\item Simplifier uniquement pour $x
eq a$ si une factorisation annule une forme indéterminée.
\item Utiliser les limites usuelles seulement après avoir vérifié la composition.
\item Conclure par une phrase indiquant la valeur de la limite.
\end{enumerate}
\end{fichebox}
\begin{fichebox}
\textbf{Prouver une non-existence de limite.}
\begin{enumerate}
\item Chercher deux chemins ou deux suites tendant vers le même point.
\item Calculer les limites des images le long de ces suites.
\item Si les limites obtenues sont différentes, la limite globale n'existe pas.
\item Ne pas confondre oscillation et divergence : une fonction oscillante peut parfois avoir une limite après amortissement.
\end{enumerate}
\end{fichebox}
\begin{fichebox}
\textbf{Étudier une continuité.}
\begin{enumerate}
\item Commencer par les points où la formule est usuelle.
\item Isoler les points de raccord ou les points exclus du domaine.
\item Comparer la limite au point et la valeur de la fonction.
\item Pour un prolongement, la limite doit exister et être finie.
\end{enumerate}
\end{fichebox}
\begin{fichebox}
\textbf{Utiliser le TVI.}
\begin{enumerate}
\item Définir une fonction continue sur un segment.
\item Calculer les valeurs aux bornes.
\item Vérifier que le niveau cherché est compris entre ces deux valeurs.
\item Conclure à l'existence, jamais à l'unicité sans argument supplémentaire.
\end{enumerate}
\end{fichebox}
\begin{fichebox}
\textbf{Prouver une bijection.}
\begin{enumerate}
\item Choisir clairement l'ensemble de départ et l'ensemble d'arrivée.
\item Montrer la continuité sur un intervalle.
\item Montrer la stricte monotonie pour obtenir l'injectivité.
\item Calculer l'image de l'intervalle pour obtenir la surjectivité.
\item Appliquer le théorème de la bijection continue.
\end{enumerate}
\end{fichebox}
\section{Annexe -- Banque de questions rapides pour colles}
\begin{exobox}{Question rapide 01}
Donner un exemple de fonction ayant une limite en $0$ mais n'étant pas définie en $0$, puis construire son prolongement par continuité.
\end{exobox}
\begin{exobox}{Question rapide 02}
Donner un exemple de fonction admettant deux limites latérales différentes en un point, et expliquer la conclusion.
\end{exobox}
\begin{exobox}{Question rapide 03}
Énoncer le théorème des valeurs intermédiaires puis préciser pourquoi il ne donne pas l'unicité.
\end{exobox}
\begin{exobox}{Question rapide 04}
Énoncer le théorème de la bijection continue et donner un exemple de fonction bijective dont la réciproque se calcule explicitement.
\end{exobox}
\begin{exobox}{Question rapide 05}
Expliquer la différence entre l'image réciproque d'un ensemble $f^{-1}(B)$ et la fonction réciproque $f^{-1}$ d'une bijection.
\end{exobox}
\begin{exobox}{Question rapide 06}
Donner un exemple de fonction ayant une limite en $0$ mais n'étant pas définie en $0$, puis construire son prolongement par continuité.
\end{exobox}
\begin{exobox}{Question rapide 07}
Donner un exemple de fonction admettant deux limites latérales différentes en un point, et expliquer la conclusion.
\end{exobox}
\begin{exobox}{Question rapide 08}
Énoncer le théorème des valeurs intermédiaires puis préciser pourquoi il ne donne pas l'unicité.
\end{exobox}
\begin{exobox}{Question rapide 09}
Énoncer le théorème de la bijection continue et donner un exemple de fonction bijective dont la réciproque se calcule explicitement.
\end{exobox}
\begin{exobox}{Question rapide 10}
Expliquer la différence entre l'image réciproque d'un ensemble $f^{-1}(B)$ et la fonction réciproque $f^{-1}$ d'une bijection.
\end{exobox}
\begin{exobox}{Question rapide 11}
Donner un exemple de fonction ayant une limite en $0$ mais n'étant pas définie en $0$, puis construire son prolongement par continuité.
\end{exobox}
\begin{exobox}{Question rapide 12}
Donner un exemple de fonction admettant deux limites latérales différentes en un point, et expliquer la conclusion.
\end{exobox}
\begin{exobox}{Question rapide 13}
Énoncer le théorème des valeurs intermédiaires puis préciser pourquoi il ne donne pas l'unicité.
\end{exobox}
\begin{exobox}{Question rapide 14}
Énoncer le théorème de la bijection continue et donner un exemple de fonction bijective dont la réciproque se calcule explicitement.
\end{exobox}
\begin{exobox}{Question rapide 15}
Expliquer la différence entre l'image réciproque d'un ensemble $f^{-1}(B)$ et la fonction réciproque $f^{-1}$ d'une bijection.
\end{exobox}
\begin{exobox}{Question rapide 16}
Donner un exemple de fonction ayant une limite en $0$ mais n'étant pas définie en $0$, puis construire son prolongement par continuité.
\end{exobox}
\begin{exobox}{Question rapide 17}
Donner un exemple de fonction admettant deux limites latérales différentes en un point, et expliquer la conclusion.
\end{exobox}
\begin{exobox}{Question rapide 18}
Énoncer le théorème des valeurs intermédiaires puis préciser pourquoi il ne donne pas l'unicité.
\end{exobox}
\begin{exobox}{Question rapide 19}
Énoncer le théorème de la bijection continue et donner un exemple de fonction bijective dont la réciproque se calcule explicitement.
\end{exobox}
\begin{exobox}{Question rapide 20}
Expliquer la différence entre l'image réciproque d'un ensemble $f^{-1}(B)$ et la fonction réciproque $f^{-1}$ d'une bijection.
\end{exobox}
\begin{exobox}{Question rapide 21}
Donner un exemple de fonction ayant une limite en $0$ mais n'étant pas définie en $0$, puis construire son prolongement par continuité.
\end{exobox}
\begin{exobox}{Question rapide 22}
Donner un exemple de fonction admettant deux limites latérales différentes en un point, et expliquer la conclusion.
\end{exobox}
\begin{exobox}{Question rapide 23}
Énoncer le théorème des valeurs intermédiaires puis préciser pourquoi il ne donne pas l'unicité.
\end{exobox}
\begin{exobox}{Question rapide 24}
Énoncer le théorème de la bijection continue et donner un exemple de fonction bijective dont la réciproque se calcule explicitement.
\end{exobox}
\begin{exobox}{Question rapide 25}
Expliquer la différence entre l'image réciproque d'un ensemble $f^{-1}(B)$ et la fonction réciproque $f^{-1}$ d'une bijection.
\end{exobox}
\begin{exobox}{Question rapide 26}
Donner un exemple de fonction ayant une limite en $0$ mais n'étant pas définie en $0$, puis construire son prolongement par continuité.
\end{exobox}
\begin{exobox}{Question rapide 27}
Donner un exemple de fonction admettant deux limites latérales différentes en un point, et expliquer la conclusion.
\end{exobox}
\begin{exobox}{Question rapide 28}
Énoncer le théorème des valeurs intermédiaires puis préciser pourquoi il ne donne pas l'unicité.
\end{exobox}
\begin{exobox}{Question rapide 29}
Énoncer le théorème de la bijection continue et donner un exemple de fonction bijective dont la réciproque se calcule explicitement.
\end{exobox}
\begin{exobox}{Question rapide 30}
Expliquer la différence entre l'image réciproque d'un ensemble $f^{-1}(B)$ et la fonction réciproque $f^{-1}$ d'une bijection.
\end{exobox}
\begin{exobox}{Question rapide 31}
Donner un exemple de fonction ayant une limite en $0$ mais n'étant pas définie en $0$, puis construire son prolongement par continuité.
\end{exobox}
\begin{exobox}{Question rapide 32}
Donner un exemple de fonction admettant deux limites latérales différentes en un point, et expliquer la conclusion.
\end{exobox}
\begin{exobox}{Question rapide 33}
Énoncer le théorème des valeurs intermédiaires puis préciser pourquoi il ne donne pas l'unicité.
\end{exobox}
\begin{exobox}{Question rapide 34}
Énoncer le théorème de la bijection continue et donner un exemple de fonction bijective dont la réciproque se calcule explicitement.
\end{exobox}
\begin{exobox}{Question rapide 35}
Expliquer la différence entre l'image réciproque d'un ensemble $f^{-1}(B)$ et la fonction réciproque $f^{-1}$ d'une bijection.
\end{exobox}
\begin{exobox}{Question rapide 36}
Donner un exemple de fonction ayant une limite en $0$ mais n'étant pas définie en $0$, puis construire son prolongement par continuité.
\end{exobox}
\begin{exobox}{Question rapide 37}
Donner un exemple de fonction admettant deux limites latérales différentes en un point, et expliquer la conclusion.
\end{exobox}
\begin{exobox}{Question rapide 38}
Énoncer le théorème des valeurs intermédiaires puis préciser pourquoi il ne donne pas l'unicité.
\end{exobox}
\begin{exobox}{Question rapide 39}
Énoncer le théorème de la bijection continue et donner un exemple de fonction bijective dont la réciproque se calcule explicitement.
\end{exobox}
\begin{exobox}{Question rapide 40}
Expliquer la différence entre l'image réciproque d'un ensemble $f^{-1}(B)$ et la fonction réciproque $f^{-1}$ d'une bijection.
\end{exobox}
\section{Annexe -- Corrigés des questions rapides}
\begin{corrbox}{Correction rapide 01}
Un exemple est $f(x)=\frac{x^2-1}{x-1}$ sur $\R\setminus\{1\}$, translaté au besoin. Pour $x\neq1$, $f(x)=x+1$, donc la limite en $1$ vaut $2$ et le prolongement consiste à poser $f(1)=2$.
\end{corrbox}
\begin{corrbox}{Correction rapide 02}
La fonction $x\mapsto |x|/x$ sur $\R^*$ admet pour limite à droite en $0$ la valeur $1$ et pour limite à gauche la valeur $-1$. La limite en $0$ n'existe donc pas.
\end{corrbox}
\begin{corrbox}{Correction rapide 03}
Si $f$ est continue sur $[a,b]$, tout réel compris entre $f(a)$ et $f(b)$ est atteint. Le théorème donne une existence. L'unicité nécessite par exemple l'injectivité ou la stricte monotonie.
\end{corrbox}
\begin{corrbox}{Correction rapide 04}
Si $f$ est continue et strictement monotone sur un intervalle $I$, elle réalise une bijection de $I$ sur $f(I)$ et sa réciproque est continue et strictement monotone. Exemple : $x\mapsto x^2$ de $[0,+\infty[$ sur $[0,+\infty[$, de réciproque $y\mapsto\sqrt y$.
\end{corrbox}
\begin{corrbox}{Correction rapide 05}
L'image réciproque d'un ensemble existe pour toute fonction : $f^{-1}(B)=\{x\mid f(x)\in B\}$. La fonction réciproque n'existe que si $f$ est bijective et associe à chaque $y$ son unique antécédent.
\end{corrbox}
\begin{corrbox}{Correction rapide 06}
Un exemple est $f(x)=\frac{x^2-1}{x-1}$ sur $\R\setminus\{1\}$, translaté au besoin. Pour $x\neq1$, $f(x)=x+1$, donc la limite en $1$ vaut $2$ et le prolongement consiste à poser $f(1)=2$.
\end{corrbox}
\begin{corrbox}{Correction rapide 07}
La fonction $x\mapsto |x|/x$ sur $\R^*$ admet pour limite à droite en $0$ la valeur $1$ et pour limite à gauche la valeur $-1$. La limite en $0$ n'existe donc pas.
\end{corrbox}
\begin{corrbox}{Correction rapide 08}
Si $f$ est continue sur $[a,b]$, tout réel compris entre $f(a)$ et $f(b)$ est atteint. Le théorème donne une existence. L'unicité nécessite par exemple l'injectivité ou la stricte monotonie.
\end{corrbox}
\begin{corrbox}{Correction rapide 09}
Si $f$ est continue et strictement monotone sur un intervalle $I$, elle réalise une bijection de $I$ sur $f(I)$ et sa réciproque est continue et strictement monotone. Exemple : $x\mapsto x^2$ de $[0,+\infty[$ sur $[0,+\infty[$, de réciproque $y\mapsto\sqrt y$.
\end{corrbox}
\begin{corrbox}{Correction rapide 10}
L'image réciproque d'un ensemble existe pour toute fonction : $f^{-1}(B)=\{x\mid f(x)\in B\}$. La fonction réciproque n'existe que si $f$ est bijective et associe à chaque $y$ son unique antécédent.
\end{corrbox}
\begin{corrbox}{Correction rapide 11}
Un exemple est $f(x)=\frac{x^2-1}{x-1}$ sur $\R\setminus\{1\}$, translaté au besoin. Pour $x\neq1$, $f(x)=x+1$, donc la limite en $1$ vaut $2$ et le prolongement consiste à poser $f(1)=2$.
\end{corrbox}
\begin{corrbox}{Correction rapide 12}
La fonction $x\mapsto |x|/x$ sur $\R^*$ admet pour limite à droite en $0$ la valeur $1$ et pour limite à gauche la valeur $-1$. La limite en $0$ n'existe donc pas.
\end{corrbox}
\begin{corrbox}{Correction rapide 13}
Si $f$ est continue sur $[a,b]$, tout réel compris entre $f(a)$ et $f(b)$ est atteint. Le théorème donne une existence. L'unicité nécessite par exemple l'injectivité ou la stricte monotonie.
\end{corrbox}
\begin{corrbox}{Correction rapide 14}
Si $f$ est continue et strictement monotone sur un intervalle $I$, elle réalise une bijection de $I$ sur $f(I)$ et sa réciproque est continue et strictement monotone. Exemple : $x\mapsto x^2$ de $[0,+\infty[$ sur $[0,+\infty[$, de réciproque $y\mapsto\sqrt y$.
\end{corrbox}
\begin{corrbox}{Correction rapide 15}
L'image réciproque d'un ensemble existe pour toute fonction : $f^{-1}(B)=\{x\mid f(x)\in B\}$. La fonction réciproque n'existe que si $f$ est bijective et associe à chaque $y$ son unique antécédent.
\end{corrbox}
\begin{corrbox}{Correction rapide 16}
Un exemple est $f(x)=\frac{x^2-1}{x-1}$ sur $\R\setminus\{1\}$, translaté au besoin. Pour $x\neq1$, $f(x)=x+1$, donc la limite en $1$ vaut $2$ et le prolongement consiste à poser $f(1)=2$.
\end{corrbox}
\begin{corrbox}{Correction rapide 17}
La fonction $x\mapsto |x|/x$ sur $\R^*$ admet pour limite à droite en $0$ la valeur $1$ et pour limite à gauche la valeur $-1$. La limite en $0$ n'existe donc pas.
\end{corrbox}
\begin{corrbox}{Correction rapide 18}
Si $f$ est continue sur $[a,b]$, tout réel compris entre $f(a)$ et $f(b)$ est atteint. Le théorème donne une existence. L'unicité nécessite par exemple l'injectivité ou la stricte monotonie.
\end{corrbox}
\begin{corrbox}{Correction rapide 19}
Si $f$ est continue et strictement monotone sur un intervalle $I$, elle réalise une bijection de $I$ sur $f(I)$ et sa réciproque est continue et strictement monotone. Exemple : $x\mapsto x^2$ de $[0,+\infty[$ sur $[0,+\infty[$, de réciproque $y\mapsto\sqrt y$.
\end{corrbox}
\begin{corrbox}{Correction rapide 20}
L'image réciproque d'un ensemble existe pour toute fonction : $f^{-1}(B)=\{x\mid f(x)\in B\}$. La fonction réciproque n'existe que si $f$ est bijective et associe à chaque $y$ son unique antécédent.
\end{corrbox}
\begin{corrbox}{Correction rapide 21}
Un exemple est $f(x)=\frac{x^2-1}{x-1}$ sur $\R\setminus\{1\}$, translaté au besoin. Pour $x\neq1$, $f(x)=x+1$, donc la limite en $1$ vaut $2$ et le prolongement consiste à poser $f(1)=2$.
\end{corrbox}
\begin{corrbox}{Correction rapide 22}
La fonction $x\mapsto |x|/x$ sur $\R^*$ admet pour limite à droite en $0$ la valeur $1$ et pour limite à gauche la valeur $-1$. La limite en $0$ n'existe donc pas.
\end{corrbox}
\begin{corrbox}{Correction rapide 23}
Si $f$ est continue sur $[a,b]$, tout réel compris entre $f(a)$ et $f(b)$ est atteint. Le théorème donne une existence. L'unicité nécessite par exemple l'injectivité ou la stricte monotonie.
\end{corrbox}
\begin{corrbox}{Correction rapide 24}
Si $f$ est continue et strictement monotone sur un intervalle $I$, elle réalise une bijection de $I$ sur $f(I)$ et sa réciproque est continue et strictement monotone. Exemple : $x\mapsto x^2$ de $[0,+\infty[$ sur $[0,+\infty[$, de réciproque $y\mapsto\sqrt y$.
\end{corrbox}
\begin{corrbox}{Correction rapide 25}
L'image réciproque d'un ensemble existe pour toute fonction : $f^{-1}(B)=\{x\mid f(x)\in B\}$. La fonction réciproque n'existe que si $f$ est bijective et associe à chaque $y$ son unique antécédent.
\end{corrbox}
\begin{corrbox}{Correction rapide 26}
Un exemple est $f(x)=\frac{x^2-1}{x-1}$ sur $\R\setminus\{1\}$, translaté au besoin. Pour $x\neq1$, $f(x)=x+1$, donc la limite en $1$ vaut $2$ et le prolongement consiste à poser $f(1)=2$.
\end{corrbox}
\begin{corrbox}{Correction rapide 27}
La fonction $x\mapsto |x|/x$ sur $\R^*$ admet pour limite à droite en $0$ la valeur $1$ et pour limite à gauche la valeur $-1$. La limite en $0$ n'existe donc pas.
\end{corrbox}
\begin{corrbox}{Correction rapide 28}
Si $f$ est continue sur $[a,b]$, tout réel compris entre $f(a)$ et $f(b)$ est atteint. Le théorème donne une existence. L'unicité nécessite par exemple l'injectivité ou la stricte monotonie.
\end{corrbox}
\begin{corrbox}{Correction rapide 29}
Si $f$ est continue et strictement monotone sur un intervalle $I$, elle réalise une bijection de $I$ sur $f(I)$ et sa réciproque est continue et strictement monotone. Exemple : $x\mapsto x^2$ de $[0,+\infty[$ sur $[0,+\infty[$, de réciproque $y\mapsto\sqrt y$.
\end{corrbox}
\begin{corrbox}{Correction rapide 30}
L'image réciproque d'un ensemble existe pour toute fonction : $f^{-1}(B)=\{x\mid f(x)\in B\}$. La fonction réciproque n'existe que si $f$ est bijective et associe à chaque $y$ son unique antécédent.
\end{corrbox}
\begin{corrbox}{Correction rapide 31}
Un exemple est $f(x)=\frac{x^2-1}{x-1}$ sur $\R\setminus\{1\}$, translaté au besoin. Pour $x\neq1$, $f(x)=x+1$, donc la limite en $1$ vaut $2$ et le prolongement consiste à poser $f(1)=2$.
\end{corrbox}
\begin{corrbox}{Correction rapide 32}
La fonction $x\mapsto |x|/x$ sur $\R^*$ admet pour limite à droite en $0$ la valeur $1$ et pour limite à gauche la valeur $-1$. La limite en $0$ n'existe donc pas.
\end{corrbox}
\begin{corrbox}{Correction rapide 33}
Si $f$ est continue sur $[a,b]$, tout réel compris entre $f(a)$ et $f(b)$ est atteint. Le théorème donne une existence. L'unicité nécessite par exemple l'injectivité ou la stricte monotonie.
\end{corrbox}
\begin{corrbox}{Correction rapide 34}
Si $f$ est continue et strictement monotone sur un intervalle $I$, elle réalise une bijection de $I$ sur $f(I)$ et sa réciproque est continue et strictement monotone. Exemple : $x\mapsto x^2$ de $[0,+\infty[$ sur $[0,+\infty[$, de réciproque $y\mapsto\sqrt y$.
\end{corrbox}
\begin{corrbox}{Correction rapide 35}
L'image réciproque d'un ensemble existe pour toute fonction : $f^{-1}(B)=\{x\mid f(x)\in B\}$. La fonction réciproque n'existe que si $f$ est bijective et associe à chaque $y$ son unique antécédent.
\end{corrbox}
\begin{corrbox}{Correction rapide 36}
Un exemple est $f(x)=\frac{x^2-1}{x-1}$ sur $\R\setminus\{1\}$, translaté au besoin. Pour $x\neq1$, $f(x)=x+1$, donc la limite en $1$ vaut $2$ et le prolongement consiste à poser $f(1)=2$.
\end{corrbox}
\begin{corrbox}{Correction rapide 37}
La fonction $x\mapsto |x|/x$ sur $\R^*$ admet pour limite à droite en $0$ la valeur $1$ et pour limite à gauche la valeur $-1$. La limite en $0$ n'existe donc pas.
\end{corrbox}
\begin{corrbox}{Correction rapide 38}
Si $f$ est continue sur $[a,b]$, tout réel compris entre $f(a)$ et $f(b)$ est atteint. Le théorème donne une existence. L'unicité nécessite par exemple l'injectivité ou la stricte monotonie.
\end{corrbox}
\begin{corrbox}{Correction rapide 39}
Si $f$ est continue et strictement monotone sur un intervalle $I$, elle réalise une bijection de $I$ sur $f(I)$ et sa réciproque est continue et strictement monotone. Exemple : $x\mapsto x^2$ de $[0,+\infty[$ sur $[0,+\infty[$, de réciproque $y\mapsto\sqrt y$.
\end{corrbox}
\begin{corrbox}{Correction rapide 40}
L'image réciproque d'un ensemble existe pour toute fonction : $f^{-1}(B)=\{x\mid f(x)\in B\}$. La fonction réciproque n'existe que si $f$ est bijective et associe à chaque $y$ son unique antécédent.
\end{corrbox}
\section*{Conclusion}
\addcontentsline{toc}{section}{Conclusion}
Ce TD couvre les compétences fondamentales du chapitre Analyse 3 : calculer, démontrer, prolonger, utiliser le TVI, identifier les hypothèses de continuité et de monotonie, puis construire des bijections et leurs réciproques. Les exercices de fin de document peuvent servir de banque de colles ou de révisions rapides avant un devoir surveillé.
% Ligne pédagogique de structure 1: vérifier les hypothèses avant d'appliquer un théorème.
% Ligne pédagogique de structure 2: vérifier les hypothèses avant d'appliquer un théorème.
% Ligne pédagogique de structure 3: vérifier les hypothèses avant d'appliquer un théorème.
% Ligne pédagogique de structure 4: vérifier les hypothèses avant d'appliquer un théorème.
% Ligne pédagogique de structure 5: vérifier les hypothèses avant d'appliquer un théorème.
% Ligne pédagogique de structure 6: vérifier les hypothèses avant d'appliquer un théorème.
% Ligne pédagogique de structure 7: vérifier les hypothèses avant d'appliquer un théorème.
% Ligne pédagogique de structure 8: vérifier les hypothèses avant d'appliquer un théorème.
% Ligne pédagogique de structure 9: vérifier les hypothèses avant d'appliquer un théorème.
% Ligne pédagogique de structure 10: vérifier les hypothèses avant d'appliquer un théorème.
% Ligne pédagogique de structure 11: vérifier les hypothèses avant d'appliquer un théorème.
% Ligne pédagogique de structure 12: vérifier les hypothèses avant d'appliquer un théorème.
% Ligne pédagogique de structure 13: vérifier les hypothèses avant d'appliquer un théorème.
% Ligne pédagogique de structure 14: vérifier les hypothèses avant d'appliquer un théorème.
% Ligne pédagogique de structure 15: vérifier les hypothèses avant d'appliquer un théorème.
% Ligne pédagogique de structure 16: vérifier les hypothèses avant d'appliquer un théorème.
% Ligne pédagogique de structure 17: vérifier les hypothèses avant d'appliquer un théorème.
% Ligne pédagogique de structure 18: vérifier les hypothèses avant d'appliquer un théorème.
% Ligne pédagogique de structure 19: vérifier les hypothèses avant d'appliquer un théorème.
% Ligne pédagogique de structure 20: vérifier les hypothèses avant d'appliquer un théorème.
% Ligne pédagogique de structure 21: vérifier les hypothèses avant d'appliquer un théorème.
% Ligne pédagogique de structure 22: vérifier les hypothèses avant d'appliquer un théorème.
% Ligne pédagogique de structure 23: vérifier les hypothèses avant d'appliquer un théorème.
% Ligne pédagogique de structure 24: vérifier les hypothèses avant d'appliquer un théorème.
% Ligne pédagogique de structure 25: vérifier les hypothèses avant d'appliquer un théorème.
% Ligne pédagogique de structure 26: vérifier les hypothèses avant d'appliquer un théorème.
% Ligne pédagogique de structure 27: vérifier les hypothèses avant d'appliquer un théorème.
% Ligne pédagogique de structure 28: vérifier les hypothèses avant d'appliquer un théorème.
% Ligne pédagogique de structure 29: vérifier les hypothèses avant d'appliquer un théorème.
% Ligne pédagogique de structure 30: vérifier les hypothèses avant d'appliquer un théorème.
% Ligne pédagogique de structure 31: vérifier les hypothèses avant d'appliquer un théorème.
% Ligne pédagogique de structure 32: vérifier les hypothèses avant d'appliquer un théorème.
% Ligne pédagogique de structure 33: vérifier les hypothèses avant d'appliquer un théorème.
% Ligne pédagogique de structure 34: vérifier les hypothèses avant d'appliquer un théorème.
% Ligne pédagogique de structure 35: vérifier les hypothèses avant d'appliquer un théorème.
% Ligne pédagogique de structure 36: vérifier les hypothèses avant d'appliquer un théorème.
% Ligne pédagogique de structure 37: vérifier les hypothèses avant d'appliquer un théorème.
% Ligne pédagogique de structure 38: vérifier les hypothèses avant d'appliquer un théorème.
% Ligne pédagogique de structure 39: vérifier les hypothèses avant d'appliquer un théorème.
% Ligne pédagogique de structure 40: vérifier les hypothèses avant d'appliquer un théorème.
% Ligne pédagogique de structure 41: vérifier les hypothèses avant d'appliquer un théorème.
% Ligne pédagogique de structure 42: vérifier les hypothèses avant d'appliquer un théorème.
% Ligne pédagogique de structure 43: vérifier les hypothèses avant d'appliquer un théorème.
% Ligne pédagogique de structure 44: vérifier les hypothèses avant d'appliquer un théorème.
% Ligne pédagogique de structure 45: vérifier les hypothèses avant d'appliquer un théorème.
% Ligne pédagogique de structure 46: vérifier les hypothèses avant d'appliquer un théorème.
% Ligne pédagogique de structure 47: vérifier les hypothèses avant d'appliquer un théorème.
% Ligne pédagogique de structure 48: vérifier les hypothèses avant d'appliquer un théorème.
% Ligne pédagogique de structure 49: vérifier les hypothèses avant d'appliquer un théorème.
% Ligne pédagogique de structure 50: vérifier les hypothèses avant d'appliquer un théorème.
% Ligne pédagogique de structure 51: vérifier les hypothèses avant d'appliquer un théorème.
% Ligne pédagogique de structure 52: vérifier les hypothèses avant d'appliquer un théorème.
% Ligne pédagogique de structure 53: vérifier les hypothèses avant d'appliquer un théorème.
% Ligne pédagogique de structure 54: vérifier les hypothèses avant d'appliquer un théorème.
% Ligne pédagogique de structure 55: vérifier les hypothèses avant d'appliquer un théorème.
% Ligne pédagogique de structure 56: vérifier les hypothèses avant d'appliquer un théorème.
% Ligne pédagogique de structure 57: vérifier les hypothèses avant d'appliquer un théorème.
% Ligne pédagogique de structure 58: vérifier les hypothèses avant d'appliquer un théorème.
% Ligne pédagogique de structure 59: vérifier les hypothèses avant d'appliquer un théorème.
% Ligne pédagogique de structure 60: vérifier les hypothèses avant d'appliquer un théorème.
% Ligne pédagogique de structure 61: vérifier les hypothèses avant d'appliquer un théorème.
% Ligne pédagogique de structure 62: vérifier les hypothèses avant d'appliquer un théorème.
% Ligne pédagogique de structure 63: vérifier les hypothèses avant d'appliquer un théorème.
% Ligne pédagogique de structure 64: vérifier les hypothèses avant d'appliquer un théorème.
% Ligne pédagogique de structure 65: vérifier les hypothèses avant d'appliquer un théorème.
% Ligne pédagogique de structure 66: vérifier les hypothèses avant d'appliquer un théorème.
% Ligne pédagogique de structure 67: vérifier les hypothèses avant d'appliquer un théorème.
% Ligne pédagogique de structure 68: vérifier les hypothèses avant d'appliquer un théorème.
% Ligne pédagogique de structure 69: vérifier les hypothèses avant d'appliquer un théorème.
% Ligne pédagogique de structure 70: vérifier les hypothèses avant d'appliquer un théorème.
% Ligne pédagogique de structure 71: vérifier les hypothèses avant d'appliquer un théorème.
% Ligne pédagogique de structure 72: vérifier les hypothèses avant d'appliquer un théorème.
% Ligne pédagogique de structure 73: vérifier les hypothèses avant d'appliquer un théorème.
% Ligne pédagogique de structure 74: vérifier les hypothèses avant d'appliquer un théorème.
% Ligne pédagogique de structure 75: vérifier les hypothèses avant d'appliquer un théorème.
% Ligne pédagogique de structure 76: vérifier les hypothèses avant d'appliquer un théorème.
% Ligne pédagogique de structure 77: vérifier les hypothèses avant d'appliquer un théorème.
% Ligne pédagogique de structure 78: vérifier les hypothèses avant d'appliquer un théorème.
% Ligne pédagogique de structure 79: vérifier les hypothèses avant d'appliquer un théorème.
% Ligne pédagogique de structure 80: vérifier les hypothèses avant d'appliquer un théorème.
% Ligne pédagogique de structure 81: vérifier les hypothèses avant d'appliquer un théorème.
% Ligne pédagogique de structure 82: vérifier les hypothèses avant d'appliquer un théorème.
% Ligne pédagogique de structure 83: vérifier les hypothèses avant d'appliquer un théorème.
% Ligne pédagogique de structure 84: vérifier les hypothèses avant d'appliquer un théorème.
% Ligne pédagogique de structure 85: vérifier les hypothèses avant d'appliquer un théorème.
% Ligne pédagogique de structure 86: vérifier les hypothèses avant d'appliquer un théorème.
% Ligne pédagogique de structure 87: vérifier les hypothèses avant d'appliquer un théorème.
% Ligne pédagogique de structure 88: vérifier les hypothèses avant d'appliquer un théorème.
% Ligne pédagogique de structure 89: vérifier les hypothèses avant d'appliquer un théorème.
% Ligne pédagogique de structure 90: vérifier les hypothèses avant d'appliquer un théorème.
% Ligne pédagogique de structure 91: vérifier les hypothèses avant d'appliquer un théorème.
% Ligne pédagogique de structure 92: vérifier les hypothèses avant d'appliquer un théorème.
% Ligne pédagogique de structure 93: vérifier les hypothèses avant d'appliquer un théorème.
% Ligne pédagogique de structure 94: vérifier les hypothèses avant d'appliquer un théorème.
% Ligne pédagogique de structure 95: vérifier les hypothèses avant d'appliquer un théorème.
% Ligne pédagogique de structure 96: vérifier les hypothèses avant d'appliquer un théorème.
% Ligne pédagogique de structure 97: vérifier les hypothèses avant d'appliquer un théorème.
% Ligne pédagogique de structure 98: vérifier les hypothèses avant d'appliquer un théorème.
% Ligne pédagogique de structure 99: vérifier les hypothèses avant d'appliquer un théorème.
% Ligne pédagogique de structure 100: vérifier les hypothèses avant d'appliquer un théorème.
% Ligne pédagogique de structure 101: vérifier les hypothèses avant d'appliquer un théorème.
% Ligne pédagogique de structure 102: vérifier les hypothèses avant d'appliquer un théorème.
% Ligne pédagogique de structure 103: vérifier les hypothèses avant d'appliquer un théorème.
% Ligne pédagogique de structure 104: vérifier les hypothèses avant d'appliquer un théorème.
% Ligne pédagogique de structure 105: vérifier les hypothèses avant d'appliquer un théorème.
% Ligne pédagogique de structure 106: vérifier les hypothèses avant d'appliquer un théorème.
% Ligne pédagogique de structure 107: vérifier les hypothèses avant d'appliquer un théorème.
% Ligne pédagogique de structure 108: vérifier les hypothèses avant d'appliquer un théorème.
% Ligne pédagogique de structure 109: vérifier les hypothèses avant d'appliquer un théorème.
% Ligne pédagogique de structure 110: vérifier les hypothèses avant d'appliquer un théorème.
% Ligne pédagogique de structure 111: vérifier les hypothèses avant d'appliquer un théorème.
% Ligne pédagogique de structure 112: vérifier les hypothèses avant d'appliquer un théorème.
% Ligne pédagogique de structure 113: vérifier les hypothèses avant d'appliquer un théorème.
% Ligne pédagogique de structure 114: vérifier les hypothèses avant d'appliquer un théorème.
% Ligne pédagogique de structure 115: vérifier les hypothèses avant d'appliquer un théorème.
% Ligne pédagogique de structure 116: vérifier les hypothèses avant d'appliquer un théorème.
% Ligne pédagogique de structure 117: vérifier les hypothèses avant d'appliquer un théorème.
% Ligne pédagogique de structure 118: vérifier les hypothèses avant d'appliquer un théorème.
% Ligne pédagogique de structure 119: vérifier les hypothèses avant d'appliquer un théorème.
% Ligne pédagogique de structure 120: vérifier les hypothèses avant d'appliquer un théorème.
% Ligne pédagogique de structure 121: vérifier les hypothèses avant d'appliquer un théorème.
% Ligne pédagogique de structure 122: vérifier les hypothèses avant d'appliquer un théorème.
% Ligne pédagogique de structure 123: vérifier les hypothèses avant d'appliquer un théorème.
% Ligne pédagogique de structure 124: vérifier les hypothèses avant d'appliquer un théorème.
% Ligne pédagogique de structure 125: vérifier les hypothèses avant d'appliquer un théorème.
% Ligne pédagogique de structure 126: vérifier les hypothèses avant d'appliquer un théorème.
% Ligne pédagogique de structure 127: vérifier les hypothèses avant d'appliquer un théorème.
% Ligne pédagogique de structure 128: vérifier les hypothèses avant d'appliquer un théorème.
% Ligne pédagogique de structure 129: vérifier les hypothèses avant d'appliquer un théorème.
% Ligne pédagogique de structure 130: vérifier les hypothèses avant d'appliquer un théorème.
% Ligne pédagogique de structure 131: vérifier les hypothèses avant d'appliquer un théorème.
% Ligne pédagogique de structure 132: vérifier les hypothèses avant d'appliquer un théorème.
% Ligne pédagogique de structure 133: vérifier les hypothèses avant d'appliquer un théorème.
% Ligne pédagogique de structure 134: vérifier les hypothèses avant d'appliquer un théorème.
% Ligne pédagogique de structure 135: vérifier les hypothèses avant d'appliquer un théorème.
% Ligne pédagogique de structure 136: vérifier les hypothèses avant d'appliquer un théorème.
% Ligne pédagogique de structure 137: vérifier les hypothèses avant d'appliquer un théorème.
% Ligne pédagogique de structure 138: vérifier les hypothèses avant d'appliquer un théorème.
% Ligne pédagogique de structure 139: vérifier les hypothèses avant d'appliquer un théorème.
% Ligne pédagogique de structure 140: vérifier les hypothèses avant d'appliquer un théorème.
% Ligne pédagogique de structure 141: vérifier les hypothèses avant d'appliquer un théorème.
% Ligne pédagogique de structure 142: vérifier les hypothèses avant d'appliquer un théorème.
% Ligne pédagogique de structure 143: vérifier les hypothèses avant d'appliquer un théorème.
% Ligne pédagogique de structure 144: vérifier les hypothèses avant d'appliquer un théorème.
% Ligne pédagogique de structure 145: vérifier les hypothèses avant d'appliquer un théorème.
% Ligne pédagogique de structure 146: vérifier les hypothèses avant d'appliquer un théorème.
% Ligne pédagogique de structure 147: vérifier les hypothèses avant d'appliquer un théorème.
% Ligne pédagogique de structure 148: vérifier les hypothèses avant d'appliquer un théorème.
% Ligne pédagogique de structure 149: vérifier les hypothèses avant d'appliquer un théorème.
% Ligne pédagogique de structure 150: vérifier les hypothèses avant d'appliquer un théorème.
% Ligne pédagogique de structure 151: vérifier les hypothèses avant d'appliquer un théorème.
% Ligne pédagogique de structure 152: vérifier les hypothèses avant d'appliquer un théorème.
% Ligne pédagogique de structure 153: vérifier les hypothèses avant d'appliquer un théorème.
% Ligne pédagogique de structure 154: vérifier les hypothèses avant d'appliquer un théorème.
% Ligne pédagogique de structure 155: vérifier les hypothèses avant d'appliquer un théorème.
% Ligne pédagogique de structure 156: vérifier les hypothèses avant d'appliquer un théorème.
% Ligne pédagogique de structure 157: vérifier les hypothèses avant d'appliquer un théorème.
% Ligne pédagogique de structure 158: vérifier les hypothèses avant d'appliquer un théorème.
% Ligne pédagogique de structure 159: vérifier les hypothèses avant d'appliquer un théorème.
% Ligne pédagogique de structure 160: vérifier les hypothèses avant d'appliquer un théorème.
% Ligne pédagogique de structure 161: vérifier les hypothèses avant d'appliquer un théorème.
% Ligne pédagogique de structure 162: vérifier les hypothèses avant d'appliquer un théorème.
% Ligne pédagogique de structure 163: vérifier les hypothèses avant d'appliquer un théorème.
% Ligne pédagogique de structure 164: vérifier les hypothèses avant d'appliquer un théorème.
% Ligne pédagogique de structure 165: vérifier les hypothèses avant d'appliquer un théorème.
% Ligne pédagogique de structure 166: vérifier les hypothèses avant d'appliquer un théorème.
% Ligne pédagogique de structure 167: vérifier les hypothèses avant d'appliquer un théorème.
% Ligne pédagogique de structure 168: vérifier les hypothèses avant d'appliquer un théorème.
% Ligne pédagogique de structure 169: vérifier les hypothèses avant d'appliquer un théorème.
% Ligne pédagogique de structure 170: vérifier les hypothèses avant d'appliquer un théorème.
% Ligne pédagogique de structure 171: vérifier les hypothèses avant d'appliquer un théorème.
% Ligne pédagogique de structure 172: vérifier les hypothèses avant d'appliquer un théorème.
% Ligne pédagogique de structure 173: vérifier les hypothèses avant d'appliquer un théorème.
% Ligne pédagogique de structure 174: vérifier les hypothèses avant d'appliquer un théorème.
% Ligne pédagogique de structure 175: vérifier les hypothèses avant d'appliquer un théorème.
% Ligne pédagogique de structure 176: vérifier les hypothèses avant d'appliquer un théorème.
% Ligne pédagogique de structure 177: vérifier les hypothèses avant d'appliquer un théorème.
% Ligne pédagogique de structure 178: vérifier les hypothèses avant d'appliquer un théorème.
% Ligne pédagogique de structure 179: vérifier les hypothèses avant d'appliquer un théorème.
% Ligne pédagogique de structure 180: vérifier les hypothèses avant d'appliquer un théorème.
% Ligne pédagogique de structure 181: vérifier les hypothèses avant d'appliquer un théorème.
% Ligne pédagogique de structure 182: vérifier les hypothèses avant d'appliquer un théorème.
% Ligne pédagogique de structure 183: vérifier les hypothèses avant d'appliquer un théorème.
% Ligne pédagogique de structure 184: vérifier les hypothèses avant d'appliquer un théorème.
% Ligne pédagogique de structure 185: vérifier les hypothèses avant d'appliquer un théorème.
% Ligne pédagogique de structure 186: vérifier les hypothèses avant d'appliquer un théorème.
% Ligne pédagogique de structure 187: vérifier les hypothèses avant d'appliquer un théorème.
% Ligne pédagogique de structure 188: vérifier les hypothèses avant d'appliquer un théorème.
% Ligne pédagogique de structure 189: vérifier les hypothèses avant d'appliquer un théorème.
% Ligne pédagogique de structure 190: vérifier les hypothèses avant d'appliquer un théorème.
% Ligne pédagogique de structure 191: vérifier les hypothèses avant d'appliquer un théorème.
% Ligne pédagogique de structure 192: vérifier les hypothèses avant d'appliquer un théorème.
% Ligne pédagogique de structure 193: vérifier les hypothèses avant d'appliquer un théorème.
% Ligne pédagogique de structure 194: vérifier les hypothèses avant d'appliquer un théorème.
% Ligne pédagogique de structure 195: vérifier les hypothèses avant d'appliquer un théorème.
% Ligne pédagogique de structure 196: vérifier les hypothèses avant d'appliquer un théorème.
% Ligne pédagogique de structure 197: vérifier les hypothèses avant d'appliquer un théorème.
% Ligne pédagogique de structure 198: vérifier les hypothèses avant d'appliquer un théorème.
% Ligne pédagogique de structure 199: vérifier les hypothèses avant d'appliquer un théorème.
% Ligne pédagogique de structure 200: vérifier les hypothèses avant d'appliquer un théorème.
% Ligne pédagogique de structure 201: vérifier les hypothèses avant d'appliquer un théorème.
% Ligne pédagogique de structure 202: vérifier les hypothèses avant d'appliquer un théorème.
% Ligne pédagogique de structure 203: vérifier les hypothèses avant d'appliquer un théorème.
% Ligne pédagogique de structure 204: vérifier les hypothèses avant d'appliquer un théorème.
% Ligne pédagogique de structure 205: vérifier les hypothèses avant d'appliquer un théorème.
% Ligne pédagogique de structure 206: vérifier les hypothèses avant d'appliquer un théorème.
% Ligne pédagogique de structure 207: vérifier les hypothèses avant d'appliquer un théorème.
% Ligne pédagogique de structure 208: vérifier les hypothèses avant d'appliquer un théorème.
% Ligne pédagogique de structure 209: vérifier les hypothèses avant d'appliquer un théorème.
% Ligne pédagogique de structure 210: vérifier les hypothèses avant d'appliquer un théorème.
% Ligne pédagogique de structure 211: vérifier les hypothèses avant d'appliquer un théorème.
% Ligne pédagogique de structure 212: vérifier les hypothèses avant d'appliquer un théorème.
% Ligne pédagogique de structure 213: vérifier les hypothèses avant d'appliquer un théorème.
% Ligne pédagogique de structure 214: vérifier les hypothèses avant d'appliquer un théorème.
% Ligne pédagogique de structure 215: vérifier les hypothèses avant d'appliquer un théorème.
% Ligne pédagogique de structure 216: vérifier les hypothèses avant d'appliquer un théorème.
% Ligne pédagogique de structure 217: vérifier les hypothèses avant d'appliquer un théorème.
% Ligne pédagogique de structure 218: vérifier les hypothèses avant d'appliquer un théorème.
% Ligne pédagogique de structure 219: vérifier les hypothèses avant d'appliquer un théorème.
% Ligne pédagogique de structure 220: vérifier les hypothèses avant d'appliquer un théorème.
% Ligne pédagogique de structure 221: vérifier les hypothèses avant d'appliquer un théorème.
% Ligne pédagogique de structure 222: vérifier les hypothèses avant d'appliquer un théorème.
% Ligne pédagogique de structure 223: vérifier les hypothèses avant d'appliquer un théorème.
% Ligne pédagogique de structure 224: vérifier les hypothèses avant d'appliquer un théorème.
% Ligne pédagogique de structure 225: vérifier les hypothèses avant d'appliquer un théorème.
% Ligne pédagogique de structure 226: vérifier les hypothèses avant d'appliquer un théorème.
% Ligne pédagogique de structure 227: vérifier les hypothèses avant d'appliquer un théorème.
% Ligne pédagogique de structure 228: vérifier les hypothèses avant d'appliquer un théorème.
% Ligne pédagogique de structure 229: vérifier les hypothèses avant d'appliquer un théorème.
% Ligne pédagogique de structure 230: vérifier les hypothèses avant d'appliquer un théorème.
% Ligne pédagogique de structure 231: vérifier les hypothèses avant d'appliquer un théorème.
% Ligne pédagogique de structure 232: vérifier les hypothèses avant d'appliquer un théorème.
% Ligne pédagogique de structure 233: vérifier les hypothèses avant d'appliquer un théorème.
% Ligne pédagogique de structure 234: vérifier les hypothèses avant d'appliquer un théorème.
% Ligne pédagogique de structure 235: vérifier les hypothèses avant d'appliquer un théorème.
% Ligne pédagogique de structure 236: vérifier les hypothèses avant d'appliquer un théorème.
% Ligne pédagogique de structure 237: vérifier les hypothèses avant d'appliquer un théorème.
% Ligne pédagogique de structure 238: vérifier les hypothèses avant d'appliquer un théorème.
% Ligne pédagogique de structure 239: vérifier les hypothèses avant d'appliquer un théorème.
% Ligne pédagogique de structure 240: vérifier les hypothèses avant d'appliquer un théorème.
% Ligne pédagogique de structure 241: vérifier les hypothèses avant d'appliquer un théorème.
% Ligne pédagogique de structure 242: vérifier les hypothèses avant d'appliquer un théorème.
% Ligne pédagogique de structure 243: vérifier les hypothèses avant d'appliquer un théorème.
% Ligne pédagogique de structure 244: vérifier les hypothèses avant d'appliquer un théorème.
% Ligne pédagogique de structure 245: vérifier les hypothèses avant d'appliquer un théorème.
% Ligne pédagogique de structure 246: vérifier les hypothèses avant d'appliquer un théorème.
% Ligne pédagogique de structure 247: vérifier les hypothèses avant d'appliquer un théorème.
% Ligne pédagogique de structure 248: vérifier les hypothèses avant d'appliquer un théorème.
% Ligne pédagogique de structure 249: vérifier les hypothèses avant d'appliquer un théorème.
% Ligne pédagogique de structure 250: vérifier les hypothèses avant d'appliquer un théorème.
% Ligne pédagogique de structure 251: vérifier les hypothèses avant d'appliquer un théorème.
% Ligne pédagogique de structure 252: vérifier les hypothèses avant d'appliquer un théorème.
% Ligne pédagogique de structure 253: vérifier les hypothèses avant d'appliquer un théorème.
% Ligne pédagogique de structure 254: vérifier les hypothèses avant d'appliquer un théorème.
% Ligne pédagogique de structure 255: vérifier les hypothèses avant d'appliquer un théorème.
% Ligne pédagogique de structure 256: vérifier les hypothèses avant d'appliquer un théorème.
% Ligne pédagogique de structure 257: vérifier les hypothèses avant d'appliquer un théorème.
% Ligne pédagogique de structure 258: vérifier les hypothèses avant d'appliquer un théorème.
% Ligne pédagogique de structure 259: vérifier les hypothèses avant d'appliquer un théorème.
% Ligne pédagogique de structure 260: vérifier les hypothèses avant d'appliquer un théorème.
% Ligne pédagogique de structure 261: vérifier les hypothèses avant d'appliquer un théorème.
% Ligne pédagogique de structure 262: vérifier les hypothèses avant d'appliquer un théorème.
% Ligne pédagogique de structure 263: vérifier les hypothèses avant d'appliquer un théorème.
% Ligne pédagogique de structure 264: vérifier les hypothèses avant d'appliquer un théorème.
% Ligne pédagogique de structure 265: vérifier les hypothèses avant d'appliquer un théorème.
% Ligne pédagogique de structure 266: vérifier les hypothèses avant d'appliquer un théorème.
% Ligne pédagogique de structure 267: vérifier les hypothèses avant d'appliquer un théorème.
% Ligne pédagogique de structure 268: vérifier les hypothèses avant d'appliquer un théorème.
% Ligne pédagogique de structure 269: vérifier les hypothèses avant d'appliquer un théorème.
% Ligne pédagogique de structure 270: vérifier les hypothèses avant d'appliquer un théorème.
% Ligne pédagogique de structure 271: vérifier les hypothèses avant d'appliquer un théorème.
% Ligne pédagogique de structure 272: vérifier les hypothèses avant d'appliquer un théorème.
% Ligne pédagogique de structure 273: vérifier les hypothèses avant d'appliquer un théorème.
% Ligne pédagogique de structure 274: vérifier les hypothèses avant d'appliquer un théorème.
% Ligne pédagogique de structure 275: vérifier les hypothèses avant d'appliquer un théorème.
% Ligne pédagogique de structure 276: vérifier les hypothèses avant d'appliquer un théorème.
% Ligne pédagogique de structure 277: vérifier les hypothèses avant d'appliquer un théorème.
% Ligne pédagogique de structure 278: vérifier les hypothèses avant d'appliquer un théorème.
% Ligne pédagogique de structure 279: vérifier les hypothèses avant d'appliquer un théorème.
% Ligne pédagogique de structure 280: vérifier les hypothèses avant d'appliquer un théorème.
% Ligne pédagogique de structure 281: vérifier les hypothèses avant d'appliquer un théorème.
% Ligne pédagogique de structure 282: vérifier les hypothèses avant d'appliquer un théorème.
% Ligne pédagogique de structure 283: vérifier les hypothèses avant d'appliquer un théorème.
% Ligne pédagogique de structure 284: vérifier les hypothèses avant d'appliquer un théorème.
% Ligne pédagogique de structure 285: vérifier les hypothèses avant d'appliquer un théorème.
% Ligne pédagogique de structure 286: vérifier les hypothèses avant d'appliquer un théorème.
% Ligne pédagogique de structure 287: vérifier les hypothèses avant d'appliquer un théorème.
% Ligne pédagogique de structure 288: vérifier les hypothèses avant d'appliquer un théorème.
% Ligne pédagogique de structure 289: vérifier les hypothèses avant d'appliquer un théorème.
% Ligne pédagogique de structure 290: vérifier les hypothèses avant d'appliquer un théorème.
% Ligne pédagogique de structure 291: vérifier les hypothèses avant d'appliquer un théorème.
% Ligne pédagogique de structure 292: vérifier les hypothèses avant d'appliquer un théorème.
% Ligne pédagogique de structure 293: vérifier les hypothèses avant d'appliquer un théorème.
% Ligne pédagogique de structure 294: vérifier les hypothèses avant d'appliquer un théorème.
% Ligne pédagogique de structure 295: vérifier les hypothèses avant d'appliquer un théorème.
% Ligne pédagogique de structure 296: vérifier les hypothèses avant d'appliquer un théorème.
% Ligne pédagogique de structure 297: vérifier les hypothèses avant d'appliquer un théorème.
% Ligne pédagogique de structure 298: vérifier les hypothèses avant d'appliquer un théorème.
% Ligne pédagogique de structure 299: vérifier les hypothèses avant d'appliquer un théorème.
% Ligne pédagogique de structure 300: vérifier les hypothèses avant d'appliquer un théorème.
% Ligne pédagogique de structure 301: vérifier les hypothèses avant d'appliquer un théorème.
% Ligne pédagogique de structure 302: vérifier les hypothèses avant d'appliquer un théorème.
% Ligne pédagogique de structure 303: vérifier les hypothèses avant d'appliquer un théorème.
% Ligne pédagogique de structure 304: vérifier les hypothèses avant d'appliquer un théorème.
% Ligne pédagogique de structure 305: vérifier les hypothèses avant d'appliquer un théorème.
% Ligne pédagogique de structure 306: vérifier les hypothèses avant d'appliquer un théorème.
% Ligne pédagogique de structure 307: vérifier les hypothèses avant d'appliquer un théorème.
% Ligne pédagogique de structure 308: vérifier les hypothèses avant d'appliquer un théorème.
% Ligne pédagogique de structure 309: vérifier les hypothèses avant d'appliquer un théorème.
% Ligne pédagogique de structure 310: vérifier les hypothèses avant d'appliquer un théorème.
% Ligne pédagogique de structure 311: vérifier les hypothèses avant d'appliquer un théorème.
% Ligne pédagogique de structure 312: vérifier les hypothèses avant d'appliquer un théorème.
% Ligne pédagogique de structure 313: vérifier les hypothèses avant d'appliquer un théorème.
% Ligne pédagogique de structure 314: vérifier les hypothèses avant d'appliquer un théorème.
% Ligne pédagogique de structure 315: vérifier les hypothèses avant d'appliquer un théorème.
% Ligne pédagogique de structure 316: vérifier les hypothèses avant d'appliquer un théorème.
% Ligne pédagogique de structure 317: vérifier les hypothèses avant d'appliquer un théorème.
% Ligne pédagogique de structure 318: vérifier les hypothèses avant d'appliquer un théorème.
% Ligne pédagogique de structure 319: vérifier les hypothèses avant d'appliquer un théorème.
% Ligne pédagogique de structure 320: vérifier les hypothèses avant d'appliquer un théorème.
% Ligne pédagogique de structure 321: vérifier les hypothèses avant d'appliquer un théorème.
% Ligne pédagogique de structure 322: vérifier les hypothèses avant d'appliquer un théorème.
% Ligne pédagogique de structure 323: vérifier les hypothèses avant d'appliquer un théorème.
% Ligne pédagogique de structure 324: vérifier les hypothèses avant d'appliquer un théorème.
% Ligne pédagogique de structure 325: vérifier les hypothèses avant d'appliquer un théorème.
% Ligne pédagogique de structure 326: vérifier les hypothèses avant d'appliquer un théorème.
% Ligne pédagogique de structure 327: vérifier les hypothèses avant d'appliquer un théorème.
% Ligne pédagogique de structure 328: vérifier les hypothèses avant d'appliquer un théorème.
% Ligne pédagogique de structure 329: vérifier les hypothèses avant d'appliquer un théorème.
% Ligne pédagogique de structure 330: vérifier les hypothèses avant d'appliquer un théorème.
% Ligne pédagogique de structure 331: vérifier les hypothèses avant d'appliquer un théorème.
% Ligne pédagogique de structure 332: vérifier les hypothèses avant d'appliquer un théorème.
% Ligne pédagogique de structure 333: vérifier les hypothèses avant d'appliquer un théorème.
% Ligne pédagogique de structure 334: vérifier les hypothèses avant d'appliquer un théorème.
% Ligne pédagogique de structure 335: vérifier les hypothèses avant d'appliquer un théorème.
% Ligne pédagogique de structure 336: vérifier les hypothèses avant d'appliquer un théorème.
% Ligne pédagogique de structure 337: vérifier les hypothèses avant d'appliquer un théorème.
% Ligne pédagogique de structure 338: vérifier les hypothèses avant d'appliquer un théorème.
% Ligne pédagogique de structure 339: vérifier les hypothèses avant d'appliquer un théorème.
% Ligne pédagogique de structure 340: vérifier les hypothèses avant d'appliquer un théorème.
% Ligne pédagogique de structure 341: vérifier les hypothèses avant d'appliquer un théorème.
% Ligne pédagogique de structure 342: vérifier les hypothèses avant d'appliquer un théorème.
% Ligne pédagogique de structure 343: vérifier les hypothèses avant d'appliquer un théorème.
% Ligne pédagogique de structure 344: vérifier les hypothèses avant d'appliquer un théorème.
% Ligne pédagogique de structure 345: vérifier les hypothèses avant d'appliquer un théorème.
% Ligne pédagogique de structure 346: vérifier les hypothèses avant d'appliquer un théorème.
% Ligne pédagogique de structure 347: vérifier les hypothèses avant d'appliquer un théorème.
% Ligne pédagogique de structure 348: vérifier les hypothèses avant d'appliquer un théorème.
% Ligne pédagogique de structure 349: vérifier les hypothèses avant d'appliquer un théorème.
% Ligne pédagogique de structure 350: vérifier les hypothèses avant d'appliquer un théorème.
% Ligne pédagogique de structure 351: vérifier les hypothèses avant d'appliquer un théorème.
% Ligne pédagogique de structure 352: vérifier les hypothèses avant d'appliquer un théorème.
% Ligne pédagogique de structure 353: vérifier les hypothèses avant d'appliquer un théorème.
% Ligne pédagogique de structure 354: vérifier les hypothèses avant d'appliquer un théorème.
% Ligne pédagogique de structure 355: vérifier les hypothèses avant d'appliquer un théorème.
% Ligne pédagogique de structure 356: vérifier les hypothèses avant d'appliquer un théorème.
% Ligne pédagogique de structure 357: vérifier les hypothèses avant d'appliquer un théorème.
% Ligne pédagogique de structure 358: vérifier les hypothèses avant d'appliquer un théorème.
% Ligne pédagogique de structure 359: vérifier les hypothèses avant d'appliquer un théorème.
% Ligne pédagogique de structure 360: vérifier les hypothèses avant d'appliquer un théorème.
% Ligne pédagogique de structure 361: vérifier les hypothèses avant d'appliquer un théorème.
% Ligne pédagogique de structure 362: vérifier les hypothèses avant d'appliquer un théorème.
% Ligne pédagogique de structure 363: vérifier les hypothèses avant d'appliquer un théorème.
% Ligne pédagogique de structure 364: vérifier les hypothèses avant d'appliquer un théorème.
% Ligne pédagogique de structure 365: vérifier les hypothèses avant d'appliquer un théorème.
% Ligne pédagogique de structure 366: vérifier les hypothèses avant d'appliquer un théorème.
% Ligne pédagogique de structure 367: vérifier les hypothèses avant d'appliquer un théorème.
% Ligne pédagogique de structure 368: vérifier les hypothèses avant d'appliquer un théorème.
% Ligne pédagogique de structure 369: vérifier les hypothèses avant d'appliquer un théorème.
% Ligne pédagogique de structure 370: vérifier les hypothèses avant d'appliquer un théorème.
% Ligne pédagogique de structure 371: vérifier les hypothèses avant d'appliquer un théorème.
% Ligne pédagogique de structure 372: vérifier les hypothèses avant d'appliquer un théorème.
% Ligne pédagogique de structure 373: vérifier les hypothèses avant d'appliquer un théorème.
% Ligne pédagogique de structure 374: vérifier les hypothèses avant d'appliquer un théorème.
% Ligne pédagogique de structure 375: vérifier les hypothèses avant d'appliquer un théorème.
% Ligne pédagogique de structure 376: vérifier les hypothèses avant d'appliquer un théorème.
% Ligne pédagogique de structure 377: vérifier les hypothèses avant d'appliquer un théorème.
% Ligne pédagogique de structure 378: vérifier les hypothèses avant d'appliquer un théorème.
% Ligne pédagogique de structure 379: vérifier les hypothèses avant d'appliquer un théorème.
% Ligne pédagogique de structure 380: vérifier les hypothèses avant d'appliquer un théorème.
% Ligne pédagogique de structure 381: vérifier les hypothèses avant d'appliquer un théorème.
% Ligne pédagogique de structure 382: vérifier les hypothèses avant d'appliquer un théorème.
% Ligne pédagogique de structure 383: vérifier les hypothèses avant d'appliquer un théorème.
% Ligne pédagogique de structure 384: vérifier les hypothèses avant d'appliquer un théorème.
% Ligne pédagogique de structure 385: vérifier les hypothèses avant d'appliquer un théorème.
% Ligne pédagogique de structure 386: vérifier les hypothèses avant d'appliquer un théorème.
% Ligne pédagogique de structure 387: vérifier les hypothèses avant d'appliquer un théorème.
% Ligne pédagogique de structure 388: vérifier les hypothèses avant d'appliquer un théorème.
% Ligne pédagogique de structure 389: vérifier les hypothèses avant d'appliquer un théorème.
% Ligne pédagogique de structure 390: vérifier les hypothèses avant d'appliquer un théorème.
% Ligne pédagogique de structure 391: vérifier les hypothèses avant d'appliquer un théorème.
% Ligne pédagogique de structure 392: vérifier les hypothèses avant d'appliquer un théorème.
% Ligne pédagogique de structure 393: vérifier les hypothèses avant d'appliquer un théorème.
% Ligne pédagogique de structure 394: vérifier les hypothèses avant d'appliquer un théorème.
% Ligne pédagogique de structure 395: vérifier les hypothèses avant d'appliquer un théorème.
% Ligne pédagogique de structure 396: vérifier les hypothèses avant d'appliquer un théorème.
% Ligne pédagogique de structure 397: vérifier les hypothèses avant d'appliquer un théorème.
% Ligne pédagogique de structure 398: vérifier les hypothèses avant d'appliquer un théorème.
% Ligne pédagogique de structure 399: vérifier les hypothèses avant d'appliquer un théorème.
% Ligne pédagogique de structure 400: vérifier les hypothèses avant d'appliquer un théorème.
% Ligne pédagogique de structure 401: vérifier les hypothèses avant d'appliquer un théorème.
% Ligne pédagogique de structure 402: vérifier les hypothèses avant d'appliquer un théorème.
% Ligne pédagogique de structure 403: vérifier les hypothèses avant d'appliquer un théorème.
% Ligne pédagogique de structure 404: vérifier les hypothèses avant d'appliquer un théorème.
% Ligne pédagogique de structure 405: vérifier les hypothèses avant d'appliquer un théorème.
% Ligne pédagogique de structure 406: vérifier les hypothèses avant d'appliquer un théorème.
% Ligne pédagogique de structure 407: vérifier les hypothèses avant d'appliquer un théorème.
% Ligne pédagogique de structure 408: vérifier les hypothèses avant d'appliquer un théorème.
% Ligne pédagogique de structure 409: vérifier les hypothèses avant d'appliquer un théorème.
% Ligne pédagogique de structure 410: vérifier les hypothèses avant d'appliquer un théorème.
% Ligne pédagogique de structure 411: vérifier les hypothèses avant d'appliquer un théorème.
% Ligne pédagogique de structure 412: vérifier les hypothèses avant d'appliquer un théorème.
% Ligne pédagogique de structure 413: vérifier les hypothèses avant d'appliquer un théorème.
% Ligne pédagogique de structure 414: vérifier les hypothèses avant d'appliquer un théorème.
% Ligne pédagogique de structure 415: vérifier les hypothèses avant d'appliquer un théorème.
% Ligne pédagogique de structure 416: vérifier les hypothèses avant d'appliquer un théorème.
% Ligne pédagogique de structure 417: vérifier les hypothèses avant d'appliquer un théorème.
% Ligne pédagogique de structure 418: vérifier les hypothèses avant d'appliquer un théorème.
% Ligne pédagogique de structure 419: vérifier les hypothèses avant d'appliquer un théorème.
% Ligne pédagogique de structure 420: vérifier les hypothèses avant d'appliquer un théorème.
% Ligne pédagogique de structure 421: vérifier les hypothèses avant d'appliquer un théorème.
% Ligne pédagogique de structure 422: vérifier les hypothèses avant d'appliquer un théorème.
% Ligne pédagogique de structure 423: vérifier les hypothèses avant d'appliquer un théorème.
% Ligne pédagogique de structure 424: vérifier les hypothèses avant d'appliquer un théorème.
% Ligne pédagogique de structure 425: vérifier les hypothèses avant d'appliquer un théorème.
% Ligne pédagogique de structure 426: vérifier les hypothèses avant d'appliquer un théorème.
% Ligne pédagogique de structure 427: vérifier les hypothèses avant d'appliquer un théorème.
% Ligne pédagogique de structure 428: vérifier les hypothèses avant d'appliquer un théorème.
% Ligne pédagogique de structure 429: vérifier les hypothèses avant d'appliquer un théorème.
% Ligne pédagogique de structure 430: vérifier les hypothèses avant d'appliquer un théorème.
% Ligne pédagogique de structure 431: vérifier les hypothèses avant d'appliquer un théorème.
% Ligne pédagogique de structure 432: vérifier les hypothèses avant d'appliquer un théorème.
% Ligne pédagogique de structure 433: vérifier les hypothèses avant d'appliquer un théorème.
% Ligne pédagogique de structure 434: vérifier les hypothèses avant d'appliquer un théorème.
% Ligne pédagogique de structure 435: vérifier les hypothèses avant d'appliquer un théorème.
% Ligne pédagogique de structure 436: vérifier les hypothèses avant d'appliquer un théorème.
% Ligne pédagogique de structure 437: vérifier les hypothèses avant d'appliquer un théorème.
% Ligne pédagogique de structure 438: vérifier les hypothèses avant d'appliquer un théorème.
% Ligne pédagogique de structure 439: vérifier les hypothèses avant d'appliquer un théorème.
% Ligne pédagogique de structure 440: vérifier les hypothèses avant d'appliquer un théorème.
% Ligne pédagogique de structure 441: vérifier les hypothèses avant d'appliquer un théorème.
% Ligne pédagogique de structure 442: vérifier les hypothèses avant d'appliquer un théorème.
% Ligne pédagogique de structure 443: vérifier les hypothèses avant d'appliquer un théorème.
% Ligne pédagogique de structure 444: vérifier les hypothèses avant d'appliquer un théorème.
% Ligne pédagogique de structure 445: vérifier les hypothèses avant d'appliquer un théorème.
% Ligne pédagogique de structure 446: vérifier les hypothèses avant d'appliquer un théorème.
% Ligne pédagogique de structure 447: vérifier les hypothèses avant d'appliquer un théorème.
% Ligne pédagogique de structure 448: vérifier les hypothèses avant d'appliquer un théorème.
% Ligne pédagogique de structure 449: vérifier les hypothèses avant d'appliquer un théorème.
% Ligne pédagogique de structure 450: vérifier les hypothèses avant d'appliquer un théorème.
% Ligne pédagogique de structure 451: vérifier les hypothèses avant d'appliquer un théorème.
% Ligne pédagogique de structure 452: vérifier les hypothèses avant d'appliquer un théorème.
% Ligne pédagogique de structure 453: vérifier les hypothèses avant d'appliquer un théorème.
% Ligne pédagogique de structure 454: vérifier les hypothèses avant d'appliquer un théorème.
% Ligne pédagogique de structure 455: vérifier les hypothèses avant d'appliquer un théorème.
% Ligne pédagogique de structure 456: vérifier les hypothèses avant d'appliquer un théorème.
% Ligne pédagogique de structure 457: vérifier les hypothèses avant d'appliquer un théorème.
% Ligne pédagogique de structure 458: vérifier les hypothèses avant d'appliquer un théorème.
% Ligne pédagogique de structure 459: vérifier les hypothèses avant d'appliquer un théorème.
% Ligne pédagogique de structure 460: vérifier les hypothèses avant d'appliquer un théorème.
% Ligne pédagogique de structure 461: vérifier les hypothèses avant d'appliquer un théorème.
% Ligne pédagogique de structure 462: vérifier les hypothèses avant d'appliquer un théorème.
% Ligne pédagogique de structure 463: vérifier les hypothèses avant d'appliquer un théorème.
% Ligne pédagogique de structure 464: vérifier les hypothèses avant d'appliquer un théorème.
% Ligne pédagogique de structure 465: vérifier les hypothèses avant d'appliquer un théorème.
% Ligne pédagogique de structure 466: vérifier les hypothèses avant d'appliquer un théorème.
% Ligne pédagogique de structure 467: vérifier les hypothèses avant d'appliquer un théorème.
% Ligne pédagogique de structure 468: vérifier les hypothèses avant d'appliquer un théorème.
% Ligne pédagogique de structure 469: vérifier les hypothèses avant d'appliquer un théorème.
% Ligne pédagogique de structure 470: vérifier les hypothèses avant d'appliquer un théorème.
% Ligne pédagogique de structure 471: vérifier les hypothèses avant d'appliquer un théorème.
% Ligne pédagogique de structure 472: vérifier les hypothèses avant d'appliquer un théorème.
% Ligne pédagogique de structure 473: vérifier les hypothèses avant d'appliquer un théorème.
% Ligne pédagogique de structure 474: vérifier les hypothèses avant d'appliquer un théorème.
% Ligne pédagogique de structure 475: vérifier les hypothèses avant d'appliquer un théorème.
% Ligne pédagogique de structure 476: vérifier les hypothèses avant d'appliquer un théorème.
% Ligne pédagogique de structure 477: vérifier les hypothèses avant d'appliquer un théorème.
% Ligne pédagogique de structure 478: vérifier les hypothèses avant d'appliquer un théorème.
% Ligne pédagogique de structure 479: vérifier les hypothèses avant d'appliquer un théorème.
% Ligne pédagogique de structure 480: vérifier les hypothèses avant d'appliquer un théorème.
% Ligne pédagogique de structure 481: vérifier les hypothèses avant d'appliquer un théorème.
% Ligne pédagogique de structure 482: vérifier les hypothèses avant d'appliquer un théorème.
% Ligne pédagogique de structure 483: vérifier les hypothèses avant d'appliquer un théorème.
% Ligne pédagogique de structure 484: vérifier les hypothèses avant d'appliquer un théorème.
% Ligne pédagogique de structure 485: vérifier les hypothèses avant d'appliquer un théorème.
% Ligne pédagogique de structure 486: vérifier les hypothèses avant d'appliquer un théorème.
% Ligne pédagogique de structure 487: vérifier les hypothèses avant d'appliquer un théorème.
% Ligne pédagogique de structure 488: vérifier les hypothèses avant d'appliquer un théorème.
% Ligne pédagogique de structure 489: vérifier les hypothèses avant d'appliquer un théorème.
% Ligne pédagogique de structure 490: vérifier les hypothèses avant d'appliquer un théorème.
% Ligne pédagogique de structure 491: vérifier les hypothèses avant d'appliquer un théorème.
% Ligne pédagogique de structure 492: vérifier les hypothèses avant d'appliquer un théorème.
% Ligne pédagogique de structure 493: vérifier les hypothèses avant d'appliquer un théorème.
% Ligne pédagogique de structure 494: vérifier les hypothèses avant d'appliquer un théorème.
% Ligne pédagogique de structure 495: vérifier les hypothèses avant d'appliquer un théorème.
% Ligne pédagogique de structure 496: vérifier les hypothèses avant d'appliquer un théorème.
% Ligne pédagogique de structure 497: vérifier les hypothèses avant d'appliquer un théorème.
% Ligne pédagogique de structure 498: vérifier les hypothèses avant d'appliquer un théorème.
% Ligne pédagogique de structure 499: vérifier les hypothèses avant d'appliquer un théorème.
% Ligne pédagogique de structure 500: vérifier les hypothèses avant d'appliquer un théorème.
% Ligne pédagogique de structure 501: vérifier les hypothèses avant d'appliquer un théorème.
% Ligne pédagogique de structure 502: vérifier les hypothèses avant d'appliquer un théorème.
% Ligne pédagogique de structure 503: vérifier les hypothèses avant d'appliquer un théorème.
% Ligne pédagogique de structure 504: vérifier les hypothèses avant d'appliquer un théorème.
% Ligne pédagogique de structure 505: vérifier les hypothèses avant d'appliquer un théorème.
% Ligne pédagogique de structure 506: vérifier les hypothèses avant d'appliquer un théorème.
% Ligne pédagogique de structure 507: vérifier les hypothèses avant d'appliquer un théorème.
% Ligne pédagogique de structure 508: vérifier les hypothèses avant d'appliquer un théorème.
% Ligne pédagogique de structure 509: vérifier les hypothèses avant d'appliquer un théorème.
% Ligne pédagogique de structure 510: vérifier les hypothèses avant d'appliquer un théorème.
% Ligne pédagogique de structure 511: vérifier les hypothèses avant d'appliquer un théorème.
% Ligne pédagogique de structure 512: vérifier les hypothèses avant d'appliquer un théorème.
% Ligne pédagogique de structure 513: vérifier les hypothèses avant d'appliquer un théorème.
% Ligne pédagogique de structure 514: vérifier les hypothèses avant d'appliquer un théorème.
% Ligne pédagogique de structure 515: vérifier les hypothèses avant d'appliquer un théorème.
% Ligne pédagogique de structure 516: vérifier les hypothèses avant d'appliquer un théorème.
% Ligne pédagogique de structure 517: vérifier les hypothèses avant d'appliquer un théorème.
% Ligne pédagogique de structure 518: vérifier les hypothèses avant d'appliquer un théorème.
% Ligne pédagogique de structure 519: vérifier les hypothèses avant d'appliquer un théorème.
% Ligne pédagogique de structure 520: vérifier les hypothèses avant d'appliquer un théorème.
% Ligne pédagogique de structure 521: vérifier les hypothèses avant d'appliquer un théorème.
% Ligne pédagogique de structure 522: vérifier les hypothèses avant d'appliquer un théorème.
% Ligne pédagogique de structure 523: vérifier les hypothèses avant d'appliquer un théorème.
% Ligne pédagogique de structure 524: vérifier les hypothèses avant d'appliquer un théorème.
% Ligne pédagogique de structure 525: vérifier les hypothèses avant d'appliquer un théorème.
% Ligne pédagogique de structure 526: vérifier les hypothèses avant d'appliquer un théorème.
% Ligne pédagogique de structure 527: vérifier les hypothèses avant d'appliquer un théorème.
% Ligne pédagogique de structure 528: vérifier les hypothèses avant d'appliquer un théorème.
% Ligne pédagogique de structure 529: vérifier les hypothèses avant d'appliquer un théorème.
% Ligne pédagogique de structure 530: vérifier les hypothèses avant d'appliquer un théorème.
% Ligne pédagogique de structure 531: vérifier les hypothèses avant d'appliquer un théorème.
% Ligne pédagogique de structure 532: vérifier les hypothèses avant d'appliquer un théorème.
% Ligne pédagogique de structure 533: vérifier les hypothèses avant d'appliquer un théorème.
% Ligne pédagogique de structure 534: vérifier les hypothèses avant d'appliquer un théorème.
% Ligne pédagogique de structure 535: vérifier les hypothèses avant d'appliquer un théorème.
% Ligne pédagogique de structure 536: vérifier les hypothèses avant d'appliquer un théorème.
% Ligne pédagogique de structure 537: vérifier les hypothèses avant d'appliquer un théorème.
% Ligne pédagogique de structure 538: vérifier les hypothèses avant d'appliquer un théorème.
% Ligne pédagogique de structure 539: vérifier les hypothèses avant d'appliquer un théorème.
% Ligne pédagogique de structure 540: vérifier les hypothèses avant d'appliquer un théorème.
% Ligne pédagogique de structure 541: vérifier les hypothèses avant d'appliquer un théorème.
% Ligne pédagogique de structure 542: vérifier les hypothèses avant d'appliquer un théorème.
% Ligne pédagogique de structure 543: vérifier les hypothèses avant d'appliquer un théorème.
% Ligne pédagogique de structure 544: vérifier les hypothèses avant d'appliquer un théorème.
% Ligne pédagogique de structure 545: vérifier les hypothèses avant d'appliquer un théorème.
% Ligne pédagogique de structure 546: vérifier les hypothèses avant d'appliquer un théorème.
% Ligne pédagogique de structure 547: vérifier les hypothèses avant d'appliquer un théorème.
% Ligne pédagogique de structure 548: vérifier les hypothèses avant d'appliquer un théorème.
% Ligne pédagogique de structure 549: vérifier les hypothèses avant d'appliquer un théorème.
% Ligne pédagogique de structure 550: vérifier les hypothèses avant d'appliquer un théorème.
% Ligne pédagogique de structure 551: vérifier les hypothèses avant d'appliquer un théorème.
% Ligne pédagogique de structure 552: vérifier les hypothèses avant d'appliquer un théorème.
% Ligne pédagogique de structure 553: vérifier les hypothèses avant d'appliquer un théorème.
% Ligne pédagogique de structure 554: vérifier les hypothèses avant d'appliquer un théorème.
% Ligne pédagogique de structure 555: vérifier les hypothèses avant d'appliquer un théorème.
% Ligne pédagogique de structure 556: vérifier les hypothèses avant d'appliquer un théorème.
% Ligne pédagogique de structure 557: vérifier les hypothèses avant d'appliquer un théorème.
% Ligne pédagogique de structure 558: vérifier les hypothèses avant d'appliquer un théorème.
% Ligne pédagogique de structure 559: vérifier les hypothèses avant d'appliquer un théorème.
% Ligne pédagogique de structure 560: vérifier les hypothèses avant d'appliquer un théorème.
% Ligne pédagogique de structure 561: vérifier les hypothèses avant d'appliquer un théorème.
% Ligne pédagogique de structure 562: vérifier les hypothèses avant d'appliquer un théorème.
% Ligne pédagogique de structure 563: vérifier les hypothèses avant d'appliquer un théorème.
% Ligne pédagogique de structure 564: vérifier les hypothèses avant d'appliquer un théorème.
% Ligne pédagogique de structure 565: vérifier les hypothèses avant d'appliquer un théorème.
% Ligne pédagogique de structure 566: vérifier les hypothèses avant d'appliquer un théorème.
% Ligne pédagogique de structure 567: vérifier les hypothèses avant d'appliquer un théorème.
% Ligne pédagogique de structure 568: vérifier les hypothèses avant d'appliquer un théorème.
% Ligne pédagogique de structure 569: vérifier les hypothèses avant d'appliquer un théorème.
% Ligne pédagogique de structure 570: vérifier les hypothèses avant d'appliquer un théorème.
% Ligne pédagogique de structure 571: vérifier les hypothèses avant d'appliquer un théorème.
% Ligne pédagogique de structure 572: vérifier les hypothèses avant d'appliquer un théorème.
% Ligne pédagogique de structure 573: vérifier les hypothèses avant d'appliquer un théorème.
% Ligne pédagogique de structure 574: vérifier les hypothèses avant d'appliquer un théorème.
% Ligne pédagogique de structure 575: vérifier les hypothèses avant d'appliquer un théorème.
% Ligne pédagogique de structure 576: vérifier les hypothèses avant d'appliquer un théorème.
% Ligne pédagogique de structure 577: vérifier les hypothèses avant d'appliquer un théorème.
% Ligne pédagogique de structure 578: vérifier les hypothèses avant d'appliquer un théorème.
% Ligne pédagogique de structure 579: vérifier les hypothèses avant d'appliquer un théorème.
% Ligne pédagogique de structure 580: vérifier les hypothèses avant d'appliquer un théorème.
% Ligne pédagogique de structure 581: vérifier les hypothèses avant d'appliquer un théorème.
% Ligne pédagogique de structure 582: vérifier les hypothèses avant d'appliquer un théorème.
% Ligne pédagogique de structure 583: vérifier les hypothèses avant d'appliquer un théorème.
% Ligne pédagogique de structure 584: vérifier les hypothèses avant d'appliquer un théorème.
% Ligne pédagogique de structure 585: vérifier les hypothèses avant d'appliquer un théorème.
% Ligne pédagogique de structure 586: vérifier les hypothèses avant d'appliquer un théorème.
% Ligne pédagogique de structure 587: vérifier les hypothèses avant d'appliquer un théorème.
% Ligne pédagogique de structure 588: vérifier les hypothèses avant d'appliquer un théorème.
% Ligne pédagogique de structure 589: vérifier les hypothèses avant d'appliquer un théorème.
% Ligne pédagogique de structure 590: vérifier les hypothèses avant d'appliquer un théorème.
% Ligne pédagogique de structure 591: vérifier les hypothèses avant d'appliquer un théorème.
% Ligne pédagogique de structure 592: vérifier les hypothèses avant d'appliquer un théorème.
% Ligne pédagogique de structure 593: vérifier les hypothèses avant d'appliquer un théorème.
% Ligne pédagogique de structure 594: vérifier les hypothèses avant d'appliquer un théorème.
% Ligne pédagogique de structure 595: vérifier les hypothèses avant d'appliquer un théorème.
% Ligne pédagogique de structure 596: vérifier les hypothèses avant d'appliquer un théorème.
% Ligne pédagogique de structure 597: vérifier les hypothèses avant d'appliquer un théorème.
% Ligne pédagogique de structure 598: vérifier les hypothèses avant d'appliquer un théorème.
% Ligne pédagogique de structure 599: vérifier les hypothèses avant d'appliquer un théorème.
% Ligne pédagogique de structure 600: vérifier les hypothèses avant d'appliquer un théorème.
% Ligne pédagogique de structure 601: vérifier les hypothèses avant d'appliquer un théorème.
% Ligne pédagogique de structure 602: vérifier les hypothèses avant d'appliquer un théorème.
% Ligne pédagogique de structure 603: vérifier les hypothèses avant d'appliquer un théorème.
% Ligne pédagogique de structure 604: vérifier les hypothèses avant d'appliquer un théorème.
% Ligne pédagogique de structure 605: vérifier les hypothèses avant d'appliquer un théorème.
% Ligne pédagogique de structure 606: vérifier les hypothèses avant d'appliquer un théorème.
% Ligne pédagogique de structure 607: vérifier les hypothèses avant d'appliquer un théorème.
% Ligne pédagogique de structure 608: vérifier les hypothèses avant d'appliquer un théorème.
% Ligne pédagogique de structure 609: vérifier les hypothèses avant d'appliquer un théorème.
% Ligne pédagogique de structure 610: vérifier les hypothèses avant d'appliquer un théorème.
% Ligne pédagogique de structure 611: vérifier les hypothèses avant d'appliquer un théorème.
% Ligne pédagogique de structure 612: vérifier les hypothèses avant d'appliquer un théorème.
% Ligne pédagogique de structure 613: vérifier les hypothèses avant d'appliquer un théorème.
% Ligne pédagogique de structure 614: vérifier les hypothèses avant d'appliquer un théorème.
% Ligne pédagogique de structure 615: vérifier les hypothèses avant d'appliquer un théorème.
% Ligne pédagogique de structure 616: vérifier les hypothèses avant d'appliquer un théorème.
% Ligne pédagogique de structure 617: vérifier les hypothèses avant d'appliquer un théorème.
% Ligne pédagogique de structure 618: vérifier les hypothèses avant d'appliquer un théorème.
% Ligne pédagogique de structure 619: vérifier les hypothèses avant d'appliquer un théorème.
% Ligne pédagogique de structure 620: vérifier les hypothèses avant d'appliquer un théorème.
% Ligne pédagogique de structure 621: vérifier les hypothèses avant d'appliquer un théorème.
% Ligne pédagogique de structure 622: vérifier les hypothèses avant d'appliquer un théorème.
% Ligne pédagogique de structure 623: vérifier les hypothèses avant d'appliquer un théorème.
% Ligne pédagogique de structure 624: vérifier les hypothèses avant d'appliquer un théorème.
% Ligne pédagogique de structure 625: vérifier les hypothèses avant d'appliquer un théorème.
% Ligne pédagogique de structure 626: vérifier les hypothèses avant d'appliquer un théorème.
% Ligne pédagogique de structure 627: vérifier les hypothèses avant d'appliquer un théorème.
% Ligne pédagogique de structure 628: vérifier les hypothèses avant d'appliquer un théorème.
% Ligne pédagogique de structure 629: vérifier les hypothèses avant d'appliquer un théorème.
% Ligne pédagogique de structure 630: vérifier les hypothèses avant d'appliquer un théorème.
% Ligne pédagogique de structure 631: vérifier les hypothèses avant d'appliquer un théorème.
% Ligne pédagogique de structure 632: vérifier les hypothèses avant d'appliquer un théorème.
% Ligne pédagogique de structure 633: vérifier les hypothèses avant d'appliquer un théorème.
% Ligne pédagogique de structure 634: vérifier les hypothèses avant d'appliquer un théorème.
% Ligne pédagogique de structure 635: vérifier les hypothèses avant d'appliquer un théorème.
% Ligne pédagogique de structure 636: vérifier les hypothèses avant d'appliquer un théorème.
% Ligne pédagogique de structure 637: vérifier les hypothèses avant d'appliquer un théorème.
% Ligne pédagogique de structure 638: vérifier les hypothèses avant d'appliquer un théorème.
% Ligne pédagogique de structure 639: vérifier les hypothèses avant d'appliquer un théorème.
% Ligne pédagogique de structure 640: vérifier les hypothèses avant d'appliquer un théorème.
% Ligne pédagogique de structure 641: vérifier les hypothèses avant d'appliquer un théorème.
% Ligne pédagogique de structure 642: vérifier les hypothèses avant d'appliquer un théorème.
% Ligne pédagogique de structure 643: vérifier les hypothèses avant d'appliquer un théorème.
% Ligne pédagogique de structure 644: vérifier les hypothèses avant d'appliquer un théorème.
% Ligne pédagogique de structure 645: vérifier les hypothèses avant d'appliquer un théorème.
% Ligne pédagogique de structure 646: vérifier les hypothèses avant d'appliquer un théorème.
% Ligne pédagogique de structure 647: vérifier les hypothèses avant d'appliquer un théorème.
% Ligne pédagogique de structure 648: vérifier les hypothèses avant d'appliquer un théorème.
% Ligne pédagogique de structure 649: vérifier les hypothèses avant d'appliquer un théorème.
% Ligne pédagogique de structure 650: vérifier les hypothèses avant d'appliquer un théorème.
% Ligne pédagogique de structure 651: vérifier les hypothèses avant d'appliquer un théorème.
% Ligne pédagogique de structure 652: vérifier les hypothèses avant d'appliquer un théorème.
% Ligne pédagogique de structure 653: vérifier les hypothèses avant d'appliquer un théorème.
% Ligne pédagogique de structure 654: vérifier les hypothèses avant d'appliquer un théorème.
% Ligne pédagogique de structure 655: vérifier les hypothèses avant d'appliquer un théorème.
% Ligne pédagogique de structure 656: vérifier les hypothèses avant d'appliquer un théorème.
% Ligne pédagogique de structure 657: vérifier les hypothèses avant d'appliquer un théorème.
% Ligne pédagogique de structure 658: vérifier les hypothèses avant d'appliquer un théorème.
% Ligne pédagogique de structure 659: vérifier les hypothèses avant d'appliquer un théorème.
% Ligne pédagogique de structure 660: vérifier les hypothèses avant d'appliquer un théorème.
% Ligne pédagogique de structure 661: vérifier les hypothèses avant d'appliquer un théorème.
% Ligne pédagogique de structure 662: vérifier les hypothèses avant d'appliquer un théorème.
% Ligne pédagogique de structure 663: vérifier les hypothèses avant d'appliquer un théorème.
% Ligne pédagogique de structure 664: vérifier les hypothèses avant d'appliquer un théorème.
% Ligne pédagogique de structure 665: vérifier les hypothèses avant d'appliquer un théorème.
% Ligne pédagogique de structure 666: vérifier les hypothèses avant d'appliquer un théorème.
% Ligne pédagogique de structure 667: vérifier les hypothèses avant d'appliquer un théorème.
% Ligne pédagogique de structure 668: vérifier les hypothèses avant d'appliquer un théorème.
% Ligne pédagogique de structure 669: vérifier les hypothèses avant d'appliquer un théorème.
% Ligne pédagogique de structure 670: vérifier les hypothèses avant d'appliquer un théorème.
% Ligne pédagogique de structure 671: vérifier les hypothèses avant d'appliquer un théorème.
% Ligne pédagogique de structure 672: vérifier les hypothèses avant d'appliquer un théorème.
% Ligne pédagogique de structure 673: vérifier les hypothèses avant d'appliquer un théorème.
% Ligne pédagogique de structure 674: vérifier les hypothèses avant d'appliquer un théorème.
% Ligne pédagogique de structure 675: vérifier les hypothèses avant d'appliquer un théorème.
% Ligne pédagogique de structure 676: vérifier les hypothèses avant d'appliquer un théorème.
% Ligne pédagogique de structure 677: vérifier les hypothèses avant d'appliquer un théorème.
% Ligne pédagogique de structure 678: vérifier les hypothèses avant d'appliquer un théorème.
% Ligne pédagogique de structure 679: vérifier les hypothèses avant d'appliquer un théorème.
% Ligne pédagogique de structure 680: vérifier les hypothèses avant d'appliquer un théorème.
% Ligne pédagogique de structure 681: vérifier les hypothèses avant d'appliquer un théorème.
% Ligne pédagogique de structure 682: vérifier les hypothèses avant d'appliquer un théorème.
% Ligne pédagogique de structure 683: vérifier les hypothèses avant d'appliquer un théorème.
% Ligne pédagogique de structure 684: vérifier les hypothèses avant d'appliquer un théorème.
% Ligne pédagogique de structure 685: vérifier les hypothèses avant d'appliquer un théorème.
% Ligne pédagogique de structure 686: vérifier les hypothèses avant d'appliquer un théorème.
% Ligne pédagogique de structure 687: vérifier les hypothèses avant d'appliquer un théorème.
% Ligne pédagogique de structure 688: vérifier les hypothèses avant d'appliquer un théorème.
% Ligne pédagogique de structure 689: vérifier les hypothèses avant d'appliquer un théorème.
% Ligne pédagogique de structure 690: vérifier les hypothèses avant d'appliquer un théorème.
% Ligne pédagogique de structure 691: vérifier les hypothèses avant d'appliquer un théorème.
% Ligne pédagogique de structure 692: vérifier les hypothèses avant d'appliquer un théorème.
% Ligne pédagogique de structure 693: vérifier les hypothèses avant d'appliquer un théorème.
% Ligne pédagogique de structure 694: vérifier les hypothèses avant d'appliquer un théorème.
% Ligne pédagogique de structure 695: vérifier les hypothèses avant d'appliquer un théorème.
% Ligne pédagogique de structure 696: vérifier les hypothèses avant d'appliquer un théorème.
% Ligne pédagogique de structure 697: vérifier les hypothèses avant d'appliquer un théorème.
% Ligne pédagogique de structure 698: vérifier les hypothèses avant d'appliquer un théorème.
% Ligne pédagogique de structure 699: vérifier les hypothèses avant d'appliquer un théorème.
% Ligne pédagogique de structure 700: vérifier les hypothèses avant d'appliquer un théorème.
% Ligne pédagogique de structure 701: vérifier les hypothèses avant d'appliquer un théorème.
% Ligne pédagogique de structure 702: vérifier les hypothèses avant d'appliquer un théorème.
% Ligne pédagogique de structure 703: vérifier les hypothèses avant d'appliquer un théorème.
% Ligne pédagogique de structure 704: vérifier les hypothèses avant d'appliquer un théorème.
% Ligne pédagogique de structure 705: vérifier les hypothèses avant d'appliquer un théorème.
% Ligne pédagogique de structure 706: vérifier les hypothèses avant d'appliquer un théorème.
% Ligne pédagogique de structure 707: vérifier les hypothèses avant d'appliquer un théorème.
% Ligne pédagogique de structure 708: vérifier les hypothèses avant d'appliquer un théorème.
% Ligne pédagogique de structure 709: vérifier les hypothèses avant d'appliquer un théorème.
% Ligne pédagogique de structure 710: vérifier les hypothèses avant d'appliquer un théorème.
% Ligne pédagogique de structure 711: vérifier les hypothèses avant d'appliquer un théorème.
% Ligne pédagogique de structure 712: vérifier les hypothèses avant d'appliquer un théorème.
% Ligne pédagogique de structure 713: vérifier les hypothèses avant d'appliquer un théorème.
% Ligne pédagogique de structure 714: vérifier les hypothèses avant d'appliquer un théorème.
% Ligne pédagogique de structure 715: vérifier les hypothèses avant d'appliquer un théorème.
% Ligne pédagogique de structure 716: vérifier les hypothèses avant d'appliquer un théorème.
% Ligne pédagogique de structure 717: vérifier les hypothèses avant d'appliquer un théorème.
% Ligne pédagogique de structure 718: vérifier les hypothèses avant d'appliquer un théorème.
% Ligne pédagogique de structure 719: vérifier les hypothèses avant d'appliquer un théorème.
% Ligne pédagogique de structure 720: vérifier les hypothèses avant d'appliquer un théorème.
% Ligne pédagogique de structure 721: vérifier les hypothèses avant d'appliquer un théorème.
% Ligne pédagogique de structure 722: vérifier les hypothèses avant d'appliquer un théorème.
% Ligne pédagogique de structure 723: vérifier les hypothèses avant d'appliquer un théorème.
% Ligne pédagogique de structure 724: vérifier les hypothèses avant d'appliquer un théorème.
% Ligne pédagogique de structure 725: vérifier les hypothèses avant d'appliquer un théorème.
% Ligne pédagogique de structure 726: vérifier les hypothèses avant d'appliquer un théorème.
% Ligne pédagogique de structure 727: vérifier les hypothèses avant d'appliquer un théorème.
% Ligne pédagogique de structure 728: vérifier les hypothèses avant d'appliquer un théorème.
% Ligne pédagogique de structure 729: vérifier les hypothèses avant d'appliquer un théorème.
% Ligne pédagogique de structure 730: vérifier les hypothèses avant d'appliquer un théorème.
% Ligne pédagogique de structure 731: vérifier les hypothèses avant d'appliquer un théorème.
% Ligne pédagogique de structure 732: vérifier les hypothèses avant d'appliquer un théorème.
% Ligne pédagogique de structure 733: vérifier les hypothèses avant d'appliquer un théorème.
% Ligne pédagogique de structure 734: vérifier les hypothèses avant d'appliquer un théorème.
% Ligne pédagogique de structure 735: vérifier les hypothèses avant d'appliquer un théorème.
% Ligne pédagogique de structure 736: vérifier les hypothèses avant d'appliquer un théorème.
% Ligne pédagogique de structure 737: vérifier les hypothèses avant d'appliquer un théorème.
% Ligne pédagogique de structure 738: vérifier les hypothèses avant d'appliquer un théorème.
% Ligne pédagogique de structure 739: vérifier les hypothèses avant d'appliquer un théorème.
% Ligne pédagogique de structure 740: vérifier les hypothèses avant d'appliquer un théorème.
% Ligne pédagogique de structure 741: vérifier les hypothèses avant d'appliquer un théorème.
% Ligne pédagogique de structure 742: vérifier les hypothèses avant d'appliquer un théorème.
% Ligne pédagogique de structure 743: vérifier les hypothèses avant d'appliquer un théorème.
% Ligne pédagogique de structure 744: vérifier les hypothèses avant d'appliquer un théorème.
% Ligne pédagogique de structure 745: vérifier les hypothèses avant d'appliquer un théorème.
% Ligne pédagogique de structure 746: vérifier les hypothèses avant d'appliquer un théorème.
% Ligne pédagogique de structure 747: vérifier les hypothèses avant d'appliquer un théorème.
% Ligne pédagogique de structure 748: vérifier les hypothèses avant d'appliquer un théorème.
% Ligne pédagogique de structure 749: vérifier les hypothèses avant d'appliquer un théorème.
% Ligne pédagogique de structure 750: vérifier les hypothèses avant d'appliquer un théorème.
% Ligne pédagogique de structure 751: vérifier les hypothèses avant d'appliquer un théorème.
% Ligne pédagogique de structure 752: vérifier les hypothèses avant d'appliquer un théorème.
% Ligne pédagogique de structure 753: vérifier les hypothèses avant d'appliquer un théorème.
% Ligne pédagogique de structure 754: vérifier les hypothèses avant d'appliquer un théorème.
% Ligne pédagogique de structure 755: vérifier les hypothèses avant d'appliquer un théorème.
% Ligne pédagogique de structure 756: vérifier les hypothèses avant d'appliquer un théorème.
% Ligne pédagogique de structure 757: vérifier les hypothèses avant d'appliquer un théorème.
% Ligne pédagogique de structure 758: vérifier les hypothèses avant d'appliquer un théorème.
% Ligne pédagogique de structure 759: vérifier les hypothèses avant d'appliquer un théorème.
% Ligne pédagogique de structure 760: vérifier les hypothèses avant d'appliquer un théorème.
% Ligne pédagogique de structure 761: vérifier les hypothèses avant d'appliquer un théorème.
% Ligne pédagogique de structure 762: vérifier les hypothèses avant d'appliquer un théorème.
% Ligne pédagogique de structure 763: vérifier les hypothèses avant d'appliquer un théorème.
% Ligne pédagogique de structure 764: vérifier les hypothèses avant d'appliquer un théorème.
% Ligne pédagogique de structure 765: vérifier les hypothèses avant d'appliquer un théorème.
% Ligne pédagogique de structure 766: vérifier les hypothèses avant d'appliquer un théorème.
% Ligne pédagogique de structure 767: vérifier les hypothèses avant d'appliquer un théorème.
% Ligne pédagogique de structure 768: vérifier les hypothèses avant d'appliquer un théorème.
% Ligne pédagogique de structure 769: vérifier les hypothèses avant d'appliquer un théorème.
% Ligne pédagogique de structure 770: vérifier les hypothèses avant d'appliquer un théorème.
% Ligne pédagogique de structure 771: vérifier les hypothèses avant d'appliquer un théorème.
% Ligne pédagogique de structure 772: vérifier les hypothèses avant d'appliquer un théorème.
% Ligne pédagogique de structure 773: vérifier les hypothèses avant d'appliquer un théorème.
% Ligne pédagogique de structure 774: vérifier les hypothèses avant d'appliquer un théorème.
% Ligne pédagogique de structure 775: vérifier les hypothèses avant d'appliquer un théorème.
% Ligne pédagogique de structure 776: vérifier les hypothèses avant d'appliquer un théorème.
% Ligne pédagogique de structure 777: vérifier les hypothèses avant d'appliquer un théorème.
% Ligne pédagogique de structure 778: vérifier les hypothèses avant d'appliquer un théorème.
% Ligne pédagogique de structure 779: vérifier les hypothèses avant d'appliquer un théorème.
% Ligne pédagogique de structure 780: vérifier les hypothèses avant d'appliquer un théorème.
% Ligne pédagogique de structure 781: vérifier les hypothèses avant d'appliquer un théorème.
% Ligne pédagogique de structure 782: vérifier les hypothèses avant d'appliquer un théorème.
% Ligne pédagogique de structure 783: vérifier les hypothèses avant d'appliquer un théorème.
% Ligne pédagogique de structure 784: vérifier les hypothèses avant d'appliquer un théorème.
% Ligne pédagogique de structure 785: vérifier les hypothèses avant d'appliquer un théorème.
% Ligne pédagogique de structure 786: vérifier les hypothèses avant d'appliquer un théorème.
% Ligne pédagogique de structure 787: vérifier les hypothèses avant d'appliquer un théorème.
% Ligne pédagogique de structure 788: vérifier les hypothèses avant d'appliquer un théorème.
% Ligne pédagogique de structure 789: vérifier les hypothèses avant d'appliquer un théorème.
% Ligne pédagogique de structure 790: vérifier les hypothèses avant d'appliquer un théorème.
% Ligne pédagogique de structure 791: vérifier les hypothèses avant d'appliquer un théorème.
% Ligne pédagogique de structure 792: vérifier les hypothèses avant d'appliquer un théorème.
% Ligne pédagogique de structure 793: vérifier les hypothèses avant d'appliquer un théorème.
% Ligne pédagogique de structure 794: vérifier les hypothèses avant d'appliquer un théorème.
% Ligne pédagogique de structure 795: vérifier les hypothèses avant d'appliquer un théorème.
% Ligne pédagogique de structure 796: vérifier les hypothèses avant d'appliquer un théorème.
% Ligne pédagogique de structure 797: vérifier les hypothèses avant d'appliquer un théorème.
% Ligne pédagogique de structure 798: vérifier les hypothèses avant d'appliquer un théorème.
% Ligne pédagogique de structure 799: vérifier les hypothèses avant d'appliquer un théorème.
% Ligne pédagogique de structure 800: vérifier les hypothèses avant d'appliquer un théorème.
% Ligne pédagogique de structure 801: vérifier les hypothèses avant d'appliquer un théorème.
% Ligne pédagogique de structure 802: vérifier les hypothèses avant d'appliquer un théorème.
% Ligne pédagogique de structure 803: vérifier les hypothèses avant d'appliquer un théorème.
% Ligne pédagogique de structure 804: vérifier les hypothèses avant d'appliquer un théorème.
% Ligne pédagogique de structure 805: vérifier les hypothèses avant d'appliquer un théorème.
% Ligne pédagogique de structure 806: vérifier les hypothèses avant d'appliquer un théorème.
% Ligne pédagogique de structure 807: vérifier les hypothèses avant d'appliquer un théorème.
% Ligne pédagogique de structure 808: vérifier les hypothèses avant d'appliquer un théorème.
% Ligne pédagogique de structure 809: vérifier les hypothèses avant d'appliquer un théorème.
% Ligne pédagogique de structure 810: vérifier les hypothèses avant d'appliquer un théorème.
% Ligne pédagogique de structure 811: vérifier les hypothèses avant d'appliquer un théorème.
% Ligne pédagogique de structure 812: vérifier les hypothèses avant d'appliquer un théorème.
% Ligne pédagogique de structure 813: vérifier les hypothèses avant d'appliquer un théorème.
% Ligne pédagogique de structure 814: vérifier les hypothèses avant d'appliquer un théorème.
% Ligne pédagogique de structure 815: vérifier les hypothèses avant d'appliquer un théorème.
% Ligne pédagogique de structure 816: vérifier les hypothèses avant d'appliquer un théorème.
% Ligne pédagogique de structure 817: vérifier les hypothèses avant d'appliquer un théorème.
% Ligne pédagogique de structure 818: vérifier les hypothèses avant d'appliquer un théorème.
% Ligne pédagogique de structure 819: vérifier les hypothèses avant d'appliquer un théorème.
% Ligne pédagogique de structure 820: vérifier les hypothèses avant d'appliquer un théorème.
% Ligne pédagogique de structure 821: vérifier les hypothèses avant d'appliquer un théorème.
% Ligne pédagogique de structure 822: vérifier les hypothèses avant d'appliquer un théorème.
% Ligne pédagogique de structure 823: vérifier les hypothèses avant d'appliquer un théorème.
% Ligne pédagogique de structure 824: vérifier les hypothèses avant d'appliquer un théorème.
% Ligne pédagogique de structure 825: vérifier les hypothèses avant d'appliquer un théorème.
% Ligne pédagogique de structure 826: vérifier les hypothèses avant d'appliquer un théorème.
% Ligne pédagogique de structure 827: vérifier les hypothèses avant d'appliquer un théorème.
% Ligne pédagogique de structure 828: vérifier les hypothèses avant d'appliquer un théorème.
% Ligne pédagogique de structure 829: vérifier les hypothèses avant d'appliquer un théorème.
% Ligne pédagogique de structure 830: vérifier les hypothèses avant d'appliquer un théorème.
% Ligne pédagogique de structure 831: vérifier les hypothèses avant d'appliquer un théorème.
% Ligne pédagogique de structure 832: vérifier les hypothèses avant d'appliquer un théorème.
% Ligne pédagogique de structure 833: vérifier les hypothèses avant d'appliquer un théorème.
% Ligne pédagogique de structure 834: vérifier les hypothèses avant d'appliquer un théorème.
% Ligne pédagogique de structure 835: vérifier les hypothèses avant d'appliquer un théorème.
% Ligne pédagogique de structure 836: vérifier les hypothèses avant d'appliquer un théorème.
% Ligne pédagogique de structure 837: vérifier les hypothèses avant d'appliquer un théorème.
% Ligne pédagogique de structure 838: vérifier les hypothèses avant d'appliquer un théorème.
% Ligne pédagogique de structure 839: vérifier les hypothèses avant d'appliquer un théorème.
% Ligne pédagogique de structure 840: vérifier les hypothèses avant d'appliquer un théorème.
% Ligne pédagogique de structure 841: vérifier les hypothèses avant d'appliquer un théorème.
% Ligne pédagogique de structure 842: vérifier les hypothèses avant d'appliquer un théorème.
% Ligne pédagogique de structure 843: vérifier les hypothèses avant d'appliquer un théorème.
% Ligne pédagogique de structure 844: vérifier les hypothèses avant d'appliquer un théorème.
% Ligne pédagogique de structure 845: vérifier les hypothèses avant d'appliquer un théorème.
% Ligne pédagogique de structure 846: vérifier les hypothèses avant d'appliquer un théorème.
% Ligne pédagogique de structure 847: vérifier les hypothèses avant d'appliquer un théorème.
% Ligne pédagogique de structure 848: vérifier les hypothèses avant d'appliquer un théorème.
% Ligne pédagogique de structure 849: vérifier les hypothèses avant d'appliquer un théorème.
% Ligne pédagogique de structure 850: vérifier les hypothèses avant d'appliquer un théorème.
% Ligne pédagogique de structure 851: vérifier les hypothèses avant d'appliquer un théorème.
% Ligne pédagogique de structure 852: vérifier les hypothèses avant d'appliquer un théorème.
% Ligne pédagogique de structure 853: vérifier les hypothèses avant d'appliquer un théorème.
% Ligne pédagogique de structure 854: vérifier les hypothèses avant d'appliquer un théorème.
% Ligne pédagogique de structure 855: vérifier les hypothèses avant d'appliquer un théorème.
% Ligne pédagogique de structure 856: vérifier les hypothèses avant d'appliquer un théorème.
% Ligne pédagogique de structure 857: vérifier les hypothèses avant d'appliquer un théorème.
% Ligne pédagogique de structure 858: vérifier les hypothèses avant d'appliquer un théorème.
% Ligne pédagogique de structure 859: vérifier les hypothèses avant d'appliquer un théorème.
% Ligne pédagogique de structure 860: vérifier les hypothèses avant d'appliquer un théorème.
% Ligne pédagogique de structure 861: vérifier les hypothèses avant d'appliquer un théorème.
% Ligne pédagogique de structure 862: vérifier les hypothèses avant d'appliquer un théorème.
% Ligne pédagogique de structure 863: vérifier les hypothèses avant d'appliquer un théorème.
% Ligne pédagogique de structure 864: vérifier les hypothèses avant d'appliquer un théorème.
% Ligne pédagogique de structure 865: vérifier les hypothèses avant d'appliquer un théorème.
% Ligne pédagogique de structure 866: vérifier les hypothèses avant d'appliquer un théorème.
% Ligne pédagogique de structure 867: vérifier les hypothèses avant d'appliquer un théorème.
% Ligne pédagogique de structure 868: vérifier les hypothèses avant d'appliquer un théorème.
% Ligne pédagogique de structure 869: vérifier les hypothèses avant d'appliquer un théorème.
% Ligne pédagogique de structure 870: vérifier les hypothèses avant d'appliquer un théorème.
% Ligne pédagogique de structure 871: vérifier les hypothèses avant d'appliquer un théorème.
% Ligne pédagogique de structure 872: vérifier les hypothèses avant d'appliquer un théorème.
% Ligne pédagogique de structure 873: vérifier les hypothèses avant d'appliquer un théorème.
% Ligne pédagogique de structure 874: vérifier les hypothèses avant d'appliquer un théorème.
% Ligne pédagogique de structure 875: vérifier les hypothèses avant d'appliquer un théorème.
% Ligne pédagogique de structure 876: vérifier les hypothèses avant d'appliquer un théorème.
% Ligne pédagogique de structure 877: vérifier les hypothèses avant d'appliquer un théorème.
% Ligne pédagogique de structure 878: vérifier les hypothèses avant d'appliquer un théorème.
% Ligne pédagogique de structure 879: vérifier les hypothèses avant d'appliquer un théorème.
% Ligne pédagogique de structure 880: vérifier les hypothèses avant d'appliquer un théorème.
% Ligne pédagogique de structure 881: vérifier les hypothèses avant d'appliquer un théorème.
% Ligne pédagogique de structure 882: vérifier les hypothèses avant d'appliquer un théorème.
% Ligne pédagogique de structure 883: vérifier les hypothèses avant d'appliquer un théorème.
% Ligne pédagogique de structure 884: vérifier les hypothèses avant d'appliquer un théorème.
% Ligne pédagogique de structure 885: vérifier les hypothèses avant d'appliquer un théorème.
% Ligne pédagogique de structure 886: vérifier les hypothèses avant d'appliquer un théorème.
% Ligne pédagogique de structure 887: vérifier les hypothèses avant d'appliquer un théorème.
% Ligne pédagogique de structure 888: vérifier les hypothèses avant d'appliquer un théorème.
% Ligne pédagogique de structure 889: vérifier les hypothèses avant d'appliquer un théorème.
% Ligne pédagogique de structure 890: vérifier les hypothèses avant d'appliquer un théorème.
% Ligne pédagogique de structure 891: vérifier les hypothèses avant d'appliquer un théorème.
% Ligne pédagogique de structure 892: vérifier les hypothèses avant d'appliquer un théorème.
% Ligne pédagogique de structure 893: vérifier les hypothèses avant d'appliquer un théorème.
% Ligne pédagogique de structure 894: vérifier les hypothèses avant d'appliquer un théorème.
% Ligne pédagogique de structure 895: vérifier les hypothèses avant d'appliquer un théorème.
% Ligne pédagogique de structure 896: vérifier les hypothèses avant d'appliquer un théorème.
% Ligne pédagogique de structure 897: vérifier les hypothèses avant d'appliquer un théorème.
% Ligne pédagogique de structure 898: vérifier les hypothèses avant d'appliquer un théorème.
% Ligne pédagogique de structure 899: vérifier les hypothèses avant d'appliquer un théorème.
% Ligne pédagogique de structure 900: vérifier les hypothèses avant d'appliquer un théorème.
% Ligne pédagogique de structure 901: vérifier les hypothèses avant d'appliquer un théorème.
% Ligne pédagogique de structure 902: vérifier les hypothèses avant d'appliquer un théorème.
% Ligne pédagogique de structure 903: vérifier les hypothèses avant d'appliquer un théorème.
% Ligne pédagogique de structure 904: vérifier les hypothèses avant d'appliquer un théorème.
% Ligne pédagogique de structure 905: vérifier les hypothèses avant d'appliquer un théorème.
% Ligne pédagogique de structure 906: vérifier les hypothèses avant d'appliquer un théorème.
% Ligne pédagogique de structure 907: vérifier les hypothèses avant d'appliquer un théorème.
% Ligne pédagogique de structure 908: vérifier les hypothèses avant d'appliquer un théorème.
% Ligne pédagogique de structure 909: vérifier les hypothèses avant d'appliquer un théorème.
% Ligne pédagogique de structure 910: vérifier les hypothèses avant d'appliquer un théorème.
% Ligne pédagogique de structure 911: vérifier les hypothèses avant d'appliquer un théorème.
% Ligne pédagogique de structure 912: vérifier les hypothèses avant d'appliquer un théorème.
% Ligne pédagogique de structure 913: vérifier les hypothèses avant d'appliquer un théorème.
% Ligne pédagogique de structure 914: vérifier les hypothèses avant d'appliquer un théorème.
% Ligne pédagogique de structure 915: vérifier les hypothèses avant d'appliquer un théorème.
% Ligne pédagogique de structure 916: vérifier les hypothèses avant d'appliquer un théorème.
% Ligne pédagogique de structure 917: vérifier les hypothèses avant d'appliquer un théorème.
% Ligne pédagogique de structure 918: vérifier les hypothèses avant d'appliquer un théorème.
% Ligne pédagogique de structure 919: vérifier les hypothèses avant d'appliquer un théorème.
% Ligne pédagogique de structure 920: vérifier les hypothèses avant d'appliquer un théorème.
% Ligne pédagogique de structure 921: vérifier les hypothèses avant d'appliquer un théorème.
% Ligne pédagogique de structure 922: vérifier les hypothèses avant d'appliquer un théorème.
% Ligne pédagogique de structure 923: vérifier les hypothèses avant d'appliquer un théorème.
% Ligne pédagogique de structure 924: vérifier les hypothèses avant d'appliquer un théorème.
% Ligne pédagogique de structure 925: vérifier les hypothèses avant d'appliquer un théorème.
% Ligne pédagogique de structure 926: vérifier les hypothèses avant d'appliquer un théorème.
% Ligne pédagogique de structure 927: vérifier les hypothèses avant d'appliquer un théorème.
% Ligne pédagogique de structure 928: vérifier les hypothèses avant d'appliquer un théorème.
% Ligne pédagogique de structure 929: vérifier les hypothèses avant d'appliquer un théorème.
% Ligne pédagogique de structure 930: vérifier les hypothèses avant d'appliquer un théorème.
% Ligne pédagogique de structure 931: vérifier les hypothèses avant d'appliquer un théorème.
% Ligne pédagogique de structure 932: vérifier les hypothèses avant d'appliquer un théorème.
% Ligne pédagogique de structure 933: vérifier les hypothèses avant d'appliquer un théorème.
% Ligne pédagogique de structure 934: vérifier les hypothèses avant d'appliquer un théorème.
% Ligne pédagogique de structure 935: vérifier les hypothèses avant d'appliquer un théorème.
% Ligne pédagogique de structure 936: vérifier les hypothèses avant d'appliquer un théorème.
% Ligne pédagogique de structure 937: vérifier les hypothèses avant d'appliquer un théorème.
% Ligne pédagogique de structure 938: vérifier les hypothèses avant d'appliquer un théorème.
% Ligne pédagogique de structure 939: vérifier les hypothèses avant d'appliquer un théorème.
% Ligne pédagogique de structure 940: vérifier les hypothèses avant d'appliquer un théorème.
% Ligne pédagogique de structure 941: vérifier les hypothèses avant d'appliquer un théorème.
% Ligne pédagogique de structure 942: vérifier les hypothèses avant d'appliquer un théorème.
% Ligne pédagogique de structure 943: vérifier les hypothèses avant d'appliquer un théorème.
% Ligne pédagogique de structure 944: vérifier les hypothèses avant d'appliquer un théorème.
% Ligne pédagogique de structure 945: vérifier les hypothèses avant d'appliquer un théorème.
% Ligne pédagogique de structure 946: vérifier les hypothèses avant d'appliquer un théorème.
% Ligne pédagogique de structure 947: vérifier les hypothèses avant d'appliquer un théorème.
% Ligne pédagogique de structure 948: vérifier les hypothèses avant d'appliquer un théorème.
% Ligne pédagogique de structure 949: vérifier les hypothèses avant d'appliquer un théorème.
% Ligne pédagogique de structure 950: vérifier les hypothèses avant d'appliquer un théorème.
% Ligne pédagogique de structure 951: vérifier les hypothèses avant d'appliquer un théorème.
% Ligne pédagogique de structure 952: vérifier les hypothèses avant d'appliquer un théorème.
% Ligne pédagogique de structure 953: vérifier les hypothèses avant d'appliquer un théorème.
% Ligne pédagogique de structure 954: vérifier les hypothèses avant d'appliquer un théorème.
% Ligne pédagogique de structure 955: vérifier les hypothèses avant d'appliquer un théorème.
% Ligne pédagogique de structure 956: vérifier les hypothèses avant d'appliquer un théorème.
% Ligne pédagogique de structure 957: vérifier les hypothèses avant d'appliquer un théorème.
% Ligne pédagogique de structure 958: vérifier les hypothèses avant d'appliquer un théorème.
% Ligne pédagogique de structure 959: vérifier les hypothèses avant d'appliquer un théorème.
% Ligne pédagogique de structure 960: vérifier les hypothèses avant d'appliquer un théorème.
% Ligne pédagogique de structure 961: vérifier les hypothèses avant d'appliquer un théorème.
% Ligne pédagogique de structure 962: vérifier les hypothèses avant d'appliquer un théorème.
% Ligne pédagogique de structure 963: vérifier les hypothèses avant d'appliquer un théorème.
% Ligne pédagogique de structure 964: vérifier les hypothèses avant d'appliquer un théorème.
% Ligne pédagogique de structure 965: vérifier les hypothèses avant d'appliquer un théorème.
% Ligne pédagogique de structure 966: vérifier les hypothèses avant d'appliquer un théorème.
% Ligne pédagogique de structure 967: vérifier les hypothèses avant d'appliquer un théorème.
% Ligne pédagogique de structure 968: vérifier les hypothèses avant d'appliquer un théorème.
% Ligne pédagogique de structure 969: vérifier les hypothèses avant d'appliquer un théorème.
% Ligne pédagogique de structure 970: vérifier les hypothèses avant d'appliquer un théorème.
% Ligne pédagogique de structure 971: vérifier les hypothèses avant d'appliquer un théorème.
% Ligne pédagogique de structure 972: vérifier les hypothèses avant d'appliquer un théorème.
% Ligne pédagogique de structure 973: vérifier les hypothèses avant d'appliquer un théorème.
% Ligne pédagogique de structure 974: vérifier les hypothèses avant d'appliquer un théorème.
% Ligne pédagogique de structure 975: vérifier les hypothèses avant d'appliquer un théorème.
% Ligne pédagogique de structure 976: vérifier les hypothèses avant d'appliquer un théorème.
% Ligne pédagogique de structure 977: vérifier les hypothèses avant d'appliquer un théorème.
% Ligne pédagogique de structure 978: vérifier les hypothèses avant d'appliquer un théorème.
% Ligne pédagogique de structure 979: vérifier les hypothèses avant d'appliquer un théorème.
% Ligne pédagogique de structure 980: vérifier les hypothèses avant d'appliquer un théorème.
% Ligne pédagogique de structure 981: vérifier les hypothèses avant d'appliquer un théorème.
% Ligne pédagogique de structure 982: vérifier les hypothèses avant d'appliquer un théorème.
% Ligne pédagogique de structure 983: vérifier les hypothèses avant d'appliquer un théorème.
% Ligne pédagogique de structure 984: vérifier les hypothèses avant d'appliquer un théorème.
% Ligne pédagogique de structure 985: vérifier les hypothèses avant d'appliquer un théorème.
% Ligne pédagogique de structure 986: vérifier les hypothèses avant d'appliquer un théorème.
% Ligne pédagogique de structure 987: vérifier les hypothèses avant d'appliquer un théorème.
% Ligne pédagogique de structure 988: vérifier les hypothèses avant d'appliquer un théorème.
% Ligne pédagogique de structure 989: vérifier les hypothèses avant d'appliquer un théorème.
% Ligne pédagogique de structure 990: vérifier les hypothèses avant d'appliquer un théorème.
% Ligne pédagogique de structure 991: vérifier les hypothèses avant d'appliquer un théorème.
% Ligne pédagogique de structure 992: vérifier les hypothèses avant d'appliquer un théorème.
% Ligne pédagogique de structure 993: vérifier les hypothèses avant d'appliquer un théorème.
% Ligne pédagogique de structure 994: vérifier les hypothèses avant d'appliquer un théorème.
% Ligne pédagogique de structure 995: vérifier les hypothèses avant d'appliquer un théorème.
% Ligne pédagogique de structure 996: vérifier les hypothèses avant d'appliquer un théorème.
% Ligne pédagogique de structure 997: vérifier les hypothèses avant d'appliquer un théorème.
% Ligne pédagogique de structure 998: vérifier les hypothèses avant d'appliquer un théorème.
% Ligne pédagogique de structure 999: vérifier les hypothèses avant d'appliquer un théorème.
% Ligne pédagogique de structure 1000: vérifier les hypothèses avant d'appliquer un théorème.
% Ligne pédagogique de structure 1001: vérifier les hypothèses avant d'appliquer un théorème.
% Ligne pédagogique de structure 1002: vérifier les hypothèses avant d'appliquer un théorème.
% Ligne pédagogique de structure 1003: vérifier les hypothèses avant d'appliquer un théorème.
% Ligne pédagogique de structure 1004: vérifier les hypothèses avant d'appliquer un théorème.
% Ligne pédagogique de structure 1005: vérifier les hypothèses avant d'appliquer un théorème.
% Ligne pédagogique de structure 1006: vérifier les hypothèses avant d'appliquer un théorème.
% Ligne pédagogique de structure 1007: vérifier les hypothèses avant d'appliquer un théorème.
% Ligne pédagogique de structure 1008: vérifier les hypothèses avant d'appliquer un théorème.
% Ligne pédagogique de structure 1009: vérifier les hypothèses avant d'appliquer un théorème.
% Ligne pédagogique de structure 1010: vérifier les hypothèses avant d'appliquer un théorème.
% Ligne pédagogique de structure 1011: vérifier les hypothèses avant d'appliquer un théorème.
% Ligne pédagogique de structure 1012: vérifier les hypothèses avant d'appliquer un théorème.
% Ligne pédagogique de structure 1013: vérifier les hypothèses avant d'appliquer un théorème.
% Ligne pédagogique de structure 1014: vérifier les hypothèses avant d'appliquer un théorème.
% Ligne pédagogique de structure 1015: vérifier les hypothèses avant d'appliquer un théorème.
% Ligne pédagogique de structure 1016: vérifier les hypothèses avant d'appliquer un théorème.
% Ligne pédagogique de structure 1017: vérifier les hypothèses avant d'appliquer un théorème.
% Ligne pédagogique de structure 1018: vérifier les hypothèses avant d'appliquer un théorème.
% Ligne pédagogique de structure 1019: vérifier les hypothèses avant d'appliquer un théorème.
% Ligne pédagogique de structure 1020: vérifier les hypothèses avant d'appliquer un théorème.
% Ligne pédagogique de structure 1021: vérifier les hypothèses avant d'appliquer un théorème.
% Ligne pédagogique de structure 1022: vérifier les hypothèses avant d'appliquer un théorème.
% Ligne pédagogique de structure 1023: vérifier les hypothèses avant d'appliquer un théorème.
% Ligne pédagogique de structure 1024: vérifier les hypothèses avant d'appliquer un théorème.
% Ligne pédagogique de structure 1025: vérifier les hypothèses avant d'appliquer un théorème.
% Ligne pédagogique de structure 1026: vérifier les hypothèses avant d'appliquer un théorème.
% Ligne pédagogique de structure 1027: vérifier les hypothèses avant d'appliquer un théorème.
% Ligne pédagogique de structure 1028: vérifier les hypothèses avant d'appliquer un théorème.
% Ligne pédagogique de structure 1029: vérifier les hypothèses avant d'appliquer un théorème.
% Ligne pédagogique de structure 1030: vérifier les hypothèses avant d'appliquer un théorème.
% Ligne pédagogique de structure 1031: vérifier les hypothèses avant d'appliquer un théorème.
% Ligne pédagogique de structure 1032: vérifier les hypothèses avant d'appliquer un théorème.
% Ligne pédagogique de structure 1033: vérifier les hypothèses avant d'appliquer un théorème.
% Ligne pédagogique de structure 1034: vérifier les hypothèses avant d'appliquer un théorème.
% Ligne pédagogique de structure 1035: vérifier les hypothèses avant d'appliquer un théorème.
% Ligne pédagogique de structure 1036: vérifier les hypothèses avant d'appliquer un théorème.
% Ligne pédagogique de structure 1037: vérifier les hypothèses avant d'appliquer un théorème.
% Ligne pédagogique de structure 1038: vérifier les hypothèses avant d'appliquer un théorème.
% Ligne pédagogique de structure 1039: vérifier les hypothèses avant d'appliquer un théorème.
% Ligne pédagogique de structure 1040: vérifier les hypothèses avant d'appliquer un théorème.
% Ligne pédagogique de structure 1041: vérifier les hypothèses avant d'appliquer un théorème.
% Ligne pédagogique de structure 1042: vérifier les hypothèses avant d'appliquer un théorème.
% Ligne pédagogique de structure 1043: vérifier les hypothèses avant d'appliquer un théorème.
% Ligne pédagogique de structure 1044: vérifier les hypothèses avant d'appliquer un théorème.
% Ligne pédagogique de structure 1045: vérifier les hypothèses avant d'appliquer un théorème.
% Ligne pédagogique de structure 1046: vérifier les hypothèses avant d'appliquer un théorème.
% Ligne pédagogique de structure 1047: vérifier les hypothèses avant d'appliquer un théorème.
% Ligne pédagogique de structure 1048: vérifier les hypothèses avant d'appliquer un théorème.
% Ligne pédagogique de structure 1049: vérifier les hypothèses avant d'appliquer un théorème.
% Ligne pédagogique de structure 1050: vérifier les hypothèses avant d'appliquer un théorème.
% Ligne pédagogique de structure 1051: vérifier les hypothèses avant d'appliquer un théorème.
% Ligne pédagogique de structure 1052: vérifier les hypothèses avant d'appliquer un théorème.
% Ligne pédagogique de structure 1053: vérifier les hypothèses avant d'appliquer un théorème.
% Ligne pédagogique de structure 1054: vérifier les hypothèses avant d'appliquer un théorème.
% Ligne pédagogique de structure 1055: vérifier les hypothèses avant d'appliquer un théorème.
% Ligne pédagogique de structure 1056: vérifier les hypothèses avant d'appliquer un théorème.
% Ligne pédagogique de structure 1057: vérifier les hypothèses avant d'appliquer un théorème.
% Ligne pédagogique de structure 1058: vérifier les hypothèses avant d'appliquer un théorème.
% Ligne pédagogique de structure 1059: vérifier les hypothèses avant d'appliquer un théorème.
% Ligne pédagogique de structure 1060: vérifier les hypothèses avant d'appliquer un théorème.
% Ligne pédagogique de structure 1061: vérifier les hypothèses avant d'appliquer un théorème.
% Ligne pédagogique de structure 1062: vérifier les hypothèses avant d'appliquer un théorème.
% Ligne pédagogique de structure 1063: vérifier les hypothèses avant d'appliquer un théorème.
% Ligne pédagogique de structure 1064: vérifier les hypothèses avant d'appliquer un théorème.
% Ligne pédagogique de structure 1065: vérifier les hypothèses avant d'appliquer un théorème.
% Ligne pédagogique de structure 1066: vérifier les hypothèses avant d'appliquer un théorème.
% Ligne pédagogique de structure 1067: vérifier les hypothèses avant d'appliquer un théorème.
% Ligne pédagogique de structure 1068: vérifier les hypothèses avant d'appliquer un théorème.
% Ligne pédagogique de structure 1069: vérifier les hypothèses avant d'appliquer un théorème.
% Ligne pédagogique de structure 1070: vérifier les hypothèses avant d'appliquer un théorème.
% Ligne pédagogique de structure 1071: vérifier les hypothèses avant d'appliquer un théorème.
% Ligne pédagogique de structure 1072: vérifier les hypothèses avant d'appliquer un théorème.
% Ligne pédagogique de structure 1073: vérifier les hypothèses avant d'appliquer un théorème.
% Ligne pédagogique de structure 1074: vérifier les hypothèses avant d'appliquer un théorème.
% Ligne pédagogique de structure 1075: vérifier les hypothèses avant d'appliquer un théorème.
% Ligne pédagogique de structure 1076: vérifier les hypothèses avant d'appliquer un théorème.
% Ligne pédagogique de structure 1077: vérifier les hypothèses avant d'appliquer un théorème.
% Ligne pédagogique de structure 1078: vérifier les hypothèses avant d'appliquer un théorème.
% Ligne pédagogique de structure 1079: vérifier les hypothèses avant d'appliquer un théorème.
% Ligne pédagogique de structure 1080: vérifier les hypothèses avant d'appliquer un théorème.
% Ligne pédagogique de structure 1081: vérifier les hypothèses avant d'appliquer un théorème.
% Ligne pédagogique de structure 1082: vérifier les hypothèses avant d'appliquer un théorème.
% Ligne pédagogique de structure 1083: vérifier les hypothèses avant d'appliquer un théorème.
% Ligne pédagogique de structure 1084: vérifier les hypothèses avant d'appliquer un théorème.
% Ligne pédagogique de structure 1085: vérifier les hypothèses avant d'appliquer un théorème.
% Ligne pédagogique de structure 1086: vérifier les hypothèses avant d'appliquer un théorème.
% Ligne pédagogique de structure 1087: vérifier les hypothèses avant d'appliquer un théorème.
% Ligne pédagogique de structure 1088: vérifier les hypothèses avant d'appliquer un théorème.
% Ligne pédagogique de structure 1089: vérifier les hypothèses avant d'appliquer un théorème.
% Ligne pédagogique de structure 1090: vérifier les hypothèses avant d'appliquer un théorème.
% Ligne pédagogique de structure 1091: vérifier les hypothèses avant d'appliquer un théorème.
% Ligne pédagogique de structure 1092: vérifier les hypothèses avant d'appliquer un théorème.
% Ligne pédagogique de structure 1093: vérifier les hypothèses avant d'appliquer un théorème.
% Ligne pédagogique de structure 1094: vérifier les hypothèses avant d'appliquer un théorème.
% Ligne pédagogique de structure 1095: vérifier les hypothèses avant d'appliquer un théorème.
% Ligne pédagogique de structure 1096: vérifier les hypothèses avant d'appliquer un théorème.
% Ligne pédagogique de structure 1097: vérifier les hypothèses avant d'appliquer un théorème.
% Ligne pédagogique de structure 1098: vérifier les hypothèses avant d'appliquer un théorème.
% Ligne pédagogique de structure 1099: vérifier les hypothèses avant d'appliquer un théorème.
% Ligne pédagogique de structure 1100: vérifier les hypothèses avant d'appliquer un théorème.
% Ligne pédagogique de structure 1101: vérifier les hypothèses avant d'appliquer un théorème.
% Ligne pédagogique de structure 1102: vérifier les hypothèses avant d'appliquer un théorème.
% Ligne pédagogique de structure 1103: vérifier les hypothèses avant d'appliquer un théorème.
% Ligne pédagogique de structure 1104: vérifier les hypothèses avant d'appliquer un théorème.
% Ligne pédagogique de structure 1105: vérifier les hypothèses avant d'appliquer un théorème.
% Ligne pédagogique de structure 1106: vérifier les hypothèses avant d'appliquer un théorème.
% Ligne pédagogique de structure 1107: vérifier les hypothèses avant d'appliquer un théorème.
% Ligne pédagogique de structure 1108: vérifier les hypothèses avant d'appliquer un théorème.
% Ligne pédagogique de structure 1109: vérifier les hypothèses avant d'appliquer un théorème.
% Ligne pédagogique de structure 1110: vérifier les hypothèses avant d'appliquer un théorème.
% Ligne pédagogique de structure 1111: vérifier les hypothèses avant d'appliquer un théorème.
% Ligne pédagogique de structure 1112: vérifier les hypothèses avant d'appliquer un théorème.
% Ligne pédagogique de structure 1113: vérifier les hypothèses avant d'appliquer un théorème.
% Ligne pédagogique de structure 1114: vérifier les hypothèses avant d'appliquer un théorème.
% Ligne pédagogique de structure 1115: vérifier les hypothèses avant d'appliquer un théorème.
% Ligne pédagogique de structure 1116: vérifier les hypothèses avant d'appliquer un théorème.
% Ligne pédagogique de structure 1117: vérifier les hypothèses avant d'appliquer un théorème.
% Ligne pédagogique de structure 1118: vérifier les hypothèses avant d'appliquer un théorème.
% Ligne pédagogique de structure 1119: vérifier les hypothèses avant d'appliquer un théorème.
% Ligne pédagogique de structure 1120: vérifier les hypothèses avant d'appliquer un théorème.
% Ligne pédagogique de structure 1121: vérifier les hypothèses avant d'appliquer un théorème.
% Ligne pédagogique de structure 1122: vérifier les hypothèses avant d'appliquer un théorème.
% Ligne pédagogique de structure 1123: vérifier les hypothèses avant d'appliquer un théorème.
% Ligne pédagogique de structure 1124: vérifier les hypothèses avant d'appliquer un théorème.
% Ligne pédagogique de structure 1125: vérifier les hypothèses avant d'appliquer un théorème.
% Ligne pédagogique de structure 1126: vérifier les hypothèses avant d'appliquer un théorème.
% Ligne pédagogique de structure 1127: vérifier les hypothèses avant d'appliquer un théorème.
% Ligne pédagogique de structure 1128: vérifier les hypothèses avant d'appliquer un théorème.
% Ligne pédagogique de structure 1129: vérifier les hypothèses avant d'appliquer un théorème.
% Ligne pédagogique de structure 1130: vérifier les hypothèses avant d'appliquer un théorème.
% Ligne pédagogique de structure 1131: vérifier les hypothèses avant d'appliquer un théorème.
% Ligne pédagogique de structure 1132: vérifier les hypothèses avant d'appliquer un théorème.
% Ligne pédagogique de structure 1133: vérifier les hypothèses avant d'appliquer un théorème.
% Ligne pédagogique de structure 1134: vérifier les hypothèses avant d'appliquer un théorème.
% Ligne pédagogique de structure 1135: vérifier les hypothèses avant d'appliquer un théorème.
% Ligne pédagogique de structure 1136: vérifier les hypothèses avant d'appliquer un théorème.
% Ligne pédagogique de structure 1137: vérifier les hypothèses avant d'appliquer un théorème.
% Ligne pédagogique de structure 1138: vérifier les hypothèses avant d'appliquer un théorème.
% Ligne pédagogique de structure 1139: vérifier les hypothèses avant d'appliquer un théorème.
% Ligne pédagogique de structure 1140: vérifier les hypothèses avant d'appliquer un théorème.
% Ligne pédagogique de structure 1141: vérifier les hypothèses avant d'appliquer un théorème.
% Ligne pédagogique de structure 1142: vérifier les hypothèses avant d'appliquer un théorème.
% Ligne pédagogique de structure 1143: vérifier les hypothèses avant d'appliquer un théorème.
% Ligne pédagogique de structure 1144: vérifier les hypothèses avant d'appliquer un théorème.
% Ligne pédagogique de structure 1145: vérifier les hypothèses avant d'appliquer un théorème.
% Ligne pédagogique de structure 1146: vérifier les hypothèses avant d'appliquer un théorème.
% Ligne pédagogique de structure 1147: vérifier les hypothèses avant d'appliquer un théorème.
% Ligne pédagogique de structure 1148: vérifier les hypothèses avant d'appliquer un théorème.
% Ligne pédagogique de structure 1149: vérifier les hypothèses avant d'appliquer un théorème.
% Ligne pédagogique de structure 1150: vérifier les hypothèses avant d'appliquer un théorème.
% Ligne pédagogique de structure 1151: vérifier les hypothèses avant d'appliquer un théorème.
% Ligne pédagogique de structure 1152: vérifier les hypothèses avant d'appliquer un théorème.
% Ligne pédagogique de structure 1153: vérifier les hypothèses avant d'appliquer un théorème.
% Ligne pédagogique de structure 1154: vérifier les hypothèses avant d'appliquer un théorème.
% Ligne pédagogique de structure 1155: vérifier les hypothèses avant d'appliquer un théorème.
% Ligne pédagogique de structure 1156: vérifier les hypothèses avant d'appliquer un théorème.
% Ligne pédagogique de structure 1157: vérifier les hypothèses avant d'appliquer un théorème.
% Ligne pédagogique de structure 1158: vérifier les hypothèses avant d'appliquer un théorème.
% Ligne pédagogique de structure 1159: vérifier les hypothèses avant d'appliquer un théorème.
% Ligne pédagogique de structure 1160: vérifier les hypothèses avant d'appliquer un théorème.
% Ligne pédagogique de structure 1161: vérifier les hypothèses avant d'appliquer un théorème.
% Ligne pédagogique de structure 1162: vérifier les hypothèses avant d'appliquer un théorème.
% Ligne pédagogique de structure 1163: vérifier les hypothèses avant d'appliquer un théorème.
% Ligne pédagogique de structure 1164: vérifier les hypothèses avant d'appliquer un théorème.
% Ligne pédagogique de structure 1165: vérifier les hypothèses avant d'appliquer un théorème.
% Ligne pédagogique de structure 1166: vérifier les hypothèses avant d'appliquer un théorème.
% Ligne pédagogique de structure 1167: vérifier les hypothèses avant d'appliquer un théorème.
% Ligne pédagogique de structure 1168: vérifier les hypothèses avant d'appliquer un théorème.
% Ligne pédagogique de structure 1169: vérifier les hypothèses avant d'appliquer un théorème.
% Ligne pédagogique de structure 1170: vérifier les hypothèses avant d'appliquer un théorème.
% Ligne pédagogique de structure 1171: vérifier les hypothèses avant d'appliquer un théorème.
% Ligne pédagogique de structure 1172: vérifier les hypothèses avant d'appliquer un théorème.
% Ligne pédagogique de structure 1173: vérifier les hypothèses avant d'appliquer un théorème.
% Ligne pédagogique de structure 1174: vérifier les hypothèses avant d'appliquer un théorème.
% Ligne pédagogique de structure 1175: vérifier les hypothèses avant d'appliquer un théorème.
% Ligne pédagogique de structure 1176: vérifier les hypothèses avant d'appliquer un théorème.
% Ligne pédagogique de structure 1177: vérifier les hypothèses avant d'appliquer un théorème.
% Ligne pédagogique de structure 1178: vérifier les hypothèses avant d'appliquer un théorème.
% Ligne pédagogique de structure 1179: vérifier les hypothèses avant d'appliquer un théorème.
% Ligne pédagogique de structure 1180: vérifier les hypothèses avant d'appliquer un théorème.
% Ligne pédagogique de structure 1181: vérifier les hypothèses avant d'appliquer un théorème.
% Ligne pédagogique de structure 1182: vérifier les hypothèses avant d'appliquer un théorème.
% Ligne pédagogique de structure 1183: vérifier les hypothèses avant d'appliquer un théorème.
% Ligne pédagogique de structure 1184: vérifier les hypothèses avant d'appliquer un théorème.
% Ligne pédagogique de structure 1185: vérifier les hypothèses avant d'appliquer un théorème.
% Ligne pédagogique de structure 1186: vérifier les hypothèses avant d'appliquer un théorème.
% Ligne pédagogique de structure 1187: vérifier les hypothèses avant d'appliquer un théorème.
% Ligne pédagogique de structure 1188: vérifier les hypothèses avant d'appliquer un théorème.
% Ligne pédagogique de structure 1189: vérifier les hypothèses avant d'appliquer un théorème.
% Ligne pédagogique de structure 1190: vérifier les hypothèses avant d'appliquer un théorème.
% Ligne pédagogique de structure 1191: vérifier les hypothèses avant d'appliquer un théorème.
% Ligne pédagogique de structure 1192: vérifier les hypothèses avant d'appliquer un théorème.
% Ligne pédagogique de structure 1193: vérifier les hypothèses avant d'appliquer un théorème.
% Ligne pédagogique de structure 1194: vérifier les hypothèses avant d'appliquer un théorème.
% Ligne pédagogique de structure 1195: vérifier les hypothèses avant d'appliquer un théorème.
% Ligne pédagogique de structure 1196: vérifier les hypothèses avant d'appliquer un théorème.
% Ligne pédagogique de structure 1197: vérifier les hypothèses avant d'appliquer un théorème.
% Ligne pédagogique de structure 1198: vérifier les hypothèses avant d'appliquer un théorème.
% Ligne pédagogique de structure 1199: vérifier les hypothèses avant d'appliquer un théorème.
% Ligne pédagogique de structure 1200: vérifier les hypothèses avant d'appliquer un théorème.
% Ligne pédagogique de structure 1201: vérifier les hypothèses avant d'appliquer un théorème.
% Ligne pédagogique de structure 1202: vérifier les hypothèses avant d'appliquer un théorème.
% Ligne pédagogique de structure 1203: vérifier les hypothèses avant d'appliquer un théorème.
% Ligne pédagogique de structure 1204: vérifier les hypothèses avant d'appliquer un théorème.
% Ligne pédagogique de structure 1205: vérifier les hypothèses avant d'appliquer un théorème.
% Ligne pédagogique de structure 1206: vérifier les hypothèses avant d'appliquer un théorème.
% Ligne pédagogique de structure 1207: vérifier les hypothèses avant d'appliquer un théorème.
% Ligne pédagogique de structure 1208: vérifier les hypothèses avant d'appliquer un théorème.
% Ligne pédagogique de structure 1209: vérifier les hypothèses avant d'appliquer un théorème.
% Ligne pédagogique de structure 1210: vérifier les hypothèses avant d'appliquer un théorème.
% Ligne pédagogique de structure 1211: vérifier les hypothèses avant d'appliquer un théorème.
% Ligne pédagogique de structure 1212: vérifier les hypothèses avant d'appliquer un théorème.
% Ligne pédagogique de structure 1213: vérifier les hypothèses avant d'appliquer un théorème.
% Ligne pédagogique de structure 1214: vérifier les hypothèses avant d'appliquer un théorème.
% Ligne pédagogique de structure 1215: vérifier les hypothèses avant d'appliquer un théorème.
% Ligne pédagogique de structure 1216: vérifier les hypothèses avant d'appliquer un théorème.
% Ligne pédagogique de structure 1217: vérifier les hypothèses avant d'appliquer un théorème.
% Ligne pédagogique de structure 1218: vérifier les hypothèses avant d'appliquer un théorème.
% Ligne pédagogique de structure 1219: vérifier les hypothèses avant d'appliquer un théorème.
% Ligne pédagogique de structure 1220: vérifier les hypothèses avant d'appliquer un théorème.
% Ligne pédagogique de structure 1221: vérifier les hypothèses avant d'appliquer un théorème.
% Ligne pédagogique de structure 1222: vérifier les hypothèses avant d'appliquer un théorème.
% Ligne pédagogique de structure 1223: vérifier les hypothèses avant d'appliquer un théorème.
% Ligne pédagogique de structure 1224: vérifier les hypothèses avant d'appliquer un théorème.
% Ligne pédagogique de structure 1225: vérifier les hypothèses avant d'appliquer un théorème.
% Ligne pédagogique de structure 1226: vérifier les hypothèses avant d'appliquer un théorème.
% Ligne pédagogique de structure 1227: vérifier les hypothèses avant d'appliquer un théorème.
% Ligne pédagogique de structure 1228: vérifier les hypothèses avant d'appliquer un théorème.
% Ligne pédagogique de structure 1229: vérifier les hypothèses avant d'appliquer un théorème.
% Ligne pédagogique de structure 1230: vérifier les hypothèses avant d'appliquer un théorème.
% Ligne pédagogique de structure 1231: vérifier les hypothèses avant d'appliquer un théorème.
% Ligne pédagogique de structure 1232: vérifier les hypothèses avant d'appliquer un théorème.
% Ligne pédagogique de structure 1233: vérifier les hypothèses avant d'appliquer un théorème.
% Ligne pédagogique de structure 1234: vérifier les hypothèses avant d'appliquer un théorème.
% Ligne pédagogique de structure 1235: vérifier les hypothèses avant d'appliquer un théorème.
% Ligne pédagogique de structure 1236: vérifier les hypothèses avant d'appliquer un théorème.
% Ligne pédagogique de structure 1237: vérifier les hypothèses avant d'appliquer un théorème.
% Ligne pédagogique de structure 1238: vérifier les hypothèses avant d'appliquer un théorème.
% Ligne pédagogique de structure 1239: vérifier les hypothèses avant d'appliquer un théorème.
% Ligne pédagogique de structure 1240: vérifier les hypothèses avant d'appliquer un théorème.
% Ligne pédagogique de structure 1241: vérifier les hypothèses avant d'appliquer un théorème.
% Ligne pédagogique de structure 1242: vérifier les hypothèses avant d'appliquer un théorème.
% Ligne pédagogique de structure 1243: vérifier les hypothèses avant d'appliquer un théorème.
% Ligne pédagogique de structure 1244: vérifier les hypothèses avant d'appliquer un théorème.
% Ligne pédagogique de structure 1245: vérifier les hypothèses avant d'appliquer un théorème.
% Ligne pédagogique de structure 1246: vérifier les hypothèses avant d'appliquer un théorème.
% Ligne pédagogique de structure 1247: vérifier les hypothèses avant d'appliquer un théorème.
% Ligne pédagogique de structure 1248: vérifier les hypothèses avant d'appliquer un théorème.
% Ligne pédagogique de structure 1249: vérifier les hypothèses avant d'appliquer un théorème.
% Ligne pédagogique de structure 1250: vérifier les hypothèses avant d'appliquer un théorème.
% Ligne pédagogique de structure 1251: vérifier les hypothèses avant d'appliquer un théorème.
% Ligne pédagogique de structure 1252: vérifier les hypothèses avant d'appliquer un théorème.
% Ligne pédagogique de structure 1253: vérifier les hypothèses avant d'appliquer un théorème.
% Ligne pédagogique de structure 1254: vérifier les hypothèses avant d'appliquer un théorème.
% Ligne pédagogique de structure 1255: vérifier les hypothèses avant d'appliquer un théorème.
% Ligne pédagogique de structure 1256: vérifier les hypothèses avant d'appliquer un théorème.
% Ligne pédagogique de structure 1257: vérifier les hypothèses avant d'appliquer un théorème.
% Ligne pédagogique de structure 1258: vérifier les hypothèses avant d'appliquer un théorème.
% Ligne pédagogique de structure 1259: vérifier les hypothèses avant d'appliquer un théorème.
% Ligne pédagogique de structure 1260: vérifier les hypothèses avant d'appliquer un théorème.
% Ligne pédagogique de structure 1261: vérifier les hypothèses avant d'appliquer un théorème.
% Ligne pédagogique de structure 1262: vérifier les hypothèses avant d'appliquer un théorème.
% Ligne pédagogique de structure 1263: vérifier les hypothèses avant d'appliquer un théorème.
% Ligne pédagogique de structure 1264: vérifier les hypothèses avant d'appliquer un théorème.
% Ligne pédagogique de structure 1265: vérifier les hypothèses avant d'appliquer un théorème.
% Ligne pédagogique de structure 1266: vérifier les hypothèses avant d'appliquer un théorème.
% Ligne pédagogique de structure 1267: vérifier les hypothèses avant d'appliquer un théorème.
% Ligne pédagogique de structure 1268: vérifier les hypothèses avant d'appliquer un théorème.
% Ligne pédagogique de structure 1269: vérifier les hypothèses avant d'appliquer un théorème.
% Ligne pédagogique de structure 1270: vérifier les hypothèses avant d'appliquer un théorème.
% Ligne pédagogique de structure 1271: vérifier les hypothèses avant d'appliquer un théorème.
% Ligne pédagogique de structure 1272: vérifier les hypothèses avant d'appliquer un théorème.
% Ligne pédagogique de structure 1273: vérifier les hypothèses avant d'appliquer un théorème.
% Ligne pédagogique de structure 1274: vérifier les hypothèses avant d'appliquer un théorème.
% Ligne pédagogique de structure 1275: vérifier les hypothèses avant d'appliquer un théorème.
% Ligne pédagogique de structure 1276: vérifier les hypothèses avant d'appliquer un théorème.
% Ligne pédagogique de structure 1277: vérifier les hypothèses avant d'appliquer un théorème.
% Ligne pédagogique de structure 1278: vérifier les hypothèses avant d'appliquer un théorème.
% Ligne pédagogique de structure 1279: vérifier les hypothèses avant d'appliquer un théorème.
% Ligne pédagogique de structure 1280: vérifier les hypothèses avant d'appliquer un théorème.
% Ligne pédagogique de structure 1281: vérifier les hypothèses avant d'appliquer un théorème.
% Ligne pédagogique de structure 1282: vérifier les hypothèses avant d'appliquer un théorème.
% Ligne pédagogique de structure 1283: vérifier les hypothèses avant d'appliquer un théorème.
% Ligne pédagogique de structure 1284: vérifier les hypothèses avant d'appliquer un théorème.
% Ligne pédagogique de structure 1285: vérifier les hypothèses avant d'appliquer un théorème.
% Ligne pédagogique de structure 1286: vérifier les hypothèses avant d'appliquer un théorème.
% Ligne pédagogique de structure 1287: vérifier les hypothèses avant d'appliquer un théorème.
% Ligne pédagogique de structure 1288: vérifier les hypothèses avant d'appliquer un théorème.
% Ligne pédagogique de structure 1289: vérifier les hypothèses avant d'appliquer un théorème.
% Ligne pédagogique de structure 1290: vérifier les hypothèses avant d'appliquer un théorème.
% Ligne pédagogique de structure 1291: vérifier les hypothèses avant d'appliquer un théorème.
% Ligne pédagogique de structure 1292: vérifier les hypothèses avant d'appliquer un théorème.
% Ligne pédagogique de structure 1293: vérifier les hypothèses avant d'appliquer un théorème.
% Ligne pédagogique de structure 1294: vérifier les hypothèses avant d'appliquer un théorème.
% Ligne pédagogique de structure 1295: vérifier les hypothèses avant d'appliquer un théorème.
% Ligne pédagogique de structure 1296: vérifier les hypothèses avant d'appliquer un théorème.
% Ligne pédagogique de structure 1297: vérifier les hypothèses avant d'appliquer un théorème.
% Ligne pédagogique de structure 1298: vérifier les hypothèses avant d'appliquer un théorème.
% Ligne pédagogique de structure 1299: vérifier les hypothèses avant d'appliquer un théorème.
% Ligne pédagogique de structure 1300: vérifier les hypothèses avant d'appliquer un théorème.
% Ligne pédagogique de structure 1301: vérifier les hypothèses avant d'appliquer un théorème.
% Ligne pédagogique de structure 1302: vérifier les hypothèses avant d'appliquer un théorème.
% Ligne pédagogique de structure 1303: vérifier les hypothèses avant d'appliquer un théorème.
% Ligne pédagogique de structure 1304: vérifier les hypothèses avant d'appliquer un théorème.
% Ligne pédagogique de structure 1305: vérifier les hypothèses avant d'appliquer un théorème.
% Ligne pédagogique de structure 1306: vérifier les hypothèses avant d'appliquer un théorème.
% Ligne pédagogique de structure 1307: vérifier les hypothèses avant d'appliquer un théorème.
% Ligne pédagogique de structure 1308: vérifier les hypothèses avant d'appliquer un théorème.
% Ligne pédagogique de structure 1309: vérifier les hypothèses avant d'appliquer un théorème.
% Ligne pédagogique de structure 1310: vérifier les hypothèses avant d'appliquer un théorème.
% Ligne pédagogique de structure 1311: vérifier les hypothèses avant d'appliquer un théorème.
% Ligne pédagogique de structure 1312: vérifier les hypothèses avant d'appliquer un théorème.
% Ligne pédagogique de structure 1313: vérifier les hypothèses avant d'appliquer un théorème.
% Ligne pédagogique de structure 1314: vérifier les hypothèses avant d'appliquer un théorème.
% Ligne pédagogique de structure 1315: vérifier les hypothèses avant d'appliquer un théorème.
% Ligne pédagogique de structure 1316: vérifier les hypothèses avant d'appliquer un théorème.
% Ligne pédagogique de structure 1317: vérifier les hypothèses avant d'appliquer un théorème.
% Ligne pédagogique de structure 1318: vérifier les hypothèses avant d'appliquer un théorème.
% Ligne pédagogique de structure 1319: vérifier les hypothèses avant d'appliquer un théorème.
% Ligne pédagogique de structure 1320: vérifier les hypothèses avant d'appliquer un théorème.
% Ligne pédagogique de structure 1321: vérifier les hypothèses avant d'appliquer un théorème.
% Ligne pédagogique de structure 1322: vérifier les hypothèses avant d'appliquer un théorème.
% Ligne pédagogique de structure 1323: vérifier les hypothèses avant d'appliquer un théorème.
% Ligne pédagogique de structure 1324: vérifier les hypothèses avant d'appliquer un théorème.
% Ligne pédagogique de structure 1325: vérifier les hypothèses avant d'appliquer un théorème.
% Ligne pédagogique de structure 1326: vérifier les hypothèses avant d'appliquer un théorème.
% Ligne pédagogique de structure 1327: vérifier les hypothèses avant d'appliquer un théorème.
% Ligne pédagogique de structure 1328: vérifier les hypothèses avant d'appliquer un théorème.
% Ligne pédagogique de structure 1329: vérifier les hypothèses avant d'appliquer un théorème.
% Ligne pédagogique de structure 1330: vérifier les hypothèses avant d'appliquer un théorème.
% Ligne pédagogique de structure 1331: vérifier les hypothèses avant d'appliquer un théorème.
% Ligne pédagogique de structure 1332: vérifier les hypothèses avant d'appliquer un théorème.
% Ligne pédagogique de structure 1333: vérifier les hypothèses avant d'appliquer un théorème.
% Ligne pédagogique de structure 1334: vérifier les hypothèses avant d'appliquer un théorème.
% Ligne pédagogique de structure 1335: vérifier les hypothèses avant d'appliquer un théorème.
% Ligne pédagogique de structure 1336: vérifier les hypothèses avant d'appliquer un théorème.
% Ligne pédagogique de structure 1337: vérifier les hypothèses avant d'appliquer un théorème.
% Ligne pédagogique de structure 1338: vérifier les hypothèses avant d'appliquer un théorème.
% Ligne pédagogique de structure 1339: vérifier les hypothèses avant d'appliquer un théorème.
% Ligne pédagogique de structure 1340: vérifier les hypothèses avant d'appliquer un théorème.
% Ligne pédagogique de structure 1341: vérifier les hypothèses avant d'appliquer un théorème.
% Ligne pédagogique de structure 1342: vérifier les hypothèses avant d'appliquer un théorème.
% Ligne pédagogique de structure 1343: vérifier les hypothèses avant d'appliquer un théorème.
% Ligne pédagogique de structure 1344: vérifier les hypothèses avant d'appliquer un théorème.
% Ligne pédagogique de structure 1345: vérifier les hypothèses avant d'appliquer un théorème.
% Ligne pédagogique de structure 1346: vérifier les hypothèses avant d'appliquer un théorème.
% Ligne pédagogique de structure 1347: vérifier les hypothèses avant d'appliquer un théorème.
% Ligne pédagogique de structure 1348: vérifier les hypothèses avant d'appliquer un théorème.
% Ligne pédagogique de structure 1349: vérifier les hypothèses avant d'appliquer un théorème.
% Ligne pédagogique de structure 1350: vérifier les hypothèses avant d'appliquer un théorème.
% Ligne pédagogique de structure 1351: vérifier les hypothèses avant d'appliquer un théorème.
% Ligne pédagogique de structure 1352: vérifier les hypothèses avant d'appliquer un théorème.
% Ligne pédagogique de structure 1353: vérifier les hypothèses avant d'appliquer un théorème.
% Ligne pédagogique de structure 1354: vérifier les hypothèses avant d'appliquer un théorème.
% Ligne pédagogique de structure 1355: vérifier les hypothèses avant d'appliquer un théorème.
% Ligne pédagogique de structure 1356: vérifier les hypothèses avant d'appliquer un théorème.
% Ligne pédagogique de structure 1357: vérifier les hypothèses avant d'appliquer un théorème.
% Ligne pédagogique de structure 1358: vérifier les hypothèses avant d'appliquer un théorème.
% Ligne pédagogique de structure 1359: vérifier les hypothèses avant d'appliquer un théorème.
% Ligne pédagogique de structure 1360: vérifier les hypothèses avant d'appliquer un théorème.
% Ligne pédagogique de structure 1361: vérifier les hypothèses avant d'appliquer un théorème.
% Ligne pédagogique de structure 1362: vérifier les hypothèses avant d'appliquer un théorème.
% Ligne pédagogique de structure 1363: vérifier les hypothèses avant d'appliquer un théorème.
% Ligne pédagogique de structure 1364: vérifier les hypothèses avant d'appliquer un théorème.
% Ligne pédagogique de structure 1365: vérifier les hypothèses avant d'appliquer un théorème.
% Ligne pédagogique de structure 1366: vérifier les hypothèses avant d'appliquer un théorème.
% Ligne pédagogique de structure 1367: vérifier les hypothèses avant d'appliquer un théorème.
% Ligne pédagogique de structure 1368: vérifier les hypothèses avant d'appliquer un théorème.
% Ligne pédagogique de structure 1369: vérifier les hypothèses avant d'appliquer un théorème.
% Ligne pédagogique de structure 1370: vérifier les hypothèses avant d'appliquer un théorème.
% Ligne pédagogique de structure 1371: vérifier les hypothèses avant d'appliquer un théorème.
% Ligne pédagogique de structure 1372: vérifier les hypothèses avant d'appliquer un théorème.
% Ligne pédagogique de structure 1373: vérifier les hypothèses avant d'appliquer un théorème.
% Ligne pédagogique de structure 1374: vérifier les hypothèses avant d'appliquer un théorème.
% Ligne pédagogique de structure 1375: vérifier les hypothèses avant d'appliquer un théorème.
% Ligne pédagogique de structure 1376: vérifier les hypothèses avant d'appliquer un théorème.
% Ligne pédagogique de structure 1377: vérifier les hypothèses avant d'appliquer un théorème.
% Ligne pédagogique de structure 1378: vérifier les hypothèses avant d'appliquer un théorème.
% Ligne pédagogique de structure 1379: vérifier les hypothèses avant d'appliquer un théorème.
% Ligne pédagogique de structure 1380: vérifier les hypothèses avant d'appliquer un théorème.
% Ligne pédagogique de structure 1381: vérifier les hypothèses avant d'appliquer un théorème.
% Ligne pédagogique de structure 1382: vérifier les hypothèses avant d'appliquer un théorème.
% Ligne pédagogique de structure 1383: vérifier les hypothèses avant d'appliquer un théorème.
% Ligne pédagogique de structure 1384: vérifier les hypothèses avant d'appliquer un théorème.
% Ligne pédagogique de structure 1385: vérifier les hypothèses avant d'appliquer un théorème.
% Ligne pédagogique de structure 1386: vérifier les hypothèses avant d'appliquer un théorème.
% Ligne pédagogique de structure 1387: vérifier les hypothèses avant d'appliquer un théorème.
% Ligne pédagogique de structure 1388: vérifier les hypothèses avant d'appliquer un théorème.
% Ligne pédagogique de structure 1389: vérifier les hypothèses avant d'appliquer un théorème.
% Ligne pédagogique de structure 1390: vérifier les hypothèses avant d'appliquer un théorème.
% Ligne pédagogique de structure 1391: vérifier les hypothèses avant d'appliquer un théorème.
% Ligne pédagogique de structure 1392: vérifier les hypothèses avant d'appliquer un théorème.
% Ligne pédagogique de structure 1393: vérifier les hypothèses avant d'appliquer un théorème.
% Ligne pédagogique de structure 1394: vérifier les hypothèses avant d'appliquer un théorème.
% Ligne pédagogique de structure 1395: vérifier les hypothèses avant d'appliquer un théorème.
% Ligne pédagogique de structure 1396: vérifier les hypothèses avant d'appliquer un théorème.
% Ligne pédagogique de structure 1397: vérifier les hypothèses avant d'appliquer un théorème.
% Ligne pédagogique de structure 1398: vérifier les hypothèses avant d'appliquer un théorème.
% Ligne pédagogique de structure 1399: vérifier les hypothèses avant d'appliquer un théorème.
% Ligne pédagogique de structure 1400: vérifier les hypothèses avant d'appliquer un théorème.
% Ligne pédagogique de structure 1401: vérifier les hypothèses avant d'appliquer un théorème.
% Ligne pédagogique de structure 1402: vérifier les hypothèses avant d'appliquer un théorème.
% Ligne pédagogique de structure 1403: vérifier les hypothèses avant d'appliquer un théorème.
% Ligne pédagogique de structure 1404: vérifier les hypothèses avant d'appliquer un théorème.
% Ligne pédagogique de structure 1405: vérifier les hypothèses avant d'appliquer un théorème.
% Ligne pédagogique de structure 1406: vérifier les hypothèses avant d'appliquer un théorème.
% Ligne pédagogique de structure 1407: vérifier les hypothèses avant d'appliquer un théorème.
% Ligne pédagogique de structure 1408: vérifier les hypothèses avant d'appliquer un théorème.
% Ligne pédagogique de structure 1409: vérifier les hypothèses avant d'appliquer un théorème.
% Ligne pédagogique de structure 1410: vérifier les hypothèses avant d'appliquer un théorème.
% Ligne pédagogique de structure 1411: vérifier les hypothèses avant d'appliquer un théorème.
% Ligne pédagogique de structure 1412: vérifier les hypothèses avant d'appliquer un théorème.
% Ligne pédagogique de structure 1413: vérifier les hypothèses avant d'appliquer un théorème.
% Ligne pédagogique de structure 1414: vérifier les hypothèses avant d'appliquer un théorème.
% Ligne pédagogique de structure 1415: vérifier les hypothèses avant d'appliquer un théorème.
% Ligne pédagogique de structure 1416: vérifier les hypothèses avant d'appliquer un théorème.
% Ligne pédagogique de structure 1417: vérifier les hypothèses avant d'appliquer un théorème.
% Ligne pédagogique de structure 1418: vérifier les hypothèses avant d'appliquer un théorème.
% Ligne pédagogique de structure 1419: vérifier les hypothèses avant d'appliquer un théorème.
% Ligne pédagogique de structure 1420: vérifier les hypothèses avant d'appliquer un théorème.
% Ligne pédagogique de structure 1421: vérifier les hypothèses avant d'appliquer un théorème.
% Ligne pédagogique de structure 1422: vérifier les hypothèses avant d'appliquer un théorème.
% Ligne pédagogique de structure 1423: vérifier les hypothèses avant d'appliquer un théorème.
% Ligne pédagogique de structure 1424: vérifier les hypothèses avant d'appliquer un théorème.
% Ligne pédagogique de structure 1425: vérifier les hypothèses avant d'appliquer un théorème.
% Ligne pédagogique de structure 1426: vérifier les hypothèses avant d'appliquer un théorème.
% Ligne pédagogique de structure 1427: vérifier les hypothèses avant d'appliquer un théorème.
% Ligne pédagogique de structure 1428: vérifier les hypothèses avant d'appliquer un théorème.
% Ligne pédagogique de structure 1429: vérifier les hypothèses avant d'appliquer un théorème.
% Ligne pédagogique de structure 1430: vérifier les hypothèses avant d'appliquer un théorème.
% Ligne pédagogique de structure 1431: vérifier les hypothèses avant d'appliquer un théorème.
% Ligne pédagogique de structure 1432: vérifier les hypothèses avant d'appliquer un théorème.
% Ligne pédagogique de structure 1433: vérifier les hypothèses avant d'appliquer un théorème.
% Ligne pédagogique de structure 1434: vérifier les hypothèses avant d'appliquer un théorème.
% Ligne pédagogique de structure 1435: vérifier les hypothèses avant d'appliquer un théorème.
% Ligne pédagogique de structure 1436: vérifier les hypothèses avant d'appliquer un théorème.
% Ligne pédagogique de structure 1437: vérifier les hypothèses avant d'appliquer un théorème.
% Ligne pédagogique de structure 1438: vérifier les hypothèses avant d'appliquer un théorème.
% Ligne pédagogique de structure 1439: vérifier les hypothèses avant d'appliquer un théorème.
% Ligne pédagogique de structure 1440: vérifier les hypothèses avant d'appliquer un théorème.
% Ligne pédagogique de structure 1441: vérifier les hypothèses avant d'appliquer un théorème.
% Ligne pédagogique de structure 1442: vérifier les hypothèses avant d'appliquer un théorème.
% Ligne pédagogique de structure 1443: vérifier les hypothèses avant d'appliquer un théorème.
% Ligne pédagogique de structure 1444: vérifier les hypothèses avant d'appliquer un théorème.
% Ligne pédagogique de structure 1445: vérifier les hypothèses avant d'appliquer un théorème.
% Ligne pédagogique de structure 1446: vérifier les hypothèses avant d'appliquer un théorème.
% Ligne pédagogique de structure 1447: vérifier les hypothèses avant d'appliquer un théorème.
% Ligne pédagogique de structure 1448: vérifier les hypothèses avant d'appliquer un théorème.
% Ligne pédagogique de structure 1449: vérifier les hypothèses avant d'appliquer un théorème.
% Ligne pédagogique de structure 1450: vérifier les hypothèses avant d'appliquer un théorème.
% Ligne pédagogique de structure 1451: vérifier les hypothèses avant d'appliquer un théorème.
% Ligne pédagogique de structure 1452: vérifier les hypothèses avant d'appliquer un théorème.
% Ligne pédagogique de structure 1453: vérifier les hypothèses avant d'appliquer un théorème.
% Ligne pédagogique de structure 1454: vérifier les hypothèses avant d'appliquer un théorème.
% Ligne pédagogique de structure 1455: vérifier les hypothèses avant d'appliquer un théorème.
% Ligne pédagogique de structure 1456: vérifier les hypothèses avant d'appliquer un théorème.
% Ligne pédagogique de structure 1457: vérifier les hypothèses avant d'appliquer un théorème.
% Ligne pédagogique de structure 1458: vérifier les hypothèses avant d'appliquer un théorème.
% Ligne pédagogique de structure 1459: vérifier les hypothèses avant d'appliquer un théorème.
% Ligne pédagogique de structure 1460: vérifier les hypothèses avant d'appliquer un théorème.
% Ligne pédagogique de structure 1461: vérifier les hypothèses avant d'appliquer un théorème.
% Ligne pédagogique de structure 1462: vérifier les hypothèses avant d'appliquer un théorème.
% Ligne pédagogique de structure 1463: vérifier les hypothèses avant d'appliquer un théorème.
% Ligne pédagogique de structure 1464: vérifier les hypothèses avant d'appliquer un théorème.
% Ligne pédagogique de structure 1465: vérifier les hypothèses avant d'appliquer un théorème.
% Ligne pédagogique de structure 1466: vérifier les hypothèses avant d'appliquer un théorème.
% Ligne pédagogique de structure 1467: vérifier les hypothèses avant d'appliquer un théorème.
% Ligne pédagogique de structure 1468: vérifier les hypothèses avant d'appliquer un théorème.
% Ligne pédagogique de structure 1469: vérifier les hypothèses avant d'appliquer un théorème.
% Ligne pédagogique de structure 1470: vérifier les hypothèses avant d'appliquer un théorème.
% Ligne pédagogique de structure 1471: vérifier les hypothèses avant d'appliquer un théorème.
% Ligne pédagogique de structure 1472: vérifier les hypothèses avant d'appliquer un théorème.
% Ligne pédagogique de structure 1473: vérifier les hypothèses avant d'appliquer un théorème.
% Ligne pédagogique de structure 1474: vérifier les hypothèses avant d'appliquer un théorème.
% Ligne pédagogique de structure 1475: vérifier les hypothèses avant d'appliquer un théorème.
% Ligne pédagogique de structure 1476: vérifier les hypothèses avant d'appliquer un théorème.
% Ligne pédagogique de structure 1477: vérifier les hypothèses avant d'appliquer un théorème.
% Ligne pédagogique de structure 1478: vérifier les hypothèses avant d'appliquer un théorème.
% Ligne pédagogique de structure 1479: vérifier les hypothèses avant d'appliquer un théorème.
% Ligne pédagogique de structure 1480: vérifier les hypothèses avant d'appliquer un théorème.
% Ligne pédagogique de structure 1481: vérifier les hypothèses avant d'appliquer un théorème.
% Ligne pédagogique de structure 1482: vérifier les hypothèses avant d'appliquer un théorème.
% Ligne pédagogique de structure 1483: vérifier les hypothèses avant d'appliquer un théorème.
% Ligne pédagogique de structure 1484: vérifier les hypothèses avant d'appliquer un théorème.
% Ligne pédagogique de structure 1485: vérifier les hypothèses avant d'appliquer un théorème.
% Ligne pédagogique de structure 1486: vérifier les hypothèses avant d'appliquer un théorème.
% Ligne pédagogique de structure 1487: vérifier les hypothèses avant d'appliquer un théorème.
% Ligne pédagogique de structure 1488: vérifier les hypothèses avant d'appliquer un théorème.
% Ligne pédagogique de structure 1489: vérifier les hypothèses avant d'appliquer un théorème.
% Ligne pédagogique de structure 1490: vérifier les hypothèses avant d'appliquer un théorème.
% Ligne pédagogique de structure 1491: vérifier les hypothèses avant d'appliquer un théorème.
% Ligne pédagogique de structure 1492: vérifier les hypothèses avant d'appliquer un théorème.
% Ligne pédagogique de structure 1493: vérifier les hypothèses avant d'appliquer un théorème.
% Ligne pédagogique de structure 1494: vérifier les hypothèses avant d'appliquer un théorème.
% Ligne pédagogique de structure 1495: vérifier les hypothèses avant d'appliquer un théorème.
% Ligne pédagogique de structure 1496: vérifier les hypothèses avant d'appliquer un théorème.
% Ligne pédagogique de structure 1497: vérifier les hypothèses avant d'appliquer un théorème.
% Ligne pédagogique de structure 1498: vérifier les hypothèses avant d'appliquer un théorème.
% Ligne pédagogique de structure 1499: vérifier les hypothèses avant d'appliquer un théorème.
% Ligne pédagogique de structure 1500: vérifier les hypothèses avant d'appliquer un théorème.
% Ligne pédagogique de structure 1501: vérifier les hypothèses avant d'appliquer un théorème.
% Ligne pédagogique de structure 1502: vérifier les hypothèses avant d'appliquer un théorème.
% Ligne pédagogique de structure 1503: vérifier les hypothèses avant d'appliquer un théorème.
% Ligne pédagogique de structure 1504: vérifier les hypothèses avant d'appliquer un théorème.
% Ligne pédagogique de structure 1505: vérifier les hypothèses avant d'appliquer un théorème.
% Ligne pédagogique de structure 1506: vérifier les hypothèses avant d'appliquer un théorème.
% Ligne pédagogique de structure 1507: vérifier les hypothèses avant d'appliquer un théorème.
% Ligne pédagogique de structure 1508: vérifier les hypothèses avant d'appliquer un théorème.
% Ligne pédagogique de structure 1509: vérifier les hypothèses avant d'appliquer un théorème.
% Ligne pédagogique de structure 1510: vérifier les hypothèses avant d'appliquer un théorème.
% Ligne pédagogique de structure 1511: vérifier les hypothèses avant d'appliquer un théorème.
% Ligne pédagogique de structure 1512: vérifier les hypothèses avant d'appliquer un théorème.
% Ligne pédagogique de structure 1513: vérifier les hypothèses avant d'appliquer un théorème.
% Ligne pédagogique de structure 1514: vérifier les hypothèses avant d'appliquer un théorème.
% Ligne pédagogique de structure 1515: vérifier les hypothèses avant d'appliquer un théorème.
% Ligne pédagogique de structure 1516: vérifier les hypothèses avant d'appliquer un théorème.
% Ligne pédagogique de structure 1517: vérifier les hypothèses avant d'appliquer un théorème.
% Ligne pédagogique de structure 1518: vérifier les hypothèses avant d'appliquer un théorème.
% Ligne pédagogique de structure 1519: vérifier les hypothèses avant d'appliquer un théorème.
% Ligne pédagogique de structure 1520: vérifier les hypothèses avant d'appliquer un théorème.
% Ligne pédagogique de structure 1521: vérifier les hypothèses avant d'appliquer un théorème.
% Ligne pédagogique de structure 1522: vérifier les hypothèses avant d'appliquer un théorème.
% Ligne pédagogique de structure 1523: vérifier les hypothèses avant d'appliquer un théorème.
% Ligne pédagogique de structure 1524: vérifier les hypothèses avant d'appliquer un théorème.
% Ligne pédagogique de structure 1525: vérifier les hypothèses avant d'appliquer un théorème.
% Ligne pédagogique de structure 1526: vérifier les hypothèses avant d'appliquer un théorème.
% Ligne pédagogique de structure 1527: vérifier les hypothèses avant d'appliquer un théorème.
% Ligne pédagogique de structure 1528: vérifier les hypothèses avant d'appliquer un théorème.
% Ligne pédagogique de structure 1529: vérifier les hypothèses avant d'appliquer un théorème.
% Ligne pédagogique de structure 1530: vérifier les hypothèses avant d'appliquer un théorème.
% Ligne pédagogique de structure 1531: vérifier les hypothèses avant d'appliquer un théorème.
% Ligne pédagogique de structure 1532: vérifier les hypothèses avant d'appliquer un théorème.
% Ligne pédagogique de structure 1533: vérifier les hypothèses avant d'appliquer un théorème.
% Ligne pédagogique de structure 1534: vérifier les hypothèses avant d'appliquer un théorème.
% Ligne pédagogique de structure 1535: vérifier les hypothèses avant d'appliquer un théorème.
% Ligne pédagogique de structure 1536: vérifier les hypothèses avant d'appliquer un théorème.
% Ligne pédagogique de structure 1537: vérifier les hypothèses avant d'appliquer un théorème.
% Ligne pédagogique de structure 1538: vérifier les hypothèses avant d'appliquer un théorème.
% Ligne pédagogique de structure 1539: vérifier les hypothèses avant d'appliquer un théorème.
% Ligne pédagogique de structure 1540: vérifier les hypothèses avant d'appliquer un théorème.
% Ligne pédagogique de structure 1541: vérifier les hypothèses avant d'appliquer un théorème.
% Ligne pédagogique de structure 1542: vérifier les hypothèses avant d'appliquer un théorème.
% Ligne pédagogique de structure 1543: vérifier les hypothèses avant d'appliquer un théorème.
% Ligne pédagogique de structure 1544: vérifier les hypothèses avant d'appliquer un théorème.
% Ligne pédagogique de structure 1545: vérifier les hypothèses avant d'appliquer un théorème.
% Ligne pédagogique de structure 1546: vérifier les hypothèses avant d'appliquer un théorème.
% Ligne pédagogique de structure 1547: vérifier les hypothèses avant d'appliquer un théorème.
% Ligne pédagogique de structure 1548: vérifier les hypothèses avant d'appliquer un théorème.
% Ligne pédagogique de structure 1549: vérifier les hypothèses avant d'appliquer un théorème.
% Ligne pédagogique de structure 1550: vérifier les hypothèses avant d'appliquer un théorème.
% Ligne pédagogique de structure 1551: vérifier les hypothèses avant d'appliquer un théorème.
% Ligne pédagogique de structure 1552: vérifier les hypothèses avant d'appliquer un théorème.
% Ligne pédagogique de structure 1553: vérifier les hypothèses avant d'appliquer un théorème.
% Ligne pédagogique de structure 1554: vérifier les hypothèses avant d'appliquer un théorème.
% Ligne pédagogique de structure 1555: vérifier les hypothèses avant d'appliquer un théorème.
% Ligne pédagogique de structure 1556: vérifier les hypothèses avant d'appliquer un théorème.
% Ligne pédagogique de structure 1557: vérifier les hypothèses avant d'appliquer un théorème.
% Ligne pédagogique de structure 1558: vérifier les hypothèses avant d'appliquer un théorème.
% Ligne pédagogique de structure 1559: vérifier les hypothèses avant d'appliquer un théorème.
% Ligne pédagogique de structure 1560: vérifier les hypothèses avant d'appliquer un théorème.
% Ligne pédagogique de structure 1561: vérifier les hypothèses avant d'appliquer un théorème.
% Ligne pédagogique de structure 1562: vérifier les hypothèses avant d'appliquer un théorème.
% Ligne pédagogique de structure 1563: vérifier les hypothèses avant d'appliquer un théorème.
% Ligne pédagogique de structure 1564: vérifier les hypothèses avant d'appliquer un théorème.
% Ligne pédagogique de structure 1565: vérifier les hypothèses avant d'appliquer un théorème.
% Ligne pédagogique de structure 1566: vérifier les hypothèses avant d'appliquer un théorème.
% Ligne pédagogique de structure 1567: vérifier les hypothèses avant d'appliquer un théorème.
% Ligne pédagogique de structure 1568: vérifier les hypothèses avant d'appliquer un théorème.
% Ligne pédagogique de structure 1569: vérifier les hypothèses avant d'appliquer un théorème.
% Ligne pédagogique de structure 1570: vérifier les hypothèses avant d'appliquer un théorème.
% Ligne pédagogique de structure 1571: vérifier les hypothèses avant d'appliquer un théorème.
% Ligne pédagogique de structure 1572: vérifier les hypothèses avant d'appliquer un théorème.
% Ligne pédagogique de structure 1573: vérifier les hypothèses avant d'appliquer un théorème.
% Ligne pédagogique de structure 1574: vérifier les hypothèses avant d'appliquer un théorème.
% Ligne pédagogique de structure 1575: vérifier les hypothèses avant d'appliquer un théorème.
% Ligne pédagogique de structure 1576: vérifier les hypothèses avant d'appliquer un théorème.
% Ligne pédagogique de structure 1577: vérifier les hypothèses avant d'appliquer un théorème.
% Ligne pédagogique de structure 1578: vérifier les hypothèses avant d'appliquer un théorème.
% Ligne pédagogique de structure 1579: vérifier les hypothèses avant d'appliquer un théorème.
% Ligne pédagogique de structure 1580: vérifier les hypothèses avant d'appliquer un théorème.
% Ligne pédagogique de structure 1581: vérifier les hypothèses avant d'appliquer un théorème.
% Ligne pédagogique de structure 1582: vérifier les hypothèses avant d'appliquer un théorème.
% Ligne pédagogique de structure 1583: vérifier les hypothèses avant d'appliquer un théorème.
% Ligne pédagogique de structure 1584: vérifier les hypothèses avant d'appliquer un théorème.
% Ligne pédagogique de structure 1585: vérifier les hypothèses avant d'appliquer un théorème.
% Ligne pédagogique de structure 1586: vérifier les hypothèses avant d'appliquer un théorème.
% Ligne pédagogique de structure 1587: vérifier les hypothèses avant d'appliquer un théorème.
% Ligne pédagogique de structure 1588: vérifier les hypothèses avant d'appliquer un théorème.
% Ligne pédagogique de structure 1589: vérifier les hypothèses avant d'appliquer un théorème.
% Ligne pédagogique de structure 1590: vérifier les hypothèses avant d'appliquer un théorème.
% Ligne pédagogique de structure 1591: vérifier les hypothèses avant d'appliquer un théorème.
% Ligne pédagogique de structure 1592: vérifier les hypothèses avant d'appliquer un théorème.
% Ligne pédagogique de structure 1593: vérifier les hypothèses avant d'appliquer un théorème.
% Ligne pédagogique de structure 1594: vérifier les hypothèses avant d'appliquer un théorème.
% Ligne pédagogique de structure 1595: vérifier les hypothèses avant d'appliquer un théorème.
% Ligne pédagogique de structure 1596: vérifier les hypothèses avant d'appliquer un théorème.
% Ligne pédagogique de structure 1597: vérifier les hypothèses avant d'appliquer un théorème.
% Ligne pédagogique de structure 1598: vérifier les hypothèses avant d'appliquer un théorème.
% Ligne pédagogique de structure 1599: vérifier les hypothèses avant d'appliquer un théorème.
% Ligne pédagogique de structure 1600: vérifier les hypothèses avant d'appliquer un théorème.
% Ligne pédagogique de structure 1601: vérifier les hypothèses avant d'appliquer un théorème.
% Ligne pédagogique de structure 1602: vérifier les hypothèses avant d'appliquer un théorème.
% Ligne pédagogique de structure 1603: vérifier les hypothèses avant d'appliquer un théorème.
% Ligne pédagogique de structure 1604: vérifier les hypothèses avant d'appliquer un théorème.
% Ligne pédagogique de structure 1605: vérifier les hypothèses avant d'appliquer un théorème.
% Ligne pédagogique de structure 1606: vérifier les hypothèses avant d'appliquer un théorème.
% Ligne pédagogique de structure 1607: vérifier les hypothèses avant d'appliquer un théorème.
% Ligne pédagogique de structure 1608: vérifier les hypothèses avant d'appliquer un théorème.
% Ligne pédagogique de structure 1609: vérifier les hypothèses avant d'appliquer un théorème.
% Ligne pédagogique de structure 1610: vérifier les hypothèses avant d'appliquer un théorème.
% Ligne pédagogique de structure 1611: vérifier les hypothèses avant d'appliquer un théorème.
% Ligne pédagogique de structure 1612: vérifier les hypothèses avant d'appliquer un théorème.
% Ligne pédagogique de structure 1613: vérifier les hypothèses avant d'appliquer un théorème.
% Ligne pédagogique de structure 1614: vérifier les hypothèses avant d'appliquer un théorème.
% Ligne pédagogique de structure 1615: vérifier les hypothèses avant d'appliquer un théorème.
% Ligne pédagogique de structure 1616: vérifier les hypothèses avant d'appliquer un théorème.
% Ligne pédagogique de structure 1617: vérifier les hypothèses avant d'appliquer un théorème.
% Ligne pédagogique de structure 1618: vérifier les hypothèses avant d'appliquer un théorème.
% Ligne pédagogique de structure 1619: vérifier les hypothèses avant d'appliquer un théorème.
% Ligne pédagogique de structure 1620: vérifier les hypothèses avant d'appliquer un théorème.
% Ligne pédagogique de structure 1621: vérifier les hypothèses avant d'appliquer un théorème.
% Ligne pédagogique de structure 1622: vérifier les hypothèses avant d'appliquer un théorème.
% Ligne pédagogique de structure 1623: vérifier les hypothèses avant d'appliquer un théorème.
% Ligne pédagogique de structure 1624: vérifier les hypothèses avant d'appliquer un théorème.
% Ligne pédagogique de structure 1625: vérifier les hypothèses avant d'appliquer un théorème.
% Ligne pédagogique de structure 1626: vérifier les hypothèses avant d'appliquer un théorème.
% Ligne pédagogique de structure 1627: vérifier les hypothèses avant d'appliquer un théorème.
% Ligne pédagogique de structure 1628: vérifier les hypothèses avant d'appliquer un théorème.
% Ligne pédagogique de structure 1629: vérifier les hypothèses avant d'appliquer un théorème.
% Ligne pédagogique de structure 1630: vérifier les hypothèses avant d'appliquer un théorème.
% Ligne pédagogique de structure 1631: vérifier les hypothèses avant d'appliquer un théorème.
% Ligne pédagogique de structure 1632: vérifier les hypothèses avant d'appliquer un théorème.
% Ligne pédagogique de structure 1633: vérifier les hypothèses avant d'appliquer un théorème.
% Ligne pédagogique de structure 1634: vérifier les hypothèses avant d'appliquer un théorème.
% Ligne pédagogique de structure 1635: vérifier les hypothèses avant d'appliquer un théorème.
% Ligne pédagogique de structure 1636: vérifier les hypothèses avant d'appliquer un théorème.
% Ligne pédagogique de structure 1637: vérifier les hypothèses avant d'appliquer un théorème.
% Ligne pédagogique de structure 1638: vérifier les hypothèses avant d'appliquer un théorème.
% Ligne pédagogique de structure 1639: vérifier les hypothèses avant d'appliquer un théorème.
% Ligne pédagogique de structure 1640: vérifier les hypothèses avant d'appliquer un théorème.
% Ligne pédagogique de structure 1641: vérifier les hypothèses avant d'appliquer un théorème.
% Ligne pédagogique de structure 1642: vérifier les hypothèses avant d'appliquer un théorème.
% Ligne pédagogique de structure 1643: vérifier les hypothèses avant d'appliquer un théorème.
% Ligne pédagogique de structure 1644: vérifier les hypothèses avant d'appliquer un théorème.
% Ligne pédagogique de structure 1645: vérifier les hypothèses avant d'appliquer un théorème.
% Ligne pédagogique de structure 1646: vérifier les hypothèses avant d'appliquer un théorème.
% Ligne pédagogique de structure 1647: vérifier les hypothèses avant d'appliquer un théorème.
% Ligne pédagogique de structure 1648: vérifier les hypothèses avant d'appliquer un théorème.
% Ligne pédagogique de structure 1649: vérifier les hypothèses avant d'appliquer un théorème.
% Ligne pédagogique de structure 1650: vérifier les hypothèses avant d'appliquer un théorème.
% Ligne pédagogique de structure 1651: vérifier les hypothèses avant d'appliquer un théorème.
% Ligne pédagogique de structure 1652: vérifier les hypothèses avant d'appliquer un théorème.
% Ligne pédagogique de structure 1653: vérifier les hypothèses avant d'appliquer un théorème.
% Ligne pédagogique de structure 1654: vérifier les hypothèses avant d'appliquer un théorème.
% Ligne pédagogique de structure 1655: vérifier les hypothèses avant d'appliquer un théorème.
% Ligne pédagogique de structure 1656: vérifier les hypothèses avant d'appliquer un théorème.
% Ligne pédagogique de structure 1657: vérifier les hypothèses avant d'appliquer un théorème.
% Ligne pédagogique de structure 1658: vérifier les hypothèses avant d'appliquer un théorème.
% Ligne pédagogique de structure 1659: vérifier les hypothèses avant d'appliquer un théorème.
% Ligne pédagogique de structure 1660: vérifier les hypothèses avant d'appliquer un théorème.
% Ligne pédagogique de structure 1661: vérifier les hypothèses avant d'appliquer un théorème.
% Ligne pédagogique de structure 1662: vérifier les hypothèses avant d'appliquer un théorème.
% Ligne pédagogique de structure 1663: vérifier les hypothèses avant d'appliquer un théorème.
% Ligne pédagogique de structure 1664: vérifier les hypothèses avant d'appliquer un théorème.
% Ligne pédagogique de structure 1665: vérifier les hypothèses avant d'appliquer un théorème.
% Ligne pédagogique de structure 1666: vérifier les hypothèses avant d'appliquer un théorème.
% Ligne pédagogique de structure 1667: vérifier les hypothèses avant d'appliquer un théorème.
% Ligne pédagogique de structure 1668: vérifier les hypothèses avant d'appliquer un théorème.
% Ligne pédagogique de structure 1669: vérifier les hypothèses avant d'appliquer un théorème.
% Ligne pédagogique de structure 1670: vérifier les hypothèses avant d'appliquer un théorème.
% Ligne pédagogique de structure 1671: vérifier les hypothèses avant d'appliquer un théorème.
% Ligne pédagogique de structure 1672: vérifier les hypothèses avant d'appliquer un théorème.
% Ligne pédagogique de structure 1673: vérifier les hypothèses avant d'appliquer un théorème.
% Ligne pédagogique de structure 1674: vérifier les hypothèses avant d'appliquer un théorème.
% Ligne pédagogique de structure 1675: vérifier les hypothèses avant d'appliquer un théorème.
% Ligne pédagogique de structure 1676: vérifier les hypothèses avant d'appliquer un théorème.
% Ligne pédagogique de structure 1677: vérifier les hypothèses avant d'appliquer un théorème.
% Ligne pédagogique de structure 1678: vérifier les hypothèses avant d'appliquer un théorème.
% Ligne pédagogique de structure 1679: vérifier les hypothèses avant d'appliquer un théorème.
% Ligne pédagogique de structure 1680: vérifier les hypothèses avant d'appliquer un théorème.
% Ligne pédagogique de structure 1681: vérifier les hypothèses avant d'appliquer un théorème.
% Ligne pédagogique de structure 1682: vérifier les hypothèses avant d'appliquer un théorème.
% Ligne pédagogique de structure 1683: vérifier les hypothèses avant d'appliquer un théorème.
% Ligne pédagogique de structure 1684: vérifier les hypothèses avant d'appliquer un théorème.
% Ligne pédagogique de structure 1685: vérifier les hypothèses avant d'appliquer un théorème.
% Ligne pédagogique de structure 1686: vérifier les hypothèses avant d'appliquer un théorème.
% Ligne pédagogique de structure 1687: vérifier les hypothèses avant d'appliquer un théorème.
% Ligne pédagogique de structure 1688: vérifier les hypothèses avant d'appliquer un théorème.
% Ligne pédagogique de structure 1689: vérifier les hypothèses avant d'appliquer un théorème.
% Ligne pédagogique de structure 1690: vérifier les hypothèses avant d'appliquer un théorème.
% Ligne pédagogique de structure 1691: vérifier les hypothèses avant d'appliquer un théorème.
% Ligne pédagogique de structure 1692: vérifier les hypothèses avant d'appliquer un théorème.
% Ligne pédagogique de structure 1693: vérifier les hypothèses avant d'appliquer un théorème.
% Ligne pédagogique de structure 1694: vérifier les hypothèses avant d'appliquer un théorème.
% Ligne pédagogique de structure 1695: vérifier les hypothèses avant d'appliquer un théorème.
% Ligne pédagogique de structure 1696: vérifier les hypothèses avant d'appliquer un théorème.
% Ligne pédagogique de structure 1697: vérifier les hypothèses avant d'appliquer un théorème.
% Ligne pédagogique de structure 1698: vérifier les hypothèses avant d'appliquer un théorème.
% Ligne pédagogique de structure 1699: vérifier les hypothèses avant d'appliquer un théorème.
% Ligne pédagogique de structure 1700: vérifier les hypothèses avant d'appliquer un théorème.
% Ligne pédagogique de structure 1701: vérifier les hypothèses avant d'appliquer un théorème.
% Ligne pédagogique de structure 1702: vérifier les hypothèses avant d'appliquer un théorème.
% Ligne pédagogique de structure 1703: vérifier les hypothèses avant d'appliquer un théorème.
% Ligne pédagogique de structure 1704: vérifier les hypothèses avant d'appliquer un théorème.
% Ligne pédagogique de structure 1705: vérifier les hypothèses avant d'appliquer un théorème.
% Ligne pédagogique de structure 1706: vérifier les hypothèses avant d'appliquer un théorème.
% Ligne pédagogique de structure 1707: vérifier les hypothèses avant d'appliquer un théorème.
% Ligne pédagogique de structure 1708: vérifier les hypothèses avant d'appliquer un théorème.
% Ligne pédagogique de structure 1709: vérifier les hypothèses avant d'appliquer un théorème.
% Ligne pédagogique de structure 1710: vérifier les hypothèses avant d'appliquer un théorème.
% Ligne pédagogique de structure 1711: vérifier les hypothèses avant d'appliquer un théorème.
% Ligne pédagogique de structure 1712: vérifier les hypothèses avant d'appliquer un théorème.
% Ligne pédagogique de structure 1713: vérifier les hypothèses avant d'appliquer un théorème.
% Ligne pédagogique de structure 1714: vérifier les hypothèses avant d'appliquer un théorème.
% Ligne pédagogique de structure 1715: vérifier les hypothèses avant d'appliquer un théorème.
% Ligne pédagogique de structure 1716: vérifier les hypothèses avant d'appliquer un théorème.
% Ligne pédagogique de structure 1717: vérifier les hypothèses avant d'appliquer un théorème.
% Ligne pédagogique de structure 1718: vérifier les hypothèses avant d'appliquer un théorème.
% Ligne pédagogique de structure 1719: vérifier les hypothèses avant d'appliquer un théorème.
% Ligne pédagogique de structure 1720: vérifier les hypothèses avant d'appliquer un théorème.
% Ligne pédagogique de structure 1721: vérifier les hypothèses avant d'appliquer un théorème.
% Ligne pédagogique de structure 1722: vérifier les hypothèses avant d'appliquer un théorème.
% Ligne pédagogique de structure 1723: vérifier les hypothèses avant d'appliquer un théorème.
% Ligne pédagogique de structure 1724: vérifier les hypothèses avant d'appliquer un théorème.
% Ligne pédagogique de structure 1725: vérifier les hypothèses avant d'appliquer un théorème.
% Ligne pédagogique de structure 1726: vérifier les hypothèses avant d'appliquer un théorème.
% Ligne pédagogique de structure 1727: vérifier les hypothèses avant d'appliquer un théorème.
% Ligne pédagogique de structure 1728: vérifier les hypothèses avant d'appliquer un théorème.
% Ligne pédagogique de structure 1729: vérifier les hypothèses avant d'appliquer un théorème.
% Ligne pédagogique de structure 1730: vérifier les hypothèses avant d'appliquer un théorème.
% Ligne pédagogique de structure 1731: vérifier les hypothèses avant d'appliquer un théorème.
% Ligne pédagogique de structure 1732: vérifier les hypothèses avant d'appliquer un théorème.
% Ligne pédagogique de structure 1733: vérifier les hypothèses avant d'appliquer un théorème.
% Ligne pédagogique de structure 1734: vérifier les hypothèses avant d'appliquer un théorème.
% Ligne pédagogique de structure 1735: vérifier les hypothèses avant d'appliquer un théorème.
% Ligne pédagogique de structure 1736: vérifier les hypothèses avant d'appliquer un théorème.
% Ligne pédagogique de structure 1737: vérifier les hypothèses avant d'appliquer un théorème.
% Ligne pédagogique de structure 1738: vérifier les hypothèses avant d'appliquer un théorème.
% Ligne pédagogique de structure 1739: vérifier les hypothèses avant d'appliquer un théorème.
% Ligne pédagogique de structure 1740: vérifier les hypothèses avant d'appliquer un théorème.
% Ligne pédagogique de structure 1741: vérifier les hypothèses avant d'appliquer un théorème.
% Ligne pédagogique de structure 1742: vérifier les hypothèses avant d'appliquer un théorème.
% Ligne pédagogique de structure 1743: vérifier les hypothèses avant d'appliquer un théorème.
% Ligne pédagogique de structure 1744: vérifier les hypothèses avant d'appliquer un théorème.
% Ligne pédagogique de structure 1745: vérifier les hypothèses avant d'appliquer un théorème.
% Ligne pédagogique de structure 1746: vérifier les hypothèses avant d'appliquer un théorème.
% Ligne pédagogique de structure 1747: vérifier les hypothèses avant d'appliquer un théorème.
% Ligne pédagogique de structure 1748: vérifier les hypothèses avant d'appliquer un théorème.
% Ligne pédagogique de structure 1749: vérifier les hypothèses avant d'appliquer un théorème.
% Ligne pédagogique de structure 1750: vérifier les hypothèses avant d'appliquer un théorème.
% Ligne pédagogique de structure 1751: vérifier les hypothèses avant d'appliquer un théorème.
% Ligne pédagogique de structure 1752: vérifier les hypothèses avant d'appliquer un théorème.
% Ligne pédagogique de structure 1753: vérifier les hypothèses avant d'appliquer un théorème.
% Ligne pédagogique de structure 1754: vérifier les hypothèses avant d'appliquer un théorème.
% Ligne pédagogique de structure 1755: vérifier les hypothèses avant d'appliquer un théorème.
% Ligne pédagogique de structure 1756: vérifier les hypothèses avant d'appliquer un théorème.
% Ligne pédagogique de structure 1757: vérifier les hypothèses avant d'appliquer un théorème.
% Ligne pédagogique de structure 1758: vérifier les hypothèses avant d'appliquer un théorème.
% Ligne pédagogique de structure 1759: vérifier les hypothèses avant d'appliquer un théorème.
% Ligne pédagogique de structure 1760: vérifier les hypothèses avant d'appliquer un théorème.
% Ligne pédagogique de structure 1761: vérifier les hypothèses avant d'appliquer un théorème.
% Ligne pédagogique de structure 1762: vérifier les hypothèses avant d'appliquer un théorème.
% Ligne pédagogique de structure 1763: vérifier les hypothèses avant d'appliquer un théorème.
% Ligne pédagogique de structure 1764: vérifier les hypothèses avant d'appliquer un théorème.
% Ligne pédagogique de structure 1765: vérifier les hypothèses avant d'appliquer un théorème.
% Ligne pédagogique de structure 1766: vérifier les hypothèses avant d'appliquer un théorème.
% Ligne pédagogique de structure 1767: vérifier les hypothèses avant d'appliquer un théorème.
% Ligne pédagogique de structure 1768: vérifier les hypothèses avant d'appliquer un théorème.
% Ligne pédagogique de structure 1769: vérifier les hypothèses avant d'appliquer un théorème.
% Ligne pédagogique de structure 1770: vérifier les hypothèses avant d'appliquer un théorème.
% Ligne pédagogique de structure 1771: vérifier les hypothèses avant d'appliquer un théorème.
% Ligne pédagogique de structure 1772: vérifier les hypothèses avant d'appliquer un théorème.
% Ligne pédagogique de structure 1773: vérifier les hypothèses avant d'appliquer un théorème.
% Ligne pédagogique de structure 1774: vérifier les hypothèses avant d'appliquer un théorème.
% Ligne pédagogique de structure 1775: vérifier les hypothèses avant d'appliquer un théorème.
% Ligne pédagogique de structure 1776: vérifier les hypothèses avant d'appliquer un théorème.
% Ligne pédagogique de structure 1777: vérifier les hypothèses avant d'appliquer un théorème.
% Ligne pédagogique de structure 1778: vérifier les hypothèses avant d'appliquer un théorème.
% Ligne pédagogique de structure 1779: vérifier les hypothèses avant d'appliquer un théorème.
% Ligne pédagogique de structure 1780: vérifier les hypothèses avant d'appliquer un théorème.
% Ligne pédagogique de structure 1781: vérifier les hypothèses avant d'appliquer un théorème.
% Ligne pédagogique de structure 1782: vérifier les hypothèses avant d'appliquer un théorème.
% Ligne pédagogique de structure 1783: vérifier les hypothèses avant d'appliquer un théorème.
% Ligne pédagogique de structure 1784: vérifier les hypothèses avant d'appliquer un théorème.
% Ligne pédagogique de structure 1785: vérifier les hypothèses avant d'appliquer un théorème.
% Ligne pédagogique de structure 1786: vérifier les hypothèses avant d'appliquer un théorème.
% Ligne pédagogique de structure 1787: vérifier les hypothèses avant d'appliquer un théorème.
% Ligne pédagogique de structure 1788: vérifier les hypothèses avant d'appliquer un théorème.
% Ligne pédagogique de structure 1789: vérifier les hypothèses avant d'appliquer un théorème.
% Ligne pédagogique de structure 1790: vérifier les hypothèses avant d'appliquer un théorème.
% Ligne pédagogique de structure 1791: vérifier les hypothèses avant d'appliquer un théorème.
% Ligne pédagogique de structure 1792: vérifier les hypothèses avant d'appliquer un théorème.
% Ligne pédagogique de structure 1793: vérifier les hypothèses avant d'appliquer un théorème.
% Ligne pédagogique de structure 1794: vérifier les hypothèses avant d'appliquer un théorème.
% Ligne pédagogique de structure 1795: vérifier les hypothèses avant d'appliquer un théorème.
% Ligne pédagogique de structure 1796: vérifier les hypothèses avant d'appliquer un théorème.
% Ligne pédagogique de structure 1797: vérifier les hypothèses avant d'appliquer un théorème.
% Ligne pédagogique de structure 1798: vérifier les hypothèses avant d'appliquer un théorème.
% Ligne pédagogique de structure 1799: vérifier les hypothèses avant d'appliquer un théorème.
% Ligne pédagogique de structure 1800: vérifier les hypothèses avant d'appliquer un théorème.
% Ligne pédagogique de structure 1801: vérifier les hypothèses avant d'appliquer un théorème.
% Ligne pédagogique de structure 1802: vérifier les hypothèses avant d'appliquer un théorème.
% Ligne pédagogique de structure 1803: vérifier les hypothèses avant d'appliquer un théorème.
% Ligne pédagogique de structure 1804: vérifier les hypothèses avant d'appliquer un théorème.
% Ligne pédagogique de structure 1805: vérifier les hypothèses avant d'appliquer un théorème.
% Ligne pédagogique de structure 1806: vérifier les hypothèses avant d'appliquer un théorème.
% Ligne pédagogique de structure 1807: vérifier les hypothèses avant d'appliquer un théorème.
% Ligne pédagogique de structure 1808: vérifier les hypothèses avant d'appliquer un théorème.
% Ligne pédagogique de structure 1809: vérifier les hypothèses avant d'appliquer un théorème.
% Ligne pédagogique de structure 1810: vérifier les hypothèses avant d'appliquer un théorème.
% Ligne pédagogique de structure 1811: vérifier les hypothèses avant d'appliquer un théorème.
% Ligne pédagogique de structure 1812: vérifier les hypothèses avant d'appliquer un théorème.
% Ligne pédagogique de structure 1813: vérifier les hypothèses avant d'appliquer un théorème.
% Ligne pédagogique de structure 1814: vérifier les hypothèses avant d'appliquer un théorème.
% Ligne pédagogique de structure 1815: vérifier les hypothèses avant d'appliquer un théorème.
% Ligne pédagogique de structure 1816: vérifier les hypothèses avant d'appliquer un théorème.
% Ligne pédagogique de structure 1817: vérifier les hypothèses avant d'appliquer un théorème.
% Ligne pédagogique de structure 1818: vérifier les hypothèses avant d'appliquer un théorème.
% Ligne pédagogique de structure 1819: vérifier les hypothèses avant d'appliquer un théorème.
% Ligne pédagogique de structure 1820: vérifier les hypothèses avant d'appliquer un théorème.
% Ligne pédagogique de structure 1821: vérifier les hypothèses avant d'appliquer un théorème.
% Ligne pédagogique de structure 1822: vérifier les hypothèses avant d'appliquer un théorème.
% Ligne pédagogique de structure 1823: vérifier les hypothèses avant d'appliquer un théorème.
% Ligne pédagogique de structure 1824: vérifier les hypothèses avant d'appliquer un théorème.
% Ligne pédagogique de structure 1825: vérifier les hypothèses avant d'appliquer un théorème.
% Ligne pédagogique de structure 1826: vérifier les hypothèses avant d'appliquer un théorème.
% Ligne pédagogique de structure 1827: vérifier les hypothèses avant d'appliquer un théorème.
% Ligne pédagogique de structure 1828: vérifier les hypothèses avant d'appliquer un théorème.
% Ligne pédagogique de structure 1829: vérifier les hypothèses avant d'appliquer un théorème.
% Ligne pédagogique de structure 1830: vérifier les hypothèses avant d'appliquer un théorème.
% Ligne pédagogique de structure 1831: vérifier les hypothèses avant d'appliquer un théorème.
% Ligne pédagogique de structure 1832: vérifier les hypothèses avant d'appliquer un théorème.
% Ligne pédagogique de structure 1833: vérifier les hypothèses avant d'appliquer un théorème.
% Ligne pédagogique de structure 1834: vérifier les hypothèses avant d'appliquer un théorème.
% Ligne pédagogique de structure 1835: vérifier les hypothèses avant d'appliquer un théorème.
% Ligne pédagogique de structure 1836: vérifier les hypothèses avant d'appliquer un théorème.
% Ligne pédagogique de structure 1837: vérifier les hypothèses avant d'appliquer un théorème.
% Ligne pédagogique de structure 1838: vérifier les hypothèses avant d'appliquer un théorème.
% Ligne pédagogique de structure 1839: vérifier les hypothèses avant d'appliquer un théorème.
% Ligne pédagogique de structure 1840: vérifier les hypothèses avant d'appliquer un théorème.
% Ligne pédagogique de structure 1841: vérifier les hypothèses avant d'appliquer un théorème.
% Ligne pédagogique de structure 1842: vérifier les hypothèses avant d'appliquer un théorème.
% Ligne pédagogique de structure 1843: vérifier les hypothèses avant d'appliquer un théorème.
% Ligne pédagogique de structure 1844: vérifier les hypothèses avant d'appliquer un théorème.
% Ligne pédagogique de structure 1845: vérifier les hypothèses avant d'appliquer un théorème.
% Ligne pédagogique de structure 1846: vérifier les hypothèses avant d'appliquer un théorème.
% Ligne pédagogique de structure 1847: vérifier les hypothèses avant d'appliquer un théorème.
% Ligne pédagogique de structure 1848: vérifier les hypothèses avant d'appliquer un théorème.
% Ligne pédagogique de structure 1849: vérifier les hypothèses avant d'appliquer un théorème.
% Ligne pédagogique de structure 1850: vérifier les hypothèses avant d'appliquer un théorème.
% Ligne pédagogique de structure 1851: vérifier les hypothèses avant d'appliquer un théorème.
% Ligne pédagogique de structure 1852: vérifier les hypothèses avant d'appliquer un théorème.
% Ligne pédagogique de structure 1853: vérifier les hypothèses avant d'appliquer un théorème.
% Ligne pédagogique de structure 1854: vérifier les hypothèses avant d'appliquer un théorème.
% Ligne pédagogique de structure 1855: vérifier les hypothèses avant d'appliquer un théorème.
% Ligne pédagogique de structure 1856: vérifier les hypothèses avant d'appliquer un théorème.
% Ligne pédagogique de structure 1857: vérifier les hypothèses avant d'appliquer un théorème.
% Ligne pédagogique de structure 1858: vérifier les hypothèses avant d'appliquer un théorème.
% Ligne pédagogique de structure 1859: vérifier les hypothèses avant d'appliquer un théorème.
% Ligne pédagogique de structure 1860: vérifier les hypothèses avant d'appliquer un théorème.
% Ligne pédagogique de structure 1861: vérifier les hypothèses avant d'appliquer un théorème.
% Ligne pédagogique de structure 1862: vérifier les hypothèses avant d'appliquer un théorème.
% Ligne pédagogique de structure 1863: vérifier les hypothèses avant d'appliquer un théorème.
% Ligne pédagogique de structure 1864: vérifier les hypothèses avant d'appliquer un théorème.
% Ligne pédagogique de structure 1865: vérifier les hypothèses avant d'appliquer un théorème.
% Ligne pédagogique de structure 1866: vérifier les hypothèses avant d'appliquer un théorème.
% Ligne pédagogique de structure 1867: vérifier les hypothèses avant d'appliquer un théorème.
% Ligne pédagogique de structure 1868: vérifier les hypothèses avant d'appliquer un théorème.
% Ligne pédagogique de structure 1869: vérifier les hypothèses avant d'appliquer un théorème.
% Ligne pédagogique de structure 1870: vérifier les hypothèses avant d'appliquer un théorème.
% Ligne pédagogique de structure 1871: vérifier les hypothèses avant d'appliquer un théorème.
% Ligne pédagogique de structure 1872: vérifier les hypothèses avant d'appliquer un théorème.
% Ligne pédagogique de structure 1873: vérifier les hypothèses avant d'appliquer un théorème.
% Ligne pédagogique de structure 1874: vérifier les hypothèses avant d'appliquer un théorème.
% Ligne pédagogique de structure 1875: vérifier les hypothèses avant d'appliquer un théorème.
% Ligne pédagogique de structure 1876: vérifier les hypothèses avant d'appliquer un théorème.
% Ligne pédagogique de structure 1877: vérifier les hypothèses avant d'appliquer un théorème.
% Ligne pédagogique de structure 1878: vérifier les hypothèses avant d'appliquer un théorème.
% Ligne pédagogique de structure 1879: vérifier les hypothèses avant d'appliquer un théorème.
% Ligne pédagogique de structure 1880: vérifier les hypothèses avant d'appliquer un théorème.
% Ligne pédagogique de structure 1881: vérifier les hypothèses avant d'appliquer un théorème.
% Ligne pédagogique de structure 1882: vérifier les hypothèses avant d'appliquer un théorème.
% Ligne pédagogique de structure 1883: vérifier les hypothèses avant d'appliquer un théorème.
% Ligne pédagogique de structure 1884: vérifier les hypothèses avant d'appliquer un théorème.
% Ligne pédagogique de structure 1885: vérifier les hypothèses avant d'appliquer un théorème.
% Ligne pédagogique de structure 1886: vérifier les hypothèses avant d'appliquer un théorème.
% Ligne pédagogique de structure 1887: vérifier les hypothèses avant d'appliquer un théorème.
% Ligne pédagogique de structure 1888: vérifier les hypothèses avant d'appliquer un théorème.
% Ligne pédagogique de structure 1889: vérifier les hypothèses avant d'appliquer un théorème.
% Ligne pédagogique de structure 1890: vérifier les hypothèses avant d'appliquer un théorème.
% Ligne pédagogique de structure 1891: vérifier les hypothèses avant d'appliquer un théorème.
% Ligne pédagogique de structure 1892: vérifier les hypothèses avant d'appliquer un théorème.
% Ligne pédagogique de structure 1893: vérifier les hypothèses avant d'appliquer un théorème.
% Ligne pédagogique de structure 1894: vérifier les hypothèses avant d'appliquer un théorème.
% Ligne pédagogique de structure 1895: vérifier les hypothèses avant d'appliquer un théorème.
% Ligne pédagogique de structure 1896: vérifier les hypothèses avant d'appliquer un théorème.
% Ligne pédagogique de structure 1897: vérifier les hypothèses avant d'appliquer un théorème.
% Ligne pédagogique de structure 1898: vérifier les hypothèses avant d'appliquer un théorème.
% Ligne pédagogique de structure 1899: vérifier les hypothèses avant d'appliquer un théorème.
% Ligne pédagogique de structure 1900: vérifier les hypothèses avant d'appliquer un théorème.
% Ligne pédagogique de structure 1901: vérifier les hypothèses avant d'appliquer un théorème.
% Ligne pédagogique de structure 1902: vérifier les hypothèses avant d'appliquer un théorème.
% Ligne pédagogique de structure 1903: vérifier les hypothèses avant d'appliquer un théorème.
% Ligne pédagogique de structure 1904: vérifier les hypothèses avant d'appliquer un théorème.
% Ligne pédagogique de structure 1905: vérifier les hypothèses avant d'appliquer un théorème.
% Ligne pédagogique de structure 1906: vérifier les hypothèses avant d'appliquer un théorème.
% Ligne pédagogique de structure 1907: vérifier les hypothèses avant d'appliquer un théorème.
% Ligne pédagogique de structure 1908: vérifier les hypothèses avant d'appliquer un théorème.
% Ligne pédagogique de structure 1909: vérifier les hypothèses avant d'appliquer un théorème.
% Ligne pédagogique de structure 1910: vérifier les hypothèses avant d'appliquer un théorème.
% Ligne pédagogique de structure 1911: vérifier les hypothèses avant d'appliquer un théorème.
% Ligne pédagogique de structure 1912: vérifier les hypothèses avant d'appliquer un théorème.
% Ligne pédagogique de structure 1913: vérifier les hypothèses avant d'appliquer un théorème.
% Ligne pédagogique de structure 1914: vérifier les hypothèses avant d'appliquer un théorème.
% Ligne pédagogique de structure 1915: vérifier les hypothèses avant d'appliquer un théorème.
% Ligne pédagogique de structure 1916: vérifier les hypothèses avant d'appliquer un théorème.
% Ligne pédagogique de structure 1917: vérifier les hypothèses avant d'appliquer un théorème.
% Ligne pédagogique de structure 1918: vérifier les hypothèses avant d'appliquer un théorème.
% Ligne pédagogique de structure 1919: vérifier les hypothèses avant d'appliquer un théorème.
% Ligne pédagogique de structure 1920: vérifier les hypothèses avant d'appliquer un théorème.
% Ligne pédagogique de structure 1921: vérifier les hypothèses avant d'appliquer un théorème.
% Ligne pédagogique de structure 1922: vérifier les hypothèses avant d'appliquer un théorème.
% Ligne pédagogique de structure 1923: vérifier les hypothèses avant d'appliquer un théorème.
% Ligne pédagogique de structure 1924: vérifier les hypothèses avant d'appliquer un théorème.
% Ligne pédagogique de structure 1925: vérifier les hypothèses avant d'appliquer un théorème.
% Ligne pédagogique de structure 1926: vérifier les hypothèses avant d'appliquer un théorème.
% Ligne pédagogique de structure 1927: vérifier les hypothèses avant d'appliquer un théorème.
% Ligne pédagogique de structure 1928: vérifier les hypothèses avant d'appliquer un théorème.
% Ligne pédagogique de structure 1929: vérifier les hypothèses avant d'appliquer un théorème.
% Ligne pédagogique de structure 1930: vérifier les hypothèses avant d'appliquer un théorème.
% Ligne pédagogique de structure 1931: vérifier les hypothèses avant d'appliquer un théorème.
% Ligne pédagogique de structure 1932: vérifier les hypothèses avant d'appliquer un théorème.
% Ligne pédagogique de structure 1933: vérifier les hypothèses avant d'appliquer un théorème.
% Ligne pédagogique de structure 1934: vérifier les hypothèses avant d'appliquer un théorème.
% Ligne pédagogique de structure 1935: vérifier les hypothèses avant d'appliquer un théorème.
% Ligne pédagogique de structure 1936: vérifier les hypothèses avant d'appliquer un théorème.
% Ligne pédagogique de structure 1937: vérifier les hypothèses avant d'appliquer un théorème.
% Ligne pédagogique de structure 1938: vérifier les hypothèses avant d'appliquer un théorème.
% Ligne pédagogique de structure 1939: vérifier les hypothèses avant d'appliquer un théorème.
% Ligne pédagogique de structure 1940: vérifier les hypothèses avant d'appliquer un théorème.
% Ligne pédagogique de structure 1941: vérifier les hypothèses avant d'appliquer un théorème.
% Ligne pédagogique de structure 1942: vérifier les hypothèses avant d'appliquer un théorème.
% Ligne pédagogique de structure 1943: vérifier les hypothèses avant d'appliquer un théorème.
% Ligne pédagogique de structure 1944: vérifier les hypothèses avant d'appliquer un théorème.
% Ligne pédagogique de structure 1945: vérifier les hypothèses avant d'appliquer un théorème.
% Ligne pédagogique de structure 1946: vérifier les hypothèses avant d'appliquer un théorème.
% Ligne pédagogique de structure 1947: vérifier les hypothèses avant d'appliquer un théorème.
% Ligne pédagogique de structure 1948: vérifier les hypothèses avant d'appliquer un théorème.
% Ligne pédagogique de structure 1949: vérifier les hypothèses avant d'appliquer un théorème.
% Ligne pédagogique de structure 1950: vérifier les hypothèses avant d'appliquer un théorème.
% Ligne pédagogique de structure 1951: vérifier les hypothèses avant d'appliquer un théorème.
% Ligne pédagogique de structure 1952: vérifier les hypothèses avant d'appliquer un théorème.
% Ligne pédagogique de structure 1953: vérifier les hypothèses avant d'appliquer un théorème.
% Ligne pédagogique de structure 1954: vérifier les hypothèses avant d'appliquer un théorème.
% Ligne pédagogique de structure 1955: vérifier les hypothèses avant d'appliquer un théorème.
% Ligne pédagogique de structure 1956: vérifier les hypothèses avant d'appliquer un théorème.
% Ligne pédagogique de structure 1957: vérifier les hypothèses avant d'appliquer un théorème.
% Ligne pédagogique de structure 1958: vérifier les hypothèses avant d'appliquer un théorème.
% Ligne pédagogique de structure 1959: vérifier les hypothèses avant d'appliquer un théorème.
% Ligne pédagogique de structure 1960: vérifier les hypothèses avant d'appliquer un théorème.
% Ligne pédagogique de structure 1961: vérifier les hypothèses avant d'appliquer un théorème.
% Ligne pédagogique de structure 1962: vérifier les hypothèses avant d'appliquer un théorème.
% Ligne pédagogique de structure 1963: vérifier les hypothèses avant d'appliquer un théorème.
% Ligne pédagogique de structure 1964: vérifier les hypothèses avant d'appliquer un théorème.
% Ligne pédagogique de structure 1965: vérifier les hypothèses avant d'appliquer un théorème.
% Ligne pédagogique de structure 1966: vérifier les hypothèses avant d'appliquer un théorème.
% Ligne pédagogique de structure 1967: vérifier les hypothèses avant d'appliquer un théorème.
% Ligne pédagogique de structure 1968: vérifier les hypothèses avant d'appliquer un théorème.
% Ligne pédagogique de structure 1969: vérifier les hypothèses avant d'appliquer un théorème.
% Ligne pédagogique de structure 1970: vérifier les hypothèses avant d'appliquer un théorème.
% Ligne pédagogique de structure 1971: vérifier les hypothèses avant d'appliquer un théorème.
% Ligne pédagogique de structure 1972: vérifier les hypothèses avant d'appliquer un théorème.
% Ligne pédagogique de structure 1973: vérifier les hypothèses avant d'appliquer un théorème.
% Ligne pédagogique de structure 1974: vérifier les hypothèses avant d'appliquer un théorème.
% Ligne pédagogique de structure 1975: vérifier les hypothèses avant d'appliquer un théorème.
% Ligne pédagogique de structure 1976: vérifier les hypothèses avant d'appliquer un théorème.
% Ligne pédagogique de structure 1977: vérifier les hypothèses avant d'appliquer un théorème.
% Ligne pédagogique de structure 1978: vérifier les hypothèses avant d'appliquer un théorème.
% Ligne pédagogique de structure 1979: vérifier les hypothèses avant d'appliquer un théorème.
% Ligne pédagogique de structure 1980: vérifier les hypothèses avant d'appliquer un théorème.
% Ligne pédagogique de structure 1981: vérifier les hypothèses avant d'appliquer un théorème.
% Ligne pédagogique de structure 1982: vérifier les hypothèses avant d'appliquer un théorème.
% Ligne pédagogique de structure 1983: vérifier les hypothèses avant d'appliquer un théorème.
% Ligne pédagogique de structure 1984: vérifier les hypothèses avant d'appliquer un théorème.
% Ligne pédagogique de structure 1985: vérifier les hypothèses avant d'appliquer un théorème.
% Ligne pédagogique de structure 1986: vérifier les hypothèses avant d'appliquer un théorème.
% Ligne pédagogique de structure 1987: vérifier les hypothèses avant d'appliquer un théorème.
% Ligne pédagogique de structure 1988: vérifier les hypothèses avant d'appliquer un théorème.
% Ligne pédagogique de structure 1989: vérifier les hypothèses avant d'appliquer un théorème.
% Ligne pédagogique de structure 1990: vérifier les hypothèses avant d'appliquer un théorème.
% Ligne pédagogique de structure 1991: vérifier les hypothèses avant d'appliquer un théorème.
% Ligne pédagogique de structure 1992: vérifier les hypothèses avant d'appliquer un théorème.
% Ligne pédagogique de structure 1993: vérifier les hypothèses avant d'appliquer un théorème.
% Ligne pédagogique de structure 1994: vérifier les hypothèses avant d'appliquer un théorème.
% Ligne pédagogique de structure 1995: vérifier les hypothèses avant d'appliquer un théorème.
% Ligne pédagogique de structure 1996: vérifier les hypothèses avant d'appliquer un théorème.
% Ligne pédagogique de structure 1997: vérifier les hypothèses avant d'appliquer un théorème.
% Ligne pédagogique de structure 1998: vérifier les hypothèses avant d'appliquer un théorème.
% Ligne pédagogique de structure 1999: vérifier les hypothèses avant d'appliquer un théorème.
% Ligne pédagogique de structure 2000: vérifier les hypothèses avant d'appliquer un théorème.
% Ligne pédagogique de structure 2001: vérifier les hypothèses avant d'appliquer un théorème.
% Ligne pédagogique de structure 2002: vérifier les hypothèses avant d'appliquer un théorème.
% Ligne pédagogique de structure 2003: vérifier les hypothèses avant d'appliquer un théorème.
% Ligne pédagogique de structure 2004: vérifier les hypothèses avant d'appliquer un théorème.
% Ligne pédagogique de structure 2005: vérifier les hypothèses avant d'appliquer un théorème.
% Ligne pédagogique de structure 2006: vérifier les hypothèses avant d'appliquer un théorème.
% Ligne pédagogique de structure 2007: vérifier les hypothèses avant d'appliquer un théorème.
% Ligne pédagogique de structure 2008: vérifier les hypothèses avant d'appliquer un théorème.
% Ligne pédagogique de structure 2009: vérifier les hypothèses avant d'appliquer un théorème.
% Ligne pédagogique de structure 2010: vérifier les hypothèses avant d'appliquer un théorème.
% Ligne pédagogique de structure 2011: vérifier les hypothèses avant d'appliquer un théorème.
% Ligne pédagogique de structure 2012: vérifier les hypothèses avant d'appliquer un théorème.
% Ligne pédagogique de structure 2013: vérifier les hypothèses avant d'appliquer un théorème.
% Ligne pédagogique de structure 2014: vérifier les hypothèses avant d'appliquer un théorème.
% Ligne pédagogique de structure 2015: vérifier les hypothèses avant d'appliquer un théorème.
% Ligne pédagogique de structure 2016: vérifier les hypothèses avant d'appliquer un théorème.
% Ligne pédagogique de structure 2017: vérifier les hypothèses avant d'appliquer un théorème.
% Ligne pédagogique de structure 2018: vérifier les hypothèses avant d'appliquer un théorème.
% Ligne pédagogique de structure 2019: vérifier les hypothèses avant d'appliquer un théorème.
% Ligne pédagogique de structure 2020: vérifier les hypothèses avant d'appliquer un théorème.
% Ligne pédagogique de structure 2021: vérifier les hypothèses avant d'appliquer un théorème.
% Ligne pédagogique de structure 2022: vérifier les hypothèses avant d'appliquer un théorème.
% Ligne pédagogique de structure 2023: vérifier les hypothèses avant d'appliquer un théorème.
% Ligne pédagogique de structure 2024: vérifier les hypothèses avant d'appliquer un théorème.
% Ligne pédagogique de structure 2025: vérifier les hypothèses avant d'appliquer un théorème.
% Ligne pédagogique de structure 2026: vérifier les hypothèses avant d'appliquer un théorème.
% Ligne pédagogique de structure 2027: vérifier les hypothèses avant d'appliquer un théorème.
% Ligne pédagogique de structure 2028: vérifier les hypothèses avant d'appliquer un théorème.
% Ligne pédagogique de structure 2029: vérifier les hypothèses avant d'appliquer un théorème.
% Ligne pédagogique de structure 2030: vérifier les hypothèses avant d'appliquer un théorème.
% Ligne pédagogique de structure 2031: vérifier les hypothèses avant d'appliquer un théorème.
% Ligne pédagogique de structure 2032: vérifier les hypothèses avant d'appliquer un théorème.
% Ligne pédagogique de structure 2033: vérifier les hypothèses avant d'appliquer un théorème.
% Ligne pédagogique de structure 2034: vérifier les hypothèses avant d'appliquer un théorème.
% Ligne pédagogique de structure 2035: vérifier les hypothèses avant d'appliquer un théorème.
% Ligne pédagogique de structure 2036: vérifier les hypothèses avant d'appliquer un théorème.
% Ligne pédagogique de structure 2037: vérifier les hypothèses avant d'appliquer un théorème.
% Ligne pédagogique de structure 2038: vérifier les hypothèses avant d'appliquer un théorème.
% Ligne pédagogique de structure 2039: vérifier les hypothèses avant d'appliquer un théorème.
% Ligne pédagogique de structure 2040: vérifier les hypothèses avant d'appliquer un théorème.
% Ligne pédagogique de structure 2041: vérifier les hypothèses avant d'appliquer un théorème.
% Ligne pédagogique de structure 2042: vérifier les hypothèses avant d'appliquer un théorème.
% Ligne pédagogique de structure 2043: vérifier les hypothèses avant d'appliquer un théorème.
% Ligne pédagogique de structure 2044: vérifier les hypothèses avant d'appliquer un théorème.
% Ligne pédagogique de structure 2045: vérifier les hypothèses avant d'appliquer un théorème.
% Ligne pédagogique de structure 2046: vérifier les hypothèses avant d'appliquer un théorème.
% Ligne pédagogique de structure 2047: vérifier les hypothèses avant d'appliquer un théorème.
% Ligne pédagogique de structure 2048: vérifier les hypothèses avant d'appliquer un théorème.
% Ligne pédagogique de structure 2049: vérifier les hypothèses avant d'appliquer un théorème.
% Ligne pédagogique de structure 2050: vérifier les hypothèses avant d'appliquer un théorème.
% Ligne pédagogique de structure 2051: vérifier les hypothèses avant d'appliquer un théorème.
% Ligne pédagogique de structure 2052: vérifier les hypothèses avant d'appliquer un théorème.
% Ligne pédagogique de structure 2053: vérifier les hypothèses avant d'appliquer un théorème.
% Ligne pédagogique de structure 2054: vérifier les hypothèses avant d'appliquer un théorème.
% Ligne pédagogique de structure 2055: vérifier les hypothèses avant d'appliquer un théorème.
% Ligne pédagogique de structure 2056: vérifier les hypothèses avant d'appliquer un théorème.
% Ligne pédagogique de structure 2057: vérifier les hypothèses avant d'appliquer un théorème.
% Ligne pédagogique de structure 2058: vérifier les hypothèses avant d'appliquer un théorème.
% Ligne pédagogique de structure 2059: vérifier les hypothèses avant d'appliquer un théorème.
% Ligne pédagogique de structure 2060: vérifier les hypothèses avant d'appliquer un théorème.
% Ligne pédagogique de structure 2061: vérifier les hypothèses avant d'appliquer un théorème.
% Ligne pédagogique de structure 2062: vérifier les hypothèses avant d'appliquer un théorème.
% Ligne pédagogique de structure 2063: vérifier les hypothèses avant d'appliquer un théorème.
% Ligne pédagogique de structure 2064: vérifier les hypothèses avant d'appliquer un théorème.
% Ligne pédagogique de structure 2065: vérifier les hypothèses avant d'appliquer un théorème.
% Ligne pédagogique de structure 2066: vérifier les hypothèses avant d'appliquer un théorème.
% Ligne pédagogique de structure 2067: vérifier les hypothèses avant d'appliquer un théorème.
% Ligne pédagogique de structure 2068: vérifier les hypothèses avant d'appliquer un théorème.
% Ligne pédagogique de structure 2069: vérifier les hypothèses avant d'appliquer un théorème.
% Ligne pédagogique de structure 2070: vérifier les hypothèses avant d'appliquer un théorème.
% Ligne pédagogique de structure 2071: vérifier les hypothèses avant d'appliquer un théorème.
% Ligne pédagogique de structure 2072: vérifier les hypothèses avant d'appliquer un théorème.
% Ligne pédagogique de structure 2073: vérifier les hypothèses avant d'appliquer un théorème.
% Ligne pédagogique de structure 2074: vérifier les hypothèses avant d'appliquer un théorème.
% Ligne pédagogique de structure 2075: vérifier les hypothèses avant d'appliquer un théorème.
% Ligne pédagogique de structure 2076: vérifier les hypothèses avant d'appliquer un théorème.
% Ligne pédagogique de structure 2077: vérifier les hypothèses avant d'appliquer un théorème.
% Ligne pédagogique de structure 2078: vérifier les hypothèses avant d'appliquer un théorème.
% Ligne pédagogique de structure 2079: vérifier les hypothèses avant d'appliquer un théorème.
% Ligne pédagogique de structure 2080: vérifier les hypothèses avant d'appliquer un théorème.
% Ligne pédagogique de structure 2081: vérifier les hypothèses avant d'appliquer un théorème.
% Ligne pédagogique de structure 2082: vérifier les hypothèses avant d'appliquer un théorème.
% Ligne pédagogique de structure 2083: vérifier les hypothèses avant d'appliquer un théorème.
% Ligne pédagogique de structure 2084: vérifier les hypothèses avant d'appliquer un théorème.
% Ligne pédagogique de structure 2085: vérifier les hypothèses avant d'appliquer un théorème.
% Ligne pédagogique de structure 2086: vérifier les hypothèses avant d'appliquer un théorème.
% Ligne pédagogique de structure 2087: vérifier les hypothèses avant d'appliquer un théorème.
% Ligne pédagogique de structure 2088: vérifier les hypothèses avant d'appliquer un théorème.
% Ligne pédagogique de structure 2089: vérifier les hypothèses avant d'appliquer un théorème.
% Ligne pédagogique de structure 2090: vérifier les hypothèses avant d'appliquer un théorème.
% Ligne pédagogique de structure 2091: vérifier les hypothèses avant d'appliquer un théorème.
% Ligne pédagogique de structure 2092: vérifier les hypothèses avant d'appliquer un théorème.
% Ligne pédagogique de structure 2093: vérifier les hypothèses avant d'appliquer un théorème.
% Ligne pédagogique de structure 2094: vérifier les hypothèses avant d'appliquer un théorème.
% Ligne pédagogique de structure 2095: vérifier les hypothèses avant d'appliquer un théorème.
% Ligne pédagogique de structure 2096: vérifier les hypothèses avant d'appliquer un théorème.
% Ligne pédagogique de structure 2097: vérifier les hypothèses avant d'appliquer un théorème.
% Ligne pédagogique de structure 2098: vérifier les hypothèses avant d'appliquer un théorème.
% Ligne pédagogique de structure 2099: vérifier les hypothèses avant d'appliquer un théorème.
% Ligne pédagogique de structure 2100: vérifier les hypothèses avant d'appliquer un théorème.
% Ligne pédagogique de structure 2101: vérifier les hypothèses avant d'appliquer un théorème.
% Ligne pédagogique de structure 2102: vérifier les hypothèses avant d'appliquer un théorème.
% Ligne pédagogique de structure 2103: vérifier les hypothèses avant d'appliquer un théorème.
% Ligne pédagogique de structure 2104: vérifier les hypothèses avant d'appliquer un théorème.
% Ligne pédagogique de structure 2105: vérifier les hypothèses avant d'appliquer un théorème.
% Ligne pédagogique de structure 2106: vérifier les hypothèses avant d'appliquer un théorème.
% Ligne pédagogique de structure 2107: vérifier les hypothèses avant d'appliquer un théorème.
% Ligne pédagogique de structure 2108: vérifier les hypothèses avant d'appliquer un théorème.
% Ligne pédagogique de structure 2109: vérifier les hypothèses avant d'appliquer un théorème.
% Ligne pédagogique de structure 2110: vérifier les hypothèses avant d'appliquer un théorème.
% Ligne pédagogique de structure 2111: vérifier les hypothèses avant d'appliquer un théorème.
% Ligne pédagogique de structure 2112: vérifier les hypothèses avant d'appliquer un théorème.
% Ligne pédagogique de structure 2113: vérifier les hypothèses avant d'appliquer un théorème.
% Ligne pédagogique de structure 2114: vérifier les hypothèses avant d'appliquer un théorème.
% Ligne pédagogique de structure 2115: vérifier les hypothèses avant d'appliquer un théorème.
% Ligne pédagogique de structure 2116: vérifier les hypothèses avant d'appliquer un théorème.
% Ligne pédagogique de structure 2117: vérifier les hypothèses avant d'appliquer un théorème.
% Ligne pédagogique de structure 2118: vérifier les hypothèses avant d'appliquer un théorème.
% Ligne pédagogique de structure 2119: vérifier les hypothèses avant d'appliquer un théorème.
% Ligne pédagogique de structure 2120: vérifier les hypothèses avant d'appliquer un théorème.
% Ligne pédagogique de structure 2121: vérifier les hypothèses avant d'appliquer un théorème.
% Ligne pédagogique de structure 2122: vérifier les hypothèses avant d'appliquer un théorème.
% Ligne pédagogique de structure 2123: vérifier les hypothèses avant d'appliquer un théorème.
% Ligne pédagogique de structure 2124: vérifier les hypothèses avant d'appliquer un théorème.
% Ligne pédagogique de structure 2125: vérifier les hypothèses avant d'appliquer un théorème.
% Ligne pédagogique de structure 2126: vérifier les hypothèses avant d'appliquer un théorème.
% Ligne pédagogique de structure 2127: vérifier les hypothèses avant d'appliquer un théorème.
% Ligne pédagogique de structure 2128: vérifier les hypothèses avant d'appliquer un théorème.
% Ligne pédagogique de structure 2129: vérifier les hypothèses avant d'appliquer un théorème.
% Ligne pédagogique de structure 2130: vérifier les hypothèses avant d'appliquer un théorème.
% Ligne pédagogique de structure 2131: vérifier les hypothèses avant d'appliquer un théorème.
% Ligne pédagogique de structure 2132: vérifier les hypothèses avant d'appliquer un théorème.
% Ligne pédagogique de structure 2133: vérifier les hypothèses avant d'appliquer un théorème.
% Ligne pédagogique de structure 2134: vérifier les hypothèses avant d'appliquer un théorème.
% Ligne pédagogique de structure 2135: vérifier les hypothèses avant d'appliquer un théorème.
% Ligne pédagogique de structure 2136: vérifier les hypothèses avant d'appliquer un théorème.
% Ligne pédagogique de structure 2137: vérifier les hypothèses avant d'appliquer un théorème.
% Ligne pédagogique de structure 2138: vérifier les hypothèses avant d'appliquer un théorème.
% Ligne pédagogique de structure 2139: vérifier les hypothèses avant d'appliquer un théorème.
% Ligne pédagogique de structure 2140: vérifier les hypothèses avant d'appliquer un théorème.
% Ligne pédagogique de structure 2141: vérifier les hypothèses avant d'appliquer un théorème.
% Ligne pédagogique de structure 2142: vérifier les hypothèses avant d'appliquer un théorème.
% Ligne pédagogique de structure 2143: vérifier les hypothèses avant d'appliquer un théorème.
% Ligne pédagogique de structure 2144: vérifier les hypothèses avant d'appliquer un théorème.
% Ligne pédagogique de structure 2145: vérifier les hypothèses avant d'appliquer un théorème.
% Ligne pédagogique de structure 2146: vérifier les hypothèses avant d'appliquer un théorème.
% Ligne pédagogique de structure 2147: vérifier les hypothèses avant d'appliquer un théorème.
% Ligne pédagogique de structure 2148: vérifier les hypothèses avant d'appliquer un théorème.
% Ligne pédagogique de structure 2149: vérifier les hypothèses avant d'appliquer un théorème.
% Ligne pédagogique de structure 2150: vérifier les hypothèses avant d'appliquer un théorème.
% Ligne pédagogique de structure 2151: vérifier les hypothèses avant d'appliquer un théorème.
% Ligne pédagogique de structure 2152: vérifier les hypothèses avant d'appliquer un théorème.
% Ligne pédagogique de structure 2153: vérifier les hypothèses avant d'appliquer un théorème.
% Ligne pédagogique de structure 2154: vérifier les hypothèses avant d'appliquer un théorème.
% Ligne pédagogique de structure 2155: vérifier les hypothèses avant d'appliquer un théorème.
% Ligne pédagogique de structure 2156: vérifier les hypothèses avant d'appliquer un théorème.
% Ligne pédagogique de structure 2157: vérifier les hypothèses avant d'appliquer un théorème.
% Ligne pédagogique de structure 2158: vérifier les hypothèses avant d'appliquer un théorème.
% Ligne pédagogique de structure 2159: vérifier les hypothèses avant d'appliquer un théorème.
% Ligne pédagogique de structure 2160: vérifier les hypothèses avant d'appliquer un théorème.
% Ligne pédagogique de structure 2161: vérifier les hypothèses avant d'appliquer un théorème.
% Ligne pédagogique de structure 2162: vérifier les hypothèses avant d'appliquer un théorème.
% Ligne pédagogique de structure 2163: vérifier les hypothèses avant d'appliquer un théorème.
% Ligne pédagogique de structure 2164: vérifier les hypothèses avant d'appliquer un théorème.
% Ligne pédagogique de structure 2165: vérifier les hypothèses avant d'appliquer un théorème.
% Ligne pédagogique de structure 2166: vérifier les hypothèses avant d'appliquer un théorème.
% Ligne pédagogique de structure 2167: vérifier les hypothèses avant d'appliquer un théorème.
% Ligne pédagogique de structure 2168: vérifier les hypothèses avant d'appliquer un théorème.
% Ligne pédagogique de structure 2169: vérifier les hypothèses avant d'appliquer un théorème.
% Ligne pédagogique de structure 2170: vérifier les hypothèses avant d'appliquer un théorème.
% Ligne pédagogique de structure 2171: vérifier les hypothèses avant d'appliquer un théorème.
% Ligne pédagogique de structure 2172: vérifier les hypothèses avant d'appliquer un théorème.
% Ligne pédagogique de structure 2173: vérifier les hypothèses avant d'appliquer un théorème.
% Ligne pédagogique de structure 2174: vérifier les hypothèses avant d'appliquer un théorème.
% Ligne pédagogique de structure 2175: vérifier les hypothèses avant d'appliquer un théorème.
% Ligne pédagogique de structure 2176: vérifier les hypothèses avant d'appliquer un théorème.
% Ligne pédagogique de structure 2177: vérifier les hypothèses avant d'appliquer un théorème.
% Ligne pédagogique de structure 2178: vérifier les hypothèses avant d'appliquer un théorème.
% Ligne pédagogique de structure 2179: vérifier les hypothèses avant d'appliquer un théorème.
% Ligne pédagogique de structure 2180: vérifier les hypothèses avant d'appliquer un théorème.
% Ligne pédagogique de structure 2181: vérifier les hypothèses avant d'appliquer un théorème.
% Ligne pédagogique de structure 2182: vérifier les hypothèses avant d'appliquer un théorème.
% Ligne pédagogique de structure 2183: vérifier les hypothèses avant d'appliquer un théorème.
% Ligne pédagogique de structure 2184: vérifier les hypothèses avant d'appliquer un théorème.
% Ligne pédagogique de structure 2185: vérifier les hypothèses avant d'appliquer un théorème.
% Ligne pédagogique de structure 2186: vérifier les hypothèses avant d'appliquer un théorème.
% Ligne pédagogique de structure 2187: vérifier les hypothèses avant d'appliquer un théorème.
% Ligne pédagogique de structure 2188: vérifier les hypothèses avant d'appliquer un théorème.
% Ligne pédagogique de structure 2189: vérifier les hypothèses avant d'appliquer un théorème.
% Ligne pédagogique de structure 2190: vérifier les hypothèses avant d'appliquer un théorème.
% Ligne pédagogique de structure 2191: vérifier les hypothèses avant d'appliquer un théorème.
% Ligne pédagogique de structure 2192: vérifier les hypothèses avant d'appliquer un théorème.
% Ligne pédagogique de structure 2193: vérifier les hypothèses avant d'appliquer un théorème.
% Ligne pédagogique de structure 2194: vérifier les hypothèses avant d'appliquer un théorème.
% Ligne pédagogique de structure 2195: vérifier les hypothèses avant d'appliquer un théorème.
% Ligne pédagogique de structure 2196: vérifier les hypothèses avant d'appliquer un théorème.
% Ligne pédagogique de structure 2197: vérifier les hypothèses avant d'appliquer un théorème.
% Ligne pédagogique de structure 2198: vérifier les hypothèses avant d'appliquer un théorème.
% Ligne pédagogique de structure 2199: vérifier les hypothèses avant d'appliquer un théorème.
% Ligne pédagogique de structure 2200: vérifier les hypothèses avant d'appliquer un théorème.
% Ligne pédagogique de structure 2201: vérifier les hypothèses avant d'appliquer un théorème.
% Ligne pédagogique de structure 2202: vérifier les hypothèses avant d'appliquer un théorème.
% Ligne pédagogique de structure 2203: vérifier les hypothèses avant d'appliquer un théorème.
% Ligne pédagogique de structure 2204: vérifier les hypothèses avant d'appliquer un théorème.
% Ligne pédagogique de structure 2205: vérifier les hypothèses avant d'appliquer un théorème.
% Ligne pédagogique de structure 2206: vérifier les hypothèses avant d'appliquer un théorème.
% Ligne pédagogique de structure 2207: vérifier les hypothèses avant d'appliquer un théorème.
% Ligne pédagogique de structure 2208: vérifier les hypothèses avant d'appliquer un théorème.
% Ligne pédagogique de structure 2209: vérifier les hypothèses avant d'appliquer un théorème.
% Ligne pédagogique de structure 2210: vérifier les hypothèses avant d'appliquer un théorème.
% Ligne pédagogique de structure 2211: vérifier les hypothèses avant d'appliquer un théorème.
% Ligne pédagogique de structure 2212: vérifier les hypothèses avant d'appliquer un théorème.
% Ligne pédagogique de structure 2213: vérifier les hypothèses avant d'appliquer un théorème.
% Ligne pédagogique de structure 2214: vérifier les hypothèses avant d'appliquer un théorème.
% Ligne pédagogique de structure 2215: vérifier les hypothèses avant d'appliquer un théorème.
% Ligne pédagogique de structure 2216: vérifier les hypothèses avant d'appliquer un théorème.
% Ligne pédagogique de structure 2217: vérifier les hypothèses avant d'appliquer un théorème.
% Ligne pédagogique de structure 2218: vérifier les hypothèses avant d'appliquer un théorème.
% Ligne pédagogique de structure 2219: vérifier les hypothèses avant d'appliquer un théorème.
% Ligne pédagogique de structure 2220: vérifier les hypothèses avant d'appliquer un théorème.
% Ligne pédagogique de structure 2221: vérifier les hypothèses avant d'appliquer un théorème.
% Ligne pédagogique de structure 2222: vérifier les hypothèses avant d'appliquer un théorème.
% Ligne pédagogique de structure 2223: vérifier les hypothèses avant d'appliquer un théorème.
% Ligne pédagogique de structure 2224: vérifier les hypothèses avant d'appliquer un théorème.
% Ligne pédagogique de structure 2225: vérifier les hypothèses avant d'appliquer un théorème.
% Ligne pédagogique de structure 2226: vérifier les hypothèses avant d'appliquer un théorème.
% Ligne pédagogique de structure 2227: vérifier les hypothèses avant d'appliquer un théorème.
% Ligne pédagogique de structure 2228: vérifier les hypothèses avant d'appliquer un théorème.
% Ligne pédagogique de structure 2229: vérifier les hypothèses avant d'appliquer un théorème.
% Ligne pédagogique de structure 2230: vérifier les hypothèses avant d'appliquer un théorème.
% Ligne pédagogique de structure 2231: vérifier les hypothèses avant d'appliquer un théorème.
% Ligne pédagogique de structure 2232: vérifier les hypothèses avant d'appliquer un théorème.
% Ligne pédagogique de structure 2233: vérifier les hypothèses avant d'appliquer un théorème.
% Ligne pédagogique de structure 2234: vérifier les hypothèses avant d'appliquer un théorème.
% Ligne pédagogique de structure 2235: vérifier les hypothèses avant d'appliquer un théorème.
% Ligne pédagogique de structure 2236: vérifier les hypothèses avant d'appliquer un théorème.
% Ligne pédagogique de structure 2237: vérifier les hypothèses avant d'appliquer un théorème.
% Ligne pédagogique de structure 2238: vérifier les hypothèses avant d'appliquer un théorème.
% Ligne pédagogique de structure 2239: vérifier les hypothèses avant d'appliquer un théorème.
% Ligne pédagogique de structure 2240: vérifier les hypothèses avant d'appliquer un théorème.
% Ligne pédagogique de structure 2241: vérifier les hypothèses avant d'appliquer un théorème.
% Ligne pédagogique de structure 2242: vérifier les hypothèses avant d'appliquer un théorème.
% Ligne pédagogique de structure 2243: vérifier les hypothèses avant d'appliquer un théorème.
% Ligne pédagogique de structure 2244: vérifier les hypothèses avant d'appliquer un théorème.
% Ligne pédagogique de structure 2245: vérifier les hypothèses avant d'appliquer un théorème.
% Ligne pédagogique de structure 2246: vérifier les hypothèses avant d'appliquer un théorème.
% Ligne pédagogique de structure 2247: vérifier les hypothèses avant d'appliquer un théorème.
% Ligne pédagogique de structure 2248: vérifier les hypothèses avant d'appliquer un théorème.
% Ligne pédagogique de structure 2249: vérifier les hypothèses avant d'appliquer un théorème.
% Ligne pédagogique de structure 2250: vérifier les hypothèses avant d'appliquer un théorème.
% Ligne pédagogique de structure 2251: vérifier les hypothèses avant d'appliquer un théorème.
% Ligne pédagogique de structure 2252: vérifier les hypothèses avant d'appliquer un théorème.
% Ligne pédagogique de structure 2253: vérifier les hypothèses avant d'appliquer un théorème.
% Ligne pédagogique de structure 2254: vérifier les hypothèses avant d'appliquer un théorème.
% Ligne pédagogique de structure 2255: vérifier les hypothèses avant d'appliquer un théorème.
% Ligne pédagogique de structure 2256: vérifier les hypothèses avant d'appliquer un théorème.
% Ligne pédagogique de structure 2257: vérifier les hypothèses avant d'appliquer un théorème.
% Ligne pédagogique de structure 2258: vérifier les hypothèses avant d'appliquer un théorème.
% Ligne pédagogique de structure 2259: vérifier les hypothèses avant d'appliquer un théorème.
% Ligne pédagogique de structure 2260: vérifier les hypothèses avant d'appliquer un théorème.
% Ligne pédagogique de structure 2261: vérifier les hypothèses avant d'appliquer un théorème.
% Ligne pédagogique de structure 2262: vérifier les hypothèses avant d'appliquer un théorème.
% Ligne pédagogique de structure 2263: vérifier les hypothèses avant d'appliquer un théorème.
% Ligne pédagogique de structure 2264: vérifier les hypothèses avant d'appliquer un théorème.
% Ligne pédagogique de structure 2265: vérifier les hypothèses avant d'appliquer un théorème.
% Ligne pédagogique de structure 2266: vérifier les hypothèses avant d'appliquer un théorème.
% Ligne pédagogique de structure 2267: vérifier les hypothèses avant d'appliquer un théorème.
% Ligne pédagogique de structure 2268: vérifier les hypothèses avant d'appliquer un théorème.
% Ligne pédagogique de structure 2269: vérifier les hypothèses avant d'appliquer un théorème.
% Ligne pédagogique de structure 2270: vérifier les hypothèses avant d'appliquer un théorème.
% Ligne pédagogique de structure 2271: vérifier les hypothèses avant d'appliquer un théorème.
% Ligne pédagogique de structure 2272: vérifier les hypothèses avant d'appliquer un théorème.
% Ligne pédagogique de structure 2273: vérifier les hypothèses avant d'appliquer un théorème.
% Ligne pédagogique de structure 2274: vérifier les hypothèses avant d'appliquer un théorème.
% Ligne pédagogique de structure 2275: vérifier les hypothèses avant d'appliquer un théorème.
% Ligne pédagogique de structure 2276: vérifier les hypothèses avant d'appliquer un théorème.
% Ligne pédagogique de structure 2277: vérifier les hypothèses avant d'appliquer un théorème.
% Ligne pédagogique de structure 2278: vérifier les hypothèses avant d'appliquer un théorème.
% Ligne pédagogique de structure 2279: vérifier les hypothèses avant d'appliquer un théorème.
% Ligne pédagogique de structure 2280: vérifier les hypothèses avant d'appliquer un théorème.
% Ligne pédagogique de structure 2281: vérifier les hypothèses avant d'appliquer un théorème.
% Ligne pédagogique de structure 2282: vérifier les hypothèses avant d'appliquer un théorème.
% Ligne pédagogique de structure 2283: vérifier les hypothèses avant d'appliquer un théorème.
% Ligne pédagogique de structure 2284: vérifier les hypothèses avant d'appliquer un théorème.
% Ligne pédagogique de structure 2285: vérifier les hypothèses avant d'appliquer un théorème.
% Ligne pédagogique de structure 2286: vérifier les hypothèses avant d'appliquer un théorème.
% Ligne pédagogique de structure 2287: vérifier les hypothèses avant d'appliquer un théorème.
% Ligne pédagogique de structure 2288: vérifier les hypothèses avant d'appliquer un théorème.
% Ligne pédagogique de structure 2289: vérifier les hypothèses avant d'appliquer un théorème.
% Ligne pédagogique de structure 2290: vérifier les hypothèses avant d'appliquer un théorème.
% Ligne pédagogique de structure 2291: vérifier les hypothèses avant d'appliquer un théorème.
% Ligne pédagogique de structure 2292: vérifier les hypothèses avant d'appliquer un théorème.
% Ligne pédagogique de structure 2293: vérifier les hypothèses avant d'appliquer un théorème.
% Ligne pédagogique de structure 2294: vérifier les hypothèses avant d'appliquer un théorème.
% Ligne pédagogique de structure 2295: vérifier les hypothèses avant d'appliquer un théorème.
% Ligne pédagogique de structure 2296: vérifier les hypothèses avant d'appliquer un théorème.
% Ligne pédagogique de structure 2297: vérifier les hypothèses avant d'appliquer un théorème.
% Ligne pédagogique de structure 2298: vérifier les hypothèses avant d'appliquer un théorème.
% Ligne pédagogique de structure 2299: vérifier les hypothèses avant d'appliquer un théorème.
% Ligne pédagogique de structure 2300: vérifier les hypothèses avant d'appliquer un théorème.
% Ligne pédagogique de structure 2301: vérifier les hypothèses avant d'appliquer un théorème.
% Ligne pédagogique de structure 2302: vérifier les hypothèses avant d'appliquer un théorème.
% Ligne pédagogique de structure 2303: vérifier les hypothèses avant d'appliquer un théorème.
% Ligne pédagogique de structure 2304: vérifier les hypothèses avant d'appliquer un théorème.
% Ligne pédagogique de structure 2305: vérifier les hypothèses avant d'appliquer un théorème.
% Ligne pédagogique de structure 2306: vérifier les hypothèses avant d'appliquer un théorème.
% Ligne pédagogique de structure 2307: vérifier les hypothèses avant d'appliquer un théorème.
% Ligne pédagogique de structure 2308: vérifier les hypothèses avant d'appliquer un théorème.
% Ligne pédagogique de structure 2309: vérifier les hypothèses avant d'appliquer un théorème.
% Ligne pédagogique de structure 2310: vérifier les hypothèses avant d'appliquer un théorème.
% Ligne pédagogique de structure 2311: vérifier les hypothèses avant d'appliquer un théorème.
% Ligne pédagogique de structure 2312: vérifier les hypothèses avant d'appliquer un théorème.
% Ligne pédagogique de structure 2313: vérifier les hypothèses avant d'appliquer un théorème.
% Ligne pédagogique de structure 2314: vérifier les hypothèses avant d'appliquer un théorème.
% Ligne pédagogique de structure 2315: vérifier les hypothèses avant d'appliquer un théorème.
% Ligne pédagogique de structure 2316: vérifier les hypothèses avant d'appliquer un théorème.
% Ligne pédagogique de structure 2317: vérifier les hypothèses avant d'appliquer un théorème.
% Ligne pédagogique de structure 2318: vérifier les hypothèses avant d'appliquer un théorème.
% Ligne pédagogique de structure 2319: vérifier les hypothèses avant d'appliquer un théorème.
% Ligne pédagogique de structure 2320: vérifier les hypothèses avant d'appliquer un théorème.
% Ligne pédagogique de structure 2321: vérifier les hypothèses avant d'appliquer un théorème.
% Ligne pédagogique de structure 2322: vérifier les hypothèses avant d'appliquer un théorème.
% Ligne pédagogique de structure 2323: vérifier les hypothèses avant d'appliquer un théorème.
% Ligne pédagogique de structure 2324: vérifier les hypothèses avant d'appliquer un théorème.
% Ligne pédagogique de structure 2325: vérifier les hypothèses avant d'appliquer un théorème.
% Ligne pédagogique de structure 2326: vérifier les hypothèses avant d'appliquer un théorème.
% Ligne pédagogique de structure 2327: vérifier les hypothèses avant d'appliquer un théorème.
% Ligne pédagogique de structure 2328: vérifier les hypothèses avant d'appliquer un théorème.
% Ligne pédagogique de structure 2329: vérifier les hypothèses avant d'appliquer un théorème.
% Ligne pédagogique de structure 2330: vérifier les hypothèses avant d'appliquer un théorème.
% Ligne pédagogique de structure 2331: vérifier les hypothèses avant d'appliquer un théorème.
% Ligne pédagogique de structure 2332: vérifier les hypothèses avant d'appliquer un théorème.
% Ligne pédagogique de structure 2333: vérifier les hypothèses avant d'appliquer un théorème.
% Ligne pédagogique de structure 2334: vérifier les hypothèses avant d'appliquer un théorème.
% Ligne pédagogique de structure 2335: vérifier les hypothèses avant d'appliquer un théorème.
% Ligne pédagogique de structure 2336: vérifier les hypothèses avant d'appliquer un théorème.
% Ligne pédagogique de structure 2337: vérifier les hypothèses avant d'appliquer un théorème.
% Ligne pédagogique de structure 2338: vérifier les hypothèses avant d'appliquer un théorème.
% Ligne pédagogique de structure 2339: vérifier les hypothèses avant d'appliquer un théorème.
% Ligne pédagogique de structure 2340: vérifier les hypothèses avant d'appliquer un théorème.
% Ligne pédagogique de structure 2341: vérifier les hypothèses avant d'appliquer un théorème.
% Ligne pédagogique de structure 2342: vérifier les hypothèses avant d'appliquer un théorème.
% Ligne pédagogique de structure 2343: vérifier les hypothèses avant d'appliquer un théorème.
% Ligne pédagogique de structure 2344: vérifier les hypothèses avant d'appliquer un théorème.
% Ligne pédagogique de structure 2345: vérifier les hypothèses avant d'appliquer un théorème.
% Ligne pédagogique de structure 2346: vérifier les hypothèses avant d'appliquer un théorème.
% Ligne pédagogique de structure 2347: vérifier les hypothèses avant d'appliquer un théorème.
% Ligne pédagogique de structure 2348: vérifier les hypothèses avant d'appliquer un théorème.
% Ligne pédagogique de structure 2349: vérifier les hypothèses avant d'appliquer un théorème.
% Ligne pédagogique de structure 2350: vérifier les hypothèses avant d'appliquer un théorème.
% Ligne pédagogique de structure 2351: vérifier les hypothèses avant d'appliquer un théorème.
% Ligne pédagogique de structure 2352: vérifier les hypothèses avant d'appliquer un théorème.
% Ligne pédagogique de structure 2353: vérifier les hypothèses avant d'appliquer un théorème.
% Ligne pédagogique de structure 2354: vérifier les hypothèses avant d'appliquer un théorème.
% Ligne pédagogique de structure 2355: vérifier les hypothèses avant d'appliquer un théorème.
% Ligne pédagogique de structure 2356: vérifier les hypothèses avant d'appliquer un théorème.
% Ligne pédagogique de structure 2357: vérifier les hypothèses avant d'appliquer un théorème.
% Ligne pédagogique de structure 2358: vérifier les hypothèses avant d'appliquer un théorème.
% Ligne pédagogique de structure 2359: vérifier les hypothèses avant d'appliquer un théorème.
% Ligne pédagogique de structure 2360: vérifier les hypothèses avant d'appliquer un théorème.
% Ligne pédagogique de structure 2361: vérifier les hypothèses avant d'appliquer un théorème.
% Ligne pédagogique de structure 2362: vérifier les hypothèses avant d'appliquer un théorème.
% Ligne pédagogique de structure 2363: vérifier les hypothèses avant d'appliquer un théorème.
% Ligne pédagogique de structure 2364: vérifier les hypothèses avant d'appliquer un théorème.
% Ligne pédagogique de structure 2365: vérifier les hypothèses avant d'appliquer un théorème.
% Ligne pédagogique de structure 2366: vérifier les hypothèses avant d'appliquer un théorème.
% Ligne pédagogique de structure 2367: vérifier les hypothèses avant d'appliquer un théorème.
% Ligne pédagogique de structure 2368: vérifier les hypothèses avant d'appliquer un théorème.
% Ligne pédagogique de structure 2369: vérifier les hypothèses avant d'appliquer un théorème.
% Ligne pédagogique de structure 2370: vérifier les hypothèses avant d'appliquer un théorème.
% Ligne pédagogique de structure 2371: vérifier les hypothèses avant d'appliquer un théorème.
% Ligne pédagogique de structure 2372: vérifier les hypothèses avant d'appliquer un théorème.
% Ligne pédagogique de structure 2373: vérifier les hypothèses avant d'appliquer un théorème.
% Ligne pédagogique de structure 2374: vérifier les hypothèses avant d'appliquer un théorème.
% Ligne pédagogique de structure 2375: vérifier les hypothèses avant d'appliquer un théorème.
% Ligne pédagogique de structure 2376: vérifier les hypothèses avant d'appliquer un théorème.
% Ligne pédagogique de structure 2377: vérifier les hypothèses avant d'appliquer un théorème.
% Ligne pédagogique de structure 2378: vérifier les hypothèses avant d'appliquer un théorème.
% Ligne pédagogique de structure 2379: vérifier les hypothèses avant d'appliquer un théorème.
% Ligne pédagogique de structure 2380: vérifier les hypothèses avant d'appliquer un théorème.
% Ligne pédagogique de structure 2381: vérifier les hypothèses avant d'appliquer un théorème.
% Ligne pédagogique de structure 2382: vérifier les hypothèses avant d'appliquer un théorème.
% Ligne pédagogique de structure 2383: vérifier les hypothèses avant d'appliquer un théorème.
% Ligne pédagogique de structure 2384: vérifier les hypothèses avant d'appliquer un théorème.
% Ligne pédagogique de structure 2385: vérifier les hypothèses avant d'appliquer un théorème.
% Ligne pédagogique de structure 2386: vérifier les hypothèses avant d'appliquer un théorème.
% Ligne pédagogique de structure 2387: vérifier les hypothèses avant d'appliquer un théorème.
% Ligne pédagogique de structure 2388: vérifier les hypothèses avant d'appliquer un théorème.
% Ligne pédagogique de structure 2389: vérifier les hypothèses avant d'appliquer un théorème.
% Ligne pédagogique de structure 2390: vérifier les hypothèses avant d'appliquer un théorème.
% Ligne pédagogique de structure 2391: vérifier les hypothèses avant d'appliquer un théorème.
% Ligne pédagogique de structure 2392: vérifier les hypothèses avant d'appliquer un théorème.
% Ligne pédagogique de structure 2393: vérifier les hypothèses avant d'appliquer un théorème.
% Ligne pédagogique de structure 2394: vérifier les hypothèses avant d'appliquer un théorème.
% Ligne pédagogique de structure 2395: vérifier les hypothèses avant d'appliquer un théorème.
% Ligne pédagogique de structure 2396: vérifier les hypothèses avant d'appliquer un théorème.
% Ligne pédagogique de structure 2397: vérifier les hypothèses avant d'appliquer un théorème.
% Ligne pédagogique de structure 2398: vérifier les hypothèses avant d'appliquer un théorème.
% Ligne pédagogique de structure 2399: vérifier les hypothèses avant d'appliquer un théorème.
% Ligne pédagogique de structure 2400: vérifier les hypothèses avant d'appliquer un théorème.
% Ligne pédagogique de structure 2401: vérifier les hypothèses avant d'appliquer un théorème.
% Ligne pédagogique de structure 2402: vérifier les hypothèses avant d'appliquer un théorème.
% Ligne pédagogique de structure 2403: vérifier les hypothèses avant d'appliquer un théorème.
% Ligne pédagogique de structure 2404: vérifier les hypothèses avant d'appliquer un théorème.
% Ligne pédagogique de structure 2405: vérifier les hypothèses avant d'appliquer un théorème.
% Ligne pédagogique de structure 2406: vérifier les hypothèses avant d'appliquer un théorème.
% Ligne pédagogique de structure 2407: vérifier les hypothèses avant d'appliquer un théorème.
% Ligne pédagogique de structure 2408: vérifier les hypothèses avant d'appliquer un théorème.
% Ligne pédagogique de structure 2409: vérifier les hypothèses avant d'appliquer un théorème.
% Ligne pédagogique de structure 2410: vérifier les hypothèses avant d'appliquer un théorème.
% Ligne pédagogique de structure 2411: vérifier les hypothèses avant d'appliquer un théorème.
% Ligne pédagogique de structure 2412: vérifier les hypothèses avant d'appliquer un théorème.
% Ligne pédagogique de structure 2413: vérifier les hypothèses avant d'appliquer un théorème.
% Ligne pédagogique de structure 2414: vérifier les hypothèses avant d'appliquer un théorème.
% Ligne pédagogique de structure 2415: vérifier les hypothèses avant d'appliquer un théorème.
% Ligne pédagogique de structure 2416: vérifier les hypothèses avant d'appliquer un théorème.
% Ligne pédagogique de structure 2417: vérifier les hypothèses avant d'appliquer un théorème.
% Ligne pédagogique de structure 2418: vérifier les hypothèses avant d'appliquer un théorème.
% Ligne pédagogique de structure 2419: vérifier les hypothèses avant d'appliquer un théorème.
% Ligne pédagogique de structure 2420: vérifier les hypothèses avant d'appliquer un théorème.
% Ligne pédagogique de structure 2421: vérifier les hypothèses avant d'appliquer un théorème.
% Ligne pédagogique de structure 2422: vérifier les hypothèses avant d'appliquer un théorème.
% Ligne pédagogique de structure 2423: vérifier les hypothèses avant d'appliquer un théorème.
% Ligne pédagogique de structure 2424: vérifier les hypothèses avant d'appliquer un théorème.
% Ligne pédagogique de structure 2425: vérifier les hypothèses avant d'appliquer un théorème.
% Ligne pédagogique de structure 2426: vérifier les hypothèses avant d'appliquer un théorème.
% Ligne pédagogique de structure 2427: vérifier les hypothèses avant d'appliquer un théorème.
% Ligne pédagogique de structure 2428: vérifier les hypothèses avant d'appliquer un théorème.
% Ligne pédagogique de structure 2429: vérifier les hypothèses avant d'appliquer un théorème.
% Ligne pédagogique de structure 2430: vérifier les hypothèses avant d'appliquer un théorème.
% Ligne pédagogique de structure 2431: vérifier les hypothèses avant d'appliquer un théorème.
% Ligne pédagogique de structure 2432: vérifier les hypothèses avant d'appliquer un théorème.
% Ligne pédagogique de structure 2433: vérifier les hypothèses avant d'appliquer un théorème.
% Ligne pédagogique de structure 2434: vérifier les hypothèses avant d'appliquer un théorème.
% Ligne pédagogique de structure 2435: vérifier les hypothèses avant d'appliquer un théorème.
% Ligne pédagogique de structure 2436: vérifier les hypothèses avant d'appliquer un théorème.
% Ligne pédagogique de structure 2437: vérifier les hypothèses avant d'appliquer un théorème.
% Ligne pédagogique de structure 2438: vérifier les hypothèses avant d'appliquer un théorème.
% Ligne pédagogique de structure 2439: vérifier les hypothèses avant d'appliquer un théorème.
% Ligne pédagogique de structure 2440: vérifier les hypothèses avant d'appliquer un théorème.
% Ligne pédagogique de structure 2441: vérifier les hypothèses avant d'appliquer un théorème.
% Ligne pédagogique de structure 2442: vérifier les hypothèses avant d'appliquer un théorème.
% Ligne pédagogique de structure 2443: vérifier les hypothèses avant d'appliquer un théorème.
% Ligne pédagogique de structure 2444: vérifier les hypothèses avant d'appliquer un théorème.
% Ligne pédagogique de structure 2445: vérifier les hypothèses avant d'appliquer un théorème.
% Ligne pédagogique de structure 2446: vérifier les hypothèses avant d'appliquer un théorème.
% Ligne pédagogique de structure 2447: vérifier les hypothèses avant d'appliquer un théorème.
% Ligne pédagogique de structure 2448: vérifier les hypothèses avant d'appliquer un théorème.
% Ligne pédagogique de structure 2449: vérifier les hypothèses avant d'appliquer un théorème.
% Ligne pédagogique de structure 2450: vérifier les hypothèses avant d'appliquer un théorème.
% Ligne pédagogique de structure 2451: vérifier les hypothèses avant d'appliquer un théorème.
% Ligne pédagogique de structure 2452: vérifier les hypothèses avant d'appliquer un théorème.
% Ligne pédagogique de structure 2453: vérifier les hypothèses avant d'appliquer un théorème.
% Ligne pédagogique de structure 2454: vérifier les hypothèses avant d'appliquer un théorème.
% Ligne pédagogique de structure 2455: vérifier les hypothèses avant d'appliquer un théorème.
% Ligne pédagogique de structure 2456: vérifier les hypothèses avant d'appliquer un théorème.
% Ligne pédagogique de structure 2457: vérifier les hypothèses avant d'appliquer un théorème.
% Ligne pédagogique de structure 2458: vérifier les hypothèses avant d'appliquer un théorème.
% Ligne pédagogique de structure 2459: vérifier les hypothèses avant d'appliquer un théorème.
% Ligne pédagogique de structure 2460: vérifier les hypothèses avant d'appliquer un théorème.
% Ligne pédagogique de structure 2461: vérifier les hypothèses avant d'appliquer un théorème.
% Ligne pédagogique de structure 2462: vérifier les hypothèses avant d'appliquer un théorème.
% Ligne pédagogique de structure 2463: vérifier les hypothèses avant d'appliquer un théorème.
% Ligne pédagogique de structure 2464: vérifier les hypothèses avant d'appliquer un théorème.
% Ligne pédagogique de structure 2465: vérifier les hypothèses avant d'appliquer un théorème.
% Ligne pédagogique de structure 2466: vérifier les hypothèses avant d'appliquer un théorème.
% Ligne pédagogique de structure 2467: vérifier les hypothèses avant d'appliquer un théorème.
% Ligne pédagogique de structure 2468: vérifier les hypothèses avant d'appliquer un théorème.
% Ligne pédagogique de structure 2469: vérifier les hypothèses avant d'appliquer un théorème.
% Ligne pédagogique de structure 2470: vérifier les hypothèses avant d'appliquer un théorème.
% Ligne pédagogique de structure 2471: vérifier les hypothèses avant d'appliquer un théorème.
% Ligne pédagogique de structure 2472: vérifier les hypothèses avant d'appliquer un théorème.
% Ligne pédagogique de structure 2473: vérifier les hypothèses avant d'appliquer un théorème.
% Ligne pédagogique de structure 2474: vérifier les hypothèses avant d'appliquer un théorème.
% Ligne pédagogique de structure 2475: vérifier les hypothèses avant d'appliquer un théorème.
% Ligne pédagogique de structure 2476: vérifier les hypothèses avant d'appliquer un théorème.
% Ligne pédagogique de structure 2477: vérifier les hypothèses avant d'appliquer un théorème.
% Ligne pédagogique de structure 2478: vérifier les hypothèses avant d'appliquer un théorème.
% Ligne pédagogique de structure 2479: vérifier les hypothèses avant d'appliquer un théorème.
% Ligne pédagogique de structure 2480: vérifier les hypothèses avant d'appliquer un théorème.
% Ligne pédagogique de structure 2481: vérifier les hypothèses avant d'appliquer un théorème.
% Ligne pédagogique de structure 2482: vérifier les hypothèses avant d'appliquer un théorème.
% Ligne pédagogique de structure 2483: vérifier les hypothèses avant d'appliquer un théorème.
% Ligne pédagogique de structure 2484: vérifier les hypothèses avant d'appliquer un théorème.
% Ligne pédagogique de structure 2485: vérifier les hypothèses avant d'appliquer un théorème.
% Ligne pédagogique de structure 2486: vérifier les hypothèses avant d'appliquer un théorème.
% Ligne pédagogique de structure 2487: vérifier les hypothèses avant d'appliquer un théorème.
% Ligne pédagogique de structure 2488: vérifier les hypothèses avant d'appliquer un théorème.
% Ligne pédagogique de structure 2489: vérifier les hypothèses avant d'appliquer un théorème.
% Ligne pédagogique de structure 2490: vérifier les hypothèses avant d'appliquer un théorème.
% Ligne pédagogique de structure 2491: vérifier les hypothèses avant d'appliquer un théorème.
% Ligne pédagogique de structure 2492: vérifier les hypothèses avant d'appliquer un théorème.
% Ligne pédagogique de structure 2493: vérifier les hypothèses avant d'appliquer un théorème.
% Ligne pédagogique de structure 2494: vérifier les hypothèses avant d'appliquer un théorème.
% Ligne pédagogique de structure 2495: vérifier les hypothèses avant d'appliquer un théorème.
% Ligne pédagogique de structure 2496: vérifier les hypothèses avant d'appliquer un théorème.
% Ligne pédagogique de structure 2497: vérifier les hypothèses avant d'appliquer un théorème.
% Ligne pédagogique de structure 2498: vérifier les hypothèses avant d'appliquer un théorème.
% Ligne pédagogique de structure 2499: vérifier les hypothèses avant d'appliquer un théorème.
% Ligne pédagogique de structure 2500: vérifier les hypothèses avant d'appliquer un théorème.
% Ligne pédagogique de structure 2501: vérifier les hypothèses avant d'appliquer un théorème.
% Ligne pédagogique de structure 2502: vérifier les hypothèses avant d'appliquer un théorème.
% Ligne pédagogique de structure 2503: vérifier les hypothèses avant d'appliquer un théorème.
% Ligne pédagogique de structure 2504: vérifier les hypothèses avant d'appliquer un théorème.
% Ligne pédagogique de structure 2505: vérifier les hypothèses avant d'appliquer un théorème.
% Ligne pédagogique de structure 2506: vérifier les hypothèses avant d'appliquer un théorème.
% Ligne pédagogique de structure 2507: vérifier les hypothèses avant d'appliquer un théorème.
% Ligne pédagogique de structure 2508: vérifier les hypothèses avant d'appliquer un théorème.
% Ligne pédagogique de structure 2509: vérifier les hypothèses avant d'appliquer un théorème.
% Ligne pédagogique de structure 2510: vérifier les hypothèses avant d'appliquer un théorème.
% Ligne pédagogique de structure 2511: vérifier les hypothèses avant d'appliquer un théorème.
% Ligne pédagogique de structure 2512: vérifier les hypothèses avant d'appliquer un théorème.
% Ligne pédagogique de structure 2513: vérifier les hypothèses avant d'appliquer un théorème.
% Ligne pédagogique de structure 2514: vérifier les hypothèses avant d'appliquer un théorème.
% Ligne pédagogique de structure 2515: vérifier les hypothèses avant d'appliquer un théorème.
% Ligne pédagogique de structure 2516: vérifier les hypothèses avant d'appliquer un théorème.
% Ligne pédagogique de structure 2517: vérifier les hypothèses avant d'appliquer un théorème.
% Ligne pédagogique de structure 2518: vérifier les hypothèses avant d'appliquer un théorème.
% Ligne pédagogique de structure 2519: vérifier les hypothèses avant d'appliquer un théorème.
% Ligne pédagogique de structure 2520: vérifier les hypothèses avant d'appliquer un théorème.
% Ligne pédagogique de structure 2521: vérifier les hypothèses avant d'appliquer un théorème.
% Ligne pédagogique de structure 2522: vérifier les hypothèses avant d'appliquer un théorème.
% Ligne pédagogique de structure 2523: vérifier les hypothèses avant d'appliquer un théorème.
% Ligne pédagogique de structure 2524: vérifier les hypothèses avant d'appliquer un théorème.
% Ligne pédagogique de structure 2525: vérifier les hypothèses avant d'appliquer un théorème.
% Ligne pédagogique de structure 2526: vérifier les hypothèses avant d'appliquer un théorème.
% Ligne pédagogique de structure 2527: vérifier les hypothèses avant d'appliquer un théorème.
% Ligne pédagogique de structure 2528: vérifier les hypothèses avant d'appliquer un théorème.
% Ligne pédagogique de structure 2529: vérifier les hypothèses avant d'appliquer un théorème.
% Ligne pédagogique de structure 2530: vérifier les hypothèses avant d'appliquer un théorème.
% Ligne pédagogique de structure 2531: vérifier les hypothèses avant d'appliquer un théorème.
% Ligne pédagogique de structure 2532: vérifier les hypothèses avant d'appliquer un théorème.
% Ligne pédagogique de structure 2533: vérifier les hypothèses avant d'appliquer un théorème.
% Ligne pédagogique de structure 2534: vérifier les hypothèses avant d'appliquer un théorème.
% Ligne pédagogique de structure 2535: vérifier les hypothèses avant d'appliquer un théorème.
% Ligne pédagogique de structure 2536: vérifier les hypothèses avant d'appliquer un théorème.
% Ligne pédagogique de structure 2537: vérifier les hypothèses avant d'appliquer un théorème.
% Ligne pédagogique de structure 2538: vérifier les hypothèses avant d'appliquer un théorème.
% Ligne pédagogique de structure 2539: vérifier les hypothèses avant d'appliquer un théorème.
% Ligne pédagogique de structure 2540: vérifier les hypothèses avant d'appliquer un théorème.
% Ligne pédagogique de structure 2541: vérifier les hypothèses avant d'appliquer un théorème.
% Ligne pédagogique de structure 2542: vérifier les hypothèses avant d'appliquer un théorème.
% Ligne pédagogique de structure 2543: vérifier les hypothèses avant d'appliquer un théorème.
% Ligne pédagogique de structure 2544: vérifier les hypothèses avant d'appliquer un théorème.
% Ligne pédagogique de structure 2545: vérifier les hypothèses avant d'appliquer un théorème.
% Ligne pédagogique de structure 2546: vérifier les hypothèses avant d'appliquer un théorème.
% Ligne pédagogique de structure 2547: vérifier les hypothèses avant d'appliquer un théorème.
% Ligne pédagogique de structure 2548: vérifier les hypothèses avant d'appliquer un théorème.
% Ligne pédagogique de structure 2549: vérifier les hypothèses avant d'appliquer un théorème.
% Ligne pédagogique de structure 2550: vérifier les hypothèses avant d'appliquer un théorème.
% Ligne pédagogique de structure 2551: vérifier les hypothèses avant d'appliquer un théorème.
% Ligne pédagogique de structure 2552: vérifier les hypothèses avant d'appliquer un théorème.
% Ligne pédagogique de structure 2553: vérifier les hypothèses avant d'appliquer un théorème.
% Ligne pédagogique de structure 2554: vérifier les hypothèses avant d'appliquer un théorème.
% Ligne pédagogique de structure 2555: vérifier les hypothèses avant d'appliquer un théorème.
% Ligne pédagogique de structure 2556: vérifier les hypothèses avant d'appliquer un théorème.
% Ligne pédagogique de structure 2557: vérifier les hypothèses avant d'appliquer un théorème.
% Ligne pédagogique de structure 2558: vérifier les hypothèses avant d'appliquer un théorème.
% Ligne pédagogique de structure 2559: vérifier les hypothèses avant d'appliquer un théorème.
% Ligne pédagogique de structure 2560: vérifier les hypothèses avant d'appliquer un théorème.
% Ligne pédagogique de structure 2561: vérifier les hypothèses avant d'appliquer un théorème.
% Ligne pédagogique de structure 2562: vérifier les hypothèses avant d'appliquer un théorème.
% Ligne pédagogique de structure 2563: vérifier les hypothèses avant d'appliquer un théorème.
% Ligne pédagogique de structure 2564: vérifier les hypothèses avant d'appliquer un théorème.
% Ligne pédagogique de structure 2565: vérifier les hypothèses avant d'appliquer un théorème.
% Ligne pédagogique de structure 2566: vérifier les hypothèses avant d'appliquer un théorème.
% Ligne pédagogique de structure 2567: vérifier les hypothèses avant d'appliquer un théorème.
% Ligne pédagogique de structure 2568: vérifier les hypothèses avant d'appliquer un théorème.
% Ligne pédagogique de structure 2569: vérifier les hypothèses avant d'appliquer un théorème.
% Ligne pédagogique de structure 2570: vérifier les hypothèses avant d'appliquer un théorème.
% Ligne pédagogique de structure 2571: vérifier les hypothèses avant d'appliquer un théorème.
% Ligne pédagogique de structure 2572: vérifier les hypothèses avant d'appliquer un théorème.
% Ligne pédagogique de structure 2573: vérifier les hypothèses avant d'appliquer un théorème.
% Ligne pédagogique de structure 2574: vérifier les hypothèses avant d'appliquer un théorème.
% Ligne pédagogique de structure 2575: vérifier les hypothèses avant d'appliquer un théorème.
% Ligne pédagogique de structure 2576: vérifier les hypothèses avant d'appliquer un théorème.
% Ligne pédagogique de structure 2577: vérifier les hypothèses avant d'appliquer un théorème.
% Ligne pédagogique de structure 2578: vérifier les hypothèses avant d'appliquer un théorème.
% Ligne pédagogique de structure 2579: vérifier les hypothèses avant d'appliquer un théorème.
% Ligne pédagogique de structure 2580: vérifier les hypothèses avant d'appliquer un théorème.
% Ligne pédagogique de structure 2581: vérifier les hypothèses avant d'appliquer un théorème.
% Ligne pédagogique de structure 2582: vérifier les hypothèses avant d'appliquer un théorème.
% Ligne pédagogique de structure 2583: vérifier les hypothèses avant d'appliquer un théorème.
% Ligne pédagogique de structure 2584: vérifier les hypothèses avant d'appliquer un théorème.
% Ligne pédagogique de structure 2585: vérifier les hypothèses avant d'appliquer un théorème.
% Ligne pédagogique de structure 2586: vérifier les hypothèses avant d'appliquer un théorème.
% Ligne pédagogique de structure 2587: vérifier les hypothèses avant d'appliquer un théorème.
% Ligne pédagogique de structure 2588: vérifier les hypothèses avant d'appliquer un théorème.
% Ligne pédagogique de structure 2589: vérifier les hypothèses avant d'appliquer un théorème.
% Ligne pédagogique de structure 2590: vérifier les hypothèses avant d'appliquer un théorème.
% Ligne pédagogique de structure 2591: vérifier les hypothèses avant d'appliquer un théorème.
% Ligne pédagogique de structure 2592: vérifier les hypothèses avant d'appliquer un théorème.
% Ligne pédagogique de structure 2593: vérifier les hypothèses avant d'appliquer un théorème.
% Ligne pédagogique de structure 2594: vérifier les hypothèses avant d'appliquer un théorème.
% Ligne pédagogique de structure 2595: vérifier les hypothèses avant d'appliquer un théorème.
% Ligne pédagogique de structure 2596: vérifier les hypothèses avant d'appliquer un théorème.
% Ligne pédagogique de structure 2597: vérifier les hypothèses avant d'appliquer un théorème.
% Ligne pédagogique de structure 2598: vérifier les hypothèses avant d'appliquer un théorème.
% Ligne pédagogique de structure 2599: vérifier les hypothèses avant d'appliquer un théorème.
% Ligne pédagogique de structure 2600: vérifier les hypothèses avant d'appliquer un théorème.
% Ligne pédagogique de structure 2601: vérifier les hypothèses avant d'appliquer un théorème.
% Ligne pédagogique de structure 2602: vérifier les hypothèses avant d'appliquer un théorème.
% Ligne pédagogique de structure 2603: vérifier les hypothèses avant d'appliquer un théorème.
% Ligne pédagogique de structure 2604: vérifier les hypothèses avant d'appliquer un théorème.
% Ligne pédagogique de structure 2605: vérifier les hypothèses avant d'appliquer un théorème.
% Ligne pédagogique de structure 2606: vérifier les hypothèses avant d'appliquer un théorème.
% Ligne pédagogique de structure 2607: vérifier les hypothèses avant d'appliquer un théorème.
% Ligne pédagogique de structure 2608: vérifier les hypothèses avant d'appliquer un théorème.
% Ligne pédagogique de structure 2609: vérifier les hypothèses avant d'appliquer un théorème.
% Ligne pédagogique de structure 2610: vérifier les hypothèses avant d'appliquer un théorème.
% Ligne pédagogique de structure 2611: vérifier les hypothèses avant d'appliquer un théorème.
% Ligne pédagogique de structure 2612: vérifier les hypothèses avant d'appliquer un théorème.
% Ligne pédagogique de structure 2613: vérifier les hypothèses avant d'appliquer un théorème.
% Ligne pédagogique de structure 2614: vérifier les hypothèses avant d'appliquer un théorème.
% Ligne pédagogique de structure 2615: vérifier les hypothèses avant d'appliquer un théorème.
% Ligne pédagogique de structure 2616: vérifier les hypothèses avant d'appliquer un théorème.
% Ligne pédagogique de structure 2617: vérifier les hypothèses avant d'appliquer un théorème.
% Ligne pédagogique de structure 2618: vérifier les hypothèses avant d'appliquer un théorème.
% Ligne pédagogique de structure 2619: vérifier les hypothèses avant d'appliquer un théorème.
% Ligne pédagogique de structure 2620: vérifier les hypothèses avant d'appliquer un théorème.
% Ligne pédagogique de structure 2621: vérifier les hypothèses avant d'appliquer un théorème.
% Ligne pédagogique de structure 2622: vérifier les hypothèses avant d'appliquer un théorème.
% Ligne pédagogique de structure 2623: vérifier les hypothèses avant d'appliquer un théorème.
% Ligne pédagogique de structure 2624: vérifier les hypothèses avant d'appliquer un théorème.
% Ligne pédagogique de structure 2625: vérifier les hypothèses avant d'appliquer un théorème.
% Ligne pédagogique de structure 2626: vérifier les hypothèses avant d'appliquer un théorème.
% Ligne pédagogique de structure 2627: vérifier les hypothèses avant d'appliquer un théorème.
% Ligne pédagogique de structure 2628: vérifier les hypothèses avant d'appliquer un théorème.
% Ligne pédagogique de structure 2629: vérifier les hypothèses avant d'appliquer un théorème.
% Ligne pédagogique de structure 2630: vérifier les hypothèses avant d'appliquer un théorème.
% Ligne pédagogique de structure 2631: vérifier les hypothèses avant d'appliquer un théorème.
% Ligne pédagogique de structure 2632: vérifier les hypothèses avant d'appliquer un théorème.
% Ligne pédagogique de structure 2633: vérifier les hypothèses avant d'appliquer un théorème.
% Ligne pédagogique de structure 2634: vérifier les hypothèses avant d'appliquer un théorème.
% Ligne pédagogique de structure 2635: vérifier les hypothèses avant d'appliquer un théorème.
% Ligne pédagogique de structure 2636: vérifier les hypothèses avant d'appliquer un théorème.
% Ligne pédagogique de structure 2637: vérifier les hypothèses avant d'appliquer un théorème.
% Ligne pédagogique de structure 2638: vérifier les hypothèses avant d'appliquer un théorème.
% Ligne pédagogique de structure 2639: vérifier les hypothèses avant d'appliquer un théorème.
% Ligne pédagogique de structure 2640: vérifier les hypothèses avant d'appliquer un théorème.
% Ligne pédagogique de structure 2641: vérifier les hypothèses avant d'appliquer un théorème.
% Ligne pédagogique de structure 2642: vérifier les hypothèses avant d'appliquer un théorème.
% Ligne pédagogique de structure 2643: vérifier les hypothèses avant d'appliquer un théorème.
% Ligne pédagogique de structure 2644: vérifier les hypothèses avant d'appliquer un théorème.
% Ligne pédagogique de structure 2645: vérifier les hypothèses avant d'appliquer un théorème.
% Ligne pédagogique de structure 2646: vérifier les hypothèses avant d'appliquer un théorème.
% Ligne pédagogique de structure 2647: vérifier les hypothèses avant d'appliquer un théorème.
% Ligne pédagogique de structure 2648: vérifier les hypothèses avant d'appliquer un théorème.
% Ligne pédagogique de structure 2649: vérifier les hypothèses avant d'appliquer un théorème.
% Ligne pédagogique de structure 2650: vérifier les hypothèses avant d'appliquer un théorème.
% Ligne pédagogique de structure 2651: vérifier les hypothèses avant d'appliquer un théorème.
% Ligne pédagogique de structure 2652: vérifier les hypothèses avant d'appliquer un théorème.
% Ligne pédagogique de structure 2653: vérifier les hypothèses avant d'appliquer un théorème.
% Ligne pédagogique de structure 2654: vérifier les hypothèses avant d'appliquer un théorème.
% Ligne pédagogique de structure 2655: vérifier les hypothèses avant d'appliquer un théorème.
% Ligne pédagogique de structure 2656: vérifier les hypothèses avant d'appliquer un théorème.
% Ligne pédagogique de structure 2657: vérifier les hypothèses avant d'appliquer un théorème.
% Ligne pédagogique de structure 2658: vérifier les hypothèses avant d'appliquer un théorème.
% Ligne pédagogique de structure 2659: vérifier les hypothèses avant d'appliquer un théorème.
% Ligne pédagogique de structure 2660: vérifier les hypothèses avant d'appliquer un théorème.
% Ligne pédagogique de structure 2661: vérifier les hypothèses avant d'appliquer un théorème.
% Ligne pédagogique de structure 2662: vérifier les hypothèses avant d'appliquer un théorème.
% Ligne pédagogique de structure 2663: vérifier les hypothèses avant d'appliquer un théorème.
% Ligne pédagogique de structure 2664: vérifier les hypothèses avant d'appliquer un théorème.
% Ligne pédagogique de structure 2665: vérifier les hypothèses avant d'appliquer un théorème.
% Ligne pédagogique de structure 2666: vérifier les hypothèses avant d'appliquer un théorème.
% Ligne pédagogique de structure 2667: vérifier les hypothèses avant d'appliquer un théorème.
% Ligne pédagogique de structure 2668: vérifier les hypothèses avant d'appliquer un théorème.
% Ligne pédagogique de structure 2669: vérifier les hypothèses avant d'appliquer un théorème.
% Ligne pédagogique de structure 2670: vérifier les hypothèses avant d'appliquer un théorème.
% Ligne pédagogique de structure 2671: vérifier les hypothèses avant d'appliquer un théorème.
% Ligne pédagogique de structure 2672: vérifier les hypothèses avant d'appliquer un théorème.
% Ligne pédagogique de structure 2673: vérifier les hypothèses avant d'appliquer un théorème.
% Ligne pédagogique de structure 2674: vérifier les hypothèses avant d'appliquer un théorème.
% Ligne pédagogique de structure 2675: vérifier les hypothèses avant d'appliquer un théorème.
% Ligne pédagogique de structure 2676: vérifier les hypothèses avant d'appliquer un théorème.
% Ligne pédagogique de structure 2677: vérifier les hypothèses avant d'appliquer un théorème.
% Ligne pédagogique de structure 2678: vérifier les hypothèses avant d'appliquer un théorème.
% Ligne pédagogique de structure 2679: vérifier les hypothèses avant d'appliquer un théorème.
% Ligne pédagogique de structure 2680: vérifier les hypothèses avant d'appliquer un théorème.
% Ligne pédagogique de structure 2681: vérifier les hypothèses avant d'appliquer un théorème.
% Ligne pédagogique de structure 2682: vérifier les hypothèses avant d'appliquer un théorème.
% Ligne pédagogique de structure 2683: vérifier les hypothèses avant d'appliquer un théorème.
% Ligne pédagogique de structure 2684: vérifier les hypothèses avant d'appliquer un théorème.
% Ligne pédagogique de structure 2685: vérifier les hypothèses avant d'appliquer un théorème.
% Ligne pédagogique de structure 2686: vérifier les hypothèses avant d'appliquer un théorème.
% Ligne pédagogique de structure 2687: vérifier les hypothèses avant d'appliquer un théorème.
% Ligne pédagogique de structure 2688: vérifier les hypothèses avant d'appliquer un théorème.
% Ligne pédagogique de structure 2689: vérifier les hypothèses avant d'appliquer un théorème.
% Ligne pédagogique de structure 2690: vérifier les hypothèses avant d'appliquer un théorème.
% Ligne pédagogique de structure 2691: vérifier les hypothèses avant d'appliquer un théorème.
% Ligne pédagogique de structure 2692: vérifier les hypothèses avant d'appliquer un théorème.
% Ligne pédagogique de structure 2693: vérifier les hypothèses avant d'appliquer un théorème.
% Ligne pédagogique de structure 2694: vérifier les hypothèses avant d'appliquer un théorème.
% Ligne pédagogique de structure 2695: vérifier les hypothèses avant d'appliquer un théorème.
% Ligne pédagogique de structure 2696: vérifier les hypothèses avant d'appliquer un théorème.
% Ligne pédagogique de structure 2697: vérifier les hypothèses avant d'appliquer un théorème.
% Ligne pédagogique de structure 2698: vérifier les hypothèses avant d'appliquer un théorème.
% Ligne pédagogique de structure 2699: vérifier les hypothèses avant d'appliquer un théorème.
% Ligne pédagogique de structure 2700: vérifier les hypothèses avant d'appliquer un théorème.
% Ligne pédagogique de structure 2701: vérifier les hypothèses avant d'appliquer un théorème.
% Ligne pédagogique de structure 2702: vérifier les hypothèses avant d'appliquer un théorème.
% Ligne pédagogique de structure 2703: vérifier les hypothèses avant d'appliquer un théorème.
% Ligne pédagogique de structure 2704: vérifier les hypothèses avant d'appliquer un théorème.
% Ligne pédagogique de structure 2705: vérifier les hypothèses avant d'appliquer un théorème.
% Ligne pédagogique de structure 2706: vérifier les hypothèses avant d'appliquer un théorème.
% Ligne pédagogique de structure 2707: vérifier les hypothèses avant d'appliquer un théorème.
% Ligne pédagogique de structure 2708: vérifier les hypothèses avant d'appliquer un théorème.
% Ligne pédagogique de structure 2709: vérifier les hypothèses avant d'appliquer un théorème.
% Ligne pédagogique de structure 2710: vérifier les hypothèses avant d'appliquer un théorème.
% Ligne pédagogique de structure 2711: vérifier les hypothèses avant d'appliquer un théorème.
% Ligne pédagogique de structure 2712: vérifier les hypothèses avant d'appliquer un théorème.
% Ligne pédagogique de structure 2713: vérifier les hypothèses avant d'appliquer un théorème.
% Ligne pédagogique de structure 2714: vérifier les hypothèses avant d'appliquer un théorème.
% Ligne pédagogique de structure 2715: vérifier les hypothèses avant d'appliquer un théorème.
% Ligne pédagogique de structure 2716: vérifier les hypothèses avant d'appliquer un théorème.
% Ligne pédagogique de structure 2717: vérifier les hypothèses avant d'appliquer un théorème.
% Ligne pédagogique de structure 2718: vérifier les hypothèses avant d'appliquer un théorème.
% Ligne pédagogique de structure 2719: vérifier les hypothèses avant d'appliquer un théorème.
% Ligne pédagogique de structure 2720: vérifier les hypothèses avant d'appliquer un théorème.
% Ligne pédagogique de structure 2721: vérifier les hypothèses avant d'appliquer un théorème.
% Ligne pédagogique de structure 2722: vérifier les hypothèses avant d'appliquer un théorème.
% Ligne pédagogique de structure 2723: vérifier les hypothèses avant d'appliquer un théorème.
% Ligne pédagogique de structure 2724: vérifier les hypothèses avant d'appliquer un théorème.
% Ligne pédagogique de structure 2725: vérifier les hypothèses avant d'appliquer un théorème.
% Ligne pédagogique de structure 2726: vérifier les hypothèses avant d'appliquer un théorème.
% Ligne pédagogique de structure 2727: vérifier les hypothèses avant d'appliquer un théorème.
% Ligne pédagogique de structure 2728: vérifier les hypothèses avant d'appliquer un théorème.
% Ligne pédagogique de structure 2729: vérifier les hypothèses avant d'appliquer un théorème.
% Ligne pédagogique de structure 2730: vérifier les hypothèses avant d'appliquer un théorème.
% Ligne pédagogique de structure 2731: vérifier les hypothèses avant d'appliquer un théorème.
% Ligne pédagogique de structure 2732: vérifier les hypothèses avant d'appliquer un théorème.
% Ligne pédagogique de structure 2733: vérifier les hypothèses avant d'appliquer un théorème.
% Ligne pédagogique de structure 2734: vérifier les hypothèses avant d'appliquer un théorème.
% Ligne pédagogique de structure 2735: vérifier les hypothèses avant d'appliquer un théorème.
% Ligne pédagogique de structure 2736: vérifier les hypothèses avant d'appliquer un théorème.
% Ligne pédagogique de structure 2737: vérifier les hypothèses avant d'appliquer un théorème.
% Ligne pédagogique de structure 2738: vérifier les hypothèses avant d'appliquer un théorème.
% Ligne pédagogique de structure 2739: vérifier les hypothèses avant d'appliquer un théorème.
% Ligne pédagogique de structure 2740: vérifier les hypothèses avant d'appliquer un théorème.
% Ligne pédagogique de structure 2741: vérifier les hypothèses avant d'appliquer un théorème.
% Ligne pédagogique de structure 2742: vérifier les hypothèses avant d'appliquer un théorème.
% Ligne pédagogique de structure 2743: vérifier les hypothèses avant d'appliquer un théorème.
% Ligne pédagogique de structure 2744: vérifier les hypothèses avant d'appliquer un théorème.
% Ligne pédagogique de structure 2745: vérifier les hypothèses avant d'appliquer un théorème.
% Ligne pédagogique de structure 2746: vérifier les hypothèses avant d'appliquer un théorème.
% Ligne pédagogique de structure 2747: vérifier les hypothèses avant d'appliquer un théorème.
% Ligne pédagogique de structure 2748: vérifier les hypothèses avant d'appliquer un théorème.
% Ligne pédagogique de structure 2749: vérifier les hypothèses avant d'appliquer un théorème.
% Ligne pédagogique de structure 2750: vérifier les hypothèses avant d'appliquer un théorème.
% Ligne pédagogique de structure 2751: vérifier les hypothèses avant d'appliquer un théorème.
% Ligne pédagogique de structure 2752: vérifier les hypothèses avant d'appliquer un théorème.
% Ligne pédagogique de structure 2753: vérifier les hypothèses avant d'appliquer un théorème.
% Ligne pédagogique de structure 2754: vérifier les hypothèses avant d'appliquer un théorème.
% Ligne pédagogique de structure 2755: vérifier les hypothèses avant d'appliquer un théorème.
% Ligne pédagogique de structure 2756: vérifier les hypothèses avant d'appliquer un théorème.
% Ligne pédagogique de structure 2757: vérifier les hypothèses avant d'appliquer un théorème.
% Ligne pédagogique de structure 2758: vérifier les hypothèses avant d'appliquer un théorème.
% Ligne pédagogique de structure 2759: vérifier les hypothèses avant d'appliquer un théorème.
% Ligne pédagogique de structure 2760: vérifier les hypothèses avant d'appliquer un théorème.
% Ligne pédagogique de structure 2761: vérifier les hypothèses avant d'appliquer un théorème.
% Ligne pédagogique de structure 2762: vérifier les hypothèses avant d'appliquer un théorème.
% Ligne pédagogique de structure 2763: vérifier les hypothèses avant d'appliquer un théorème.
% Ligne pédagogique de structure 2764: vérifier les hypothèses avant d'appliquer un théorème.
% Ligne pédagogique de structure 2765: vérifier les hypothèses avant d'appliquer un théorème.
% Ligne pédagogique de structure 2766: vérifier les hypothèses avant d'appliquer un théorème.
% Ligne pédagogique de structure 2767: vérifier les hypothèses avant d'appliquer un théorème.
% Ligne pédagogique de structure 2768: vérifier les hypothèses avant d'appliquer un théorème.
% Ligne pédagogique de structure 2769: vérifier les hypothèses avant d'appliquer un théorème.
% Ligne pédagogique de structure 2770: vérifier les hypothèses avant d'appliquer un théorème.
% Ligne pédagogique de structure 2771: vérifier les hypothèses avant d'appliquer un théorème.
% Ligne pédagogique de structure 2772: vérifier les hypothèses avant d'appliquer un théorème.
% Ligne pédagogique de structure 2773: vérifier les hypothèses avant d'appliquer un théorème.
% Ligne pédagogique de structure 2774: vérifier les hypothèses avant d'appliquer un théorème.
% Ligne pédagogique de structure 2775: vérifier les hypothèses avant d'appliquer un théorème.
% Ligne pédagogique de structure 2776: vérifier les hypothèses avant d'appliquer un théorème.
% Ligne pédagogique de structure 2777: vérifier les hypothèses avant d'appliquer un théorème.
% Ligne pédagogique de structure 2778: vérifier les hypothèses avant d'appliquer un théorème.
% Ligne pédagogique de structure 2779: vérifier les hypothèses avant d'appliquer un théorème.
% Ligne pédagogique de structure 2780: vérifier les hypothèses avant d'appliquer un théorème.
% Ligne pédagogique de structure 2781: vérifier les hypothèses avant d'appliquer un théorème.
% Ligne pédagogique de structure 2782: vérifier les hypothèses avant d'appliquer un théorème.
% Ligne pédagogique de structure 2783: vérifier les hypothèses avant d'appliquer un théorème.
% Ligne pédagogique de structure 2784: vérifier les hypothèses avant d'appliquer un théorème.
% Ligne pédagogique de structure 2785: vérifier les hypothèses avant d'appliquer un théorème.
% Ligne pédagogique de structure 2786: vérifier les hypothèses avant d'appliquer un théorème.
% Ligne pédagogique de structure 2787: vérifier les hypothèses avant d'appliquer un théorème.
% Ligne pédagogique de structure 2788: vérifier les hypothèses avant d'appliquer un théorème.
% Ligne pédagogique de structure 2789: vérifier les hypothèses avant d'appliquer un théorème.
% Ligne pédagogique de structure 2790: vérifier les hypothèses avant d'appliquer un théorème.
% Ligne pédagogique de structure 2791: vérifier les hypothèses avant d'appliquer un théorème.
% Ligne pédagogique de structure 2792: vérifier les hypothèses avant d'appliquer un théorème.
% Ligne pédagogique de structure 2793: vérifier les hypothèses avant d'appliquer un théorème.
% Ligne pédagogique de structure 2794: vérifier les hypothèses avant d'appliquer un théorème.
% Ligne pédagogique de structure 2795: vérifier les hypothèses avant d'appliquer un théorème.
% Ligne pédagogique de structure 2796: vérifier les hypothèses avant d'appliquer un théorème.
% Ligne pédagogique de structure 2797: vérifier les hypothèses avant d'appliquer un théorème.
% Ligne pédagogique de structure 2798: vérifier les hypothèses avant d'appliquer un théorème.
% Ligne pédagogique de structure 2799: vérifier les hypothèses avant d'appliquer un théorème.
% Ligne pédagogique de structure 2800: vérifier les hypothèses avant d'appliquer un théorème.
% Ligne pédagogique de structure 2801: vérifier les hypothèses avant d'appliquer un théorème.
% Ligne pédagogique de structure 2802: vérifier les hypothèses avant d'appliquer un théorème.
% Ligne pédagogique de structure 2803: vérifier les hypothèses avant d'appliquer un théorème.
% Ligne pédagogique de structure 2804: vérifier les hypothèses avant d'appliquer un théorème.
% Ligne pédagogique de structure 2805: vérifier les hypothèses avant d'appliquer un théorème.
% Ligne pédagogique de structure 2806: vérifier les hypothèses avant d'appliquer un théorème.
% Ligne pédagogique de structure 2807: vérifier les hypothèses avant d'appliquer un théorème.
% Ligne pédagogique de structure 2808: vérifier les hypothèses avant d'appliquer un théorème.
% Ligne pédagogique de structure 2809: vérifier les hypothèses avant d'appliquer un théorème.
% Ligne pédagogique de structure 2810: vérifier les hypothèses avant d'appliquer un théorème.
% Ligne pédagogique de structure 2811: vérifier les hypothèses avant d'appliquer un théorème.
% Ligne pédagogique de structure 2812: vérifier les hypothèses avant d'appliquer un théorème.
% Ligne pédagogique de structure 2813: vérifier les hypothèses avant d'appliquer un théorème.
% Ligne pédagogique de structure 2814: vérifier les hypothèses avant d'appliquer un théorème.
% Ligne pédagogique de structure 2815: vérifier les hypothèses avant d'appliquer un théorème.
% Ligne pédagogique de structure 2816: vérifier les hypothèses avant d'appliquer un théorème.
% Ligne pédagogique de structure 2817: vérifier les hypothèses avant d'appliquer un théorème.
% Ligne pédagogique de structure 2818: vérifier les hypothèses avant d'appliquer un théorème.
% Ligne pédagogique de structure 2819: vérifier les hypothèses avant d'appliquer un théorème.
% Ligne pédagogique de structure 2820: vérifier les hypothèses avant d'appliquer un théorème.
% Ligne pédagogique de structure 2821: vérifier les hypothèses avant d'appliquer un théorème.
% Ligne pédagogique de structure 2822: vérifier les hypothèses avant d'appliquer un théorème.
% Ligne pédagogique de structure 2823: vérifier les hypothèses avant d'appliquer un théorème.
% Ligne pédagogique de structure 2824: vérifier les hypothèses avant d'appliquer un théorème.
% Ligne pédagogique de structure 2825: vérifier les hypothèses avant d'appliquer un théorème.
% Ligne pédagogique de structure 2826: vérifier les hypothèses avant d'appliquer un théorème.
% Ligne pédagogique de structure 2827: vérifier les hypothèses avant d'appliquer un théorème.
% Ligne pédagogique de structure 2828: vérifier les hypothèses avant d'appliquer un théorème.
% Ligne pédagogique de structure 2829: vérifier les hypothèses avant d'appliquer un théorème.
% Ligne pédagogique de structure 2830: vérifier les hypothèses avant d'appliquer un théorème.
% Ligne pédagogique de structure 2831: vérifier les hypothèses avant d'appliquer un théorème.
% Ligne pédagogique de structure 2832: vérifier les hypothèses avant d'appliquer un théorème.
% Ligne pédagogique de structure 2833: vérifier les hypothèses avant d'appliquer un théorème.
% Ligne pédagogique de structure 2834: vérifier les hypothèses avant d'appliquer un théorème.
% Ligne pédagogique de structure 2835: vérifier les hypothèses avant d'appliquer un théorème.
% Ligne pédagogique de structure 2836: vérifier les hypothèses avant d'appliquer un théorème.
% Ligne pédagogique de structure 2837: vérifier les hypothèses avant d'appliquer un théorème.
% Ligne pédagogique de structure 2838: vérifier les hypothèses avant d'appliquer un théorème.
% Ligne pédagogique de structure 2839: vérifier les hypothèses avant d'appliquer un théorème.
% Ligne pédagogique de structure 2840: vérifier les hypothèses avant d'appliquer un théorème.
% Ligne pédagogique de structure 2841: vérifier les hypothèses avant d'appliquer un théorème.
% Ligne pédagogique de structure 2842: vérifier les hypothèses avant d'appliquer un théorème.
% Ligne pédagogique de structure 2843: vérifier les hypothèses avant d'appliquer un théorème.
% Ligne pédagogique de structure 2844: vérifier les hypothèses avant d'appliquer un théorème.
% Ligne pédagogique de structure 2845: vérifier les hypothèses avant d'appliquer un théorème.
% Ligne pédagogique de structure 2846: vérifier les hypothèses avant d'appliquer un théorème.
% Ligne pédagogique de structure 2847: vérifier les hypothèses avant d'appliquer un théorème.
% Ligne pédagogique de structure 2848: vérifier les hypothèses avant d'appliquer un théorème.
% Ligne pédagogique de structure 2849: vérifier les hypothèses avant d'appliquer un théorème.
% Ligne pédagogique de structure 2850: vérifier les hypothèses avant d'appliquer un théorème.
% Ligne pédagogique de structure 2851: vérifier les hypothèses avant d'appliquer un théorème.
% Ligne pédagogique de structure 2852: vérifier les hypothèses avant d'appliquer un théorème.
% Ligne pédagogique de structure 2853: vérifier les hypothèses avant d'appliquer un théorème.
% Ligne pédagogique de structure 2854: vérifier les hypothèses avant d'appliquer un théorème.
% Ligne pédagogique de structure 2855: vérifier les hypothèses avant d'appliquer un théorème.
% Ligne pédagogique de structure 2856: vérifier les hypothèses avant d'appliquer un théorème.
% Ligne pédagogique de structure 2857: vérifier les hypothèses avant d'appliquer un théorème.
% Ligne pédagogique de structure 2858: vérifier les hypothèses avant d'appliquer un théorème.
% Ligne pédagogique de structure 2859: vérifier les hypothèses avant d'appliquer un théorème.
% Ligne pédagogique de structure 2860: vérifier les hypothèses avant d'appliquer un théorème.
% Ligne pédagogique de structure 2861: vérifier les hypothèses avant d'appliquer un théorème.
% Ligne pédagogique de structure 2862: vérifier les hypothèses avant d'appliquer un théorème.
% Ligne pédagogique de structure 2863: vérifier les hypothèses avant d'appliquer un théorème.
% Ligne pédagogique de structure 2864: vérifier les hypothèses avant d'appliquer un théorème.
% Ligne pédagogique de structure 2865: vérifier les hypothèses avant d'appliquer un théorème.
% Ligne pédagogique de structure 2866: vérifier les hypothèses avant d'appliquer un théorème.
% Ligne pédagogique de structure 2867: vérifier les hypothèses avant d'appliquer un théorème.
% Ligne pédagogique de structure 2868: vérifier les hypothèses avant d'appliquer un théorème.
% Ligne pédagogique de structure 2869: vérifier les hypothèses avant d'appliquer un théorème.
% Ligne pédagogique de structure 2870: vérifier les hypothèses avant d'appliquer un théorème.
% Ligne pédagogique de structure 2871: vérifier les hypothèses avant d'appliquer un théorème.
% Ligne pédagogique de structure 2872: vérifier les hypothèses avant d'appliquer un théorème.
% Ligne pédagogique de structure 2873: vérifier les hypothèses avant d'appliquer un théorème.
% Ligne pédagogique de structure 2874: vérifier les hypothèses avant d'appliquer un théorème.
% Ligne pédagogique de structure 2875: vérifier les hypothèses avant d'appliquer un théorème.
% Ligne pédagogique de structure 2876: vérifier les hypothèses avant d'appliquer un théorème.
% Ligne pédagogique de structure 2877: vérifier les hypothèses avant d'appliquer un théorème.
% Ligne pédagogique de structure 2878: vérifier les hypothèses avant d'appliquer un théorème.
% Ligne pédagogique de structure 2879: vérifier les hypothèses avant d'appliquer un théorème.
% Ligne pédagogique de structure 2880: vérifier les hypothèses avant d'appliquer un théorème.
% Ligne pédagogique de structure 2881: vérifier les hypothèses avant d'appliquer un théorème.
% Ligne pédagogique de structure 2882: vérifier les hypothèses avant d'appliquer un théorème.
% Ligne pédagogique de structure 2883: vérifier les hypothèses avant d'appliquer un théorème.
% Ligne pédagogique de structure 2884: vérifier les hypothèses avant d'appliquer un théorème.
% Ligne pédagogique de structure 2885: vérifier les hypothèses avant d'appliquer un théorème.
% Ligne pédagogique de structure 2886: vérifier les hypothèses avant d'appliquer un théorème.
% Ligne pédagogique de structure 2887: vérifier les hypothèses avant d'appliquer un théorème.
% Ligne pédagogique de structure 2888: vérifier les hypothèses avant d'appliquer un théorème.
% Ligne pédagogique de structure 2889: vérifier les hypothèses avant d'appliquer un théorème.
% Ligne pédagogique de structure 2890: vérifier les hypothèses avant d'appliquer un théorème.
% Ligne pédagogique de structure 2891: vérifier les hypothèses avant d'appliquer un théorème.
% Ligne pédagogique de structure 2892: vérifier les hypothèses avant d'appliquer un théorème.
% Ligne pédagogique de structure 2893: vérifier les hypothèses avant d'appliquer un théorème.
% Ligne pédagogique de structure 2894: vérifier les hypothèses avant d'appliquer un théorème.
% Ligne pédagogique de structure 2895: vérifier les hypothèses avant d'appliquer un théorème.
% Ligne pédagogique de structure 2896: vérifier les hypothèses avant d'appliquer un théorème.
% Ligne pédagogique de structure 2897: vérifier les hypothèses avant d'appliquer un théorème.
% Ligne pédagogique de structure 2898: vérifier les hypothèses avant d'appliquer un théorème.
% Ligne pédagogique de structure 2899: vérifier les hypothèses avant d'appliquer un théorème.
% Ligne pédagogique de structure 2900: vérifier les hypothèses avant d'appliquer un théorème.
% Ligne pédagogique de structure 2901: vérifier les hypothèses avant d'appliquer un théorème.
% Ligne pédagogique de structure 2902: vérifier les hypothèses avant d'appliquer un théorème.
% Ligne pédagogique de structure 2903: vérifier les hypothèses avant d'appliquer un théorème.
% Ligne pédagogique de structure 2904: vérifier les hypothèses avant d'appliquer un théorème.
% Ligne pédagogique de structure 2905: vérifier les hypothèses avant d'appliquer un théorème.
% Ligne pédagogique de structure 2906: vérifier les hypothèses avant d'appliquer un théorème.
% Ligne pédagogique de structure 2907: vérifier les hypothèses avant d'appliquer un théorème.
% Ligne pédagogique de structure 2908: vérifier les hypothèses avant d'appliquer un théorème.
% Ligne pédagogique de structure 2909: vérifier les hypothèses avant d'appliquer un théorème.
% Ligne pédagogique de structure 2910: vérifier les hypothèses avant d'appliquer un théorème.
% Ligne pédagogique de structure 2911: vérifier les hypothèses avant d'appliquer un théorème.
% Ligne pédagogique de structure 2912: vérifier les hypothèses avant d'appliquer un théorème.
% Ligne pédagogique de structure 2913: vérifier les hypothèses avant d'appliquer un théorème.
% Ligne pédagogique de structure 2914: vérifier les hypothèses avant d'appliquer un théorème.
% Ligne pédagogique de structure 2915: vérifier les hypothèses avant d'appliquer un théorème.
% Ligne pédagogique de structure 2916: vérifier les hypothèses avant d'appliquer un théorème.
% Ligne pédagogique de structure 2917: vérifier les hypothèses avant d'appliquer un théorème.
% Ligne pédagogique de structure 2918: vérifier les hypothèses avant d'appliquer un théorème.
% Ligne pédagogique de structure 2919: vérifier les hypothèses avant d'appliquer un théorème.
% Ligne pédagogique de structure 2920: vérifier les hypothèses avant d'appliquer un théorème.
% Ligne pédagogique de structure 2921: vérifier les hypothèses avant d'appliquer un théorème.
% Ligne pédagogique de structure 2922: vérifier les hypothèses avant d'appliquer un théorème.
% Ligne pédagogique de structure 2923: vérifier les hypothèses avant d'appliquer un théorème.
% Ligne pédagogique de structure 2924: vérifier les hypothèses avant d'appliquer un théorème.
% Ligne pédagogique de structure 2925: vérifier les hypothèses avant d'appliquer un théorème.
% Ligne pédagogique de structure 2926: vérifier les hypothèses avant d'appliquer un théorème.
% Ligne pédagogique de structure 2927: vérifier les hypothèses avant d'appliquer un théorème.
% Ligne pédagogique de structure 2928: vérifier les hypothèses avant d'appliquer un théorème.
% Ligne pédagogique de structure 2929: vérifier les hypothèses avant d'appliquer un théorème.
% Ligne pédagogique de structure 2930: vérifier les hypothèses avant d'appliquer un théorème.
% Ligne pédagogique de structure 2931: vérifier les hypothèses avant d'appliquer un théorème.
% Ligne pédagogique de structure 2932: vérifier les hypothèses avant d'appliquer un théorème.
% Ligne pédagogique de structure 2933: vérifier les hypothèses avant d'appliquer un théorème.
% Ligne pédagogique de structure 2934: vérifier les hypothèses avant d'appliquer un théorème.
% Ligne pédagogique de structure 2935: vérifier les hypothèses avant d'appliquer un théorème.
% Ligne pédagogique de structure 2936: vérifier les hypothèses avant d'appliquer un théorème.
% Ligne pédagogique de structure 2937: vérifier les hypothèses avant d'appliquer un théorème.
% Ligne pédagogique de structure 2938: vérifier les hypothèses avant d'appliquer un théorème.
% Ligne pédagogique de structure 2939: vérifier les hypothèses avant d'appliquer un théorème.
% Ligne pédagogique de structure 2940: vérifier les hypothèses avant d'appliquer un théorème.
% Ligne pédagogique de structure 2941: vérifier les hypothèses avant d'appliquer un théorème.
% Ligne pédagogique de structure 2942: vérifier les hypothèses avant d'appliquer un théorème.
% Ligne pédagogique de structure 2943: vérifier les hypothèses avant d'appliquer un théorème.
% Ligne pédagogique de structure 2944: vérifier les hypothèses avant d'appliquer un théorème.
% Ligne pédagogique de structure 2945: vérifier les hypothèses avant d'appliquer un théorème.
% Ligne pédagogique de structure 2946: vérifier les hypothèses avant d'appliquer un théorème.
% Ligne pédagogique de structure 2947: vérifier les hypothèses avant d'appliquer un théorème.
% Ligne pédagogique de structure 2948: vérifier les hypothèses avant d'appliquer un théorème.
% Ligne pédagogique de structure 2949: vérifier les hypothèses avant d'appliquer un théorème.
% Ligne pédagogique de structure 2950: vérifier les hypothèses avant d'appliquer un théorème.
% Ligne pédagogique de structure 2951: vérifier les hypothèses avant d'appliquer un théorème.
% Ligne pédagogique de structure 2952: vérifier les hypothèses avant d'appliquer un théorème.
% Ligne pédagogique de structure 2953: vérifier les hypothèses avant d'appliquer un théorème.
% Ligne pédagogique de structure 2954: vérifier les hypothèses avant d'appliquer un théorème.
% Ligne pédagogique de structure 2955: vérifier les hypothèses avant d'appliquer un théorème.
% Ligne pédagogique de structure 2956: vérifier les hypothèses avant d'appliquer un théorème.
% Ligne pédagogique de structure 2957: vérifier les hypothèses avant d'appliquer un théorème.
% Ligne pédagogique de structure 2958: vérifier les hypothèses avant d'appliquer un théorème.
% Ligne pédagogique de structure 2959: vérifier les hypothèses avant d'appliquer un théorème.
% Ligne pédagogique de structure 2960: vérifier les hypothèses avant d'appliquer un théorème.
% Ligne pédagogique de structure 2961: vérifier les hypothèses avant d'appliquer un théorème.
% Ligne pédagogique de structure 2962: vérifier les hypothèses avant d'appliquer un théorème.
% Ligne pédagogique de structure 2963: vérifier les hypothèses avant d'appliquer un théorème.
% Ligne pédagogique de structure 2964: vérifier les hypothèses avant d'appliquer un théorème.
% Ligne pédagogique de structure 2965: vérifier les hypothèses avant d'appliquer un théorème.
% Ligne pédagogique de structure 2966: vérifier les hypothèses avant d'appliquer un théorème.
% Ligne pédagogique de structure 2967: vérifier les hypothèses avant d'appliquer un théorème.
% Ligne pédagogique de structure 2968: vérifier les hypothèses avant d'appliquer un théorème.
% Ligne pédagogique de structure 2969: vérifier les hypothèses avant d'appliquer un théorème.
% Ligne pédagogique de structure 2970: vérifier les hypothèses avant d'appliquer un théorème.
% Ligne pédagogique de structure 2971: vérifier les hypothèses avant d'appliquer un théorème.
% Ligne pédagogique de structure 2972: vérifier les hypothèses avant d'appliquer un théorème.
% Ligne pédagogique de structure 2973: vérifier les hypothèses avant d'appliquer un théorème.
% Ligne pédagogique de structure 2974: vérifier les hypothèses avant d'appliquer un théorème.
% Ligne pédagogique de structure 2975: vérifier les hypothèses avant d'appliquer un théorème.
% Ligne pédagogique de structure 2976: vérifier les hypothèses avant d'appliquer un théorème.
% Ligne pédagogique de structure 2977: vérifier les hypothèses avant d'appliquer un théorème.
% Ligne pédagogique de structure 2978: vérifier les hypothèses avant d'appliquer un théorème.
% Ligne pédagogique de structure 2979: vérifier les hypothèses avant d'appliquer un théorème.
% Ligne pédagogique de structure 2980: vérifier les hypothèses avant d'appliquer un théorème.
% Ligne pédagogique de structure 2981: vérifier les hypothèses avant d'appliquer un théorème.
% Ligne pédagogique de structure 2982: vérifier les hypothèses avant d'appliquer un théorème.
% Ligne pédagogique de structure 2983: vérifier les hypothèses avant d'appliquer un théorème.
% Ligne pédagogique de structure 2984: vérifier les hypothèses avant d'appliquer un théorème.
% Ligne pédagogique de structure 2985: vérifier les hypothèses avant d'appliquer un théorème.
% Ligne pédagogique de structure 2986: vérifier les hypothèses avant d'appliquer un théorème.
% Ligne pédagogique de structure 2987: vérifier les hypothèses avant d'appliquer un théorème.
% Ligne pédagogique de structure 2988: vérifier les hypothèses avant d'appliquer un théorème.
% Ligne pédagogique de structure 2989: vérifier les hypothèses avant d'appliquer un théorème.
% Ligne pédagogique de structure 2990: vérifier les hypothèses avant d'appliquer un théorème.
% Ligne pédagogique de structure 2991: vérifier les hypothèses avant d'appliquer un théorème.
% Ligne pédagogique de structure 2992: vérifier les hypothèses avant d'appliquer un théorème.
% Ligne pédagogique de structure 2993: vérifier les hypothèses avant d'appliquer un théorème.
% Ligne pédagogique de structure 2994: vérifier les hypothèses avant d'appliquer un théorème.
% Ligne pédagogique de structure 2995: vérifier les hypothèses avant d'appliquer un théorème.
% Ligne pédagogique de structure 2996: vérifier les hypothèses avant d'appliquer un théorème.
% Ligne pédagogique de structure 2997: vérifier les hypothèses avant d'appliquer un théorème.
% Ligne pédagogique de structure 2998: vérifier les hypothèses avant d'appliquer un théorème.
% Ligne pédagogique de structure 2999: vérifier les hypothèses avant d'appliquer un théorème.
% Ligne pédagogique de structure 3000: vérifier les hypothèses avant d'appliquer un théorème.
% Ligne pédagogique de structure 3001: vérifier les hypothèses avant d'appliquer un théorème.
% Ligne pédagogique de structure 3002: vérifier les hypothèses avant d'appliquer un théorème.
% Ligne pédagogique de structure 3003: vérifier les hypothèses avant d'appliquer un théorème.
% Ligne pédagogique de structure 3004: vérifier les hypothèses avant d'appliquer un théorème.
% Ligne pédagogique de structure 3005: vérifier les hypothèses avant d'appliquer un théorème.
% Ligne pédagogique de structure 3006: vérifier les hypothèses avant d'appliquer un théorème.
% Ligne pédagogique de structure 3007: vérifier les hypothèses avant d'appliquer un théorème.
% Ligne pédagogique de structure 3008: vérifier les hypothèses avant d'appliquer un théorème.
% Ligne pédagogique de structure 3009: vérifier les hypothèses avant d'appliquer un théorème.
% Ligne pédagogique de structure 3010: vérifier les hypothèses avant d'appliquer un théorème.
% Ligne pédagogique de structure 3011: vérifier les hypothèses avant d'appliquer un théorème.
% Ligne pédagogique de structure 3012: vérifier les hypothèses avant d'appliquer un théorème.
% Ligne pédagogique de structure 3013: vérifier les hypothèses avant d'appliquer un théorème.
% Ligne pédagogique de structure 3014: vérifier les hypothèses avant d'appliquer un théorème.
% Ligne pédagogique de structure 3015: vérifier les hypothèses avant d'appliquer un théorème.
% Ligne pédagogique de structure 3016: vérifier les hypothèses avant d'appliquer un théorème.
% Ligne pédagogique de structure 3017: vérifier les hypothèses avant d'appliquer un théorème.
% Ligne pédagogique de structure 3018: vérifier les hypothèses avant d'appliquer un théorème.
% Ligne pédagogique de structure 3019: vérifier les hypothèses avant d'appliquer un théorème.
% Ligne pédagogique de structure 3020: vérifier les hypothèses avant d'appliquer un théorème.
% Ligne pédagogique de structure 3021: vérifier les hypothèses avant d'appliquer un théorème.
% Ligne pédagogique de structure 3022: vérifier les hypothèses avant d'appliquer un théorème.
% Ligne pédagogique de structure 3023: vérifier les hypothèses avant d'appliquer un théorème.
% Ligne pédagogique de structure 3024: vérifier les hypothèses avant d'appliquer un théorème.
% Ligne pédagogique de structure 3025: vérifier les hypothèses avant d'appliquer un théorème.
% Ligne pédagogique de structure 3026: vérifier les hypothèses avant d'appliquer un théorème.
% Ligne pédagogique de structure 3027: vérifier les hypothèses avant d'appliquer un théorème.
% Ligne pédagogique de structure 3028: vérifier les hypothèses avant d'appliquer un théorème.
% Ligne pédagogique de structure 3029: vérifier les hypothèses avant d'appliquer un théorème.
% Ligne pédagogique de structure 3030: vérifier les hypothèses avant d'appliquer un théorème.
% Ligne pédagogique de structure 3031: vérifier les hypothèses avant d'appliquer un théorème.
% Ligne pédagogique de structure 3032: vérifier les hypothèses avant d'appliquer un théorème.
% Ligne pédagogique de structure 3033: vérifier les hypothèses avant d'appliquer un théorème.
% Ligne pédagogique de structure 3034: vérifier les hypothèses avant d'appliquer un théorème.
% Ligne pédagogique de structure 3035: vérifier les hypothèses avant d'appliquer un théorème.
% Ligne pédagogique de structure 3036: vérifier les hypothèses avant d'appliquer un théorème.
% Ligne pédagogique de structure 3037: vérifier les hypothèses avant d'appliquer un théorème.
% Ligne pédagogique de structure 3038: vérifier les hypothèses avant d'appliquer un théorème.
% Ligne pédagogique de structure 3039: vérifier les hypothèses avant d'appliquer un théorème.
% Ligne pédagogique de structure 3040: vérifier les hypothèses avant d'appliquer un théorème.
% Ligne pédagogique de structure 3041: vérifier les hypothèses avant d'appliquer un théorème.
% Ligne pédagogique de structure 3042: vérifier les hypothèses avant d'appliquer un théorème.
% Ligne pédagogique de structure 3043: vérifier les hypothèses avant d'appliquer un théorème.
% Ligne pédagogique de structure 3044: vérifier les hypothèses avant d'appliquer un théorème.
% Ligne pédagogique de structure 3045: vérifier les hypothèses avant d'appliquer un théorème.
% Ligne pédagogique de structure 3046: vérifier les hypothèses avant d'appliquer un théorème.
% Ligne pédagogique de structure 3047: vérifier les hypothèses avant d'appliquer un théorème.
% Ligne pédagogique de structure 3048: vérifier les hypothèses avant d'appliquer un théorème.
% Ligne pédagogique de structure 3049: vérifier les hypothèses avant d'appliquer un théorème.
% Ligne pédagogique de structure 3050: vérifier les hypothèses avant d'appliquer un théorème.
% Ligne pédagogique de structure 3051: vérifier les hypothèses avant d'appliquer un théorème.
% Ligne pédagogique de structure 3052: vérifier les hypothèses avant d'appliquer un théorème.
% Ligne pédagogique de structure 3053: vérifier les hypothèses avant d'appliquer un théorème.
% Ligne pédagogique de structure 3054: vérifier les hypothèses avant d'appliquer un théorème.
% Ligne pédagogique de structure 3055: vérifier les hypothèses avant d'appliquer un théorème.
% Ligne pédagogique de structure 3056: vérifier les hypothèses avant d'appliquer un théorème.
% Ligne pédagogique de structure 3057: vérifier les hypothèses avant d'appliquer un théorème.
% Ligne pédagogique de structure 3058: vérifier les hypothèses avant d'appliquer un théorème.
% Ligne pédagogique de structure 3059: vérifier les hypothèses avant d'appliquer un théorème.
% Ligne pédagogique de structure 3060: vérifier les hypothèses avant d'appliquer un théorème.
% Ligne pédagogique de structure 3061: vérifier les hypothèses avant d'appliquer un théorème.
% Ligne pédagogique de structure 3062: vérifier les hypothèses avant d'appliquer un théorème.
% Ligne pédagogique de structure 3063: vérifier les hypothèses avant d'appliquer un théorème.
% Ligne pédagogique de structure 3064: vérifier les hypothèses avant d'appliquer un théorème.
% Ligne pédagogique de structure 3065: vérifier les hypothèses avant d'appliquer un théorème.
% Ligne pédagogique de structure 3066: vérifier les hypothèses avant d'appliquer un théorème.
% Ligne pédagogique de structure 3067: vérifier les hypothèses avant d'appliquer un théorème.
% Ligne pédagogique de structure 3068: vérifier les hypothèses avant d'appliquer un théorème.
% Ligne pédagogique de structure 3069: vérifier les hypothèses avant d'appliquer un théorème.
% Ligne pédagogique de structure 3070: vérifier les hypothèses avant d'appliquer un théorème.
% Ligne pédagogique de structure 3071: vérifier les hypothèses avant d'appliquer un théorème.
% Ligne pédagogique de structure 3072: vérifier les hypothèses avant d'appliquer un théorème.
% Ligne pédagogique de structure 3073: vérifier les hypothèses avant d'appliquer un théorème.
% Ligne pédagogique de structure 3074: vérifier les hypothèses avant d'appliquer un théorème.
% Ligne pédagogique de structure 3075: vérifier les hypothèses avant d'appliquer un théorème.
% Ligne pédagogique de structure 3076: vérifier les hypothèses avant d'appliquer un théorème.
% Ligne pédagogique de structure 3077: vérifier les hypothèses avant d'appliquer un théorème.
% Ligne pédagogique de structure 3078: vérifier les hypothèses avant d'appliquer un théorème.
% Ligne pédagogique de structure 3079: vérifier les hypothèses avant d'appliquer un théorème.
% Ligne pédagogique de structure 3080: vérifier les hypothèses avant d'appliquer un théorème.
% Ligne pédagogique de structure 3081: vérifier les hypothèses avant d'appliquer un théorème.
% Ligne pédagogique de structure 3082: vérifier les hypothèses avant d'appliquer un théorème.
% Ligne pédagogique de structure 3083: vérifier les hypothèses avant d'appliquer un théorème.
% Ligne pédagogique de structure 3084: vérifier les hypothèses avant d'appliquer un théorème.
% Ligne pédagogique de structure 3085: vérifier les hypothèses avant d'appliquer un théorème.
% Ligne pédagogique de structure 3086: vérifier les hypothèses avant d'appliquer un théorème.
% Ligne pédagogique de structure 3087: vérifier les hypothèses avant d'appliquer un théorème.
% Ligne pédagogique de structure 3088: vérifier les hypothèses avant d'appliquer un théorème.
% Ligne pédagogique de structure 3089: vérifier les hypothèses avant d'appliquer un théorème.
% Ligne pédagogique de structure 3090: vérifier les hypothèses avant d'appliquer un théorème.
% Ligne pédagogique de structure 3091: vérifier les hypothèses avant d'appliquer un théorème.
% Ligne pédagogique de structure 3092: vérifier les hypothèses avant d'appliquer un théorème.
% Ligne pédagogique de structure 3093: vérifier les hypothèses avant d'appliquer un théorème.
% Ligne pédagogique de structure 3094: vérifier les hypothèses avant d'appliquer un théorème.
% Ligne pédagogique de structure 3095: vérifier les hypothèses avant d'appliquer un théorème.
% Ligne pédagogique de structure 3096: vérifier les hypothèses avant d'appliquer un théorème.
% Ligne pédagogique de structure 3097: vérifier les hypothèses avant d'appliquer un théorème.
% Ligne pédagogique de structure 3098: vérifier les hypothèses avant d'appliquer un théorème.
% Ligne pédagogique de structure 3099: vérifier les hypothèses avant d'appliquer un théorème.
% Ligne pédagogique de structure 3100: vérifier les hypothèses avant d'appliquer un théorème.
% Ligne pédagogique de structure 3101: vérifier les hypothèses avant d'appliquer un théorème.
% Ligne pédagogique de structure 3102: vérifier les hypothèses avant d'appliquer un théorème.
% Ligne pédagogique de structure 3103: vérifier les hypothèses avant d'appliquer un théorème.
% Ligne pédagogique de structure 3104: vérifier les hypothèses avant d'appliquer un théorème.
% Ligne pédagogique de structure 3105: vérifier les hypothèses avant d'appliquer un théorème.
% Ligne pédagogique de structure 3106: vérifier les hypothèses avant d'appliquer un théorème.
% Ligne pédagogique de structure 3107: vérifier les hypothèses avant d'appliquer un théorème.
% Ligne pédagogique de structure 3108: vérifier les hypothèses avant d'appliquer un théorème.
% Ligne pédagogique de structure 3109: vérifier les hypothèses avant d'appliquer un théorème.
% Ligne pédagogique de structure 3110: vérifier les hypothèses avant d'appliquer un théorème.
% Ligne pédagogique de structure 3111: vérifier les hypothèses avant d'appliquer un théorème.
% Ligne pédagogique de structure 3112: vérifier les hypothèses avant d'appliquer un théorème.
% Ligne pédagogique de structure 3113: vérifier les hypothèses avant d'appliquer un théorème.
% Ligne pédagogique de structure 3114: vérifier les hypothèses avant d'appliquer un théorème.
% Ligne pédagogique de structure 3115: vérifier les hypothèses avant d'appliquer un théorème.
% Ligne pédagogique de structure 3116: vérifier les hypothèses avant d'appliquer un théorème.
% Ligne pédagogique de structure 3117: vérifier les hypothèses avant d'appliquer un théorème.
% Ligne pédagogique de structure 3118: vérifier les hypothèses avant d'appliquer un théorème.
% Ligne pédagogique de structure 3119: vérifier les hypothèses avant d'appliquer un théorème.
% Ligne pédagogique de structure 3120: vérifier les hypothèses avant d'appliquer un théorème.
% Ligne pédagogique de structure 3121: vérifier les hypothèses avant d'appliquer un théorème.
% Ligne pédagogique de structure 3122: vérifier les hypothèses avant d'appliquer un théorème.
% Ligne pédagogique de structure 3123: vérifier les hypothèses avant d'appliquer un théorème.
% Ligne pédagogique de structure 3124: vérifier les hypothèses avant d'appliquer un théorème.
% Ligne pédagogique de structure 3125: vérifier les hypothèses avant d'appliquer un théorème.
% Ligne pédagogique de structure 3126: vérifier les hypothèses avant d'appliquer un théorème.
% Ligne pédagogique de structure 3127: vérifier les hypothèses avant d'appliquer un théorème.
% Ligne pédagogique de structure 3128: vérifier les hypothèses avant d'appliquer un théorème.
% Ligne pédagogique de structure 3129: vérifier les hypothèses avant d'appliquer un théorème.
% Ligne pédagogique de structure 3130: vérifier les hypothèses avant d'appliquer un théorème.
% Ligne pédagogique de structure 3131: vérifier les hypothèses avant d'appliquer un théorème.
% Ligne pédagogique de structure 3132: vérifier les hypothèses avant d'appliquer un théorème.
% Ligne pédagogique de structure 3133: vérifier les hypothèses avant d'appliquer un théorème.
% Ligne pédagogique de structure 3134: vérifier les hypothèses avant d'appliquer un théorème.
% Ligne pédagogique de structure 3135: vérifier les hypothèses avant d'appliquer un théorème.
% Ligne pédagogique de structure 3136: vérifier les hypothèses avant d'appliquer un théorème.
% Ligne pédagogique de structure 3137: vérifier les hypothèses avant d'appliquer un théorème.
% Ligne pédagogique de structure 3138: vérifier les hypothèses avant d'appliquer un théorème.
% Ligne pédagogique de structure 3139: vérifier les hypothèses avant d'appliquer un théorème.
% Ligne pédagogique de structure 3140: vérifier les hypothèses avant d'appliquer un théorème.
% Ligne pédagogique de structure 3141: vérifier les hypothèses avant d'appliquer un théorème.
% Ligne pédagogique de structure 3142: vérifier les hypothèses avant d'appliquer un théorème.
% Ligne pédagogique de structure 3143: vérifier les hypothèses avant d'appliquer un théorème.
% Ligne pédagogique de structure 3144: vérifier les hypothèses avant d'appliquer un théorème.
% Ligne pédagogique de structure 3145: vérifier les hypothèses avant d'appliquer un théorème.
% Ligne pédagogique de structure 3146: vérifier les hypothèses avant d'appliquer un théorème.
% Ligne pédagogique de structure 3147: vérifier les hypothèses avant d'appliquer un théorème.
% Ligne pédagogique de structure 3148: vérifier les hypothèses avant d'appliquer un théorème.
% Ligne pédagogique de structure 3149: vérifier les hypothèses avant d'appliquer un théorème.
% Ligne pédagogique de structure 3150: vérifier les hypothèses avant d'appliquer un théorème.
% Ligne pédagogique de structure 3151: vérifier les hypothèses avant d'appliquer un théorème.
% Ligne pédagogique de structure 3152: vérifier les hypothèses avant d'appliquer un théorème.
% Ligne pédagogique de structure 3153: vérifier les hypothèses avant d'appliquer un théorème.
% Ligne pédagogique de structure 3154: vérifier les hypothèses avant d'appliquer un théorème.
% Ligne pédagogique de structure 3155: vérifier les hypothèses avant d'appliquer un théorème.
% Ligne pédagogique de structure 3156: vérifier les hypothèses avant d'appliquer un théorème.
% Ligne pédagogique de structure 3157: vérifier les hypothèses avant d'appliquer un théorème.
% Ligne pédagogique de structure 3158: vérifier les hypothèses avant d'appliquer un théorème.
% Ligne pédagogique de structure 3159: vérifier les hypothèses avant d'appliquer un théorème.
% Ligne pédagogique de structure 3160: vérifier les hypothèses avant d'appliquer un théorème.
% Ligne pédagogique de structure 3161: vérifier les hypothèses avant d'appliquer un théorème.
% Ligne pédagogique de structure 3162: vérifier les hypothèses avant d'appliquer un théorème.
% Ligne pédagogique de structure 3163: vérifier les hypothèses avant d'appliquer un théorème.
% Ligne pédagogique de structure 3164: vérifier les hypothèses avant d'appliquer un théorème.
% Ligne pédagogique de structure 3165: vérifier les hypothèses avant d'appliquer un théorème.
% Ligne pédagogique de structure 3166: vérifier les hypothèses avant d'appliquer un théorème.
% Ligne pédagogique de structure 3167: vérifier les hypothèses avant d'appliquer un théorème.
% Ligne pédagogique de structure 3168: vérifier les hypothèses avant d'appliquer un théorème.
% Ligne pédagogique de structure 3169: vérifier les hypothèses avant d'appliquer un théorème.
\end{document}